% !TeX program = xelatex
% 《Bochner-Riesz 平均讲义》中文译本
% 原书: K. M. Davis & Y.-C. Chang, Lectures on Bochner-Riesz Means,
%       LMS Lecture Note Series 114, Cambridge University Press, 1987.
\documentclass[
    ThemeColor=MidnightBlue,
    PaperSize=a4paper,
    BookTitleCN={Bochner--Riesz 平均讲义},
    BookTitleEN={Lectures on Bochner-Riesz Means},
    OriginalAuthor={Katherine Michelle Davis, Yang-Chun Chang},
    OriginalEdition={LMS Lecture Note Series 114},
    OriginalPublisher={Cambridge University Press},
    OriginalYear={1987},
    Translator={数译},
    ModelUsed={GLM-5.3-Flash},
    TranslationDate={2026年9月2日},
    BibStyle={numeric-comp}
]{mathtranslation}

% ---------- 参考文献bibLaTeX ----------
% biblatex 由 mathtranslation.cls v3.1 统一加载
%   (backend=biber, style=\MT@BibStyle, sorting=nyt,
%    giveninits=true, maxalphanames=2, maxbibnames=99)
% 本书只需补上 backref 并挂上资源文件；此处再 \usepackage 会触发选项冲突。
\ExecuteBibliographyOptions{backref=true, backrefstyle=three}
\addbibresource{refs.bib}

% ---------- 索引 ----------
\usepackage[xindy]{imakeidx}
\makeindex[
    columns=2,
    program=xindy,
    intoc=true,
    options=-M texindy -I xelatex -C utf8,
    title={名词索引暨英译}]

% ---------- 定理编号: 与原书一致, 各环境共用每章一个计数器 ----------
% 原书编号形如 0.1, 2.14, 8.5 (章.序号), 定理/引理/命题/定义/注/例共用一个序号列.
\makeatletter
\newcounter{mtthm}[chapter]
\renewcommand\themtthm{\thechapter.\arabic{mtthm}}
% v3.1 把定理类环境注册为 [section] 计数（旧版为 [chapter]）。本书要求
% 各环境共用“章.序号”的同一个序号列，故先解除按节重置，再把显示切到
% \themtthm；寄存器本身仍沿用类文件声明，amsthm 样式与 cref 名称照旧。
\@removefromreset{theorem}{section}
\@removefromreset{definition}{section}
\@removefromreset{lemma}{section}
\@removefromreset{corollary}{section}
\@removefromreset{proposition}{section}
\@removefromreset{remark}{section}
\@removefromreset{example}{section}
\let\c@theorem\c@mtthm
\let\c@definition\c@mtthm
\let\c@lemma\c@mtthm
\let\c@corollary\c@mtthm
\let\c@proposition\c@mtthm
\let\c@remark\c@mtthm
\let\c@example\c@mtthm
\renewcommand{\thetheorem}{\themtthm}
\renewcommand{\thedefinition}{\themtthm}
\renewcommand{\thelemma}{\themtthm}
\renewcommand{\thecorollary}{\themtthm}
\renewcommand{\theproposition}{\themtthm}
\renewcommand{\theremark}{\themtthm}
\renewcommand{\theexample}{\themtthm}
\makeatother

% 原书第 5 章使用的 FACT 环境 (编号 5.0 起)
\newtheorem{fact}{事实}[chapter]
\makeatletter
\let\c@fact\c@mtthm
\makeatother

% 原书第 8 章使用的 NOTATION 环境
\newtheorem{ntn}{记号}[chapter]
\makeatletter
\let\c@ntn\c@mtthm
\makeatother
\newenvironment{notation}{\begin{ntn}}{\end{ntn}}
\crefname{ntn}{记号}{记号}

% ---------- 原书带编号的短论标题 (如 2.5, 2.21), 占用条目计数器 ----------
\newcommand{\essayhead}[2]{%
  \par\vspace{1.2em}%
  \refstepcounter{mtthm}\label{#1}%
  {\noindent{\headfont \themtthm\quad #2}\par}%
  \vspace{0.4em}\nobreak}

% 原书正文中出现的引用编号在本书参考文献表中按字母序重新编号,
% 引用均由 \cite 自动生成, 并附 backref.

% ---------- 封面 ----------
% v3.1 的 \makecover 已按制作规范在出版社信息下方印出“译者 xx 翻译及重排”，
% 旧版针对 v1.2 封面宏的 patchcmd 已无标的，故整体移除。

% ---------- 其他 ----------
% headheight 由 mathtranslation.cls 的 geometry 统一设为 15pt，不再单独覆盖。
\usetikzlibrary{patterns}
\usepackage{enumitem}
\allowdisplaybreaks

% 图编号与原书一致: 原书全书按出现顺序连续编号 (图 1-16), 不随章重置.
% \counterwithout 现由 LaTeX 内核提供，无需再加载 chngcntr。
\counterwithout{figure}{chapter}
% 中文正文段落在个别行内公式/西文连字符处缺少可行断点时会轻微超出版心,
% \emergencystretch 仅在需要时给段落提供额外伸张, 不改动文字与标点.
\emergencystretch=1.5em

\begin{document}

\makecover
% 封面之后的空白页 (封面背面) 若沿用页面样式会显示孤立的页码,
% 此页不设页眉页脚, 使其真正留白.
\thispagestyle{empty}

\frontmatter
% 前言 (原书 Introduction, 第 ix 页)
\chapter*{前言}
\markboth{前言}{前言}
\addcontentsline{toc}{chapter}{前言}

\begin{center}
\itshape 献给 John——他始终不离不弃, 陪伴我们历经变迁.
\end{center}

\vspace{2em}

本书旨在对 Bochner--Riesz 平均的几何理论作一个自足的阐述. 这一主题处理 Fourier 分析中最基本的问题: Fourier 级数何时收敛到原来的函数? 二十世纪七十年代中期, 这一问题取得了实质性的进展, 但相关技巧至今仍只能见于专业文献. 我们的意图是给出一个研究生能够读懂的论述. 为了保持叙述的一致性, 我们对某些重要课题作了舍弃. 我们假定读者熟悉研究生水平的实分析, 并熟悉分布理论与 Fourier 变换的基本事实. 一本基本的参考书是 Stein 与 Weiss 的 \emph{Introduction to Fourier Analysis on Euclidean Spaces}~\cite{steinweiss}, 以及 Rudin 的两部著作\footnote{译者注: 原书此处引用编号重复作 [43] 与 [43], 应为 [42] 与 [43].}~\cite{rudinrca,rudinfa}.

在写作本书的过程中, 多年来我们从与同事们的广泛交谈中获益匪浅. 我们感谢 E.~Fabes, R.~Fefferman 与 E.~M.~Stein 教授对第~1、2、3 章材料的帮助. 第~4、5 章的内容受到了与 A.~W.~Knapp 教授交谈的影响. 对于第~7、8 章的一般思想, 我们感激 A.~Córdoba 与 C.~Fefferman 教授. 本书初稿写作期间, 第一作者得到了美国国家科学基金会资助 (NSF grants MCS 8202165 与 8001799). 第一作者感谢 W.~Beiglböck 教授与海德堡大学应用数学研究所, 使他有机会完成本书的倒数第二稿. 第二作者感谢中华人民共和国教育部在整个项目期间给予的支持. 两位作者都感谢 K.~Merryfield 教授对本书早期版本所作的细致审读.

两位作者感谢 University Research Institute 在本书成书过程中给予的经费支持.

最后, 第一作者感谢 Margaret Combs 与 Jan Duffy, 她们总是乐于提供帮助.


\mainmatter
% 第 0 章 Preliminaries (原书第 1-5 页)
\setcounter{chapter}{-1}
\chapter{预备知识}
\label{chap:0}

\section{引言}
\label{sec:0.1}
\index{Fourier 级数!收敛 (convergence of Fourier series)}

在本节中我们简要讨论本书的主题: Fourier 级数何时收敛? 我们令 $\mathbb{T}^n$ 表示 $n$ 维环面 $[0,1)^n$, 其元素为 $\theta=(\theta_1,\theta_2,\ldots,\theta_n)\in\mathbb{T}^n$, 而 $k=(k_1,k_2,\ldots,k_n)$ 为整数的 $n$ 元组. 若 $f\in L^p(\mathbb{T}^n)$, 则与 $f$ 相伴的 Fourier 变换和形式 Fourier 级数定义为
\begin{align*}
\hat{f}(k) &= \int f(\theta)e^{-2\pi ik\theta}\,d\theta, \\
f &= \sum_{k\in\mathbb{Z}^n}\hat{f}(k)e^{2\pi ik\theta}.
\end{align*}

本书的核心主题是分析 $f$ 的形式 Fourier 级数在何种意义下确实收敛于 $f$. 当然, 我们必须先取有限和, 再取极限. 由于求和取遍整数的所有 $n$ 元组, 即取遍所有 $k\in\mathbb{Z}^n$, 我们需要某种方法给 $\mathbb{Z}^n$ 的格点排序. 最简单的办法是把它们全部纳入不断扩张的球中: 定义 $|k|^2=k_1^2+k_2^2+\cdots+k_n^2$, 然后把 $f$ 的 Fourier 级数的第 $R$ 个球形部分和 (spherical partial sum) 定义为
\[
S_Rf(\theta)=\sum_{|k|<R}\hat{f}(k)e^{2\pi ik\theta}.
\]

在这个定义下, 我们要研究的基本问题是: $S_Rf$ 何时在 $L^p$ 中收敛到 $f$?

在一维情形, 答案是经典的, 由 M.~Riesz 于 $1910$ 年给出 (见 Zygmund~\cite{zygmund}, 第~7 章~2.4). 收敛对 $L^p$ 中所有的 $f$ 成立, 当且仅当 $1<p<\infty$.

在更高的维数, 这个问题直到最近才得到解决. Carl Herz 在 $1954$ 年证明, 收敛的一个必要条件是 $\frac{2n}{n+1}<p<\frac{2n}{n-1}$. Charles Fefferman 在 $1972$ 年证明, $L^p$ 中的收敛当且仅当 $p=2$ 时成立. Fefferman 的证明及其推论正是本书的主要内容.

Fefferman 的结果表明, 这不是给 Fourier 级数求和的正确方式. 那么还有别的方法吗? $\mathbb{Z}^n$ 的点本可以按球以外的某种次序分组. 同心多边形是一个显而易见的尝试, 但结果表明它并不比重复几个一维结果更有趣: 它不给出任何新数学, 还让人得以回避对 Fefferman 结果的深入思考. 回避思考一个课题几乎总是错误的; 往最好的方面说, 你日后也免不了要大吃几惊.

一位经典分析学家会立刻告诉你球形部分和的替代品: 使用求和法 (summability method). 我们将分析 S.~Bochner 引入的一种方法, 它本身是 Riesz 某个求和法的变形: 设 $(\xi)_+=\xi$ 当 $\xi>0$; 当 $\xi<0$ 时令它为 $0$. 于是我们定义 $\alpha$ 阶 \emph{Bochner--Riesz 平均} (Bochner--Riesz means) 为\index{Bochner--Riesz 平均 (Bochner--Riesz means)|emph}
\[
S_R^{\alpha}f(\theta)=\sum_{|k|<R}\left[\left(1-\frac{|k|^2}{R^2}\right)_{+}\right]^{\alpha}\hat{f}(k)e^{2\pi ik\theta}.
\]

关键在于: 若 $\alpha=0$, 则 $S_R^{\alpha}f=S_Rf$, 因此通过研究 $\alpha$ 趋于零时的极限行为, 我们有望理解 $S_R$ 的问题究竟出在哪里.

很好; 那么关于 $\alpha$ 接近零时的 $S_R^{\alpha}$, 已知的结果是什么? 高维时的行为极其复杂, 事实上完整的答案仅在二维时是已知的. 我们不打算写出答案的公式, 而是画一张图. 下面的图~\ref{fig:1} 与图~\ref{fig:2} 展示了 $S_R^{\alpha}$ 在 $1\leq p\leq\infty$ 上的 $L^p$ 有界性: 纵轴以 $\alpha$ 标度, 横轴以 $\frac{1}{p}$ 标度. 虚线与空心圆表示已知无界的点, 阴影区域表示已知有界的部分.

本书的目的, 就是给出图中各结果的详细证明.

\begin{figure}[htbp]
\centering
\begin{tikzpicture}[scale=0.9]
  % 阴影区域 (竖线填充)
  \fill[pattern=vertical lines, pattern color=black!70]
    (0,2) -- (1.25,0) -- (3.75,0) -- (5,2) -- (5,4) -- (0,4) -- cycle;
  % 横轴向外延伸的虚线
  \draw[dashed] (-0.5,0) -- (0,0);
  \draw[dashed] (5,0) -- (5.5,0);
  % 边框
  \draw (0,0) rectangle (5,4);
  % 虚线 V 形边界
  \draw[dashed] (0,2) -- (1.25,0);
  \draw[dashed] (3.75,0) -- (5,2);
  % 标记点
  \draw (0,2) circle (2.4pt);
  \draw (5,2) circle (2.4pt);
  \draw (1.25,0) circle (2.4pt);
  \draw (3.75,0) circle (2.4pt);
  \fill (2.5,0) circle (2.4pt);
  % 标注
  \node[left] at (-0.1,2) {$\alpha=\frac{1}{2}$};
  \node[below] at (1.25,-0.15) {$\frac{1}{4}$};
  \node[below] at (3.75,-0.15) {$\frac{3}{4}$};
\end{tikzpicture}
\caption{二维情形下 $S_R^\alpha$ 的 $L^p$ 有界性区域}
\label{fig:1}
\end{figure}

\begin{figure}[htbp]
\centering
\begin{tikzpicture}[scale=0.9]
  % 阴影区域 (竖线填充): 虚线 V 之上、实线 V 之上
  \fill[pattern=vertical lines, pattern color=black!70]
    (0,3.4) -- (1.5,0) -- (2.1,0.9) -- (2.6,0) -- (3.1,0.9) -- (3.7,0) -- (5,2.6) -- (5,4.6) -- (0,4.6) -- cycle;
  % 边框
  \draw (0,0) rectangle (5,4.6);
  % 虚线 V 形边界
  \draw[dashed] (0,3.4) -- (1.5,0);
  \draw[dashed] (3.7,0) -- (5,2.6);
  % 实线内边界
  \draw[thick] (2.1,0.9) -- (2.6,0) -- (3.1,0.9);
  % 标记点
  \draw (0,3.4) circle (2.4pt);
  \draw (5,2.6) circle (2.4pt);
  \draw (1.5,0) circle (2.4pt);
  \draw (3.7,0) circle (2.4pt);
  % 标注
  \node[left] at (-0.1,3.4) {$\alpha=\frac{n-1}{2}$};
  \node[left] at (-0.1,1.5) {$\alpha=\frac{n-1}{2(n+1)}$};
  \node[below] at (1.5,-0.15) {$\frac{n-1}{2n}$};
  \node[below] at (3.7,-0.15) {$\frac{n+1}{2n}$};
\end{tikzpicture}
\caption{高维情形下 $S_R^\alpha$ 的 $L^p$ 有界性区域}
\label{fig:2}
\end{figure}

\section{泛函分析的作用}
\label{sec:0.2}

本节给出攻克收敛问题的一个思路. 最好的办法是快速回顾一个标准结果: $\mathbb{T}$ 上存在连续函数, 其 Fourier 级数在 $\theta=0$ 处发散. 当然, 这个函数可以显式地构造出来, 但在一般维数下, 显式的计算是难以驾驭的. 我们转而求助于泛函分析的方法.

抽象的问题是这样的. 若我取一个有限 Fourier 级数, 其 Fourier 级数的部分和只是在进行到某处后停止, 我便收回原来的函数. 于是, Fourier 级数的部分和在所谓的三角多项式上收敛. 三角多项式在连续函数中稠密 (Stone--Weierstrass 定理), 而连续函数在 $L^p$ 中稠密. 我需要把``部分和在一个稠密子集上收敛''的信息过渡为``在整个函数空间上收敛'': 我需要交换极限. 标准的实分析告诉我, 某种一致性允许交换极限; 泛函分析则告诉我, 一致性是必要条件. 这正是共鸣定理 (uniform boundedness theorem) 的要义.

\begin{theorem}\label{thm:0.1}
\index{共鸣定理 (uniform boundedness theorem)}
存在 $f\in C(\mathbb{T})$, 使得 $\sup_N|S_Nf(0)|=\infty$.
\end{theorem}

\begin{proof}
我们先把这个问题转化为函数空间上的算子问题. 定义算子 $S_N:C(\mathbb{T})\to\mathbb{C}$ 为 $S_Nf=\sum_{|k|\leq N}\hat{f}(k)$. 由于和是有限的, 每个 $S_N$ 都是 $C(\mathbb{T})$ 上的连续线性泛函. 每个这样的线性泛函都由对一个有限 Borel 测度的积分给出; 在此情形,
\[
\sum_{|k|\leq N}\hat{f}(k)=\sum\int e^{-2\pi ik\theta}f(\theta)\,d\theta=\int K_N(\theta)f(\theta)\,d\theta,
\]
其中
\[
K_N(\theta)=\sum_{|k|\leq N}e^{-2\pi ik\theta}=\frac{\sin 2\pi\bigl(N+\frac{1}{2}\bigr)\theta}{\sin(\pi\theta)}.
\]
于是在此情形, 给出 $S_N$ 的测度是绝对连续的, 其全变差范数恰为 $\|K_N\|_1$. 这甚至不难计算; Zygmund~\cite{zygmund} 中给出了精确的计算:
\[
\|K_N\|_1=\frac{2}{\pi}\log N+O(1).
\]
真正的问题是把 $S_N$ 在个别函数上的作用, 与它作为空间 $C(\mathbb{T})$ 上的算子的行为联系起来. 这正是共鸣定理所告诉我们的: 要么 i)~$S_N$ 在算子范数下一致有界; 要么 ii)~对一个稠密集中的 $f$, $\sup_N|S_Nf|=\infty$. 由于第一个结论不成立, 第二个必然成立.
\end{proof}

这个结果绝佳地代表了本书余下部分要做的事情: 先用一些泛函分析的技巧, 把收敛问题化为函数空间上算子的有界性问题; 泛函分析之后, 是对具体算子的计算; 最终分析以一个详细的计算收尾.

\section{背景知识}
\label{sec:0.3}
\index{速降函数 (rapidly decreasing functions)}
\index{缓增分布 (tempered distributions)}

下一章将给出分析 Fourier 级数收敛所需的泛函分析. 最好的技巧是把问题从 Fourier 级数转移到 Fourier 积分. 这样做有好处: 显式计算一个积分, 比把一个级数求和成封闭形式要容易. 但也有坏处: 定义 Fourier 变换的那些积分并不绝对收敛. 克服这一困难的现代做法是, 先在一个万事皆收敛的稠密子集上工作, 然后再取极限. 这正是速降函数理论的要点, 本节对此作一概述. 更详细的论述见 Stein 与 Weiss~\cite{steinweiss}.

对 $f\in L^1(\mathbb{R}^n)$, Fourier 变换 $\hat{f}$ 定义为
\[
\hat{f}(\xi)=\int f(x)e^{-2\pi ix\xi}\,dx.
\]
该积分绝对收敛, 且 $\|\hat{f}\|_{\infty}\leq\|f\|_1$. 而要对 $f\in L^2$ 定义 $\hat{f}$, 就得耍些花招、用点骗术了.

速降函数空间或 Schwartz 函数空间记作 $\mathcal{S}$, 定义为 $\mathbb{R}^n$ 上所有使半范数 $\sup_x\allowbreak|x^{\alpha}D^{\beta}\phi(x)|$ 有限的光滑函数 $f$ 构成的类. 这里我们使用了多重指标记号: 若令 $\alpha$ 表示多重指标 $\alpha=(\alpha_1,\alpha_2,\ldots,\alpha_n)$, 则 $x^{\alpha}=x_1^{\alpha_1}\cdots x_n^{\alpha_n}$. 当这些半范数随 $\alpha$ 变动时, 它们在 $\mathcal{S}$ 上诱导出一个拓扑. $\mathcal{S}$ 的对偶空间称为缓增分布 (tempered distribution) 空间. 于是不难证明:

\begin{theorem}\label{thm:0.2}
\index{Fourier 变换!在 $\mathcal{S}$ 上 (Fourier transform on $\mathcal{S}$)}
\begin{enumerate}[label=\alph*)]
\item 若 $1\leq p<\infty$, 则 $\mathcal{S}$ 在 $L^p$ 中稠密.
\item Fourier 变换是 $\mathcal{S}$ 到 $\mathcal{S}$ 的连续单射满映射.
\item 若 $\mathcal{S}$ 中两个函数 $f$ 与 $g$ 的卷积定义为
\[
f\star g=\int f(x-y)g(y)\,dy,
\]
则该函数仍在 $\mathcal{S}$ 中, 并且 $\widehat{f\star g}=\hat{f}\,\hat{g}$.
\item 若 $f$ 与 $g$ 属于 $\mathcal{S}$, 则 $\int f\bar{g}=\int\hat{f}\,\overline{\hat{g}}$.
\item $\|f\|_2=\|\hat{f}\|_2$.
\item 分布 $u$ 的 Fourier 变换定义为 $\hat{u}(f)=u(\hat{f})$. 则 $\hat{u}$ 仍是一个分布.
\item 平移算子 $\tau_y$ 定义为 $(\tau_y f)(x)=f(x-y)$. 若把分布 $u$ 与 $\mathcal{S}$ 中函数的卷积定义为 $(u\star f)(x)=u(\tau_x\tilde{f})$, 其中 $\tilde{f}(y)=f(-y)$, 则 $u\star f$ 是光滑函数.
\end{enumerate}
\end{theorem}

空间 $\mathcal{S}$ 是一个技术性工具, 它使得我们能够方便地对那些原本可能发散的积分进行计算. 本书的其余部分还需要一些记号. 若 $T$ 是从 $L^p$ 到 $L^p$ 的有界线性算子, 我们把它的算子范数记作 ${}_p\|T\|$. 回忆
\[
{}_p\|T\|=\sup\{\|Tf\|_p \mid \|f\|_p=1\},
\]
以及
\[
{}_p\|T\|=\sup\left\{\left|\int Tfg\,\right| \mid \|f\|_p=\|g\|_{p'}=1\right\},
\]
其中 $\frac{1}{p}+\frac{1}{p'}=1$.

\section{第\texorpdfstring{$0$}{0}章评注}
\label{sec:0.4}
\index{Lebesgue 常数 (Lebesgue constants)}
\index{球形求和 (spherical summability)}

\noindent (0.1):\quad Bochner--Riesz 平均是 Riesz 平均 $(1-|n|^2)_+^{\alpha}$ 的变形; 我们将会看到, 这两种求和法在 $L^p$ 中给出相同的结果, 但 Bochner--Riesz 平均要容易计算得多. 对 $S_R^{\alpha}$ 最初的兴趣来自它与逐点求和及 Fourier 级数局部化问题的联系 (这个课题如今因其困难而名声不佳). Bochner 的论文~\cite{bochner} 表明, $n$ 维问题比一维问题更为深刻: 局部化并非对所有的 $\alpha$ 都成立, 这与一维的结果形成对照. 其中的基本思想也可见 Stein 与 Weiss~\cite{steinweiss} 第~7 章定理~4.2.

把 Fourier 系数按同心球而不是同心多边形分组, 是有充分理由的. 用 Bochner 自己的话说 (\cite{bochner}, 第~179--180 页):

\begin{quote}
``初等指数函数 $u(x)=e^{i(n_1x_1+\cdots+n_kx_k)}$ ($n_1,\ldots,n_k$ 均为整数) 是特征值问题
\[
\triangle u(x)=-\lambda u(x)
\]
的一组完备的正则解, 只要把这个系统视为定义在 (闭) 环面
\[
0\leq x_1\leq 2\pi,\ \ldots,\ 0\leq x_k\leq 2\pi
\]
上, 且 $\Delta$ 是关于环面上 Euclid 度量的 Laplace 算子, 即
\[
\frac{\partial^2 u}{\partial x_1^2}+\cdots+\frac{\partial^2 u}{\partial x_k^2}.
\]
由于 $\lambda=n_1^2+\cdots+n_k^2$, 我们书写级数的方式……满足一个十分自然的原则: 按特征值 $\lambda$ 的大小给级数的项排序.''
\end{quote}

用 Bochner 的话来复述: 在考察球形求和法时, 我们实际上是在考察紧 Riemann 流形上某个微分算子的特征函数展开的收敛性问题. C.~Fefferman 的结果已被用来证明: 当 $n>1$ 时, 这样的展开永不收敛. 这正是我们的主要论点: 球形平均的发散并非选择了古怪求和法所造成的反常, 而是高维几何的一个基本事实.

图~\ref{fig:1} 与图~\ref{fig:2} 所总结的结果归功于许多人. 概述可参见 Herz~\cite{herz}、C.~Fefferman~\cite{fefferman-spherical}、Carleson 与 Sjölin~\cite{christ}\footnote{译者注: 原书此处引作 Carleson and Sjölin [5], 但原书参考文献 [5] 实为 M.~Christ 的博士论文, 且原书参考文献缺失序号 6; 此处当指 L.~Carleson 与 P.~Sjölin 1972 年在 Studia Math. 上关于圆盘乘子问题的著名论文.} 以及 Tomas~\cite{tomas1}.

\noindent (0.2):\quad $\|K_R\|_1$ (称为 Lebesgue 常数) 的计算归功于 Fejér ($1910$ 年). 这类技巧在高维时一败涂地, 因为 $K_R(\theta)$ 的级数无法求和. 可以证明, 当 $n\geq 2$ 时 $\|K_R\|_1=O\bigl(R^{\frac{n-1}{2}}\bigr)$, 这已是最佳估计; 证明及更早结果的文献见 Shapiro~\cite{shapiro}. 在试图为部分和 $K_R$ 寻找显式公式时, 你很快会意识到: 一个真正好的公式将使你能够精确算出 $\mathbb{R}^n$ 中半径为 $R$ 的球内整点的个数. 不幸的是, 对这个量给出好的估计是数论中最困难的问题之一. 这是个严肃的问题, 因为对某些 $R$, 不存在整数对 $(m,n)$ 使 $m^2+n^2=R^2$; 而对稍微不同的 $R$, 这样的整数对又比比皆是. 你不能指望从 $K_R$ 的一个好公式得到任何真正规整的表达式.

关于一点处 Fourier 级数发散的经典详细证明, 最好查阅 Zygmund~\cite{zygmund} 第~7 章第一卷.

% 第 1 章 Multiplier Theory (原书第 7-26 页)
\chapter{乘子理论}
\label{chap:1}

\section{\texorpdfstring{$L^p$}{Lp} 上的乘子}
\label{sec:1.1}

在对 Fourier 积分``求和''时会遇到一些问题; 本节的目的是看清这些问题是什么, 以及如何克服它们.

若 $f\in L^2$, 预期的逆 Fourier 变换
\[
\check{f}=\int f(\xi)e^{2\pi ix\xi}\,d\xi
\]
根本不收敛. 传统的补救办法是转而考察``部分和'', 即有界区域上的积分; 这迫使积分收敛. 人们随后希望取极限并得到极限的收敛性. 我们定义部分和算子 $S_R$ 为
\[
S_Rf(x)=\int_{|\xi|\leq R}\hat{f}(\xi)e^{2\pi ix\xi}\,d\xi,
\]
它是绝对收敛的. 我们希望对所有的 $f\in L^2$, $\lim_{R\to\infty}S_Rf$ 在 $L^2$ 中存在, 并且极限就是 $f$. 第一个问题是, 我们甚至不知道 $S_Rf$ 是否属于 $L^2$, 更谈不上在 $L^2$ 中收敛. 于是我们从算子 $S_R$ 的泛函分析入手: 它们把 $L^2$ 函数映为 $L^2$ 函数吗? 我们令 $B$ 表示 $\mathbb{R}^n$ 中的单位球, 并令 $\chi_B\bigl(\frac{\xi}{R}\bigr)$ 为半径 $R$ 的球的特征函数. 由此可得
\[
S_Rf(x)=\int\chi_B\Bigl(\frac{\xi}{R}\Bigr)\hat{f}(\xi)e^{2\pi ix\xi}\,d\xi.
\]
这些就是我们要研究的算子.

\begin{definition}\label{def:1.1}
\index{Fourier 乘子 (Fourier multiplier)|emph}
$L^p(\mathbb{R}^n)$ 上的 \emph{Fourier 乘子算子} (Fourier multiplier operator) 是指 $L^p$ 上的有界线性算子 $T$, 对它存在 $\mu\in L^\infty(\mathbb{R}^n)$, 使得对一切 Schwartz 函数 $f$ 有
\[
Tf(x)=\int\mu(\xi)\hat{f}(\xi)e^{2\pi ix\xi}\,d\xi.
\]
此时, $\mu$ 称为与算子 $T$ 相伴的 \emph{Fourier 乘子} (Fourier multiplier).
\end{definition}

\noindent\textbf{注:} 这个算子实际上是由一个泛函方程定义的: $\widehat{Tf}=\mu\hat{f}$, 这是十分间接的. 要真正使用乘子算子, 我们需要能停留在其中一侧的刻画: 要么停留在函数一侧, 要么停留在 Fourier 变换一侧.

\begin{proposition}\label{prop:1.2}
\index{卷积核 (convolution kernel)}
设 $T$ 是 $L^p$ 的 Fourier 乘子. 则存在一个分布 $K$, 使得对一切 $f\in\mathcal{S}$ 有
\[
Tf(x)=K\star f(x).
\]
$K$ 称为 $T$ 的\emph{卷积核} (convolution kernel).
\end{proposition}

\begin{proof}
为了看清发生了什么, 先设 $\mu\in\mathcal{S}$. 则
\begin{align*}
Tf(x)&=\int\mu(\xi)e^{2\pi ix\xi}\hat{f}(\xi)\,d\xi
=\int\mu(\xi)e^{2\pi ix\xi}\int e^{-2\pi i\xi y}f(y)\,dy\,d\xi \\
&=\iint\mu(\xi)e^{2\pi i\xi(x-y)}\,d\xi\, f(y)\,dy
=\int\check{\mu}(x-y)f(y)\,dy=\check{\mu}\star f(x).
\end{align*}
由于所涉对象皆为 Schwartz 函数, 我们可以交换积分次序, 并取分布 $K$ 为函数 $\check{\mu}$. 要给出一般情形的证明, 需要许多分布方面的事实. 首先, 若 $f\in\mathcal{S}$, 则 $\hat{f}\in\mathcal{S}$, 又因 $\mu\in L^\infty$, 故 $\mu\hat{f}\in L^1$, 于是
\[
Tf(x)=\int e^{2\pi ix\xi}\mu(\xi)\hat{f}(\xi)\,d\xi
\]
对一切 $x$ 存在.

分布 $K$ 定义为 $K(f)=(T\tilde{f})(0)$. $f$ 的积分表示表明这是一个分布: $f\mapsto\hat{f}$ 是 $\mathcal{S}$ 到 $\mathcal{S}$ 的连续映射, 而用 $L^\infty$ 函数作积分是众所周知的分布. 最后, 注意到
\begin{align*}
(K\star f)(x)&\equiv K(\tau_x\tilde{f})\equiv T\bigl(\widetilde{\tau_x\tilde{f}}\bigr)(0)
=T(\tau_{-x}f)(0) \\
&=\int\mu(\xi)\widehat{\tau_{-x}f}(\xi)\,d\xi
=\int\mu(\xi)e^{2\pi ix\xi}\hat{f}(\xi)\,d\xi\equiv Tf(x);
\end{align*}
倒数第二个等号成立是因为
\[
\widehat{\tau_{-x}f}(\xi)=\int e^{-2\pi iy\xi}f(y+x)\,dy=e^{2\pi ix\xi}\hat{f}(\xi). \qedhere
\]
\end{proof}

\begin{remark}\label{rem:1.3}
\begin{enumerate}[label=\alph*)]
\item 作为分布, $\tilde{K}$ 就是 $\mu$, 而 $\tilde{\mu}$ 就是 $K$.
\item 分布 $K$ 可能非常糟糕; 看算子 $Tf=f$. 这里 Fourier 反演公式告诉我们 $\mu\equiv 1$. 给出这个算子的分布是什么? 与 $T$ 对应的分布是 Dirac $\delta$ 测度, $\delta(f)=f(0)$. 把形式步骤做一遍,
\[
(\delta\star f)(x)\equiv\delta(\tau_x\tilde{f})=(\tau_x\tilde{f})(0)=\tilde{f}(-x)=f(x).
\]
寓意在于: 即使是平凡的算子, 也可能产生出甚至不能用对某个函数的积分来表示的分布. 更糟的是, 下面我们勾画一个例子, 其中 $K$ 甚至不能用对某个有限测度的积分来表示. 设 $\mu$ 是 $L^\infty$ 函数 $i\,\mathrm{sign}(\xi)$, 那么我们将在注~\ref{rem:2.13} 中算出 $\check{\mu}(x)=\frac{1}{\pi x}$. 但必须承认, $\check{\mu}(f)=\frac{1}{\pi}\int\frac{1}{x}f(x)\,dx$ 通常并不收敛. 严格地说, 我们必须取主值积分
\[
\lim_{\epsilon\to0}\frac{1}{\pi}\int_{|x|\geq\epsilon}\frac{1}{x}f(x)\,dx.
\]
当乘子不是可积函数时, 就应当预期这类分布的出现.
\end{enumerate}
\end{remark}

\begin{proposition}\label{prop:1.4}
若 $T$ 是 $L^p$ 的乘子且 $f\in\mathcal{S}$, 则
\[
\widehat{Tf}(\xi)=\mu(\xi)\hat{f}(\xi).
\]
\end{proposition}

\begin{proof}
为此, 我们必须看到: 作为分布, $\widehat{Tf}$ 与 $\mu(\xi)\hat{f}(\xi)$ 的作用相同. 逐点的相等是立刻得到的, 因为这两个分布都由对函数的积分给出; 这正是取 $f\in\mathcal{S}$ 的全部意义所在: 此时 $\hat{f}\in\mathcal{S}$ 且 $\mu\hat{f}\in L^1$. 由此可知 $Tf$ 甚至是连续的. 好, 通过对 $g\in\mathcal{S}$ 积分来验证相等: 由定理~\ref{thm:0.2} 扩张到整个 $L^2$ 的结果,
\[
\int\widehat{Tf}\,g=\int Tf\,\hat{g}.
\]
但是
\[
\int\mu\hat{f}\,g=\iint\mu\hat{f}\,g=\int(\mu\hat{f})\check{g}=\int Tf\,\hat{g}. \qedhere
\]
\end{proof}

三元组 $(T,K,\mu)$ 称为\emph{乘子三元组} (multiplier triple).\index{乘子三元组 (multiplier triple)} $T$ 是乘子算子, $K$ 是卷积核, $\mu$ 是乘子. 在本书的其余部分, 我们会经常使用诸如``设 $T$ 是以 $\mu$ 为乘子的算子''或``若 $\mu$ 是以 $K$ 为核的乘子''这样的措辞.

\begin{proposition}\label{prop:1.5}
设 $T$ 是 $L^2$ 上的乘子算子, 乘子为 $\mu$. 则
\[
{}_2\|T\|=\|\mu\|_{\infty}.
\]
\end{proposition}

\begin{proof}
由 Plancherel 定理,
\[
\|Tf\|_2=\|\widehat{Tf}\|_2=\|\mu\hat{f}\|_2\leq\|\mu\|_{\infty}\|\hat{f}\|_2=\|\mu\|_{\infty}\|f\|_2,
\]
故 ${}_2\|T\|\leq\|\mu\|_{\infty}$. 现在证明: 若 $\epsilon>0$, 则 ${}_2\|T\|\geq\|\mu\|_{\infty}-\epsilon$.

由 $L^\infty$ 的定义, 集合 $E=\{\xi \mid |\mu(\xi)|\geq\|\mu\|_{\infty}-\epsilon\}$ 具有正测度; 我们可以选取 $E$ 的一个测度为正但有限的子集 $S$. 于是可以计算这个函数的范数: $\check{\chi}_S\in L^2$, 且 $\|\check{\chi}_S\|_2=\|\hat{\chi}_S\|_2=|S|^{\frac{1}{2}}$. 我们还能计算 $T$ 作用在这个函数上所得的范数:
\begin{align*}
{}_2\|T\|\,\|\check{\chi}_S\|_2={}_2\|T\|\,|S|^{\frac{1}{2}}&\geq\|T(\check{\chi}_S)\|_2=\|\mu\chi_S\|_2 \\
&\geq\inf|\mu(\xi)\chi_S(\xi)|\,\|\chi_S\|_2
\geq(\|\mu\|_{\infty}-\epsilon)|S|^{\frac{1}{2}}
=(\|\mu\|_{\infty}-\epsilon)\|\check{\chi}_S\|_2. \qedhere
\end{align*}
\end{proof}

\begin{remark}\label{rem:1.6}
我们用到了 $T(\check{\chi}_S)=\mu\chi_S$ 这一结果. 我们实际上只在 $f\in\mathcal{S}$ 时证明了它, 但集合 $S$ 的特征函数当然不在 $\mathcal{S}$ 中. 这是折磨着本课题的那些琐碎技术问题的典型, 它们全都源于这样一个事实: 我们的乘子只定义在 $L^p$ 的一个稠密集上, 再由连续性延拓到其余部分. 在这条注记中, 我们想展示如何处理一个典型的琐碎问题.

对 $f\in L^2$, 选取 $f_j\in\mathcal{S}$ 使 $f_j$ 在 $L^2$ 中收敛于 $f$. 由于 Fourier 变换在 $L^2$ 上连续, $\hat{f}_j$ 在 $L^2$ 中收敛于 $\hat{f}$. 由于 $T$ 是 $L^2$ 上的有界 (即连续) 算子, $Tf_j$ 在 $L^2$ 中收敛于 $Tf$, 且 $\widehat{Tf_j}$ 在 $L^2$ 中收敛于 $\widehat{Tf}$. 我们可以取子列 $j_k$, 使 $\widehat{Tf_{j_k}}$ 与 $\hat{f}_{j_k}$ 分别几乎处处逐点收敛于 $\widehat{Tf}$ 与 $\hat{f}$. 于是对几乎所有的 $\xi$,
\[
\widehat{Tf}(\xi)=\lim_{k\to\infty}\widehat{Tf_{j_k}}(\xi)
=\lim_{k\to\infty}\mu(\xi)\hat{f}_{j_k}(\xi)=\mu(\xi)\hat{f}(\xi). \qedhere
\]
\end{remark}

\begin{proposition}\label{prop:1.7}
$T$ 是 $L^1$ 的乘子当且仅当卷积核 $K$ 是一个有限测度 $dm$; 此时
\[
{}_1\|T\|=\|dm\|.
\]
\end{proposition}

\begin{proof}
这个结果有一个简单的直观解释. Dirac $\delta$ 函数几乎是一个 $L^1$ 函数, 而且它的作用如同恒等算子. 于是 ${}_1\|T\|$ 应当类似于 $\|K\star\delta\|_1=\|K\|_1$. 这个证明展示了如何处理``$\delta$ 不是 $L^1$ 函数''这一``小小的''技术麻烦. 若我们假设 $K$ 是一个有限测度, 则
\begin{align*}
\|K\star f\|_1&=\|K(\tau_x\tilde{f})\|_1
=\left\|\int(\tau_x\tilde{f})(y)\,dm(y)\right\|_1 \\
&\leq\int\|(\tau_x\tilde{f})(\cdot)\|_1\,d|m|(y)=\|f\|_1\,\|dm\|.
\end{align*}
由此得 ${}_1\|T\|\leq\|dm\|$.

对于另一个方向, 我需要触及 $\delta$ 函数, 尽管它不是 $L^1$ 函数. 我们将使用一个绝对标准的技术: 不直接使用 $\delta$, 而是用属于 $L^1$ 的函数逼近它. 我们选取 Gauss 核:
\[
\omega_t(x)=(2\pi t)^{-n}\exp\bigl[-(|x|/2t)^2\bigr];
\]
我们把它安排得使 $\|\omega_t\|_1=1$, $\omega_t\in\mathcal{S}$, 并且对每个 $f\in\mathcal{S}$, 当 $t$ 趋于零时 $\omega_t\star f$ 在 $\mathcal{S}$ 中收敛于 $f$. 当然我们得到
\[
\|T(\omega_t)\|_1\leq{}_1\|T\|\,\|\omega_t\|_1={}_1\|T\|;
\]
问题是我们能否得到 $T(\omega_t)=K\star\omega_t$ 收敛到 $dm$. 若 $K$ 是 $\mathcal{S}$ 中的函数, 这不成问题; 这里我们只知道 $T(\omega_t)$ 在 $L^1$ 中一致有界. 要从这类有界性信息得到收敛性信息, 显然需要某种紧性条件. 不幸的是, $L^1$ 中的单位球不是序列紧的, 于是我们只能耍花招、用骗术.

$T(\omega_t)$ 属于 $L^1$, 并且可以把它们视为全变差范数被 $1$ 控制的有限测度. 令人惊奇的是, 有限测度空间 $\mathcal{M}$ 上存在一个拓扑, 使其中的闭有界集是序列紧的. 我们把 $\mathcal{M}$ 视为 $C_0$ (在无穷远处消失的连续函数空间) 的对偶. 我们在 $\mathcal{M}$ 上赋予的拓扑称为弱$^\star$拓扑: $\mu_j$ 收敛到 $\mu$ 意味着对所有的 $f\in C_0$, $\mu_j(f)$ 收敛到 $\mu(f)$. 于是 Alaoglu 定理 (Rudin~\cite{rudinfa}) 给出了所需的紧性: 存在子列 $T(\omega_{t_j})$ 弱$^\star$收敛到半径 ${}_1\|T\|$ 的球中的某个有限测度 $dm$. 我们已经知道了 $dm$ 是什么, 并且知道 $dm$ 的全变差范数被 ${}_1\|T\|$ 控制. 这幅美好图景中唯一欠缺的小细节, 是``分布 $K$ 由对 $dm$ 的积分给出''这一论断. 这正是弱$^\star$收敛的作用. 由于 $\mathcal{S}\subset C_0$, 对每个 $f\in\mathcal{S}$,
\begin{align*}
dm(f)&\equiv\int f\,dm=\lim_{j\to\infty}\int f\,T(\omega_{t_j})
=\lim_{j\to\infty}\int f\,K\star\omega_{t_j} \\
&=\lim_{j\to\infty}\int f(x)K\bigl(\tau_x\tilde{\omega}_{t_j}(y)\bigr)\,dx
=\lim_{j\to\infty}K\left(\int f\,\omega_{t_j}(y-x)\,dx\right).
\end{align*}
最后一个等号成立是因为该积分可以用 Riemann 和逼近, 且逼近在 $\mathcal{S}$ 中收敛, 于是我们交换了逼近与分布的作用. 最后,
\[
\lim_{j\to\infty}K\bigl(\omega_{t_j}\bigr)\star f=K\Bigl(\lim_{j\to\infty}\omega_{t_j}\star f\Bigr)=K(f).
\]
这被称为一个``正则性''结果: 起初我们只知道 $K$ 是某个可能很糟糕的分布, 而事实上它不得不是一个好得多的对象——一个有限测度. 一个分布怎么会被突然强制成为一个测度呢? $L^1$ 范数的有界性迫使一列光滑逼近收敛. 这是我们看到的第一个``算子的 $L^p$ 有界性强制正则性''的例子; 请再次留意泛函分析所扮演的角色.
\end{proof}

\begin{lemma}\label{lem:1.8}
若 $T$ 是 $L^p$ 的乘子, 则它也是 $L^{p'}$ 的乘子, 其中 $\frac{1}{p}+\frac{1}{p'}=1$.
\end{lemma}

\begin{proof}
为了看清这里的想法, 回忆 $L^{p'}$ 是 $L^p$ 的对偶空间, 因此 $T$ 在 $L^p$ 上的有界性蕴含伴随算子 $T^{\ast}$ 在 $L^{p'}$ 上的有界性. 我们的目标是建立 $T$ 与 $T^{\ast}$ 的联系. 若 $T$ 由与某个足够好的函数 (比如说一个漂亮的 Schwartz 函数 $K$) 的卷积给出, 则我们可以计算:
\begin{align*}
\int Tf\,\bar{g}&=\int K\star f\,\bar{g}
=\iint K(x-y)f(y)\bar{g}(x)\,dy\,dx \\
&=\iint K(x-y)\bar{g}(x)f(y)\,dy\,dx
=\int\bigl[K\star(\widetilde{\bar{g}})\bigr]f(y)\,dy \\
&=\int f(y)\,\widetilde{K\star(\widetilde{\bar{g}})}(y)\,dy.
\end{align*}
这个计算意味着: 若我们用 $Jf=\bar{f}$ 定义 $L^p$ 上的等距算子 $J$, 则 $T^{\ast}=JTJ$. 由于 $J$ 是等距算子, 我们得到
\[
{}_p\|T\|={}_{p'}\|T^{\ast}\|={}_{p'}\|T\|.
\]
真正的困难在于当 $K$ 不是漂亮的 Schwartz 函数时交换积分次序. 我们将依照刚才给出的直观来做一些不愉快的分布计算.
\begin{align*}
\int Tf\,\bar{g}&=\int K\star f(x)\,\bar{g}(x)\,dx
=\int K\bigl(\tau_x(\tilde{f})(y)\bigr)\bar{g}(x)\,dx \\
&=K\left(\int(\tau_x\tilde{f})(y)\bar{g}(x)\,dx\right)
=K\left(\int f(x-y)\bar{g}(x)\,dx\right) \\
&=K\left(\int f(x)\bar{g}(x+y)\,dx\right)
=K\left(\int f(x)(\tau_{-x}\bar{g})(y)\,dx\right) \\
&=\int f(x)\bigl[K\star\tilde{\bar{g}}(x)\bigr]\,dx.
\end{align*}
等等, 等等.
\end{proof}

% 原书编号由 1.8 直接跳至 1.10, 此处跳过 1.9 以与原书一致.
\addtocounter{mtthm}{1}

\begin{lemma}\label{lem:1.10}
若 $K\in L^1$, 则 $\|K\star f\|_p\leq\|K\|_1\|f\|_p$.
\end{lemma}

\begin{proof}
与命题~\ref{prop:1.7} 的证明一样,
\[
\|K\star f\|_p=\left\|\int K(y)f(x-y)\,dy\right\|_p
\leq\int|K(y)|\,\|f(x-y)\|_{p,x}\,dy=\|K\|_1\|f\|_p. \qedhere
\]
\end{proof}

\begin{remark}\label{rem:1.11}
现在我们对 $L^1$ 与 $L^2$ 上的乘子已经有了相当完整的了解; 我们希望对 $1<p<2$ 时的 $L^p$ 有更好的理解. 由对偶性, 即由引理~\ref{lem:1.8}, 我们将能把握 $2<p<\infty$ 时的 $L^p$ 乘子.

想想恒等算子, 它在所有 $L^p$ 上有界. 如果我们了解了一个函数在 $L^1$ 与 $L^2$ 中的情形, 为什么这能告诉我们它在 $1<p<2$ 时的 $L^p$ 中的任何信息? 这里需要的是某种凸性: 若以图形的方式思考, $T$ 在 $L^p$ 上的算子范数应当位于 $L^1$ 与 $L^2$ 两点的连线之下. 而这恰好不成立. 看看 H\"older 不等式:
\[
\int|f|^{tp_0+(1-t)p_1}\leq\left(\int|f|^{p_0}\right)^{t}\left(\int|f|^{p_1}\right)^{1-t};
\]
它告诉我们 $\log\bigl(\int|f|^p\bigr)$ 是 $p$ 的凸函数.
\end{remark}

\section{\texorpdfstring{$L^{p'}$}{Lp'} 上的乘子}
\label{sec:1.2}

\begin{theorem}[三线定理]\label{thm:1.12}
\index{三线定理 (three-lines theorem)}
设 $S=\{z\in\mathbb{C} \mid 0\leq\operatorname{Re}(z)\leq 1\}$. 设 $F$ 在 $S$ 上有界连续, 在 $S$ 的内部解析. 令 $k_x=\sup_y|F(x+iy)|$. 则 $\log k_x$ 是 $x$ 的凸函数.
\end{theorem}

\textbf{证明思路}: 这是复变理论中的标准结果; 第一步是把 $F$ 重新调整到在 $0+iy$ 与 $1+iy$ 上以 $1$ 为界; 调整后的函数是 $G(z)=F(z)k_0^{z-1}k_1^{-z}$. 然后我们设法证明 $|G|\leq 1$ 处处成立. 若 $G$ 在 $|y|\to\infty$ 时消失, 我们构造一个矩形, 使 $|G|\leq 1$ 在其边界上成立, 在其内部应用最大模原理. 若 $G$ 在无穷远处不消失, 我们利用 $G$ 的有界性, 乘以一个趋于零的指数函数. 让这个指数函数越来越慢地趋于无穷, 便可恢复 $G$. 由此得出 $G(z)$ 处处以 $1$ 为界, 从而 $|F(z)|\leq k_0^{1-z}k_1^{z}$. 取上确界再取对数便得到所要的结果.

\begin{theorem}\label{thm:1.13}
设 $T$ 是在 $L^{p_i}$ ($i=0,1$) 上有界的线性算子. 则 $\log{}_p\|T\|$ 是 $\left[\frac{1}{p_0},\frac{1}{p_1}\right]$ 上的凸函数. 也就是说, 若 $\frac{1}{p}=\frac{t}{p_0}+\frac{1-t}{p_1}$, 其中 $0\leq t\leq 1$, 则
\[
{}_p\|T\|\leq{}_{p_0}\|T\|^{t}\,{}_{p_1}\|T\|^{1-t}.
\]
\end{theorem}

\begin{proof}
显然我们想使用定理~\ref{thm:1.12}, 并取 $F(z)=\int|Tf|^z$. 这在 $Tf(x)=0$ 的那些 $x$ 处会遇到技术困难, 而且我无法在不破坏积分有限性的前提下通过加上 $\epsilon$ 把 $T(f)$ 提升到 $0$ 以上. 再者, 这样也看不到 ${}_p\|T\|$. 我用一步解决所有这些问题: 我将用对一个 $L^{p'}$ 函数的积分来检验 $T(f)$ 的大小, 并通过对简单函数计算 $T$ 来避开零点问题. 于是, 设 $f,g$ 为满足
\[
\|f\|_p=\|g\|_{p'}=1
\]
的简单函数,
\[
f(x)=\sum|a_j|e^{i\theta_j}\chi_{E_j};\quad g(x)=\sum|b_k|e^{i\phi_k}\chi_{F_k}.
\]
这里 $E_j$, $F_k$ 是有限测度集; 我们将证明
\[
\int Tfg\leq{}_{p_0}\|T\|^{t}\,{}_{p_1}\|T\|^{1-t}.
\]
先令 $\alpha(z)=\frac{z}{p_0}+\frac{1-z}{p_1}$. 它把 $\frac{1}{p}$ 区间映入 $[0,1]$. 然后令
\[
F(z)=\int T\left(\sum|a_j|^{p\alpha}e^{i\theta_j}\chi_{E_j}\right)\left(\sum|b_k|^{p'(1-\alpha)}e^{i\phi_k}\chi_{F_k}\right).
\]
这并不是把 $Tf$ 升到 $z$ 次幂, 但它是带状区域 $S$ 内的 $z$ 的解析函数, 且有界连续. 注意到
\[
|F(x+iy)|\leq{}_{p_0}\|T\|\,\left\|\sum|a_j|^{p\alpha}e^{i\theta_j}\chi_{E_j}\right\|_{p_0}\left\|\sum|b_k|^{p'(1-\alpha)}e^{i\phi_k}\chi_{F_k}\right\|_{p'}.
\]
但是
\[
\left\|\sum|a_j|^{p\alpha}e^{i\theta_j}\chi_{E_j}\right\|_{p_0}
=\left\|\sum|a_j|^{\operatorname{Re}(p\alpha)}\chi_{E_j}\right\|_{p_0},
\]
这是因为集合 $E_j$ 两两不交. 现在
\[
\operatorname{Re}(\alpha p)=p\operatorname{Re}\left[\frac{1-iy}{p_0}+\frac{iy}{p_1}\right];
\]
由此可得
\begin{align*}
\left\|\sum|a_j|^{\operatorname{Re}(p\alpha)}\chi_{E_j}\right\|_{p_0}
&=\left\|\sum|a_j|^{\frac{p}{p_0}}\chi_{E_j}\right\|_{p_0} \\
&=\left\|\sum|a_j|\chi_{E_j}\right\|_p^{\frac{p}{p_0}}=\|f\|_p^{\frac{p}{p_0}}=1.
\end{align*}
对第二项有类似的估计, 于是
\[
|F(0+iy)|\leq{}_{p_0}\|T\|;\quad |F(1+iy)|\leq{}_{p_1}\|T\|.
\]
定理现在由定理~\ref{thm:1.12} 得出.
\end{proof}

\noindent\textbf{注:} 这个结果称为 Riesz--Thorin 内插定理.\index{内插!Riesz--Thorin (Riesz--Thorin interpolation)}

\begin{theorem}\label{thm:1.14}
若 $1\leq p\leq 2$, 则 $\|\hat{f}\|_{p'}\leq\|f\|_p$.
\end{theorem}

\begin{proof}
\begin{align*}
\|Tf\|_{\infty}&=\|\hat{f}\|_{\infty}\leq\|f\|_1 \\
\|Tf\|_2&=\|\hat{f}\|_2=\|f\|_2
\end{align*}
这两个情形表明, 在 $p=1$ 与 $p=2$ 这两个特殊情形, $T$ 把 $L^p$ 映入 $L^{p'}$. $p$ 与 $p'$ 的关系是线性的; 内插保持这一关系. 该映射在每个端点空间上范数以 $1$ 为界, 内插同样保持这一范数.
\end{proof}

\begin{remark}\label{rem:1.15}
\begin{enumerate}[label=\alph*)]
\item 最初, Fourier 变换只对 $L^1$ 函数有定义, 因为只有在那里积分才绝对收敛. Plancherel 定理使我们能把该算子扩张到 $L^2$, 但若 $f\in L^2$, $\hat{f}$ 只能几乎处处定义. 上面的结果称为 Hausdorff--Young 定理\index{Hausdorff--Young 定理 (Hausdorff--Young theorem)}, 它使我们能把定义扩张到 $1\leq p\leq 2$ 的 $L^p$. 这样, $L^p$ 中函数的 Fourier 变换就是 Schwartz 函数的 Fourier 变换在 $L^{p'}$ 中的极限.
\item Hausdorff--Young 不等式差不多已是我们能做到的最好结果, 至少当大小用 $L^p$ 范数来度量时如此. 下面的技巧演示了这一点: 令 $f_{\delta}(x)=f(\delta x)$. 则对固定的 $f$,
\begin{align*}
\|f_{\delta}\|_p&=\|f\|_p\,\delta^{-\frac{n}{p}}=c\,\delta^{-\frac{n}{p}} \\
\hat{f}_{\delta}(\xi)&=\delta^{-n}\hat{f}(\delta^{-1}\xi) \\
\|\hat{f}_{\delta}\|_q&=c'\,\delta^{-\frac{n}{q'}}.
\end{align*}
因此, $\|\hat{f}\|_q\leq C\|f\|_p$ 要对某个 $f$ 成立, 必须
\[
c'\,\delta^{-\frac{n}{q'}}\leq c\,\delta^{-\frac{n}{p}}
\]
对所有的 $\delta>0$ 成立. 于是 $q'=p$.
\end{enumerate}
\end{remark}

\begin{lemma}\label{lem:1.16}
设 $(T,K,\mu)$ 是 $L^p$ 的乘子. 则
\[
\|\mu\|_{\infty}\leq{}_p\|T\|.
\]
\end{lemma}

\begin{proof}
注意 ${}_p\|T\|={}_{p'}\|T\|$, 于是我们可以在 $L^p$ 与 $L^{p'}$ 之间内插. $L^2$ 总是介于两者之间, 所以该算子在 $L^2$ 上总是有界的. 而且 $\frac{1}{2}=\frac{1}{2}\bigl(\frac{1}{p}+\frac{1}{p'}\bigr)$, 故
\[
\|\mu\|_{\infty}={}_2\|T\|\leq{}_p\|T\|^{\frac{1}{2}}\,{}_{p'}\|T\|^{\frac{1}{2}}={}_p\|T\|. \qedhere
\]
\end{proof}

\begin{remark}\label{rem:1.17}
注意到 ${}_p\|T\|\leq\|K\|_1$, 正如在引理~\ref{lem:1.10} 中所指出的, 尽管对 $L^p$ 的乘子而言卷积核的 $L^1$ 范数很可能不是有限的. 但我们现在确实得到了乘子三元组 $(T,K,\mu)$ 的三个分量之间的关系:
\[
\|\mu\|_{\infty}\leq{}_p\|T\|\leq\|K\|_1.
\]
当然, 这个关系对大多数用途而言毫无价值: 只有当两端相等, 即 $\|\mu\|_{\infty}=\|K\|_1$ 时, 我们才能指望这个估计是好的. 但这意味着 $\|\hat{K}\|_{\infty}=\|K\|_1$. 而一般地 $\|\hat{K}\|_{\infty}\leq\|K\|_1$. 另一方面, 若 $K$ 是正函数,
\[
\|K\|_1=\int K=\int K(x)e^{2\pi i0x}\,dx=\hat{K}(0)\leq\|\hat{K}\|_{\infty};
\]
于是对正的 $K$,
\[
\|\hat{K}\|_{\infty}=\|K\|_1={}_p\|T\|.
\]
因此, 对正的卷积核, 乘子的全部情况是一清二楚的: 它给出某个 $L^p$ 的乘子当且仅当它给出所有 $L^p$ 的乘子, 而且此时它的 $L^1$ 范数恰是它的 $L^p$ 算子范数. 另一方面, 我们在引言中看到, 真正最有意思的乘子算子有着极不正的卷积核. 对这类乘子, $\|K\|_1$ 与 $\|\hat{K}\|_{\infty}$ 之间的差距很大, 但为了控制 $T$ 的算子范数, 我们需要对二者同时加以控制.

这是 Fourier 分析的中心问题之一: 需要对 $K$ 与 $\hat{K}$ 的同时控制.
\end{remark}

\begin{remark}\label{rem:1.18}
对 $1<p<2$ 的 $L^p$ 函数还有别的控制手段. 最简单的就是把一个 $L^p$ 函数切成小块, 使每一块属于别的某个 Lebesgue 类. 其基本例子是取
\[
G=\{x \mid |f(x)|\leq 1\};\quad B=\{x \mid |f(x)|>1\}.
\]
则 $f=f\chi_G+f\chi_B=g+b$. 这是把 $f$ 分解为一个``好''函数 $g$ 与一个``坏''函数 $b$; 好的部分在于
\begin{gather*}
\|g\|_2^2=\int_G|f|^2=\int_G|f|^{2-p}|f|^p\leq\int_G|f|^p\leq\|f\|_p^p; \\
\|b\|_1=\int|b|=\int_B|f|\leq|B|^{\frac{1}{p}}\|f\|_p.
\end{gather*}
于是我们看到, $1$ 与 $2$ 之间任一 $p$ 的 $L^p$ 中的任意函数都可以写成一个 $L^1$ 函数与一个 $L^2$ 函数之和. 若 $T$ 是线性算子, 那么仅知道 $T$ 在 $L^1$ 与 $L^2$ 上的行为, 我们便有希望追踪 $T$ 的 $L^p$ 有界性. 现在我们把这一切更仔细地做一遍.
\end{remark}

\begin{definition}\label{def:1.19}
\index{弱 $L^p$ (weak $L^p$)|emph}
可测函数 $f$ 属于\emph{弱 $L^p$} (weak $L^p$), 如果
\[
|\{x \mid |f(x)|>\lambda\}|\leq C\lambda^{-p},
\]
其中 $C$ 与 $\lambda$ 无关, $0<\lambda<\infty$.

算子 $T$ 称为\emph{弱 $(p,p)$} 的, 如果
\[
|\{x \mid |Tf(x)|>\lambda\}|\leq C\frac{\|f\|_p^p}{\lambda^p},
\]
其中 $C$ 与 $\lambda$ 和 $f$ 无关.
\end{definition}

\begin{remark}\label{rem:1.20}
\begin{enumerate}[label=\alph*)]
\item 通常我们把 $|\{f>\lambda\}|$ 作为 $|\{x \mid |f(x)|>\lambda\}|$ 的缩写; 由于我们常可取正函数, 这种记法实际上从不引起问题. 这称为 $f$ 的\emph{分布函数} (distribution function)\index{分布函数 (distribution function)}, 常记作 $\lambda(f)$.
\item 一个典型的属于弱 $L^1$ 但不属于 $L^1$ 的函数是 $f(x)=\frac{1}{x}$. 一个标准计算表明 $L^p$ 函数总是弱 $L^p$ 的:
\begin{align*}
\lambda^p|\{x \mid |f(x)|>\lambda\}|&=\int_{\{f>\lambda\}}\lambda^p\,dx \\
&\leq\int_{\{f>\lambda\}}|f(x)|^p\,dx\leq\int|f|^p=\|f\|_p^p.
\end{align*}
\item 若弱 $L^p$ 想要在分析正经的 $L^p$ 函数时派上用场, 分布函数就必须与 $L^p$ 范数联系起来. 其关系为:
\[
\|f\|_p^p=p\int_0^{\infty}|\{x \mid |f(x)|>\lambda\}|\;\lambda^{p-1}\,d\lambda.
\]
最简单的证明使用简单函数, 这把整件事化为 $f=\chi_I$ (其中 $I$ 是一个区间) 的情形. 此时 $\|f\|_p^p=|I|$, 且当 $\lambda>1$ 时 $|\{x \mid |f(x)|>\lambda\}|=0$, 当 $0<\lambda\leq 1$ 时它等于 $|I|$. 于是该积分为
\[
p|I|\int_0^1\lambda^{p-1}\,d\lambda=|I|. \qedhere
\]
\end{enumerate}
\end{remark}

\begin{theorem}\label{thm:1.21}
\index{内插!Marcinkiewicz (Marcinkiewicz interpolation)}
设 $T$ 是弱 $(p_0,p_0)$ 且弱 $(p_1,p_1)$ 的线性算子. 则 $T$ 在介于 $p_0$ 与 $p_1$ 之间的 $p$ 上是有界算子.
\end{theorem}

\begin{proof}
为了给作者们省点力气, 我们只做 $1\leq p_0<p<p_1<\infty$ 的情形. 如同在注~\ref{rem:1.18} 中那样, 令:
\begin{gather*}
G=\Bigl\{|f|\leq\frac{\lambda}{2}\Bigr\},\quad B=\Bigl\{|f|>\frac{\lambda}{2}\Bigr\}, \\
g=f\chi_G;\quad b=f\chi_B,
\end{gather*}
并注意到 $|Tf|=|Tg+Tb|\leq|Tg|+|Tb|$. 不幸的是, 分布函数不是线性的, 所以 $Tf$ 的这个分解不能立即派上用场. 但线性有一个替代品:
\begin{align*}
|\{|Tf|>\lambda\}|&\leq|\{|Tg|+|Tb|>\lambda\}| \\
&\leq\Bigl|\Bigl\{|Tg|>\frac{\lambda}{2}\Bigr\}\Bigr|+\Bigl|\Bigl\{|Tb|>\frac{\lambda}{2}\Bigr\}\Bigr| \\
&\leq C_0\lambda^{-p_0}\int|b|^{p_0}+C_1\lambda^{-p_1}\int|g|^{p_1}.
\end{align*}
现在计算:
\[
\|Tf\|_p^p=p\int|\{|Tf|>\lambda\}|\,\lambda^{p-1}\,d\lambda,
\]
发现它被两项之和控制; 第一项是:
\begin{align*}
C_0p&\int_0^{\infty}\lambda^{p-p_0-1}\int_{\{|f|\geq\frac{\lambda}{2}\}}|f|^{p_0}\,dx\,d\lambda \\
&=C_0p\iint_0^{2|f|}\lambda^{p-p_0-1}\,d\lambda\,|f(x)|^{p_0}\,dx
=\frac{2^{p-1}p\,C_0}{p-p_0}\|f\|_p^p.
\end{align*}
类似地, 第二项被下式控制:
\[
\frac{2^{p-1}p\,C_1}{p_1-p}\|f\|_p^p.
\]
\end{proof}

\noindent\textbf{注:} 这个结果归功于 Marcinkiewicz; 注意到当 $p$ 趋向 $p_0$ 或 $p_1$ 时, 对乘子的界变得非常糟糕. 它们变坏的速度在端点处提供了额外的信息. 要记住的基本例子是 $\frac{1}{x}$: 这里 $L^1$ 范数不是有限的, 但发散性由 $\log x$ 度量. 一个类似的结果在 Marcinkiewicz 内插定理中成立; 参见章末的评注.

\begin{lemma}\label{lem:1.22}
若 $f$, $g$ 属于 $\mathcal{S}$, 则
\begin{enumerate}[label=\alph*)]
\item $\widehat{\bar{f}}(\xi)=\overline{\hat{f}(-\xi)}$.
\item 若 $f_{\delta}(x)=f(\delta x)$, 则 $\hat{f}_{\delta}(\xi)=\delta^{-n}\hat{f}(\delta^{-1}\xi)$.
\item 若 $R$ 表示 $\mathbb{R}^n$ 的一个旋转, 且 $f_R(x)=f(Rx)$, 则 $\hat{f}_R(\xi)=\hat{f}(R^{-1}\xi)$.
\item $\widehat{\tau_y f}(\xi)=e^{2\pi iy\xi}\hat{f}(\xi)$.
\item $\widehat{fg}(\xi)=\hat{f}\star\hat{g}(\xi)$.
\item $\widehat{f\star g}(\xi)=\hat{f}(\xi)\hat{g}(\xi)$.
\item $D_{\xi_j}\hat{f}(\xi)=-2\pi i\,\widehat{x_j f}(\xi)$.
\end{enumerate}
\end{lemma}

\begin{proof}
所有这些都可以由变量替换与定义直接得出. 以前没有见过这些结果的读者, 请务必亲手算出细节. 这就是在和 Fourier 变换交朋友.
\end{proof}

\begin{theorem}\label{thm:1.23}
\index{de Leeuw 定理 (de Leeuw theorem)}
\index{Lebesgue 点 (Lebesgue point)}
设 $(T,\mu,K)$ 是 $L^p(\mathbb{R}^n)$ 的 Fourier 乘子三元组. 则:
\begin{enumerate}[label=\alph*)]
\item $\operatorname{Re}(\mu)$ 与 $\operatorname{Im}(\mu)$ 都是 $L^p$ 乘子.
\item 设 $T_{\delta}$ 是以 $\mu(\delta\xi)$ 为乘子的算子. 则 ${}_p\|T_{\delta}\|={}_p\|T\|$.
\item 设 $T_{\eta}$ 是以 $\mu(\xi-\eta)$ 为乘子的算子. 则 ${}_p\|T_{\eta}\|={}_p\|T\|$.
\item 若 $\tilde{\psi}\in L^1$, 设 $T_{\psi}$ 是以 $\psi\mu$ 为乘子的算子. 则 ${}_p\|T_{\psi}\|\leq\|\psi\|_1\,{}_p\|T\|$.
\item 若 $\psi\in L^1$, 设 $T_{\psi}$ 是以 $\psi\star\mu$ 为乘子的算子. 则 ${}_p\|T_{\psi}\|\leq\|\psi\|_1\,{}_p\|T\|$.
\item 设 $\xi_1,\xi_2,\ldots,\xi_{n-k}$ 固定, 并定义
\[
\mathbb{R}^{k_0}=\{\xi\in\mathbb{R}^n \mid \xi_1\cdot\xi+\cdots+\xi_{n-k}\cdot\xi=\lambda\}.
\]
设 $\mu_0=\mu|_{\mathbb{R}^{k_0}}$, 并设 $T_0$ 是 $L^p(\mathbb{R}^{k_0})$ 上以 $\mu_0$ 为乘子的算子. 则对几乎所有的 $\lambda$, ${}_p\|T_0\|\leq{}_p\|T\|$.
\end{enumerate}
\end{theorem}

\begin{proof}
a)--d) 的证明都是容易的, 而且都利用了这样一个事实: Fourier 变换在共轭、伸缩、旋转、平移等等变换下行为是可以预料的. 我们做一个典型的证明, 以期读者获得同样的体验:

d)
\begin{align*}
T_{\eta}f(x)&=\int\mu(\xi-\eta)e^{2\pi ix\xi}\hat{f}(\xi)\,d\xi \\
&=\int\mu(\xi)e^{2\pi ix(\xi+\eta)}\hat{f}(\xi+\eta)\,d\xi \\
&=e^{2\pi ix\eta}\int\mu(\xi)e^{2\pi ix\xi}\,\widehat{\bigl(e^{-2\pi ix\eta}f(x)\bigr)}(\xi)\,d\xi=JTJ^{-1}(f),
\end{align*}
其中 $J$ 是乘以 $e^{2\pi ix\eta}$ 的等距算子.

e) 我们将使用对偶性, 这是 Fourier 乘子游戏中一个非常标准的技巧. 选取 $\mathcal{S}$ 中的 $f$, $g$, 使 $\|f\|_p=\|g\|_{p'}=1$, 并计算:
\begin{align*}
\int T_{\psi}f\,\bar{g}&=\int\widehat{T_{\psi}f}\;\widehat{\bar{g}}
=\int\psi\star\mu\,\hat{f}\;\widehat{\bar{g}} \\
&=\iint\psi(\eta)\mu(\xi-\eta)\hat{f}(\xi)\widehat{\bar{g}}(\xi)\,d\eta\,d\xi
=\int\psi(\eta)\int\mu(\xi-\eta)\hat{f}(\xi)\widehat{\bar{g}}(\xi)\,d\xi\,d\eta \\
&\leq\|\psi\|_1\left\|\int\mu(\xi-\eta)\hat{f}(\xi)\widehat{\bar{g}}(\xi)\,d\xi\right\|_{\infty,\eta}\|g\|_{p'}
\leq\|\psi\|_1\,{}_p\|T\|,
\end{align*}
最后一步用了上面的 d).

f) 我们只对 $n=2$, $k=1$ 的情形给出证明, 并且甚至只限于 $\xi_1=(0,1)$. 更高维的情形非常类似, 只是记号要多得多; 而 $\mathbb{R}^2$ 中更一般的情形可以由这一情形经旋转与平移得到.

在我们的特殊情形, $\mathbb{R}^{1_0}=\{(\cdot,\lambda)\}$, 我们必须证明: 对几乎所有的 $\lambda$, $\mu(\xi_1,\lambda)$ 是 $L^p(\mathbb{R}^1)$ 的乘子. 由于这个 $\lambda$ 可能会与先前的记号发生实实在在的摩擦, 我们干脆证明: 对几乎所有的 $\xi_2$, $\mu(\xi_1,\xi_2)$ 是 $L^p(\mathbb{R}^1)$ 的乘子.

第一个痛点是 $\mu$ 只是可测的, 而 $\mu$ 在一维超平面上的限制并没有很好的定义. 我们将使用乘子理论中另一个标准技巧. 先假设 $\mu$ 连续. 设 $f_i$, $g_i$ ($i=1,2$) 满足 $\|f_i\|_p=\|g_i\|_{p'}=1$; 令 $f(x_1,x_2)=f_1(x_1)f_2(x_2)$, $g$ 类似. 则:
\begin{align*}
&\int\left[\int\mu(\xi_1,\xi_2)\hat{f}_1(\xi_1)\hat{g}_1(\xi_1)\,d\xi_1\right]\hat{f}_2(\xi_2)\hat{g}_2(\xi_2)\,d\xi_2 \\
&\quad=\int\mu(\xi)\hat{f}(\xi)\hat{g}(\xi)\,d\xi \\
&\quad=\int\widehat{Tf}\,\hat{g}=\int Tfg\leq\|Tf\|_p\|g\|_{p'}\leq{}_p\|T\|.
\end{align*}
定义
\[
\sigma(\xi_2)=\int\mu(\xi_1,\xi_2)\hat{f}_1(\xi_1)\hat{g}_1(\xi_1)\,d\xi_1.
\]
把上面的计算改写, 得到:
\[
\int\sigma(\xi_2)\hat{f}_2(\xi_2)\hat{g}_2(\xi_2)\,d\xi_2\leq{}_p\|T\|.
\]
因此, $\sigma$ 是 $L^p(\mathbb{R}^1)$ 的乘子. 但引理~\ref{lem:1.16} 告诉我们 $\|\sigma(\xi_2)\|_{\infty}\leq{}_p\|T\|$. 由于 $\mu$ 连续, $\sigma$ 也连续, 故对所有的 $\xi_2$, $\sigma(\xi_2)\leq{}_p\|T\|$. 改写一下,
\begin{align*}
\left|\int\mu(\xi_1,\xi_2)\hat{f}_1(\xi_1)\hat{g}_1(\xi_1)\,d\xi_1\right|&\leq{}_p\|T\| \\
\left|\int\mu_0(\xi_1)\hat{f}_1(\xi_1)\hat{g}_1(\xi_1)\,d\xi_1\right|&\leq{}_p\|T\| \\
\left|\int T_0f_1\,g_1\right|&\leq{}_p\|T\| \\
{}_p\|T_0\|&\leq{}_p\|T\|.
\end{align*}
我们假设了 $\mu$ 连续; 现在用一个标准技巧把连续性假设换成几乎处处的条件. 设 $S$ 表示单位正方形, 并令 $\psi_{\delta}(\xi)=\delta^{-2}\chi_S(\delta^{-1}\xi)$. 则 $\psi_{\delta}\in L^1$, $\mu\in L^{\infty}$, 故 $\psi_{\delta}\star\mu$ 连续, 于是对每个 $\xi_2$,
\begin{align*}
\int\psi_{\delta}\star\mu\,(\xi_1,\xi_2)\hat{f}_1(\xi_1)\hat{g}_1(\xi_1)\,d\xi_1
&\leq{}_p\|T\psi_{\delta}\| \\
&\leq{}_p\|T\|\,\|\psi_{\delta}\|_1={}_p\|T\|.
\end{align*}
现在, 对几乎每个 $\xi_2$, 几乎每个 $\xi_1$ 都是 $\mu_0$ 的 Lebesgue 点. 定理中所说的 $\xi_2$ 正是这些; 取定一个. Lebesgue 点 $\xi_1$ 的定义是
\begin{align*}
\lim_{\delta\to0}\psi_{\delta}\star\mu(\xi_1,\xi_2)
&=\lim_{\delta\to0}\delta^{-2}\int_{-\delta}^{\delta}\int_{-\delta}^{\delta}\mu(\eta_1-\xi_1,\eta_2-\xi_2)\,d\eta_1\,d\eta_2 \\
&=\mu(\xi_1,\xi_2).
\end{align*}
对这个取定的 $\xi_2$, $\psi_{\delta}\star\mu(\xi_1,\xi_2)$ 在 $\xi_1$ 中几乎处处收敛于 $\mu(\xi_1,\xi_2)=\mu_0(\xi_1)$. 当然 $(\psi_{\delta}\star\mu)\hat{f}_1\hat{g}_1$ 也收敛; 要使积分收敛, 我需要控制收敛定理. 但是
\begin{align*}
|(\psi_{\delta}\star\mu)\hat{f}_1\hat{g}_1|&\leq\|\psi_{\delta}\star\mu\|_{\infty}\,|\hat{f}_1|\,|\hat{g}_1| \\
&\leq\|\psi_{\delta}\|_1\|\mu\|_{\infty}\,|\hat{f}_1\hat{g}_1|.
\end{align*}
因此,
\[
\int\mu_0(\xi_1)\hat{f}_1(\xi_1)\hat{g}_1(\xi_1)\,d\xi_1
=\lim_{\delta\to0}\int\psi_{\delta}\star\mu\,\hat{f}_1\hat{g}_1\leq{}_p\|T\|,
\]
从而对几乎所有的 $\xi_2$, ${}_p\|T_0\|\leq{}_p\|T\|$.
\end{proof}

\begin{remark}\label{rem:1.24}
\index{乘子!的变形 (deformation of multipliers)}
这些结果之所以重要, 是因为它们让我们能把乘子变形为我们能够处理的东西. 例如, 若 $\psi$ 属于 $\mathcal{S}$, $\mu$ 是一个乘子, 则 $\psi\mu$ 与 $\psi\star\mu$ 都是乘子. 我们在上面定理的证明中已经用过这一点: 我们把一个任意的乘子毫无副作用地变成了连续函数. 一般地, $\mu$ 只属于 $L^{\infty}$, $K$ 是某个分布; 但若我们用一个紧支集函数 $\psi$ 去修改它, 则 $K_{\psi}$ (即 $\psi\mu$ 的逆 Fourier 变换) 是一个连续函数. 在下面的讨论中, 把一个糟糕的分布变成一个漂亮正则的函数, 将是我们最喜爱的技巧.

真正的重要性是哲学层面的. $(T,K,\mu)$ 的 $L^p$ 有界性受 $K$ 的大小影响, 而 $K$ 反映了 $\check{\mu}=K$ 的光滑性. 但对 $\hat{\mu}$ 的计算是 $\mu$ 关于所有指数函数的整体平均, 因此很难理解 $\mu$ 的局部性质如何能影响 Fourier 变换 $K$. 这正是要有 $\psi$ 或 $\psi\star\mu$ 在场的意义: 前者把 $\mu$ 变成紧支集函数, 后者把 $\mu$ 变成更光滑的函数.
\end{remark}

\section{乘子与级数的求和性}
\label{sec:1.3}

\begin{theorem}\label{thm:1.25}
设 $\mu\in L^{\infty}(\mathbb{R}^n)$ 连续, 并令
\[
S_Rf(\theta)=\sum\mu\Bigl(\frac{n}{R}\Bigr)\hat{f}(n)e^{2\pi in\theta}.
\]
若算子 $S_R$ 在 $L^p(\mathbb{T}^n)$ ($R>0$) 上一致有界, 则 $\mu$ 是 $L^p(\mathbb{R}^n)$ 的乘子, 其算子为 $T$, 且
\[
{}_p\|T\|\leq\sup_R\,{}_p\|S_R\|.
\]
\end{theorem}

\begin{proof}
我们首先必须能从 $L^p(\mathbb{T}^n)$ 上的函数过渡到 $L^p(\mathbb{R}^n)$ 上的函数, 然后转移算子. 设 $f$ 与 $g$ 是 $\mathbb{R}^n$ 上光滑的紧支集函数, 注意到当 $R$ 充分大时 $f_R(x)=f(Rx)$ 支在 $\mathbb{R}^n$ 的单位立方体中, 因此可以把它看作生活在 $n$ 维环面上. 为今后使用, 我们计算这对 $L^p$ 范数的影响:
\[
\|f_R\|_p\,\|g_R\|_{p'}=R^{-n}\|f\|_p\,\|g\|_{p'}.
\]
要从 Fourier 级数过渡到 Fourier 积分, 我们使用一个古老的想法: Fourier 积分不过是周期非常大的 Fourier 级数; 关键在于
\[
R^{-n}\sum\mu\Bigl(\frac{n}{R}\Bigr)\hat{f}\Bigl(\frac{n}{R}\Bigr)\hat{g}\Bigl(\frac{n}{R}\Bigr)
\]
是 $\int\mu(\xi)\hat{f}(\xi)\hat{g}(\xi)\,d\xi$ 的 Riemann 和逼近. 我们要做的只是安排好参数. 具体安排如下:
\[
\int S_Rf_R(\theta)g_R(\theta)\,d\theta\leq{}_p\|S_R\|\,\|f_R\|_p\,\|g_R\|_{p'},
\]
因此
\[
R^n\int S_R(f_R)g_R\leq{}_p\|S_R\|\,\|f\|_p\,\|g\|_{p'};
\]
另一方面,
\[
R^n\int S_R(f_R)g_R=R^n\sum\mu\Bigl(\frac{n}{R}\Bigr)\hat{f}_R(n)\hat{g}_R(n);
\]
而且我们可以用 Fourier 积分把 $\hat{f}_R(n)$ 表示为
\[
\hat{f}_R(n)=R^{-n}\hat{f}\Bigl(\frac{n}{R}\Bigr).
\]
于是得到
\[
R^{-n}\sum\mu\Bigl(\frac{n}{R}\Bigr)\hat{f}\Bigl(\frac{n}{R}\Bigr)\hat{g}\Bigl(\frac{n}{R}\Bigr)
\leq{}_p\|S_R\|\,\|f\|_p\,\|g\|_{p'}.
\]
鉴于左端收敛,
\[
\int\mu(\xi)\hat{f}(\xi)\hat{g}(\xi)\,d\xi\leq{}_p\|S_R\|\,\|f\|_p\,\|g\|_{p'},
\]
以及
\[
\int Tfg\leq\sup_R\,{}_p\|S_R\|\,\|f\|_p\,\|g\|_{p'}.
\]
这蕴含 $T$ 在 $L^p$ 上有界, 其范数不超过 ${}_p\|S_R\|$.
\end{proof}

\begin{theorem}\label{thm:1.26}
设 $\mu$ 连续, 且 $T$ 是 $L^p$ 的乘子. 则
\[
Sf(\theta)=\sum\mu(n)\hat{f}(n)e^{2\pi in\theta}
\]
是 $L^p(\mathbb{T}^n)$ 上的有界算子.
\end{theorem}

\begin{lemma}\label{lem:1.27}
设 $\omega_{\delta}(y)=e^{-\pi\delta|y|^2}$. 若 $f$ 连续且周期, 则
\[
\lim_{\epsilon\to0}\epsilon^{\frac{n}{2}}\int f(x)e^{-\epsilon\pi|x|^2}\,dx=\int f(x)\,dx.
\]
\end{lemma}

\begin{lemma}\label{lem:1.28}
设 $\mu$, $T$, $S$ 如定理~\ref{thm:1.26} 所述. 设 $P$ 与 $Q$ 为三角多项式. 若 $\alpha+\beta=1$, 则
\[
\lim_{\epsilon\to0}\epsilon^{\frac{n}{2}}\int T(P\omega_{\epsilon\alpha})\,\overline{[Q(x)\omega_{\epsilon\beta}]}(x)\,dx=\int(SP)(x)\bar{Q}(x)\,dx.
\]
\end{lemma}

\begin{proof}[引理~\ref{lem:1.27} 的证明]
由于结果对 $f$ 是线性的, 我们可以假设 $f$ 是三角多项式, 进而可以假设 $f(x)=e^{2\pi imx}$. $m=0$ 的情形是平凡的, 于是可以假设 $\int f(x)\,dx=0$. 但是
\[
\int e^{2\pi imx}e^{-\epsilon\pi|x|^2}\,dx=\epsilon^{-\frac{n}{2}}e^{-\pi\frac{|m|^2}{\epsilon}},
\]
故
\[
\lim_{\epsilon\to0}\epsilon^{\frac{n}{2}}\int f(x)e^{-\pi\epsilon|x|^2}\,dx=0. \qedhere
\]
\end{proof}

\begin{proof}[引理~\ref{lem:1.28} 的证明]
由于结果是线性的, 我们可以取
\[
P(x)=e^{2\pi imx};\quad Q(x)=e^{2\pi ikx}.
\]
则:
\begin{align*}
&\epsilon^{\frac{n}{2}}\int T(P\omega_{\epsilon\alpha})\,\overline{[Q(x)\omega_{\epsilon\beta}]}(x)\,dx \\
&=\epsilon^{\frac{n}{2}}\int\mu(\xi)\,\widehat{(P\omega_{\epsilon\alpha})}(\xi)\,\overline{\widehat{Q\omega_{\epsilon\beta}}}(\xi)\,d\xi \\
&=\epsilon^{\frac{n}{2}}\int\mu(\xi)e^{-\pi\frac{|\xi-m|^2}{\epsilon\alpha}}(\epsilon\alpha)^{-\frac{n}{2}}\,e^{-\pi\frac{|\xi-k|^2}{\epsilon\beta}}(\epsilon\beta)^{-\frac{n}{2}}\,d\xi.
\end{align*}
若 $m\neq k$, 这个极限可以算出为零, 于是我们得到
\[
\int S(P)\bar{Q}=\int\mu(m)e^{2\pi i(m-k)\theta}\,d\theta=0.
\]
若 $m=k$, 极限为
\[
\lim_{\epsilon\to0}(\epsilon\alpha\beta)^{-\frac{n}{2}}\int\mu(\xi)e^{-\pi|\xi-m|^2\bigl(\frac{1}{\alpha}+\frac{1}{\beta}\bigr)\frac{1}{\epsilon}}\,d\xi.
\]
由于我们把事情安排得使 $\alpha+\beta=1$, 故 $\frac{1}{\alpha}+\frac{1}{\beta}=\frac{1}{\alpha\beta}$, 于是极限为 $\mu(m)$. 但是
\[
\int S(P)\bar{Q}=\int\mu(m)e^{0}\,d\theta=\mu(m). \qedhere
\]
\end{proof}

\begin{proof}[定理~\ref{thm:1.26} 的证明]
三角多项式在 $L^p(\mathbb{T}^n)$ 中稠密, 因此我们可以用它们来计算算子范数. 于是
\begin{align*}
\int S(P)\bar{Q}&=\lim_{\epsilon\to0}\epsilon^{\frac{n}{2}}\int T(P\omega_{\epsilon\alpha})\,\bar{Q}\omega_{\epsilon\beta} \\
&\leq\limsup\,{}_p\|T\|\,\epsilon^{\frac{n}{2}}\|P\omega_{\epsilon\alpha}\|_p\,\|Q\omega_{\epsilon\beta}\|_{p'} \\
&=\limsup\,{}_p\|T\|\left(\epsilon^{\frac{n}{2}}\int|P(x)|^pe^{-\pi p\alpha|x|^2}\right)^{\frac{1}{p}}\left(\epsilon^{\frac{n}{2}}\int|Q|^{p'}e^{-\pi p'\beta|x|^2}\right)^{\frac{1}{p'}}.
\end{align*}
现在我们可以选取 $\alpha=\frac{1}{p}$, $\beta=\frac{1}{p'}$; 于是引理~\ref{lem:1.27} 告诉我们
\[
\int S(P)\bar{Q}\leq{}_p\|T\|\left(\int|P|^p\right)^{\frac{1}{p}}\left(\int|Q|^{p'}\right)^{\frac{1}{p'}}. \qedhere
\]
\end{proof}

\begin{theorem}\label{thm:1.29}
设 $\mu$ 连续, 且 $S_R$, $T$ 如上. 则 $S_R(f)$ 对所有 $f\in L^p(\mathbb{T}^n)$ 在 $L^p(\mathbb{T}^n)$ 中收敛, 当且仅当 $T$ 在 $L^p(\mathbb{R}^n)$ 上有界.
\end{theorem}

\begin{proof}
共鸣定理正是为了告诉我们``一个过程在一大堆元素上收敛蕴含算子一致有界''而设计的. 在此情形, 假设 $S_R(f)$ 总是收敛. 则算子 $f\mapsto S_R(f)$ 在 $L^p$ 上一致有界. 由前面的引理, $T$ 在 $L^p(\mathbb{R}^n)$ 上有界.

反之, 假设 $T$ 在 $L^p(\mathbb{R}^n)$ 上有界. 则如定理~\ref{thm:1.23}c) 中定义的 $T_{\epsilon}$ 也有界, 范数为 ${}_p\|T\|$; 于是由定理~\ref{thm:1.26}, 算子 $T_{\epsilon}$ 在 $L^p(\mathbb{T}^n)$ 上一致有界. $S_R(f)$ 在 $L^p(\mathbb{T}^n)$ 中的收敛现在可以由稠密集上的收敛得出. 设 $P$ 为三角多项式;
\[
\lim_{R\to\infty}S_RP(\theta)=\lim_{R\to\infty}\sum\mu\Bigl(\frac{n}{R}\Bigr)a_ne^{2\pi in\theta}
=\mu(0)\sum a_ne^{2\pi in\theta}=\mu(0)P(\theta). \qedhere
\]
\end{proof}

\begin{remark}\label{rem:1.30}
至此我们完成了对收敛结果的分析; 我们心里应当想着这样的情形: 当 $|\xi|<1$ 时 $\mu(\xi)=1$, 当 $|\xi|\geq 1$ 时 $\mu(\xi)=0$. 此时 $S_R$ 表示 $f$ 的 Fourier 级数的第 $R$ 个部分和. Fourier 级数的 $L^p$ 收敛等价于一个 Fourier 乘子的 $L^p$ 有界性. 我们将在接下来几章中研究这种有界性.
\end{remark}

\section{第 1 章评注}
\label{sec:1.4}

\noindent (1.1):\quad 传统上, Fourier 乘子算子是以另一种方式定义的. 人们先分析在 $L^p$ 上有界且与平移交换的线性算子 $T$: $\tau(y)T=T\tau(y)$, 并证明这样的 $T$ 由与一个分布的卷积给出. 然后证明 $T$ 在 $L^{p'}$ 上也有界, 从而在 $L^2$ 上有界, 并且在 Fourier 变换一侧由一个乘法给出. 这种做法的一个优点是, 不必像我们这样一开始就假设 $\mu\in L^{\infty}$. 这一思路在 Stein 与 Weiss~\cite{steinweiss} 第~I 章第~3 节中被系统地采用.

这引出了一个问题: 为什么与平移交换的算子应当由卷积给出? 原因是平移本身就是一个与平移交换的算子, 因为 $\mathbb{R}^n$ 是 Abel 群. 若我们用 $\delta_y(f)=f(y)$ 定义有限测度 $\delta_y$, 则
\[
(\delta_y\star f)(x)=\delta_y(\tau_x\tilde{f})(y)=\tilde{f}(y-x)=f(x-y).
\]
所以平移实际上由与一个有限测度的卷积给出. 对平移取和以及取和的极限, 我可以得到更多这类算子. 平移之和的极限可能恰好收敛于一个积分; 例如 $\int K(y)\delta_y\,dy$, 这个分布在函数上的作用是 $\int K(x-y)f(y)\,dy=(K\star f)(x)$. 寓意是: 卷积不过是平移的加权平均. 这是贯穿本书的一个重要直观想法; 一个典型的例子是关于区间特征函数的平均. 例如, 在证明微积分基本定理时, 必须计算
\[
\frac{1}{h}\int_x^{x+h}f(y)\,dy
\]
的极限. 若 $I$ 表示区间 $[0,h]$, 那么变量替换表明这等于 $\frac{1}{h}\int\chi_I(x-y)f(y)\,dy$. 所以微积分基本定理实际上分析的是卷积算子的极限.

这种哲学的问题在于, 没有理由指望 Fourier 变换会进入这幅图景. 思考这些问题的最简单的场所是 $L^2(\mathbb{T}^n)$. 平移是 $L^2$ 上一族交换的有界自伴算子, 因此谱定理说它们可以被同时对角化. 由于
\begin{align*}
\tau_{\psi}\bigl(e^{2\pi in\theta}\bigr)&=e^{2\pi in(\theta-\psi)} \\
&=e^{-2\pi in\psi}e^{2\pi in\theta},
\end{align*}
可知 $\mathbb{C}e^{2\pi in\theta}$ 是所有平移的联合特征空间. 由于指数函数在 $L^2$ 中完备, 这就是 $L^2$ 按特征空间的分解.

沿着这条思路继续下去, 谱定理由 $f\in L^2$ 在各特征空间上的投影恢复出 $f$. 到由单位向量 $\Phi$ 张成的一维特征空间的投影一般由 $\langle f,\Phi\rangle\Phi$ 给出; 在此情形即为
\[
\left[\int f(\theta)\,\overline{\bigl(e^{2\pi in\theta}\bigr)}\,d\theta\right]e^{2\pi in\theta}=\hat{f}(n)e^{2\pi in\theta}.
\]
这意味着投影之和在 $L^2$ 中收敛于 $f$; 也就是说,
\[
\sum\hat{f}(n)e^{2\pi in\theta}\to f
\]
在 $L^2$ 中成立. 谱定理``解释''了 Fourier 级数收敛到原函数这一事实.

仍在这个抽象的框架中, 取一个在 $L^2$ 上有界且与平移交换的算子 $T$. 那么 $T$ 保持平移算子的特征空间, 因此 $T$ 也被指数函数对角化. 对角算子由其特征值决定; 这些特征值由公式
\[
T\bigl(e^{2\pi in\theta}\bigr)=\mu(n)e^{2\pi in\theta}
\]
给出. $T$ 在一般函数 $f$ 上的作用可以借助对角化表示为
\[
T(f)(\theta)=\sum\mu(n)e^{2\pi in\theta},
\]
于是 $T$ 显露出是一个 Fourier 乘子算子. 我们会预期 $T$ 的 $L^2$ 范数是 $T$ 的最大特征值, 即诸数 $|\mu(n)|$ 中的最大者; 简言之: ${}_2\|T\|=\|\mu\|_{\infty}$.

\noindent (1.12):\quad 三线定理的证明见 Stein 与 Weiss~\cite{steinweiss}, 第~V 章引理~1.4.

\noindent (1.13):\quad 这个结果归功于 M.~Riesz; 这里的证明归功于 Thorin 与 Tamarkin 以及 Zygmund~\cite{zygmund}. 该结果通常称为 Riesz--Thorin 内插定理. 它可以扩张为从 $L^{p_i}$ 到 $L^{q_i}$ 的映射, 但带有限制 $p_i\leq q_i$.

\noindent (1.19):\quad 弱 $L^p$ 空间是 Lorentz 空间的特例; 我们不是对分布函数的增长要求一致界, 而是施加可积性条件. 为 $L^p$ 与弱 $L^p$ 空间发展的大量理论都可以在 Lorentz 空间的框架中进行. 细致的讨论见 Stein 与 Weiss~\cite{steinweiss} 第~V 章, 或 Hunt~\cite{hunt}.

\noindent (1.21):\quad 这个结果归功于 Marcinkiewicz~\cite{marcinkiewicz-multipliers}. 与 Riesz--Thorin 定理一样, 它可以扩张为从 $L^{p_i}$ 到 $L^{q_i}$ 的映射. 与 Riesz--Thorin 定理不同的是, 限制 $p_i\leq q_i$ 是不必要的.

\noindent (1.23):\quad 这个结果归功于 K.~de Leeuw~\cite{deleeuw}; 这里的证明归功于 Jodeit~\cite{jodeit}. 整个结果在 $\mathbb{T}^n$ 上要清楚得多; 也就是,
\begin{align*}
\bigl(T_0f\bigr)(\theta_1)=\sum_{n_1}\mu(n_1,n_2)e^{2\pi in\theta_1}\hat{f}(\theta_1)
  \\
=\sum_{n_1}\sum_{n_2}\mu(n_1,n_2)e^{2\pi in\theta_1}\hat{f}(n_1)\bigl[e^{2\pi in_2\theta_2}\bigr](n_2)=T\bigl[f(\theta_1)e^{2\pi in_2\theta_2}\bigr].
\end{align*}
由于环面的测度为 $1$, 容易把一个变量的函数的 $L^p$ 范数与两个变量的函数的 $L^p$ 范数联系起来.

\noindent (1.25), (1.26):\quad 这些结果最早出现在 Stein 与 Weiss 的工作中; 见~\cite{steinweiss} 第~VII 章. 注意, 如下陈述的结果并不适用于由 Fourier 级数求和产生的那些乘子; 在那种情形, 乘子 $\mu$ 是一个集合的特征函数. 更细心一点的表述表明, 它们适用于使 $\mathbb{Z}^n$ 中每个格点都是 Lebesgue 点的乘子.

这整个一节引出了一个问题: 为什么要费心考察 $S_R$ 在 $L^p$ 中的收敛? 毕竟, 我真正想知道的是 $S_Rf(\theta)$ 对每个 $\theta$ 收敛到 $f(\theta)$. $L^p$ 与现实有什么关系?

首先, 第~\ref{chap:0} 章中那个 Fourier 级数在一点不收敛的连续函数的例子表明, 逐点收敛是奢望. Georg Cantor 为构造其 Fourier 级数在越来越复杂的集合上发散的连续函数, 发展了他的可数集与不可数集理论. 该理论中基本的现代结果是: 给定一个零测集, 存在 $L^1$ 中的函数, 其 Fourier 级数在该集合上发散.

有两条路可走. 第一条路是分析使函数的 Fourier 级数在一点收敛所需的最小正则性. 关于这个问题已有大量重要工作; 答案是复杂的, 正如与之类似的问题``一个实数列满足什么条件才能保证相应的级数收敛''的答案一样复杂. 正如学过收敛判别法的学生所知, 收敛级数与发散级数之间并没有一条清晰的分界线.

我们选择另一条路: 既然几乎处处收敛已是我们所能期望的最好结果, 那么对 $f$ 施加什么样的大小条件才能保证 $f$ 的 Fourier 级数几乎处处收敛到 $f$? 这里有一组出人意料的简单答案. 第一个结果归功于 Calderón (见 Zygmund~\cite{zygmund}) 与 Stein~\cite{stein-limits}. 若 $1\leq p\leq 2$, 且 $S_Rf(\theta)$ 对 $L^p(\mathbb{T}^n)$ 中所有的 $f$ 对几乎每个 $\theta$ 收敛到 $f$, 则算子 $S^{\ast}f(\theta)=\sup_R|S_Rf(\theta)|$ 是弱 $(p,p)$ 的. 这正是你希望泛函分析为你做的事情: 你知道几乎处处收敛对某个稠密的函数集成立; 要扩张到所有函数, 就需要对某个量的一致估计. 由于我们想要的是逐点收敛, 正确的一致性不是算子在 $L^p$ 上的一致有界性, 而是某个逐点一致算子 $S^{\ast}$ 在 $L^p$ 中的有界性.

但这有几个推论. 若收敛对某个 $p_0<2$ 的 $L^{p_0}$ 中所有函数成立, 则对满足 $p_0\leq p\leq\infty$ 的所有 $p$, 逐点收敛都会发生, 因为所有这些 $L^p$ 都包含于 $L^{p_0}$. 由 Marcinkiewicz 内插定理, $S^{\ast}$ 在 $p_0<p<2$ 时是强有界的. 但 $|S_Rf(\theta)|\leq S^{\ast}f(\theta)$, 于是算子 $S_R$ 在这个 $p$ 范围上也一致有界. 这种分析是 C.~Fefferman 证明多元 Fourier 级数发散的一个关键组成部分.

寓意在于: $L^p$ 的结果不仅有助于得到逐点结果, 而且是必不可少的. 关于 Fourier 级数之外其他类似结果成立的领域的综述, 见 Gilbert~\cite{gilbert} 的文章.

% 第 2 章 The Hilbert Transform (原书第 27-60 页)
\chapter{Hilbert 变换}
\label{chap:2}

本章是全书最长也最重要的一章. 我们只讨论一个主题: Fourier 级数的一维求和理论. 我们要做的, 是建立 $(T,K,\mu)$ 如何相互作用的基本模型. 这意味着我们需要研究一个函数与其 Fourier 变换之间的关系.

\section{引言}
\label{sec:2.1}

我们先来分析一类 Fourier 乘子, 即 Bochner--Riesz 平均. 它们的优点在于提供了一种受到精确控制而又可以变化的光滑性. 它们构成一个模型, 用于研究光滑性如何影响 Fourier 变换.

\begin{definition}\label{def:2.1}
若 $\lambda>-1$, 定义
\[
\mu_{\lambda}(\xi)=\bigl(1-|\xi|^{2}\bigr)_{+}^{\lambda}.
\]
也就是说, 若 $|\xi|<1$, 则 $\mu_\lambda(\xi)=(1-|\xi|^2)^{\lambda}$; 若 $|\xi|\geq 1$, 则 $\mu_\lambda(\xi)=0$. 我们还定义 $K_\lambda=\check{\mu}_\lambda$ 与 $T_\lambda f=K_\lambda\star f$.

三元组 $(K_{\lambda},T_{\lambda},\mu_{\lambda})$ 称为 $\lambda$ 阶 Bochner--Riesz 平均.
\end{definition}

\begin{remark}\label{rem:2.2}
若 $0\leq\lambda\leq 1$, 则 $\mu_{\lambda}$ 是 $\lambda$ 阶 H\"older 连续的, 即
\[
|\mu_{\lambda}(r)-\mu_{\lambda}(s)|\leq C|r-s|^{\lambda},
\]
而对任何 $\lambda_0<\lambda$, $\lambda_0$ 阶 H\"older 连续性都不成立. 若 $\lambda\geq 1$, 则 $\mu_\lambda$ 有 $[\lambda]$ 阶导数, 而且这些导数是 $\lambda-[\lambda]$ 阶 H\"older 连续的.

这里的要点是: 乘子 $\mu_\lambda$ 具有可变而受控的光滑性. 此外, 若 $\lambda=0$, 则 $\mu_0$ 是 $\mathbb{R}^n$ 中单位球的特征函数, 于是第~\ref{chap:1} 章的工作告诉我们, $\mu_0$ 的 $L^p$ 有界性等价于 Fourier 级数的 $L^p$ 收敛性.
\end{remark}

\begin{proposition}\label{prop:2.3}
卷积核 $K_{\lambda}(x)$ 由下式给出:
\[
K_{\lambda}(x)=c_{\lambda}\frac{J_{\frac{1}{2}+\lambda}(2\pi|x|)}{|x|^{\frac{1}{2}+\lambda}},
\qquad |K_{\lambda}(x)|\leq C|x|^{-(1+\lambda)}.
\]
\end{proposition}

\begin{remark}\label{rem:2.4}
\index{Bessel 函数 (Bessel functions)}
记号 $J_{\nu}$ 表示 $\nu$ 阶 Bessel 函数, 即微分方程
\[
z^{2}D^{2}J_{\nu}+zD J_{\nu}+(z^{2}-\nu^{2})J_{\nu}=0
\]
的解, 我们关心的是在 $z=0$ 处解析的那个解.

有了对 Bessel 函数的这一认定, 该命题容易证明; 我们只需计算
$z^{\nu}\int_{-1}^{+1}e^{izt}(1-t^2)^{\nu-\frac{1}{2}}\,dt$
满足该微分方程且解析. 我们可以通过改变积分围道并作渐近展开来得到 $|K_\lambda|$ 大小的估计. 事实上, 我们将在第~\ref{chap:6} 章中详细地做这件事. 我们现在不做这个计算, 因为本书这一部分的目的, 是看清 $\mu_\lambda$ 的光滑性如何影响 $K_\lambda$ 的大小. 上述命题确切地告诉我们: $K_\lambda$ 在无穷远处的衰减速度与乘子 $\mu_\lambda$ 的光滑度 $\lambda$ 成正比. 这些 Bessel 函数公式的问题在于它们是整体的: $K_\lambda$ 的显式公式是对整个 $\mu_\lambda$ 作平均得到的. 但 $\mu_\lambda$ 是一个除 $|\xi|=1$ 外处处光滑的函数; 既然我们想理解 $|\xi|=1$ 处奇性的影响, 把 $\mu_\lambda$ 在整个 $[-1,+1]$ 上平均在概念上就是错误的. 我们其实更愿意在 $\mu_\lambda$ 于 $|\xi|=1$ 处的奇性附近局部化.
\end{remark}

\essayhead{ess:2.5}{Fourier 变换的一个直观计算}

我们想重新思考计算 Fourier 变换的整个问题; 我们要考察的问题是: 函数的光滑性如何影响其 Fourier 变换的衰减速度. 我们打算先理解一些简单的例子, 然后把它推广为一般的理解.

关于``光滑性--衰减''问题, 我们有两个基本模型. 第一个是在分析 Hausdorff--Young 定理时做过的一个计算. 我们在那里算得 $\widehat{f\bigl(\frac{x}{\delta}\bigr)}(\xi)=\delta\hat{f}(\delta\xi)$. 为什么这是一个关于光滑性的结果? 因为当 $\delta$ 变小时, $f\bigl(\frac{x}{\delta}\bigr)$ 的支集变小, 而 $f$ 的大小不变. 收缩 $\delta$ 迫使 $f$ 的大小在很短的范围内发生巨大变化: 导数增大了. 它对 Fourier 变换的影响是乘以因子 $\delta$ (小), 并把 $\hat{f}$ 伸缩一个因子 $\delta$. 若把 $\hat{f}$ 想成紧支集的, 那么 $\hat{f}(\delta\xi)$ 的支集随 $\delta$ 变小而扩大. 这意味着 Fourier 变换不如 $\delta$ 大时衰减得那么快.

这一切都是定性的; 定量的分析包含在不等式
\[
|x^{2}\check{\psi}(x)|\leq\|x^{2}\check{\psi}\|_{\infty}\leq C\|D_{\xi}^{2}\psi\|_{1}.
\]
之中. 在这里应用它有一个技术上的困难: 对 $0\leq\lambda<1$, 函数 $\mu_\lambda$ 根本没有导数. 我们只好自力更生. 我们将为处理这类奇性建立一个标准模型. 我们要一点一点地逼近 $\mu_\lambda$ 的奇性; 把 $\mu_\lambda$ 写成一些标准函数之和, 其中每一个的奇性都是我们能控制的. 设 $\mu_\lambda=\sum\mu_\lambda\phi_j$, 其中每个 $\phi_j$ 光滑且支集远离 $|\xi|=1$. 这一安排使得 $\mu_\lambda\phi_j$ 光滑, 从而可以用标准的光滑性估计求出其 Fourier 变换的大小. 当然, 我们必须把所有单个的 Fourier 变换加起来, 而当 $j\to\infty$ 时, 我们预计 $\mu_\lambda\phi_j$ 越来越不光滑, 于是它的 Fourier 变换越来越大: 潜在地, 我们有了一个求不出和的级数. 为了摆脱困境, 我们需要控制 $\mu_\lambda\phi_j$ 在其支集上的变化. 若支集在 $[a_j,a_{j+1}]$ 中, 我们需要看到 $\mu_\lambda$ 的大小在区间的长度上没有显著变化. 于是, 我们需要常数 $c_1,c_2$ 使得
\[
c_{1}\leq\frac{\mu_{\lambda}(a_{j+1})}{\mu_{\lambda}(a_{j})}\leq c_{2}.
\]
对接近奇点 $1$ 的 $a_j$, 这提示
\[
c_{1}\leq\left|\frac{(1-a_{j+1})}{(1-a_{j})}\right|\leq c_{2},
\]
而 $a_{j}=1-2^{-j}$ 就很合适.

我们偷点懒, 假装奇点在 $0$ 而不在 $1$, 于是让所有的 $\phi_j$ 都是同一个函数的伸缩 (利用``伸缩使导数变大''这一直观). 设 $\phi_j(\xi)=\phi(2^j\xi)$, 并假装 $\phi_j$ 的支集包含在 $I_j\cup-I_j$ 中, 这里 $I_j=(1-2^{-j},\,1-2^{-(j+1)})$. 那么 $\mu_\lambda\phi_j$ 的大小约为 $2^{-\lambda j}$, 于是近似地,
\[
\widetilde{\mu_{\lambda}}=\sum_{j=0}^{\infty}2^{-\lambda j}\check{\phi}_{j}\,\widetilde{\mu}_{\lambda}(x)=\sum_{j=0}^{\infty}2^{-\lambda j}\check{\phi}_{j}(x)=\sum_{j=0}^{\infty}2^{-\lambda j}2^{-j}\check{\phi}(2^{-j}x).
\]
近似地, 于是,
\[
|K_{\lambda}(x)|\leq\sum_{j=0}^{\infty}2^{-(1+\lambda)j}|\check{\phi}(2^{-j}x)|.
\]
若选取 $\phi$ 非常光滑, 我们就能容易地控制其大小; 对 $\phi\in\mathcal{S}$, 当 $|y|\leq 1$ 时 $|\check{\phi}(y)|\leq C$, 当 $|y|\geq 1$ 时 $|\check{\phi}(y)|\leq C|y|^{-2}$. 利用这个精确的估计,
\begin{align*}
|K_{\lambda}(x)|&\leq C_{1}\sum_{|2^{-j}x|\leq1}2^{-(1+\lambda)j}+C_{2}\sum_{|2^{-j}x|\geq1}2^{-(1+\lambda)j}|2^{-j}x|^{-2} \\
&\leq C_{1}\sum_{[\log|x|]-1}^{\infty}2^{-(1+\lambda)j}+C_{2}\sum_{0}^{[\log|x|]+1}2^{-(1-\lambda)j}|x|^{-2} \\
&\leq C_{1}|x|^{-(1+\lambda)}+C_{2}|x|^{-(\lambda+1)}|x|^{-2}.
\end{align*}
这就证明了 $0\leq\lambda<1$ 时的估计; 若 $\lambda\geq 1$, 则需要对 $\check{\phi}$ 利用其在无穷远处更快衰减的估计.

\essayhead{ess:2.6}{Fourier 变换的一个精确计算}

在~\ref*{ess:2.5} 中计算 Fourier 变换时, 我们第一次偷懒是在集合 $I_j$ 上用估计 $\mu_{\lambda}\phi_j=2^{-\lambda j}$. 第二次偷懒是假定诸 $\phi_j$ 都是同一个 $\phi$ 的伸缩; 我们需要用诸如
\[
|x^{2}\check{\psi}(x)|\leq\|x^{2}\check{\psi}\|_{\infty}\leq C\|D_{\xi}^{2}\psi\|_{1}.
\]
这样更精确的估计来取代伸缩直观.

我们先来精确地构造 $\phi_j$. 设 $\psi(x)=\exp(-x^{-2})$; 则 $\psi$ 是 $C^\infty$ 的且支在 $[0,\infty)$ 中. 于是函数 $\psi(x)\psi(1-x)$ 光滑且支在 $[0,1]$ 中. 取这个函数的适当伸缩与平移, 记为 $\phi$, 使其支在 $[-\frac{5}{4},-\frac{1}{4}]$ 中. 令 $\phi_j(\xi)=\phi\bigl(2^j(|\xi|-1)\bigr)$. 注意, 若把 $I_j$ 定义为区间 $[1-5\cdot 2^{-(j+2)},\,1-2^{-(j+2)}]$, 则 $\phi_j$ 的支集包含在 $J_j=I_j\cup-I_j$ 中. 现在令
\[
\chi_{j}(\xi)=\frac{\phi_{j}(\xi)}{\sum_{k=0}^{\infty}\phi_{k}(\xi)}
\]
当 $\xi\in J_j$; 而当 $\xi$ 在 $J_j$ 之外时 $\chi_j(\xi)=0$. 注意, 若 $\xi\in J_j$, 则 $\bigl|D_\xi^\alpha\chi_j(\xi)\bigr|\leq C_\alpha 2^{\alpha j}$.

精确化的一个义务是必须核查许多琐碎的细节; 例如对 $\xi\in J_j$ 有 $0<\sum\phi_k<\infty$. 这两个不等式中, 第一个是平凡的: $\phi_k\geq 0$ 且 $\sum\phi_k\geq\phi_j$. 而在 $J_j$ 的内部 $\phi_j>0$, 而 $J_j$ 的端点位于 $J_{j-1}$ 或 $J_{j+1}$ 的内部, 故和为正. 此外, $J_j$ 只与 $J_{j-1}$ 或 $J_{j+1}$ 相交, 于是对固定的 $\xi$, 这个和至多有三个非零项, 因而自然是有限的.

最后, 令 $\psi_j=\mu_\lambda\phi_j$. 由构造, $\psi_j\in\mathcal{S}$, $\operatorname{supp}\psi_j\subset J_j$, 于是 $\sum\psi_j=\mu_\lambda$. 现在估计 $\check{\psi}_j$: 利用 $\mu_\lambda$ 在 $J_j$ 上光滑这一事实,
\[
|D_{\xi}^{\alpha}\mu_{\lambda}(\xi)|\;\chi_{J_{j}}(\xi)\leq C_{\alpha}2^{-(\lambda-\alpha)j}.
\]
于是
\begin{align*}
|\check{\psi}_{j}(x)|&\leq\|\check{\psi}_{j}\|_{\infty}\leq\|\phi_{j}\|_{1}
=\int_{J_{j}}|\mu_{\lambda}|\,\chi_{J_{j}}\leq C\,2^{\lambda j}|J_{j}|, \\
&=C\,2^{-(1+\lambda)j};\qquad
|x^{2}\check{\psi}_{j}(x)|\leq\|x^{2}\check{\psi}_{j}\|_{\infty}\leq\|D_{\xi}^{2}\psi_{j}\|_{1} \\
&\leq C\|D_{\xi}^{2}\mu_{\lambda}\,\chi_{J_{j}}\|_{1}
+C\|D_{\xi}\mu_{\lambda}\,D_{\xi}\chi_{J_{j}}\|_{1}
+\|\mu_{\lambda}D_{\xi}^{2}\chi_{J_{j}}\|_{1} \\
&\leq C|J_{j}|\left(2^{-(\lambda-2)j}+2^{-(\lambda-1)j}2^{j}+2^{-\lambda j}2^{2j}\right).
\end{align*}
总之, $|\check{\psi}_{j}(x)|\leq 2^{-\lambda j}|2^{-j}x|^{-2}$.

好, 现在的证明与~\ref*{ess:2.5} 中完全一样. 再一次, 若 $\lambda\geq 1$, 我们想要的估计就是出不来: $\sum_{j\geq 0}2^{(1-\lambda)j}|x|^{-2}$ 给出 $|x|^{-2}$ 量级的项, 当 $1-\lambda\leq 0$ 时这根本不够. 为了对付这一点, 我们可以估计 $|x|^3|\check{\psi}_j|$ 与 $D_{\xi}^{3}(\mu_{\lambda}\chi_{J_j})$.

\essayhead{ess:2.7}{Fourier 变换与标准模型}

``analysis''一词可以拆成它的希腊语词根, 意为``向上''与``分解''. 在西方思想史上, 分析就是把一个对象拆成更小的部分并试图理解每一部分的过程. 这里有两个希望: 其一, 尽管整体可能无法研究, 但各个部分是可以研究的; 其二, 人们希望通过对部分的研究能够以某种方式重建整体. 后一个希望常常是个问题; 注意到, 我们并没有恢复出关于 $\check{\mu}_{\lambda}$ 的完全信息, 而只得到了 $\check{\mu}_{\lambda}$ 在无穷远处大小的信息.

我们在本节想做的是让整个分析过程在直观上更有吸引力. 我们看到, 孤立 $|\xi|=1$ 处的奇性有很多好处; 我们让 $\mu_\lambda$ 在小区间上表现得如同常数. 于是, 为控制 $\mu_\lambda=\sum\mu_\lambda\chi_{I_j}=\sum 2^{-\lambda j}\chi_{I_j}$ 的 Fourier 变换, 我们只需控制 $|\check{\chi}_{I_j}|$. 现在我们只关心 $|\check{\chi}_{I_j}|$ 的大小; 这意味着我们可以把 $I_j$ 平移到以原点为中心, 而只使 Fourier 变换改变一个绝对值为 $1$ 的相位因子. 于是所有的 $\chi_{I_j}$ 都是区间 $I=[-1,1]$ 的特征函数的真正伸缩. 而伸缩对 Fourier 变换的影响是已知的: $|\check{\chi}_{I_j}(x)|\leq 2^{-j}|\check{\chi}_{I}(2^{-j}x)|$. 起初我们需要理解一整族复杂函数的 Fourier 变换; 现在我们要做的只是理解单个区间的特征函数的 Fourier 变换.

直观上,
\[
\check{\chi}_{I}(x)=\int_{I}e^{2\pi ix\xi}\,d\xi
\]
是一个振荡指数函数在区间上的平均. 若频率很大, 振荡就快, 整个区间上的平均值接近零. 能够对 $\check{\chi}_I$ 作出贡献的频率必须与区间的大小相匹配: 频率必须等于区间的长度. 我们必须有 $|2^{-j}x|\leq 1$. 对这样的 $x$,
\[
|\check{\mu}_{\lambda}(x)|=2^{-(1+\lambda)j}\leq|x|^{-(1+\lambda)}.
\]
坏消息是这个计算是错的; 若波的频率不是区间长度的整数倍, Fourier 变换在区间上就不会平均为零. 误差将是指数函数一个波长的量级, 即 $|2^{-j}x|^{-1}$. 若 $x$ 很大, 这就给 $|\check{\mu}_\lambda(x)|$ 带来大小为 $\sum_{|2^{-j}x|>1}2^{-(1+\lambda)j}|2^{-j}x|^{-1}$ 的误差, 对大的 $x$ 它是 $|x|^{-1}$ 量级的. 这个估计糟透了.

我们可以用光滑函数 $\psi$ (支在 $I_j$ 中) 替换 $\chi_I$ 来得到更好的估计, 然后再伸缩与平移. 推理与之前相同, 只是若频率不能恰好装进区间 $I_j$, 计算的误差将与指数函数的一个波长乘以函数 $\psi$ 在区间端点处的高度成比例. 由于 $\psi$ 在区间端点附近无穷阶消失, 这一贡献会非常小. 若 $\psi$ 是线性的, 高度就是 $0+(\text{斜率})(\xi \text{ 的改变量})=(D_\xi\psi)(|\text{波长}|)$. 总误差贡献将是 $|\text{波长}|^2=|2^{-j}x|^2$; 它给总和贡献一项大小为 $|x|^{-(1+\lambda)}$ 的量, 这正是我们想要的. 若 $\lambda\geq 1$, 这一切都不成立, 我们需要利用 $\psi$ 无穷阶消失的事实, 于是 $\psi$ 在端点附近的高度如同 $|D_\xi^k\psi|\,|\text{波长}|^k$, 等等. 我们稍后会做一个精确的版本; 直观是清楚的.

这里有一个基本的要点, 在量子力学中称为不确定性原理: 我们不能简单地忽略那些很快但不完全匹配区间的频率. 这些频率贡献大小与频率成比例的误差. 若函数 $\phi$ 的支集被限制在一个小集合中, 则 $\hat{\phi}$ 不是紧支集的; 它是实解析的, 所以肯定不是紧支集的. 不可能把 $\phi$ 与 $\hat{\phi}$ 同时局部化在一个小集合中. 既然控制 Fourier 乘子需要同时控制 $\mu_{\lambda}$ 与 $K_{\lambda}=\check{\mu}_{\lambda}$, 我们可以预期我们的结果来之不易.

\begin{corollary}\label{cor:2.8}
若 $\lambda>0$, 则 $T_\lambda$ 在 $L^p(\mathbb{R}^1)$ ($1\leq p\leq\infty$) 上有界.
\end{corollary}

\begin{proof}
若 $\lambda>0$, 则 $K_\lambda\in L^1$, 因此
\[
{}_{p}\|T_{\lambda}\|\leq\|K_{\lambda}\|_{1}<\infty. \qedhere
\]
\end{proof}

\begin{lemma}\label{lem:2.9}
若 $\lambda=0$, 则 $T_0$ 在 $L^1(\mathbb{R}^1)$ 上无界.
\end{lemma}

\begin{proof}
若 $\lambda=0$, 我们对 $T_\lambda$ 的估计并不能表明 $K_\lambda\in L^1$. 但话说回来, 我们的估计本来就很糟糕. 另一方面, 若对所有的 $f\in\mathcal{S}$ 都有 $T_0 f\in L^1$, 则 $\widehat{T_0 f}(\xi)$ 对所有 $\xi$ 连续. 这意味着 $\chi_{[-1,1]}\hat{f}(\xi)$ 对所有使 $\hat{f}\in\mathcal{S}$ 的 $\hat{f}$ 连续. 这是荒谬的.
\end{proof}

\begin{remark}\label{rem:2.10}
\begin{enumerate}[label=\alph*)]
\item 引理~\ref{lem:2.9} 的证明几乎是平凡的; 它比紧环面 $\mathbb{T}^{n}$ 上的类似结果容易得多. 但~\ref*{lem:2.9} 与 $\mathbb{T}^{n}$ 上的那个结果是等价的.
\item 我们的标准模型能否帮助解释 $\lambda>0$ 与 $\lambda=0$ 之间的差别? 直观上, $\mu_{\lambda}=\sum 2^{-\lambda j}\phi(2^j\xi)$. 由于 $\phi\in\mathcal{S}$, 所有的 $\phi(2^j\xi)$ 都给出 $L^p$ 乘子, 其范数为 $\|\check{\phi}\|_1$. 于是我们用一个非常平凡的估计: ${}_p\|T_{\lambda}\|\leq\sum 2^{-\lambda j}\|\check{\phi}\|_1=C\sum 2^{-\lambda j}$. $\lambda>0$ 与 $\lambda=0$ 之间的差别一目了然. 但这没有解释 $T_0$ 将在 $p>1$ 时于 $L^p$ 上有界这一事实; 这表明, 对 $\lambda=0$, 我们需要粗估计之外的手段来把 $\mu_\lambda$ 从它的碎片重新拼合起来. 不过我们是能够做到的.
\end{enumerate}
\end{remark}

\section{Hilbert 变换}
\label{sec:2.2}

\begin{definition}\label{def:2.11}
\index{Hilbert 变换 (Hilbert transform)|emph}
Hilbert 变换 $H$ 是如下给出的 Fourier 乘子算子, 定义在 $L^2(\mathbb{R}^1)$ 上:
\[
\widehat{Hf}(\xi)=i\,\mathrm{sign}(\xi)\,\hat{f}(\xi).
\]
\end{definition}

\begin{lemma}\label{lem:2.12}
$T_{0}$ 在 $L^{p}$ 上有界当且仅当 $H$ 在 $L^{p}$ 上有界.
\end{lemma}

\begin{proof}
两个算子都以一个尖锐的跳跃间断为特征; 由 $H$ 制造 $T_0$ 并不困难. 注意:
\begin{gather*}
\chi_{(-\infty,0)}(\xi)=\frac{1+i\,(i\,\mathrm{sign}\,\xi)}{2};\quad
\chi_{(0,\infty)}(\xi)=\frac{1-i\,(i\,\mathrm{sign}\,\xi)}{2}; \\
\chi_{(-1,1)}=\chi_{(-\infty,1)}(\xi)\chi_{(-1,\infty)}(\xi)
=\chi_{(-\infty,0)}(\xi-1)\chi_{(0,\infty)}(\xi+1).
\end{gather*}
于是定理~\ref{thm:1.23} 表明, Hilbert 变换的这些改动不会产生显著影响: ${}_p\|T_0\|\leq\frac{1}{4}\bigl(1+{}_p\|H\|\bigr)^2$.

反之, 若 $T_0$ 在 $L^p$ 上有界, 则把 $\chi_{(-1,1)}$ 平移就得到 $\chi_{(0,2)}$, 而乘子范数不变. 再伸缩得 $\chi_{(0,R)}$, 然后对 Schwartz 函数取极限得 $\chi_{(0,\infty)}$. 对负的一侧作类似的推理, 然后把它们相加得到 $i\,\mathrm{sign}\,\xi$ 并无困难.
\end{proof}

\begin{remark}\label{rem:2.13}
把 $T_0$ 换成 $H$ 看起来是件特别愚蠢的事. 毕竟, 区间的特征函数属于 $L^1$, 于是与它相伴的卷积核属于 $L^1$. 但 $H$ 的乘子太糟糕了, 它的卷积核不可能属于 $L^1$ (否则其 Fourier 变换就会在无穷远处衰减). 我刚刚把一个局部化的、卷积核可积的乘子, 换成了一个整体的乘子, 其核甚至不是局部可积的.
\end{remark}

\begin{proposition}\label{prop:2.14}
若 $f\in\mathcal{S}$, 则
\[
Hf(x)=\lim_{\epsilon\to0}\frac{1}{\pi}\int_{\epsilon\le|y|\le\frac{1}{\epsilon}}\frac{1}{y}f(x-y)\,dy.
\]
\end{proposition}

\begin{proof}
我们在分布意义下的计算中会稍微偷点懒:
\begin{align*}
\lim_{\epsilon}\frac{1}{\pi}\int_{\epsilon\leq|y|\leq\frac{1}{\epsilon}}\frac{1}{y}e^{-2\pi iy\xi}\,dy
&=\lim\frac{2i}{\pi}\int_{\epsilon}^{\frac{1}{\epsilon}}\frac{1}{y}\sin 2\pi\xi y\,dy \\
&=\frac{2i}{\pi}\,\mathrm{sign}\,\xi\,\lim_{\epsilon}\int_{\epsilon}^{\frac{1}{\epsilon}}\frac{1}{y}\sin y\,dy
=\left(\frac{2i}{\pi}\,\mathrm{sign}\,\xi\right)\left(\frac{\pi}{2}\right).
\end{align*}
没有哪个真正的分析学家能抗拒不用围道积分来求积分
\[
\lim_{\epsilon}\int_{\epsilon}^{\frac{1}{\epsilon}}\frac{1}{y}\sin y\,dy
\]
的诱惑. 下面的技巧是 K.~Merryfield 告诉我们的; 这些积分都经过精心安排使之绝对收敛, 因此我们所做的形式运算都是合法的.
\[
\int_{0}^{\infty}\frac{\sin y}{y}\,dy
=\left[\frac{(1-\cos y)}{y}\right]_{0}^{\infty}+\int_{0}^{\infty}\frac{1-\cos x}{x^{2}}\,dx.
\]
但 $x^{-2}=\int_{0}^{\infty}s\,e^{-sx}\,ds$, 于是上式右端的积分为
\begin{align*}
\int_{0}^{\infty}\!\!\int_{0}^{\infty}&(1-\cos x)\,s\,e^{-sx}\,dx\,ds
=\int_{0}^{\infty}\left(\frac{1}{s}-\frac{s}{s^{2}+1}\right)s\,ds \\
&=\int_{0}^{\infty}\frac{1}{s^{2}+1}\,ds=\arctan\infty-\arctan 0=\frac{\pi}{2}. \qedhere
\end{align*}
\end{proof}

\begin{remark}\label{rem:2.15}
\begin{enumerate}[label=\alph*)]
\item 非正式地, 我们将写 $Hf=\frac{1}{x}\star f$. 当然, 由于 $x=0$ 处的奇性, 这个卷积只对非常特殊的 $f$ 收敛. 这正是我们把 $T_0$ 换成 $H$ 时所预期的. 我们有所收获吗? 与 $K_0$ 相应的卷积核是 $\frac{\sin 2\pi x}{x}$, 它在零附近是收敛的. 但 $K_0$ 与 $\frac{1}{x}$ 都不在 $L^1$ 中, 于是这些算子的 $L^p$ 有界性不依赖于大小, 而是依赖于核的正部与负部之间的相消 (对照注~\ref{rem:1.17}: 无论卷积核的绝对值还是其正部与负部, 都不给出在任何 $p$ 于 $L^p$ 上有界的乘子). 而一旦不带绝对值, 我们至少得到一个在 $L^2$ 上有界的算子.

结论是: 相消 (cancellation) 对理解 $T_0$ 或 $H$ 都是本质性的. 但两者的卷积核分别是 $\frac{\sin 2\pi x}{x}$ 与 $\frac{1}{x}$. Hilbert 变换的相消性质要容易控制得多.
\item 对 $(\mathrm{sign}\,\xi)$ 的逆变换作一个直观计算是容易的. 由于 $\mathrm{sign}\,\xi=\frac{\xi}{|\xi|}$ 在伸缩 $\xi\mapsto\delta\xi$ ($\delta>0$) 下不变, 其卷积核 $K$ 必须满足 $\delta K(\delta x)=K(x)$, 于是 $K(x)=\frac{\omega(x/|x|)}{|x|}$, 其中 $\omega$ 是某个在伸缩下不变的函数. 但 $\frac{1}{x}$ 是奇函数, 而 Fourier 变换保持奇性, 所以 $\omega$ 必须是奇函数中在伸缩下不变的那一个的倍数, 即 $\frac{x}{|x|}$. 这意味着 $K$ 是 $\frac{c}{x}$ 的形式.
\item 我们知道 $H$ 在 $L^{1}$ 上无界, 而在 $L^{2}$ 上有界. 这两个事实我们都有证明, 而且都是用 Fourier 变换技术证明的. 但这还不够好, 因为 Fourier 变换平均得太多, 而 $H$ 在 $L^{1}$ 上的无界性源于乘子在 $\xi=0$ 处的一个纯局部的间断. 对于卷积核 $\frac{1}{x}$ 与 $L^p$ 有界性之间的关系, 我们没有局部层面的理解; 作为一个分布, 它对任何 $L^p$ 中的任何函数都同样糟糕. 下面的结果, 以及再下一节, 将构造出对``卷积核如何影响算子的 $L^p$ 有界性''的更好理解.
\end{enumerate}
\end{remark}

\begin{lemma}\label{lem:2.16}
设 $I=[0,1]$, $f=\chi_I$. 则 $Hf\in L^p$ 当且仅当 $1<p<\infty$.
\end{lemma}

\begin{proof}
为证 $Hf\notin L^{1}$, 我们必须设法避免相消; 在积分 $\int_{-1}^{1}\frac{1}{x-y}\,dy$ 中我们专看 $x-y>0$ 的部分. 于是若 $y\in(-1,1)$, $x\in(n,n+1)$, 则 $\frac{1}{x-y}\geq\frac{1}{n}$, 因此
\[
Hf(x)\geq\sum\frac{2}{n}\chi_{(n,n+1)}(x),\qquad \|Hf\|_{1}=\infty.
\]
为证正面的结果, 我们把靠近奇性的点与远离奇性的点区分开: 对 $x$ 分成两个区间 $(0,2)$ 与 $(2,\infty)$. 在后者上, 我们在 $(n,n+1)$ 上用估计 $\frac{1}{x-y}\leq\frac{1}{n-1}$; 于是对 $p>1$, $\int_{2}^{\infty}\leq 2^{p}\sum\frac{1}{n^{p}}<\infty$.

在区间 $(0,2)$ 上, 我们利用相消得到
\[
Hf(x)=\log|x+1|-\log|x-1|.
\]
第一项属于 $L^\infty(0,2)$; 第二项可以在 $(0,1)$ 与 $(1,2)$ 上对称地分析; 我们将一点一点地逼近奇性以度量其影响. 在 $(0,1)$ 上,
$\int_0^1|Hf|^p=\sum_k\int_{1-2^{-k}}^{1-2^{-(k+1)}}|Hf|^p\leq\sum k^{p}2^{-k}<\infty$.
\end{proof}

\begin{remark}\label{rem:2.17}
计算的目的是洞察; 我们学到了什么? 首先, 我们得到一个并不十分意外的结果: 远离 $f$ 的支集处, $Hf$ 的行为如同 $\frac{1}{x}$. 这部分地控制了 $Hf$ 的 $L^p$ 行为. 毕竟在无穷远处, $\frac{1}{x}\in L^p$ 当且仅当 $p>1$. 但还有另一个区域, 即 $x$ 靠近 $f$ 的支集的区域, 在那里 $\frac{1}{x}$ 的大小变得无关紧要, 而相消变得重要. 这个区域需要用到 $\int_{\epsilon\leq|y|\leq\frac{1}{\epsilon}}\frac{1}{x}\,dy=0$ 这一事实.
\end{remark}

\section{Cotlar 定理}
\label{sec:2.3}
\index{殆正交性 (almost orthogonality)}

在本节中, 我们将在不使用 Fourier 变换的情况下证明 $H$ 在 $L^2(\mathbb{R}^1)$ 上的有界性. 我们尤其想看到核 $\frac{1}{x}$ 的相消是如何被用来控制 $H$ 的.

我们将求助于标准模型; 我们应该一点点逼近 $i\,\mathrm{sign}(\xi)$ 在 $\xi=0$ 处的奇性, 把它写成
\[
i\,\mathrm{sign}(\xi)=\sum 2^{-0j}\phi(2^{j}\xi).
\]
我们想在卷积核一侧给出它的解释:
\[
K(x)=\sum 2^{-j}\check{\phi}(2^{-j}x).
\]
直观上, $\phi$ 的大部分支集在 $[1,2]$ 中; 这是不对的, 但它使我们猜到: 逼近卷积核奇性的正确方式是写
\[
\frac{1}{x}=\sum_{j=-\infty}^{\infty}K_{j}(x);
\]
这里
\[
K_{j}(x)=\frac{1}{x}\chi_{(2^{j},2^{j+1})}(|x|).
\]
正如所料, $\sum{}_2\|K_j\|=\infty$, 于是 $H$ 的 $L^2$ 有界性取决于诸 $K_j$ 之间如何相互作用.

使我们能从 $K_j$ 重建 $H$ 的关键直观来自 Fourier 变换分析: 在 Fourier 变换一侧, $K_j$ 的作用如同乘以 $\phi_j$. 这些 $\phi_j$ 非常像区间 $[2^{-(j+1)},2^{-j}]$ 的特征函数. 它们是投影算子; 而且是到 $L^2$ 的相互正交的子空间上的投影 (我甚至可以把它们想成沿对角线为 $1$ 的对角算子). 度量这类算子大小的正确方法, 不是把所有单个部分的大小加起来, 而是取各部分范数中的最大者, 在这里它是有限的.

这一讨论是``$H$ 在 $L^2$ 上有界''的 Fourier 变换证明的算子理论版本. 我们要做的是把这个关于投影的想法翻译到卷积核一侧. 从算子理论的角度, ``是投影''这一条件即是: 与核 $K_j$ 相应的算子 $T_j$ 满足正交性条件 $T_kT_j=0$. 现在我们要做两件事: i) 核实正交性对我们的算子确实成立 (它们并不是被定义为投影的); ii) 核实正交性足以给出 $L^2$ 有界性.

为了核查正交性, 我们不允许使用估计
\[
{}_{2}\|T_{k}T_{j}\|=\|\hat{K}_{k}\hat{K}_{j}\|_{\infty},
\]
因为我们不使用 Fourier 变换. 我们仅有的把卷积核与算子范数联系起来的估计是
\[
{}_{2}\|T_{k}T_{j}\|=\|K_{k}\star K_{j}\|_{1},
\]
所以我们需要理解 $K_j\star K_k$ 的大小. 在 $k=1$、$j$ 很大的特殊情形, 我们看图~\ref{fig:2}. 卷积 $K_j\star K_k$ 应当被视为一个平均过程; 在区域 $(2^j,2^{j+1})$ 中, $K_j$ 几乎不变化, 而在此区域之外它是零. 另一方面, $K_1$ 在长度为 $2$ 的区间中在正负之间跳跃, 于是我们预计它对缓慢变化的函数 $K_j$ 的平均值为零.

这是我们的希望; 不幸的是, 这个直观是错的. 有两个问题: $K_j$ 缓慢变化但非常数, 于是 $K_1$ 的正部与负部``看到''的 $K_j$ 的区域不同, 不会平均为零; $K_j$ 的光滑度控制着平均值到底有多大. 其次, 当 $x\in(2^j-1,2^j)$ 时, $K_1$ 的负部对 $0$ 作平均, 而正部对 $K_j$ 作平均. 这给出一个在测度为 $1$ 的集合上大小为 $2^{-j}$ 的量:
\[
\|K_{j}\star K_{1}\|_{1}>2^{-j}.
\]
卷积确实很小, 但不是零, 于是我们关于投影的整套哲学体系土崩瓦解.

还有一条生路: 拟正交投影 (almost orthogonal projections) 理论. 这一理论利用了 $T_jT_k$ 非常小这一事实, 并巧妙地反复使用这个估计, 使各项之间的相互作用非常小. 该理论使用谱半径定理: 算子 $T$ 的范数由
\[
\|T\|=\lim_{n\to\infty}\|T^{n}\|^{\frac{1}{n}}
\]
给出. 取 $T$ 的 $n$ 次幂使我能迭代任何已有的估计; 另一方面, 我必须记录多个 $T_{i_j}$ 与 $T_{i_k}$ 相互平均的组合学.

\begin{theorem}\label{thm:2.18}
\index{Cotlar 定理 (Cotlar's lemma)}
设 $\{T_{k}\}$ 是 Hilbert 空间上的一族有界算子, 并设
\begin{align*}
&\left\|T_{k}\right\|\leq C, \\
&\|T_{i}^{\star}T_{j}\|\leq\phi^{2}(i-j), \\
&\|T_{i}T_{j}^{\star}\|\leq\phi^{2}(i-j).
\end{align*}
则
\[
\Bigl\|\sum_{j=-N}^{N}T_{j}\Bigr\|\leq\sum_{j=-\infty}^{\infty}\phi(j).
\]
\end{theorem}

\begin{proof}
设 $T$ 表示 $\sum_{j=-N}^{N}T_{j}$; 则 $T^{\star}T$ 是自伴的, 谱半径定理适用于它. 我们真正想知道的是 $T$ 的范数; 其关系为:
\[
\|T f\|^{2}=(T f,T f)=(f,T^{\star}T f)\leq\|f\|^{2}\|T^{\star}T\|;
\]
因此
\[
\left\|T\right\|\leq\left\|T^{\star}T\right\|^{\frac{1}{2}}.
\]
于是我们可以满足于估计这个算子的范数. 做一下组合,
\[
(T^{\star}T)^{n}=\sum T_{i_{1}}^{\star}T_{i_{2}}T_{i_{3}}^{\star}\ldots T_{i_{2n}}.
\]
我们预计能控制成对的迭代, 尽管对复杂的和的关系也许不行:
\begin{align*}
\|T_{i_{1}}^{\star}\ldots T_{i_{2n}}\|&\leq\phi^{2}(i_{1}-i_{2})\ldots\phi^{2}(i_{2n-1}-i_{2n}), \\
\|T_{i_{1}}^{\star}\ldots T_{i_{2n}}\|&\leq C^{2}\|T_{i_{2}}\ldots T_{i_{2n-1}}^{\star}\| \\
&\leq\phi^{2}(i_{2}-i_{3})\ldots\phi^{2}(i_{2n-2}-i_{2n-1}).
\end{align*}
把这些估计相乘再开方, 得到
\[
\left\|T_{i_{1}}^{\star}\ldots T_{i_{2n}}\right\|\leq C\,\phi(i_{1}-i_{2})\ldots\phi(i_{2n-1}-i_{2n}).
\]
最后,
\begin{align*}
\|(T^{\star}T)^{n}\|&\leq C\sum\phi(i_{1}-i_{2})\phi(i_{2}-i_{3})\cdots\phi(i_{2n-1}-i_{2n}) \\
&\leq\left[\sum_{j=-\infty}^{\infty}\phi(j)\right]\left[\sum\phi(i_{1}-i_{2})\phi(i_{2}-i_{3})\cdots\phi(i_{2n-2}-i_{2n-1})\right] \\
&\leq\left[\sum_{j=-\infty}^{\infty}\phi(j)\right]^{2n-2}\left[\sum\phi(i_{1}-i_{2})\right] \\
&\leq\left[\sum_{j=-\infty}^{\infty}\phi(j)\right]^{2n-1}(2N).
\end{align*}
总之,
\[
\|(T^{\star}T)^{n}\|^{\frac{1}{n}}\leq\lim_{n\to\infty}[2CN]^{\frac{1}{n}}\left[\sum\phi(j)\right]^{2-\frac{1}{n}},
\]
结论随之得证.
\end{proof}

\begin{theorem}\label{thm:2.19}
设 $I_j=(2^j,2^{j+1})$, $K_j=\frac{1}{x}\chi_{I_j}(|x|)$, $T_j f(x)=K_j\star f(x)$. 则
\[
{}_{2}\|T_{j}\|\leq 3,
\qquad
{}_{2}\|T_{j}^{\star}T_{i}\|={}_{2}\|T_{j}T_{i}^{\star}\|\leq 2^{-|i-j|}.
\]
因此, Hilbert 变换在 $L^{2}$ 上有界.
\end{theorem}

\begin{proof}
首先, ${}_{2}\|T_{j}\|\leq\|K_{j}\|_{1}=\log 4\leq 3$.

此外, 回忆 $T_i^{\star}=J_0J_1T_iJ_0J_1$, 其中诸 $J$ 是共轭与反射等距算子, 我们得到估计:
\[
{}_{2}\|T_{j}^{\star}T_{i}\|={}_{2}\|T_{j}T_{i}^{\star}\|={}_{2}\|T_{j}T_{i}\|\leq\|K_{i}\star K_{j}\|_{1}.
\]
而且, 若 $i>j$, 一个小小的变量替换表明
\[
K_{i}\star K_{j}(x)=2^{-j}\bigl(K_{i-j}\star K_{0}\bigr)\bigl(2^{-j}x\bigr),
\]
而这个伸缩不改变 $K_i\star K_j$ 的 $L^1$ 范数. 所以只需证明当 $m$ 很大时的估计
\[
\|K_{m}\star K_{0}\|_{1}\leq C2^{-m}.
\]

我们对卷积的直观理解是: 它表示一个函数对另一个函数的平均; 当函数的不同支集相互重叠时, 计算就变得棘手. 我们先把这种重叠理清楚: 令
\[
S_{m}=\{y \mid 1<|y|<2,\; 2^{m}<|x-y|<2^{m+1}\}.
\]
则 $S_{m}$ 是
\[
K_{m}\star K_{0}(x)=\int K_{m}(x-y)K_{0}(y)\,dy
\]
中被积函数非零的 $y$ 的集合. 我们的直观告诉我们这个积分很小, 因为 $K_{0}$ 的平均值为零而 $K_{m}$ 缓慢变化. 当我们注意到
\[
\int_{S_{m}}\frac{1}{x}\,\frac{1}{y}\,dy=0
\]
时, 这两个事实都可以加以利用; 于是
\begin{align*}
K_{m}\star K_{0}(x)&=\int K_{m}(x-y)K_{0}(y)\,dy
=\int_{S_{m}}\frac{1}{x-y}\,\frac{1}{y}\,dy \\
&=\int_{S_{m}}\left[\frac{1}{x-y}-\frac{1}{x}\right]\frac{1}{y}\,dy.
\end{align*}
由``缓慢变化''的技巧, 这两个分式之差应当很小; 我们可以用 Taylor 定理来估计它. 但既然我们可以真的把减法做出来, 何乐而不为:
\[
\left[\frac{1}{x-y}-\frac{1}{x}\right]\frac{1}{y}=\frac{1}{x^{2}}\left[1-\frac{y}{x}\right]^{-1}.
\]
而若 $m$ 很大,
\[
\left|\frac{y}{x}\right|\leq\frac{2}{2^{m}}<\frac{1}{2},
\]
因此
\[
\left|K_{m}\star K_{0}(x)\right|\leq 2\,\frac{1}{x^{2}}\int_{S_{m}}dy\leq\frac{4}{x^{2}},
\]
因为 $|S_m|\leq 2$. 最后, 我们要计算 $K_m\star K_0$ 的 $L^1$ 范数. 其支集为
\[
\{x \mid \text{存在 } y \text{ 使 } 1<|y|<2,\; 2^{m}<|x-y|<2^{m+1}\}
\subset\{x \mid 2^{m}-2<|x|<2^{m}+2\},
\]
于是
\[
\left\|K_{m}\star K_{0}\right\|_{1}<4\int\frac{1}{x^{2}}\,dx\leq C2^{-m}.
\]
现在对 $\phi(j)=2^{-|j|}$ 应用定理~\ref{thm:2.18}. 我们看到
\[
S_{N}f(x)=\sum_{j=-N}^{N}T_{j}f(x)
\]
在 $L^{2}$ 上一致有界. 于是
\[
\lim_{N\to\infty}S_{N}f(x)=\lim_{\epsilon\to0}\int_{\epsilon<|y|<\frac{1}{\epsilon}}\frac{1}{x-y}f(y)\,dy
\]
在 $L^{2}$ 中收敛. 在 Schwartz 函数上, 其极限为 $Hf$, 因此
\[
{}_{2}\|H\|\leq\sup_{N}\,{}_2\Bigl\|\sum_{j=-N}^{N}T_{j}\Bigr\|<\infty. \qedhere
\]
\end{proof}

\begin{remark}\label{rem:2.20}
看清这个证明为什么成功是重要的. 在早先的直观中, 我们得到过简单的估计
\[
Hf\sim\frac{1}{x};
\]
而这里, 不知怎的, 我们得到了
\[
H f\sim\frac{1}{x^{2}}.
\]
改进从何而来? 它来自两个来源: 我们计算了 $K_m\star K_0$, 由于 $m$ 那么大, 这就像计算 $H(K_0)$; 然后我们利用 $K_m$ 的光滑性与 $\int K_0=0$ 来做估计. 这个计算的一个显式版本如下. 设 $f$ 支在一个区间 $I$ 上, 其中心为 $y_0$; 设 $\int f=0$. 设 $\tilde{I}$ 是与 $I$ 同心但长度三倍的区间. 若 $x\notin\tilde{I}$ 且 $y\in I$, 则:
\begin{align*}
Hf(x)&=\int\frac{1}{x-y}f(y)\,dy
=\int\left[\frac{1}{x-y}-\frac{1}{x-y_{0}}\right]f(y)\,dy \\
&=\int_{I}\frac{y-y_{0}}{(x-y)(x-y_{0})}\,f(y)\,dy.
\end{align*}
但是
\[
\left|\frac{1}{x-y}\right|\leq\frac{C}{|x|};\qquad
\frac{|y-y_{0}|}{|x-y_{0}|}\leq\frac{C|I|}{|x|},
\]
于是
\[
|H f(x)|\leq\frac{C|I|}{|x|}\int_{I}|f|.
\]
\end{remark}

\essayhead{ess:2.21}{攻克 Hilbert 变换的 \texorpdfstring{$L^{p}$}{Lp} 有界性}

我们知道 $H$ 在 $L^1$ 上无界, 在 $L^2$ 上有界. 此外, 我们的 $L^1$ 反例给出 $Hf\sim\frac{1}{x}$, 这在某种程度上暗示 $H$ 是弱 $(1,1)$ 的. 于是, 要证明 $L^p$ 有界性, 在弱 $L^1$ 与 $L^2$ 之间内插就够了.

为了执行这个纲领, 我需要知道 $L^1$ 与 $L^2$ 有何不同. Marcinkiewicz 内插的证明提示了 $f\in L^1$ 的一个分裂. 令
\[
G=\Bigl\{x \,\Big|\, |f(x)|\leq\frac{\lambda}{2}\Bigr\};\quad
B=\Bigl\{x \,\Big|\, |f(x)|>\frac{\lambda}{2}\Bigr\};\quad g=\chi_{G},\ \text{等等.}
\]
那么 $g$ 实在是很好:
\begin{align*}
\lambda^{2}\left|\left\{\left|H g\right|>\lambda\right\}\right|
&\leq\int_{\left\{\left|H g\right|>\lambda\right\}}\left|H g\right|^{2} \\
&\leq\left\|H g\right\|_{2}^{2}\leq\left\|g\right\|_{2}^{2}
=\int_{G}\left|f\right|^{2}\leq\lambda\int_{G}\left|f\right| \\
&\leq\lambda\|f\|_{1};
\end{align*}
这给出了 $H$ 在 $f$ 的好部分上的弱 $L^{1}$ 行为. 为了控制 $H$ 在坏部分上的作用, 我们把前面的两个观察合在一起: 若 $f$ 在一个区间上为常数, 我们能控制 $Hf$; 若 $f$ 在一个区间上平均值为零, 我们能控制 $Hf$. 当然 $f$ 并不支在一个区间上, 但我们可以强迫它如此. 注意, 若 $f$ 属于 $\mathcal{S}$, 则 $B$ 是开集, 从而可以写成互不相交的区间 $I$ 的并. 现在令 $f_I$ 表示 $f$ 在 $I$ 上的平均:
\[
f_{I}=\frac{1}{\left|I\right|}\int_{I}f.
\]
则
\[
f\chi_{I}=\bigl[f_{I}+(f-f_{I})\bigr]\chi_{I}=f_{I}\chi_{I}+b_{I};
\]
现在 $f_{I}\chi_{I}$ 在每个 $I$ 上为常数, 而 $b_{I}$ 在 $I$ 上的平均为零. 总而言之,
\[
b=f\chi_{B}=\sum b_{I}+\sum f_{I}\chi_{I}.
\]
我们有了分解, 也能控制 Hilbert 变换在各部分上的作用; 还有什么会阻止我们完成证明吗? 有的:

\begin{enumerate}[label=\arabic*)]
\item 在 $\tilde{I}$ 之外, 我们能控制 $Hb_{I}$; 但在 $\tilde{I}$ 上无法控制.
\item Minkowski 不等式不成立; 我们不能简单地写
\[
|\{\textstyle\sum|f_{I}|>\lambda\}|\leq C\sum|\{|f_{I}|>\lambda\}|;
\]
常数 $C$ 将依赖于和式中项数的多少.
\end{enumerate}

第一个困难意味着我们必须放弃控制 $Hb$ 在 $B^{\star}=\bigcup\tilde{I}$ 上的一切希望. 这倒不是什么大问题:
\begin{align*}
\left|\left\{x\in B^{\star}\,\middle|\,|H b(x)|>\lambda\right\}\right|
&\leq|B^{\star}|\leq\sum|\tilde{I}|=3\sum|I|=3|B|=3\Bigl|\Bigl\{|f|>\frac{\lambda}{2}\Bigr\}\Bigr| \\
&\leq\frac{6}{\lambda}\int_{\{|f|>\frac{\lambda}{2}\}}|f|
\leq\frac{6}{\lambda}\|f\|_{1}.
\end{align*}
这给出了 $b$ 在 $B^{\star}$ 上的弱 $L^1$ 估计. 第二个困难更难对付; 我们将被迫在某个使三角不等式成立的空间中做估计; 否则我们无从把 $Hf$ 从它的碎片重建起来. 现在, $b$ 是 $f$ 的 $L^1$ 部分, 我们预计 $Hb$ 会表现得像 $\frac{1}{x^{2}}$, 于是 $L^1$ 应当是控制 $Hb$ 的合适空间:
\[
\lambda\left|\left\{x\in(B^{\star})^{c}\,\middle|\,\left|H\Bigl(\sum b_{I}\Bigr)\right|>\lambda\right\}\right|
\leq\int_{(B^{\star})^{c}}\left|H\Bigl(\sum b_{I}\Bigr)\right|
\leq\sum\int_{(B^{\star})^{c}}|H(b_{I})|.
\]
我们无法指望控制 $H$ 在诸 $\tilde{I}$ 相交处的作用, 而且除了在 $(\tilde{I})^{c}$ 上, 我们对 $H(b_I)$ 也没有特别好的控制. 这意味着我们必须扩大积分区域; $(B^{\star})^{c}\subset(\tilde{I})^{c}$, 且:
\[
\sum\int_{(B^{\star})^{c}}|H(b_{I})|\leq\sum\int_{(\tilde{I})^{c}}|H(b_{I})|.
\]
现在逐项分析. 设 $y_{I}$ 为 $I$ 的中心, 则
\[
\left(H b_{I}\right)(x)=\int_{I}\left[\frac{1}{x-y}-\frac{1}{x-y_{I}}\right]b_{I}(y)\,dy,
\]
\begin{align*}
\int_{(\tilde{I})^{c}}|H(b_{I})(x)|\,dx
&\leq\int_{(\tilde{I})^{c}}\int_{I}\left|\frac{1}{x-y}-\frac{1}{x-y_{I}}\right||b_{I}(y)|\,dy\,dx \\
&=\int_{I}|b_{I}(y)|\int_{(\tilde{I})^{c}}\left|\frac{1}{x-y}-\frac{1}{x-y_{I}}\right|dx\,dy.
\end{align*}
内层积分可以显式算出; 它被
\[
4\log\left|1-\frac{|y-y_{I}|}{|I|}\right|
\]
控制. 由于 $y\in I$, $x\in(\tilde{I})^{c}$, 有 $\frac{|y-y_{I}|}{|I|}\leq\frac{1}{2}$, 于是我们得到估计
\[
\int_{(\tilde{I})^{c}}|H(b_{I})(x)|\,dx\leq 8\int_{I}|b_{I}(y)|\,dy.
\]
从而
\begin{align*}
\sum\int_{(B^{\star})^{c}}|H(b_{I})|
&\leq 8\sum\int_{I}|b_{I}(y)|\,dy \\
&\leq 8\sum\int_{I}\bigl(|f(y)|+|f_{I}(y)|\bigr)\,dy
=8\sum\int_{I}|f(y)|\,dy+8\sum|I|\,|f_{I}| \\
&=16\sum\int_{I}|f(y)|\,dy
=16\int_{B}|f(y)|\,dy\leq 16\|f\|_{1}.
\end{align*}
这一项就这样处理了. 现在剩下的只有常数项 $f_{I}\chi_{I}$. 处理常数项这种事想必不必费什么力气. 和之前一样, 我们需要一个使三角不等式成立的空间; 但也和之前一样, 我们的直观告诉我们 $H(f)(x)\sim\frac{1}{x}$, 所以 $L^{1}$ 并不是合适的工作空间. $L^{2}$ 是令人满意的; 于是,
\begin{align*}
&\left|\left\{\left|H\Bigl(\sum f_{I}\chi_{I}\Bigr)\right|>\lambda\right\}\right| \\
&\leq\frac{1}{\lambda^{2}}\int\left|H\Bigl(\sum f_{I}\chi_{I}\Bigr)\right|^{2} \\
&\leq C\frac{1}{\lambda^{2}}\int\left|\sum f_{I}\chi_{I}\right|^{2} \\
&\leq C\frac{1}{\lambda^{2}}\sum\int_{I}\left|f_{I}\chi_{I}\right|^{2};
\end{align*}
最后一步是因为诸 $I$ 互不相交. 要完成证明, 我们还需要去掉一个 $\lambda$ 的幂和一个 $|f_I|$ 的幂: 我们需要估计
\[
\frac{1}{\left|I\right|}\int_{I}\left|f\right|\leq C\lambda,
\]
即 $|f_{I}|\leq C\lambda$; 这将给我们最终估计
\begin{align*}
C\frac{1}{\lambda}\sum\int_{I}|f_{I}|
&=C\frac{1}{\lambda}\sum|I|\,|f_{I}| \\
&=C\frac{1}{\lambda}\sum\int_{I}|f|
=C\frac{1}{\lambda}\int_{B}|f|\leq C\frac{1}{\lambda}\|f\|_{1}.
\end{align*}
当然, 若和式中只有一项, 我们真正需要的是估计
\[
\frac{1}{\left|I\right|}\int_{I}\left|f\right|\leq C\lambda.
\]
唉, 正确的估计全都朝相反的方向: 由于 $I\subset B$,
\[
|f_{I}|\geq\frac{\lambda}{2}.
\]
这一切意味着, 常数项 $f_{I}\chi_{I}$ 才是控制 $H$ 的真正症结所在. 这本不该让我们吃惊: 正是它们在 $\infty$ 处表现得像 $\frac{1}{x}$; 正是它们贡献了非 $L^{1}$ 的行为.

我们能做什么? 已知的不等式
\[
|f_{I}|\geq\frac{\lambda}{2}
\]
与我们想要的不等式
\[
|f_{I}|\leq C\lambda
\]
提示我们必须小心地选取区间 $I$. 它们必须与函数 $f$ 紧密匹配; 直观上说, 选取它们时必须使 $f$ 在 $I$ 上变化不太大. 形式上, 我们需要的是以下四条性质: 我们必须把 $\mathbb{R}$ 分解成互不相交的集合
\begin{enumerate}[label=\arabic*)]
\item $\mathbb{R}=G\cup B$;
\item $B=\bigcup I$, 诸 $I$ 互不相交且满足 $f_I\leq C\lambda$;
\item $|B|\leq C\frac{1}{\lambda}\|f\|_{1}$;
\item 若 $x\in G$, 则 $|f(x)|\leq\lambda$.
\end{enumerate}

若能对任何 $L^1$ 函数做到这一点, 我们就可以用我们的锦囊妙计证明 $H$ 在 $L^1$ 上弱有界; 然后内插并用对偶性得到它在 $1$ 与 $\infty$ 之间的所有 $L^p$ 上有界. 我们所获得的, 是 $H$ 的完全消除. 现在的问题成了一个关于 $L^1$ 函数的问题: 它们能否被切成在小区间上变化不太大的碎块. 问题是: 平均而言, 一个 $L^1$ 函数到底有多少``光滑性''?

\section{Hardy--Littlewood 极大函数}
\label{sec:2.4}

如何控制一个函数在区间上的变化? 让 $f_I>\lambda$ 是容易的部分; 对任意的区间和任意的 $\lambda$, 别指望还能同时得到 $f_I<C\lambda$. 我们先这样选区间使得 $f_I<C\lambda$; 所有满足这一点的区间都被扔进 $B$; 然后我们再操心其余三条性质.

取定一个 $\lambda$ 和一个 $f$, 不妨设 $f$ 非负. 第一个小细节是: 到底有没有好区间. 由于 $\int f<\infty$, 我们总能找到 $J$, $N$ 使得当 $|J|\geq N$ 时
\[
\frac{1}{\left|J\right|}\int_{J}f<\lambda.
\]
另一方面, 我们不能把所有这样的区间都拿来用, 因为我们打算用 $\frac{1}{\lambda}\|f\|_{1}$ 控制 $|B|$. 有一个简单但重要的绕过办法: 取一个满足
\[
\frac{1}{\left|J\right|}\int_{J}f<\lambda
\]
的区间 $J$. 现在设 $I$ 表示 $J$ 的左半或右半, 并假设
\[
\lambda<\frac{1}{|I|}\int_{I}f.
\]
这样的区间留作 $B$ 的一部分是好的, 因为它具有我们所需的平均值的下界性质. 当然困难的部分是上界性质. 诀窍在这里:
\[
\frac{1}{\left|I\right|}\int_{I}f\leq\frac{1}{\left|I\right|}\int_{J}f
=\frac{2}{\left|J\right|}\int_{J}f<2\lambda.
\]
于是我们没有丢掉平均值在上面的界. 现在继续这个剖分过程. 区间的左右两半总是互不相交的, 而若我们找到一个特别好的半区间 $I$ 满足
\[
\lambda<\frac{1}{|I|}\int_{I}f,
\]
就把它留给 $B$. 这一过程产生一族互不相交的区间, 其平均值下界为 $\lambda$、上界为 $2\lambda$. 由于诸区间互不相交, 我们甚至得到
\[
|B|=\sum|I|\leq\sum\frac{1}{\lambda}\int_{I}f\leq\frac{1}{\lambda}\|f\|_{1}.
\]
那么, 这一次什么地方出错了? 首先, 我们甚至还没有尝试控制 $B^{c}$. 其次, 凭什么说, 若
\[
\frac{1}{|J|}\int_{J}f<\lambda,
\]
我就总能找到一个子区间 $I$ 使
\[
\lambda\leq\frac{1}{|J|}\int_{I}f\,?
\]
当然, 没有什么能阻止我把 $J$ 的两半再对分, 一再对分下去. 这个过程凭什么会停止? 嗯, 我们确实还有一些选择余地. 若这个过程永不停止, 我们显然不想要 $J$ 进 $B$; 另一方面, $G$ 中的点必须满足: 若 $x\in G$, 则 $|f(x)|\leq\lambda$. 而关于 $x\in G$ 我们所知道的, 只是存在一个包含 $x$ 的无穷区间序列 $I_j$, 满足 $|I_{j+1}|=\frac{1}{2}|I_j|$, 且始终没有 $\lambda<f_{I_j}$, 也就是说始终有 $f_{I_j}\leq\lambda$. 那么凭什么下结论说 $f\leq\lambda$?

这是一个基本问题: 我们能由 $f$ 的平均值的信息恢复关于 $f$ 的信息吗? 若 $f$ 连续, 答案是简单的:
\[
\lim_{h\to0}\frac{1}{h}\int_{x}^{x+h}f(y)\,dy=f(x),
\]
因为这正是对积分的求导. 若令 $I_{j}=[x,x+2^{-j}]$, 我们得到
\[
\lim_{j\to\infty}f_{I_{j}}=\lim_{j\to\infty}\frac{1}{|I_{j}|}\int_{I_{j}}f=f(x),
\]
而 $f_{I_j}\leq\lambda$ 蕴含 $f(x)\leq\lambda$.

然而, 使用它仍有技术上的困难. 第一, $I_j$ 并不是区间 $[x,x+2^{-j}]$; 除了 $x\in I_j$ 之外我们对 $I_j$ 毫无控制. 所以我们不能直接用对积分求导的论证.

第二, 众所周知, 积分的导数依赖于函数. Darboux 定理告诉我, 函数的导数满足介值性质. 这意味着若 $f$ 跳过某个值, 就不存在可微的 $F$ 使 $F'=f$. 我们并非总能由平均值恢复 $f$.

要得到正确的积分微分学, 我们必须忽略零测集上发生的事, 只要求几乎处处收敛:
\[
\lim_{\substack{|I|\to0\\ x\in I}}f_{I}=f(x)
\]
是否对几乎所有的 $x$ 成立? 若 $f$ 连续可微, 这是一个容易的 $\epsilon$--$\delta$ 论证. 问题一如既往是从稠密集过渡到整个 $L^{1}$. 这是一个标准问题, 也有一个标准解法: 取一列在 $L^{1}$ 中收敛于 $f$ 的光滑函数 $f_n$. 我想说的是
\[
f(x)=\lim_{n}f_{n}(x)=\lim_{n}\lim_{I}(f_{n})_{I}=\lim_{I}\lim_{n}(f_{n})_{I}=\lim_{I}f_{I}.
\]
所以我真正想做的是交换两个极限, 但没有某种一致收敛我永远做不到. 我们在给 Fourier 级数求和时遇到过同样的问题; 当时需要对算子范数的一致估计. 这里也一样, 只不过这是一个逐点的算子, 所以我们需要逐点的一致性; 我们需要对
\[
\sup_{\substack{I\\ x\in I}}\frac{1}{|I|}\int_{I}f
\]
的估计.

\begin{definition}\label{def:2.22}
\index{Hardy--Littlewood 极大函数 (Hardy--Littlewood maximal function)|emph}
设 $f$ 局部可积. $f$ 的 Hardy--Littlewood 极大函数记作 $f^{\ast}$, 定义为
\[
f^{\ast}(x)=\sup_{\substack{I\\ x\in I}}\frac{1}{|I|}\int_{I}|f(y)|\,dy.
\]
\end{definition}

\begin{lemma}\label{lem:2.23}
设 $f\in L^1$ 且 $f$ 不几乎处处为零. 则 $f^{\star}\notin L^1$.
\end{lemma}

\begin{proof}
选取 $N$ 使 $\int_{-N}^{N}|f|\geq\frac{1}{2}\|f\|_1$. 现在任取 $|x|\geq N$ 的 $x$. 令 $I=[-|x|,|x|]$;
\begin{align*}
f^{\star}(x)&\geq\frac{1}{|I|}\int_{I}|f|
=\frac{1}{2|x|}\int_{-|x|}^{|x|}|f|\geq\frac{1}{2|x|}\int_{-N}^{N}|f| \\
&\geq\frac{1}{4|x|}\,\|f\|_{1}.
\end{align*}
于是对大的 $x$,
\[
f^{\star}(x)\geq\frac{C}{x},\qquad C=\|f\|_{1}\geq 0. \qedhere
\]
\end{proof}

\begin{theorem}\label{thm:2.24}
映射 $f\mapsto f^{\star}$ 是弱 $(1,1)$ 的.
\end{theorem}

\begin{proof}
我们可以假设 $f$ 非负. 令
\[
E=\{x \mid f^{\star}(x)>\lambda\}.
\]
关于如何估计 $E$ 的测度, 我们有一套实际的哲学, 即: 对每个 $x\in E$, 存在 $I$ 使 $x\in I$ 且 $f_I>\lambda$. 于是对任何 $x\in I$, 有 $f^{\star}>\lambda$, 从而 $x\in E$. 这意味着 $E$ 被一大堆区间覆盖. 若这些区间互不相交, 我们就得到
\begin{align*}
|E|&=|\bigcup I|=\sum|I| \\
&\leq\sum\frac{1}{\lambda}\int_{I}f=\frac{1}{\lambda}\int_{E}f\leq\frac{1}{\lambda}\|f\|_{1}.
\end{align*}
糟糕的是, 按我们定义它们的方式, 这些区间 $I$ 不可能互不相交.

我们的策略是丢掉大部分区间, 但保留足够多, 使我们不至扔掉该集合的大部分质量. 做到这一点的能力称为覆盖引理 (covering lemma)\index{覆盖引理 (covering lemma)}, 它与 $L^{1}$ 函数无关, 只与区间的几何及其相交方式有关.

在决定保留哪些区间、丢弃哪些区间时, 我们实际上没有多少选择. 若保留区间 $I_1$, 就必须丢掉所有与它相交的区间. 我们如何确定被丢掉的区间没有占 $E$ 的大部分测度? 有一个简单的窍门: 若有两个互相重叠的区间, 我们总保留较大的那个, 记为 $I_\star$. 假设有 $20$ 个区间原本与 $I_\star$ 相交, 而我们把它们全丢了. 我们是否丢掉了 $E$ 的大部分测度? 没有! 若 $x\in I\cap I_\star$, 则 $I\subset(x-|I|,\,x+|I|)\subset(x-|I_\star|,\,x+|I_\star|)$, 因为 $I_\star$ 是诸区间中最大的. 由于 $x\in I_\star$, 整个 $(x-|I_\star|,\,x+|I_\star|)$ 包含在一个与 $I_\star$ 同心、长度为其 $5$ 倍的区间中. 在丢掉所有与 $I_\star$ 相交的区间时, 我们至多把合起来的测度缩小一个 $5$ 的因子, 这是一个固定的因子, 与区间的个数、$I_\star$ 的大小或任何东西都无关. 我们的窍门将会奏效. 当然, 还有技术细节要处理.

我们首先假设 $E$ 有有限测度, 然后对每个 $x\in E$ 取定一个 $I$ 使 $x\in I$ 且 $f_I>\lambda$. 这个取定的集合族记为 $\mathcal{B}_\infty$. 我们在 $\mathcal{B}_\infty$ 中找不到最大元, 但可以接近. 设 $L_1=\sup\{|I| \mid I\in\mathcal{B}_\infty\}$. 若 $L_1=\infty$, 取 $I$ 使 $|I|>|E|$, 就完事了. 否则, 取 $I_1$ 使
\[
|I_{1}|>\frac{1}{2}L_{1};\qquad
\mathcal{B}_{2}=\{I\in\mathcal{B}_{1} \mid I\cap I_{1}=\varnothing\},\qquad
L_{2}=\sup\{|I| \mid I\in\mathcal{B}_{2}\},
\]
依此类推. 当然若 $\sum|I_j|=\infty$, 我们又完事了; 否则就得干点活.

我们一直没有试着证明 $\bigcup I\subset\bigcup I_j^{\star}$; 为此需要: 若 $x\in E$, 则 $x\in I$ 对某个 $I$ 成立, 而我们必须找到一个 $k$ 使 $x\in I_k^{\star}$. 这个策略对诸 $\mathcal{B}$ 意味着什么? 我们需要找到一个 $k$ 使得 $I\in\mathcal{B}_k$ 但 $I\notin\mathcal{B}_{k+1}$. 这并不太难, 因为诸 $\mathcal{B}$ 中区间的长度越来越小. 设 $I\in\mathcal{B}_j$ 对所有的 $j$ 成立; 则 $\frac{1}{2}|I|<\frac{1}{2}L_j<|I_j|$. 另一方面, $\sum|I_j|<\infty$, 于是要么 $|I|=0$, 要么该级数在某处终止. 但级数只有当 $L_j=0$ 时才终止, 而那时又有 $|I|=0$. 这表明存在 $k$ 使 $I\in\mathcal{B}_k$ 而 $I\notin\mathcal{B}_{k+1}$.

这对我们有利, 因为 $I\in\mathcal{B}_k$ 蕴含
\[
\frac{1}{2}|I|\leq\frac{1}{2}L_{k}\leq|I_{k}|,
\]
而若 $I\notin\mathcal{B}_{k+1}$, 则存在 $x\in I\cap I_k$. 正如直观的引入所言, $I\subset I_k^{\star}$. 于是, 对每个 $I\in\mathcal{B}_1$, 存在 $k$ 使 $I\subset I_k^{\star}$. 从而 $E\subset\bigcup I_k^{\star}$, 且
\[
\left|E\right|\leq\left|\bigcup I_{k}^{\star}\right|\leq\sum\left|I_{k}^{\star}\right|=5\sum\left|I_{j}\right|,
\]
故 $|\{f^{\star}>\lambda\}|\leq\frac{5}{\lambda}\|f\|_{1}$.
\end{proof}

\begin{corollary}\label{cor:2.25}
若 $1<p\leq\infty$, 则 $\|f^{\star}\|_p\leq C_p\|f\|_p$.
\end{corollary}

\begin{proof}
注意到
\[
|f_{I}|\leq\frac{1}{|I|}\int_{I}|f|\leq\|f\|_{\infty},
\]
故 $f^{\star}(x)\leq\|f\|_{\infty}$, 从而
\[
\|f^{\star}\|_{\infty}\leq\|f\|_{\infty}.
\]
于是可以用 Marcinkiewicz 内插定理, 即定理~\ref{thm:1.21}, 得到结论. 有一点小问题: 我们只对线性算子证明了 Marcinkiewicz 定理, 但给出的证明对次线性算子, 即满足 $T(f+g)\leq T(f)+T(g)$ 的算子, 同样适用.
\end{proof}

\begin{corollary}\label{cor:2.26}
若 $f\in L^{1}$, 则对几乎每个 $x$,
\[
\lim_{\substack{|I|\to0\\ x\in I}}\frac{1}{|I|}\int_{I}f(y)\,dy=f(x).
\]
\end{corollary}

\begin{proof}
若 $f$ 连续, 这是平凡的. 要从稠密集过渡到整个 $L^{1}$, 我们利用收敛的一致性. 证明会有点特别, 因为极大算子事实上在 $L^{1}$ 上并不有界. 令
\[
T f(x)=\left|\limsup_{I}f_{I}(x)-\liminf_{I}f_{I}(x)\right|.
\]
若 $g$ 连续, 则 $Tg=0$; 取任意的 $f\in L^1$ 与任意连续的 $g$, 注意到
\[
T f=T(f-g+g)\leq T(f-g)+Tg=T(f-g)\leq 2(f-g)^{\star}.
\]
于是
\[
|\{x \mid T f(x)>\lambda\}|\leq\Bigl|\Bigl\{x \,\Big|\, (f-g)^{\star}(x)>\frac{\lambda}{2}\Bigr\}\Bigr|\leq\frac{10}{\lambda}\|f-g\|_{1}.
\]
由于 $g$ 任意, 而左端与 $g$ 无关,
\[
|\{x \mid T f(x)>\lambda\}|=0.
\]
这意味着对几乎所有的 $x$, $Tf=0$.
\end{proof}

\begin{corollary}\label{cor:2.27}
若 $p<1$ 且 $I$ 是任意区间, 则
\[
\|f^{\star}\chi_{I}\|_{p}\leq C_{p}|I|^{\frac{1}{p}-1}\|f\|_{1}.
\]
\end{corollary}

\begin{proof}
\begin{align*}
\int_{I}|f^{\star}|^{p}&=p\int_{0}^{\infty}\lambda^{p-1}|I\cap\{f^{\star}\geq\lambda\}|\,d\lambda \\
&\leq p\int_{0}^{\infty}\lambda^{p-1}\min\Bigl\{|I|,\frac{c}{\lambda}\|f\|_{1}\Bigr\}\,d\lambda \\
&=p\int_{0}^{c|I|^{-1}\|f\|_{1}}|I|\,\lambda^{p-1}\,d\lambda
+c\,p\int_{c|I|^{-1}\|f\|_{1}}^{\infty}\lambda^{p-1}\lambda^{-1}\|f\|_{1}\,d\lambda \\
&=c|I|^{1-p}\|f\|_{1}^{p}+\frac{C_{p}}{1-p}|I|^{1-p}\|f\|_{1}^{p}. \qedhere
\end{align*}
\end{proof}

\begin{theorem}\label{thm:2.28}
\index{Calderon--Zygmund 分解 (Calder'on--Zygmund decomposition)}
设 $f\in L^1$, $f>0$, $\lambda\geq 0$ 取定. 则存在可测集 $G$, $B$ 与区间 $I$ 使得:
\begin{enumerate}[label=\arabic*)]
\item $\mathbb{R}=G\cup B$; $G\cap B=\varnothing$;
\item $B=\bigcup I$; 诸 $I$ 内部互不相交且满足 $\lambda<f_I\leq 2\lambda$;
\item $|B|\leq\frac{c}{\lambda}\|f\|_1$;
\item 对几乎所有的 $x\in G$, $|f(x)|\leq\lambda$.
\end{enumerate}
\end{theorem}

\begin{proof}
我们从本节开头的做法开始, 把 $\mathbb{R}$ 分成等长区间 $J$, 满足
\[
|J|\geq\frac{1}{\lambda}\|f\|_{1},
\]
于是自动地有
\[
f_{J}=\frac{1}{|J|}\int_{J}f\leq\frac{1}{|J|}\|f\|_{1}\leq\lambda.
\]
若 $\lambda<f_J$ 我们就保留 $J$; 否则把 $J$ 对分成 $J_0$, $J_1$. 若 $\lambda<f_{J_i}$, 就把 $J_i$ 放进保留的集合族. 若我们决定保留某个区间, 注意到
\[
f_{J_{i}}=\frac{2}{\left|J\right|}\int_{J_{i}}f
\leq\frac{2}{\left|J\right|}\int_{J}f\leq 2\lambda,
\]
于是对保留的区间自动有 $\lambda<f_{I_i}\leq 2\lambda$. 此外, 保留的区间内部互不相交, 因为它们是内部互不相交的区间的子集. 保留区间的并记为 $B$, 并令 $G$ 为其补集. 定理中唯一不能从构造直接得出的是 4). 若 $x\in G$, 则 $x$ 属于某个未被保留的初始区间 $J$, 也属于某个未被保留的剖分, 等等. 于是有一列区间 $I_k$, 使 $x\in I_k$ 且 $|I_{k+1}|=\frac{1}{2}|I_k|$. 由于没有哪个 $I_k$ 被留进 $B$, 故 $f_{I_k}\leq\lambda$; 于是由推论~\ref{cor:2.26}, 对几乎所有 $x\in G$, $f(x)\leq\lambda$.
\end{proof}

\begin{remark}\label{rem:2.29}
我们其实更愿意取 $B=\{x \mid f^{\star}(x)\geq\lambda\}$, 那样就可以避免上面的选取过程. 但我们怎能构造出区间 $I$? 条件 $f_I\leq 2\lambda$ 提示 $I$ 与 $g$ 相伴, 但事实上 $I\cap G=\varnothing$. 这是怎么回事? 设 $y\in\tilde{I}\cap G$. 则
\[
f_{I}=\frac{3}{|\tilde{I}|}\int_{I}f\leq 3f_{\tilde{I}}\leq 3f^{\star}(y)\leq 3\lambda.
\]
所以这种形式的界其实与``属于 $G$''大有关系. 要把它说清楚, 我们需要知道: 若 $x\in B$, 则存在区间 $I$ 使 $x\in I$ 且 $G\cap\tilde{I}\neq\varnothing$. 换一种说法: $x$ 到 $B^c$ 的距离与 $I$ 的长度是可比的.
\end{remark}

\begin{theorem}\label{thm:2.30}
\index{Whitney 分解 (Whitney decomposition)}
\index{二进区间 (dyadic intervals)}
设 $B$ 是 $\mathbb{R}$ 的真开子集. 则 $B$ 可以写成内部互不相交的区间 $I$ 的并, 满足
\[
|I|\leq d(B^{c},I)\leq 4|I|.
\]
\end{theorem}

\begin{proof}
最省事的窍门是令
\[
\mathcal{B}_{k}=\{I \mid 2^{k}<|I|<2^{k+1}\ \text{且}\ 2^{k}<d(I,B^{c})\leq 2^{k+1}\}.
\]
$\mathcal{B}_{k}$ 中的区间将覆盖 $B$, 且它们自动具有正确的距离性质. 当然它们并不互不相交, 所以还有活要干. 我们将像之前一样, 从一个标准的区间族出发不断取半. 令
\[
\mathcal{D}_{j}=\{I \mid I=[k2^{-j},(k+1)2^{-j}],\ k\in\mathbb{Z}\}.
\]
$\mathcal{D}_j$ 中的区间称为二进区间 (dyadic intervals); 它们自动具有我们需要的互不相交性质. 而且 $\mathcal{D}_{j+1}$ 中的区间是 $\mathcal{D}_j$ 中区间的一半. 这个互不相交性质可以换个说法: 若两个区间 $I$ 与 $J$ 的内部有非平凡的交, 则其中一个包含另一个.

现在令
\[
\mathcal{B}_{j}=\{J\in\mathcal{D}_{j} \mid \text{存在 } x\in J \text{ 使 } 2\cdot 2^{-j}<d(x,B^{c})\leq 2\cdot 2^{-(j-1)}\}.
\]
我们断言: 若 $J\in\mathcal{B}_{j}$, 则 $|J|\leq d(J,B^{c})\leq 4|J|$. 这应当是平凡的, 因为 $J$ 的长度约为 $2^{-j}$, 而 $J$ 与 $B^c$ 的距离和 $x$ 与 $B^c$ 的距离可比. 细节如下: 我们知道 $J$ 中有 $x$ 使
\[
2\cdot 2^{-j}<d(x,B^{c})\leq 2\cdot 2^{-(j-1)};
\]
因此
\[
2\left|J\right|<d(x,B^{c})\leq 4\left|J\right|,
\]
进而
\[
d(J,B^{c})=\inf\{|z-y| \mid z\in J,\,y\in B^{c}\}
\leq\inf\{|x-y| \mid y\in B^{c}\}=d(x,B^{c}).
\]
类似地,
\begin{align*}
d(J,B^{c})&=\inf\{|z-y| \mid z\in J,\,y\in B^{c}\} \\
&=\inf\{|z-x+x-y| \mid z\in J,\,y\in B^{c}\} \\
&\geq\inf\{|x-y| \mid y\in B^{c}\}-\sup\{|z-x| \mid z\in J\}
=d(x,B^{c})-|J|\geq 2|J|-|J|.
\end{align*}
为完成定理的证明, 我们需要证明 $B$ 是互不相交的区间 (大概是诸 $\mathcal{B}_j$ 中的 $I$) 的并. 若 $J\in\mathcal{B}_j$, 则 $d(J,B^c)>0$, 由于 $B^c$ 是闭的, 故 $J\subset B$. 另一方面, 若 $x\in B$, 则由 $B$ 是开集知 $d(x,B^c)>0$. 于是存在 $J\in\mathcal{D}_j$ 使 $x\in J$, 因为 $\bigcup\mathcal{D}_j=\mathbb{R}$. 为了得到互不相交性, 我们将使用最大尺寸的区间并利用二进性质. 此处的事情有点神秘, 因为许多区间相互包含, 我们容易把它们重复计数. 我们需要证明的是: 若 $x\in B$, 则 $x$ 包含在一个、且只有一个最大尺寸的区间中. 现在, 若区间 $J$ 包含 $x$, 则
\[
|J|\leq d(J,B^{c})\leq d(x,B^{c}),
\]
于是包含 $x$ 的区间的长度有界. 而且, 包含 $x$ 的二进区间按包含关系是全序的, 所以它们的并是包含 $x$ 的最大尺寸的二进区间. 极大性现在给出了互不相交性.
\end{proof}

\begin{corollary}\label{cor:2.31}
设 $f\in L^1$, $f\geq 0$ 与 $\lambda$ 取定. 则存在集合 $G$, $B$ 与内部互不相交的区间 $I$, 满足:
\begin{enumerate}[label=\arabic*)]
\item $\mathbb{R}=G\cup B$;
\item $B=\bigcup I$, $f_I\leq c\lambda$;
\item $|B|\leq\frac{c}{\lambda}\|f\|_1$;
\item 对几乎所有 $x\in G$, $f(x)\leq\lambda$.
\end{enumerate}
\end{corollary}

\begin{proof}
令
\[
G=\{x \mid f^{\star}(x)\leq\lambda\};\qquad B=\{x \mid f^{\star}(x)>\lambda\}.
\]
则第 1) 条显然成立, 而积分的微分学自动给出 4). 由于极大算子是弱 $(1,1)$ 的, 3) 也是自动的. 为得到诸区间, 注意 $B$ 是开集, 于是定理~\ref{thm:2.30} 给出具有良好距离性质的互不相交区间 $I$. 这些性质必须翻译成 2) 中那种平均型的条件.

设 $I$ 是我们的区间之一, 并设 $\tilde{I}$ 是与 $I$ 同心但长度十倍的区间. 由于 $d(I,B^c)\leq 4|I|$, 必有 $x$ 使 $x\in\tilde{I}\cap B^c$. 当然
\[
f_{\tilde{I}}\leq f^{\star}(x)\leq\lambda.
\]
但随后
\[
f_{I}\leq 10f_{\tilde{I}}\leq 10\lambda. \qedhere
\]
\end{proof}

现在我们可以证明 Hilbert 变换的有界性了. 但我们不如把本质抽象出来, 证明一个更一般的定理.

\begin{theorem}\label{thm:2.32}
\index{H\"ormander 条件 (H\"ormander condition)}
\index{奇异积分 (singular integrals)}
设 $(T,K,\mu)$ 是 $L^{2}$ 的乘子, 且 $K$ 满足
\[
\sup_{y}\int_{|x|>2|y|}|K(x-y)-K(x)|\,dx\leq C.
\]
则 $T$ 是弱 $(1,1)$ 的, 从而在所有满足 $1<p<\infty$ 的 $p$ 上有界.
\end{theorem}

\begin{proof}
我们真正要做的, 只是重复上一节末尾的计算. 取定 $\lambda>0$, 对 $f>0$ 与 $\frac{\lambda}{2}$ 应用定理~\ref{thm:2.28}. 令
\[
g=f\chi_{G},\qquad b=f\chi_{B},\qquad h_{I}=(b-b_{I})\chi_{I}.
\]
由于 $|Tf|\leq|Tg|+|Tb|$, 我们可以对每一部分分别做弱 $(1,1)$ 估计.
\begin{align*}
\Bigl|\Bigl\{|T g|>\frac{\lambda}{2}\Bigr\}\Bigr|
&\leq\frac{4}{\lambda^{2}}\|T g\|_{2}^{2}
\leq\frac{C}{\lambda^{2}}\|g\|_{2}^{2}=\frac{C}{\lambda^{2}}\int_{G}|g|^{2} \\
&=\frac{C}{\lambda^{2}}\int_{G}|f|\,|f|\leq\frac{C}{\lambda}\int_{G}|f|
\leq\frac{C}{\lambda}\|f\|_{1}.
\end{align*}
为处理坏的部分, 我们写 $b=\sum(b-b_I)\chi_I+b_I\chi_I$. 注意诸 $I$ 的互不相交性蕴含这个和几乎处处收敛. 则
\begin{align*}
\Bigl|\Bigl\{\,\Bigl|T\Bigl(\sum b_{I}\chi_{I}\Bigr)\Bigr|>\frac{\lambda}{4}\Bigr\}\Bigr|
&\leq\frac{16}{\lambda^{2}}\Bigl\|T\Bigl(\sum b_{I}\chi_{I}\Bigr)\Bigr\|_{2}^{2}
\leq\frac{C}{\lambda^{2}}\Bigl\|\sum b_{I}\chi_{I}\Bigr\|_{2}^{2} \\
&=\frac{C}{\lambda^{2}}\int\Bigl|\sum b_{I}\chi_{I}\Bigr|^{2}
=\frac{C}{\lambda^{2}}\sum\int|b_{I}\chi_{I}|^{2}
=\frac{C}{\lambda^{2}}\sum\int_{I}|b_{I}|^{2}.
\end{align*}
这个和能从平方中提出来, 是因为和中的区间互不相交; 也就是说, 逐点地,
\[
\Bigl|\sum b_{I}\chi_{I}\Bigr|^{2}=\sum|b_{I}\chi_{I}|^{2}.
\]
现在
\[
b_{I}=(f\chi_{B})_{I}=f_{I}\leq C\lambda,
\]
于是我们又有
\begin{align*}
\frac{C}{\lambda^{2}}\sum\int_{I}|b_{I}|^{2}
&\leq\frac{C}{\lambda}\sum\int_{I}|b_{I}|
=\frac{C}{\lambda}\sum|b_{I}|\,|I|=\frac{C}{\lambda}\sum\int_{I}|b| \\
&=\frac{C}{\lambda}\int_{B}|f|\leq\frac{C}{\lambda}\|f\|_{1}.
\end{align*}
如我们所知, 式子中麻烦的部分是 $L^1$ 的部分, 即诸 $h_I$. 设 $B^\star=\bigcup\tilde{I}$; 则
\[
\Bigl|\Bigl\{x\in B^{\star}\,\Big|\,\Bigl|T\Bigl(\sum h_{I}\Bigr)\Bigr|>\frac{\lambda}{4}\Bigr\}\Bigr|
\leq|B^{\star}|\leq 4|B|\leq\frac{C}{\lambda}\|f\|_{1}.
\]
另一方面,
\begin{align*}
\Bigl|\Bigl\{x\in(B^{\star})^{c}\,\Big|\,\Bigl|T\Bigl(\sum h_{I}\Bigr)\Bigr|>\frac{\lambda}{4}\Bigr\}\Bigr|
&\leq\frac{4}{\lambda}\int_{(B^{\star})^{c}}\Bigl|T\Bigl(\sum h_{I}\Bigr)\Bigr|
\leq\frac{4}{\lambda}\sum\int_{(\tilde{I})^{c}}|T(h_{I})| \\
&=\frac{4}{\lambda}\sum\int_{(\tilde{I})^{c}}\left|\int K(x-y)h_{I}(y)\,dy\right|dx.
\end{align*}
现在 $h_I=(b-b_I)\chi_I$, 于是 $\int_I h_I=0$. 若 $y_I$ 表示 $I$ 的中心, 那么我们可以随意插入一项 $-\int_I K(x-y_I)h_I\,dy$; 这给出
\begin{align*}
&\frac{4}{\lambda}\sum\int_{(\tilde{I})^{c}}\left|\int_{I}[K(x-y)-K(x-y_{I})]\,h_{I}(y)\,dy\right|dx \\
&\leq\frac{4}{\lambda}\sum\int_{I}\int_{(\tilde{I})^{c}}|K(x-y)-K(x-y_{I})|\,dx\,|h_{I}(y)|\,dy \\
&=\frac{4}{\lambda}\sum\int_{I}\int_{(\tilde{I}-y_{I})^{c}}|K(z+y-y_{I})-K(z)|\,dz\,|h_{I}(y)|\,dy.
\end{align*}
为分析这个积分, 注意到若令 $I=[-L+y_I,\,L+y_I]$, 并设 $y\in I$, $z\in(\tilde{I}-y_I)^c$, 则 $|z|>2L>2|y-y_I|$, 于是对此类 $y$,
\begin{align*}
&\int_{(\tilde{I}-y_{I})^{c}}|K(z+y-y_{I})-K(z)|\,dz\;|h_{I}(y)| \\
&\quad \leq\int_{|z|>2|y-y_{I}|}|K(z+y-y_{I})-K(z)|\,dz\;|h_{I}(y)| \\
&\leq C,
\end{align*}
最后一步用了 $K$ 满足的假设. 这意味着最后一项被
\[
\frac{C}{\lambda}\sum\int_{I}|h_{I}|
\leq\frac{C}{\lambda}\sum\int_{I}|b|+\frac{C}{\lambda}\sum\int_{I}|b_{I}|
\leq\frac{C}{\lambda}\|f\|_{1}
\]
控制. 至此证明完成, 因为所有四项都被 $\frac{C}{\lambda}\|f\|_1$ 控制.
\end{proof}

\begin{remark}\label{rem:2.33}
该定理的假设适用于核
\[
K_{\epsilon}(x)=\frac{1}{x}\chi_{\{\epsilon<|x|<\frac{1}{\epsilon}\}}(x),
\]
而且假设中的所有估计对 $\epsilon$ 一致. 由此, 算子 $T_{\epsilon}f=K_{\epsilon}\star f$ 在 $L^p$ 上一致有界, 因此
\[
\lim_{\epsilon\to0}K_{\epsilon}\star f
\]
在 $L^p$ 中收敛. 在 $L^2$ 的一个稠密集上, 它收敛到 Hilbert 变换, 于是 $T_\epsilon f$ 在 $L^p$ 中收敛到 $Hf$. 我们无法对逐点收敛下任何结论; 这还是那个交换极限的老问题, 需要某种一致的估计. 我们需要理解 $\sup_\epsilon|K_\epsilon\star f|$.
\end{remark}

\section{极大 Hilbert 变换}
\label{sec:2.5}

上一节中, 我们证明了 Hilbert 变换在 $L^{p}$ 上有界; 不幸的是, 论证是通过一个 $L^{p}$ 极限过程进行的. 缺陷在于: 对 $L^{p}$ 中的任意函数, 我们其实并不知道 Hilbert 变换是否逐点存在. 本节的目的是通过控制正确的极大函数来获得逐点信息.

\begin{definition}\label{def:2.34}
\index{极大 Hilbert 变换 (maximal Hilbert transform)|emph}
函数 $f$ 的极大 Hilbert 变换记作 $H^{\star}f$, 定义为
\[
H^{\star}f(x)=\sup_{\epsilon>0}\left|\int_{|y|>\epsilon}f(x-y)\frac{1}{y}\,dy\right|.
\]
\end{definition}

\begin{remark}\label{rem:2.35}
我想在一个更一般的背景下考察极大 Hilbert 变换与 Hardy--Littlewood 极大函数, 以找出一个统一的主题. 对 $f^{\ast}$, 设 $\psi(x)=\frac{1}{2}\chi_{I}(x)$, 这里 $I=[-1,1]$. 若 $\psi_{\epsilon}(x)=\frac{1}{\epsilon}\psi\bigl(\frac{x}{\epsilon}\bigr)$, 且 $I_{\epsilon}=[-\epsilon,\epsilon]$, 则
\begin{align*}
\psi_{\epsilon}\star f(x)&=\frac{1}{2\epsilon}\int_{x-\epsilon}^{x+\epsilon}f(y)\,dy
=\frac{1}{|I_{\epsilon}|}\int_{I_{\epsilon}}f(y)\,dy \\
&\leq\sup_{x\in I}\frac{1}{|I|}\int_{I}f(y)\,dy\leq f^{\star}(x).
\end{align*}
对 $H^{\star}f$ 有类似的公式: 设
\[
\psi(x)=\frac{1}{x}\chi_{(1,\infty)}(|x|),
\]
并注意到
\[
\psi_{\epsilon}\star f(x)=\int_{|y|>\epsilon}f(x-y)\frac{1}{y}\,dy.
\]
同样地,
\[
H^{\star}f=\sup_{\epsilon>0}|\psi_{\epsilon}\star f|.
\]
\end{remark}

\begin{definition}\label{def:2.36}
\index{恒等逼近 (approximate identity)|emph}
设 $\psi\in L^1$, $\psi\geq 0$, $\int\psi=1$; 令
\[
\psi_{\epsilon}(x)=\frac{1}{\epsilon}\psi\Bigl(\frac{x}{\epsilon}\Bigr).
\]
则 $\psi_{\epsilon}$ 称为一个\emph{恒等逼近} (approximate identity).
\end{definition}

\begin{remark}\label{rem:2.37}
``恒等逼近''这个名称来自我们即将证明的一个结果: 若 $\psi$ 足够好, 则当 $\epsilon\to 0$ 时 $\psi_{\epsilon}\star f$ 收敛到 $f$. 也就是说, 就卷积而言 $\psi_{\epsilon}$ 的作用近似于恒等算子.

我们想做的是度量 $\psi_{\epsilon}\star f$ 收敛到 $f$ 的速度. 利用 $\int\psi_{\epsilon}=1$,
\[
|\psi_{\epsilon}\star f(x)-f(x)|=\left|\int\psi_{\epsilon}(y)\bigl[f(x-y)-f(x)\bigr]dy\right|.
\]
现在假设 $f$ 有某种连续性; 比如
\[
|f(x-y)-f(x)|\leq C|y|^{\lambda}.
\]
那么上面的积分被
\[
C\int|\psi_{\epsilon}(y)|\,|y|^{\lambda}\,dy
=C\epsilon^{\lambda}\int|\psi(y)|\,|y|^{\lambda}\,dy=C\epsilon^{\lambda}
\]
控制. 于是 $\psi_{\epsilon}\star f$ 收敛到 $f$, 且收敛速度由 $f$ 的光滑性控制. 这个论证在比如 $f\in\mathcal{S}$ 时有效, 而要过渡到 $L^{p}$, 我们照例需要一致估计.
\end{remark}

\begin{proposition}\label{prop:2.38}
设 $\psi$ 是偶的、递减的且属于 $L^{1}$. 则存在与 $\psi$ 无关的常数 $C$, 使得
\[
\sup_{\epsilon>0}|\psi_{\epsilon}\star f(x)|\leq C\|\psi\|_{1}f^{\star}(x).
\]
\end{proposition}

\begin{proof}
我们将证明: 对所有满足假设的 $\psi$,
\[
|\psi\star f(0)|\leq C\|\psi\|_{1}f^{\star}(0).
\]
这使我们可以耍点小聪明. 首先, 把 $f$ 换成 $\tau_{x}f$, 并注意到
\begin{align*}
|\psi\star f(x)|&=|(\tau_{x}\psi\star f)(0)|
=|\psi\star(\tau_{x} f)(0)|\leq C\|\psi\|_{1}(\tau_{x} f)^{\star}(0) \\
&=C\|\psi\|_{1}f^{\star}(x).
\end{align*}
这些等式成立是因为卷积与平移交换, 并且估计
\[
|\psi\star f(0)|\leq C\|\psi\|_{1}f^{\star}(0)
\]
对 $\tau_x f$ 与对 $f$ 同样有效. 下一个花招: 我们注意到, 若 $\psi$ 满足假设, 则 $\psi_\epsilon$ 也满足, 而且这两个函数的 $L^1$ 范数相等. 因此
\[
|\psi_{\epsilon}\star f(x)|\leq C\|\psi\|_{1}f^{\star}(x).
\]
现在对两边取上确界即得命题的结论.

为估计
\[
|\psi\star f(0)|=\left|\int\psi(-x)f(x)\,dx\right|=\left|\int\psi(x)f(x)\,dx\right|,
\]
我们注意到 $\psi$ 递减且可积, 于是
\[
\lim_{R\to\infty}\psi(R)=0.
\]
局部可积性还蕴含
\[
\lim_{\epsilon\to0}\int_{0}^{\epsilon}f=0,
\]
于是我们可以分部积分并忽略边界项:
\[
\left|\int_{0}^{\infty}\psi(x)f(x)\,dx\right|
=\left|\int_{0}^{\infty}\int_{0}^{y}f(x)\,dx\,d\psi(y)\right|.
\]
由于 $\psi$ 递减, $d\psi\leq 0$, 于是最后的积分被
\begin{align*}
-\int_{0}^{\infty}\int_{0}^{y}|f(x)|\,dx\,d\psi(y)
&\leq-\int_{0}^{\infty}2y\,f^{\star}(0)\,d\psi(y) \\
&=-2f^{\star}(0)\int_{0}^{\infty}y\,d\psi(y)
=2f^{\star}(0)\int_{0}^{\infty}\psi(y)\,dy=f^{\star}(0)\|\psi\|_{1}
\end{align*}
控制.
\end{proof}

\begin{proposition}\label{prop:2.39}
\[
H^{\star}f(x)\leq[H f]^{\star}(x)+C f^{\star}(x).
\]
\end{proposition}

\begin{proof}
尽管我们知道
\[
H^{\star}f=\sup_{\epsilon>0}|\psi_{\epsilon}\star f(x)|,
\]
$\psi_{\epsilon}$ 却不可积. 这里的窍门是把它写成一个可积部分加上一个奇异部分, 然后把奇异部分的作用看作一个 Hilbert 变换; 可积部分贡献 $Cf^\star(x)$, 奇异部分给我们 $[Hf]^\star(x)$. 我们关于 Hilbert 变换的直观是: 可积部分来自核的光滑性, 大小约为 $\frac{1}{x^{2}}$; 为得到这个估计, 设 $\phi$ 是光滑、偶且递减的函数, 满足 $\phi\geq 0$, $\int\phi=1$, $\operatorname{supp}\phi\subset[-1,1]$. 令 $\varphi(x)=\psi(x)-\phi\star\frac{1}{x}$. 则:
\begin{align*}
|\varphi(x)|=\left|\psi(x)-\phi\star\frac{1}{x}\right|
&=\left|\psi(x)\int\phi(y)\,dy-\phi\star\frac{1}{x}\right| \\
&\leq\int_{|y|\leq1}\left|\frac{1}{x}-\frac{1}{x-y}\right|\phi(y)\,dy \\
&\leq\frac{1}{x^{2}}\int_{|y|\leq1}\left|\frac{y}{1-\frac{y}{x}}\right|\phi(y)\,dy.
\end{align*}
若 $|x|\geq 2$, 我们可以利用 $|y|\leq 1$ 来估计被积函数, 得到
\[
|\varphi(x)|\leq\frac{C}{x^{2}},
\]
正如我们所预料的. 于是
\[
\psi=\phi\star\frac{1}{x}+\varphi,
\qquad
\psi_{\epsilon}=\Bigl(\phi\star\frac{1}{x}\Bigr)_{\epsilon}+\varphi_{\epsilon}
=\phi_{\epsilon}\star\frac{1}{x}+\varphi_{\epsilon},
\]
且
\begin{align*}
H^{\star}f(x)&\leq\sup_{\epsilon}\left|\phi_{\epsilon}\star\frac{1}{x}\star f(x)\right|+\left|\varphi_{\epsilon}\star f(x)\right| \\
&\leq\sup_{\epsilon}\bigl[\phi_{\epsilon}\star H f(x)\bigr]+\sup_{\epsilon}\left|\varphi_{\epsilon}\star f(x)\right|.
\end{align*}
由于 $\phi$ 被一个可积递减函数控制, 命题~\ref{prop:2.38} 适用. 我们已经看到, 在 $|x|\geq 2$ 的范围内 $\varphi$ 被一个偶的可积递减函数控制; 若 $|x|\leq 2$, 则要看得更仔细. 在 $|x|\leq 1$ 的范围内, $\psi(x)=0$, 于是 $\varphi=-\phi\star\frac{1}{x}$. 由于 $\phi$ 光滑, $-\phi\star\frac{1}{x}$ 连续从而有界. 在 $1\leq|x|\leq 2$ 的范围内, $\phi$ 光滑, $\phi\star\frac{1}{x}$ 也光滑. 于是 $\varphi$ 被一个偶的递减可积函数控制, 命题~\ref{prop:2.38} 适用, 第二项也被 Hardy--Littlewood 极大函数控制.
\end{proof}

\begin{theorem}\label{thm:2.40}
\index{积分的微分 (differentiation of integrals)}
$H^{\star}$ 是弱 $(1,1)$ 的.
\end{theorem}

\begin{proof}
首先注意, 命题~\ref{prop:2.39} 蕴含 $H^{\star}$ 在 $1<p<\infty$ 时于 $L^p$ 上有界, 因为 Hilbert 变换与 Hardy--Littlewood 极大函数在那里都有界. 我们不能得出该算子在 $L^1$ 上有任何有界性的结论, 因为它是两个都不保持 $L^1$ 的算子的复合. 我们要做的是使用定理~\ref{thm:2.32} (Hilbert 变换有界性之证明) 中的技术. 我们可以沿用之前那些区间 $I$, 因为它们与 $\epsilon$ 无关. 技术困难出现在估计形如
\[
\int_{I^{c}}\int_{I}|\psi_{\epsilon}(x-y)-\psi_{\epsilon}(x)|\,dx\,dy
\]
的项时. 我们用 $\psi_{\epsilon}$ 截断破坏了 $\frac{1}{x}$ 的光滑性, 所以不能用惯常的光滑性论证. 另一方面, 若 $I$ 不靠近奇点, 我们就是安全的. 这意味着我们将不得不让区间与 $\epsilon$ 相适配.

如同定理~\ref{thm:2.32} 的证明, 我们只须努力得到 $H^{\star}\bigl(\sum h_I\bigr)$ 的估计, 因为所有其他估计都把问题送回 $L^2$, 而 Hilbert 变换与极大函数都是从 $L^2$ 到 $L^2$ 有界的. 对常数项, 取定一个 $\epsilon$ 并把所有的 $I$ 分成三类:

\begin{enumerate}[label=\alph*)]
\item 考察所有那些对一切 $y\in I$ 都有 $|x-y|<\epsilon$ 的 $I$. 则
\[
\int\psi_{\epsilon}(x-y)h_{I}(y)\,dy=\int_{|x-y|>\epsilon}\frac{1}{x-y}h_{I}(y)\chi_{I}(y)\,dy=0.
\]
这一类我们可以处理.
\item 那些对一切 $y\in I$ 都有 $|x-y|>\epsilon$ 的 $I$. 这些区间看不到我们引入的粗糙性, 所以可以像以前那样估计:
\[
\int_{|x-y|>\epsilon}K(x-y)h_{I}(y)\chi_{I}(y)\,dy
=\int_{I}\bigl[K(x-y)-K(x-y_{I})\bigr]h_{I}(y)\,dy,
\]
估计照旧进行.
\item 那些存在 $y_0\in I$ 使 $|x-y_0|=\epsilon$ 的 $I$. 此时卷积核中有一个从 $0$ 到 $\frac{1}{x-y}$ 的跳跃. 幸运的是, $|x-y_0|$ 与 $\epsilon$ 之间的关系使我们能用 $\epsilon$ 控制 $x-y$, 于是 $\int_I\frac{1}{x-y}h_I(y)\,dy$ 可以联系到 $h_I$ 在一个半径为 $\epsilon$ 的区间上的平均. 后者由 Hardy--Littlewood 极大函数控制.
\end{enumerate}

要真正完成这个证明, 我们必须先把条件 $x\in\tilde{I}^{c}$, $y\in I$ 与条件 $|x-y|<\epsilon$ 联系起来. 我们断言: 对这样的 $x,y$,
\[
\frac{\epsilon}{2}<|x-y|<\frac{3\epsilon}{2}.
\]
事实上, $|x-y|\geq|x-y_{0}|-|y-y_{0}|\geq\epsilon-|I|$. 但
\[
\epsilon=|x-y_{0}|\geq 2|I|;\qquad -|I|>-\frac{\epsilon}{2};\qquad |x-y|>\frac{\epsilon}{2}.
\]
另一方面, $y\in I$ 蕴含
\[
|x-y|\leq|x-y_{0}|+|y-y_{0}|=\epsilon+|y-y_{0}|\leq\epsilon+|I|<\frac{3\epsilon}{2}.
\]
于是:
\begin{align*}
&\left|\int_{|x-y|<\epsilon}\frac{1}{x-y}\,h_{I}(y)\chi_{I}(y)\,dy\right| \\
&\leq\int_{|x-y|<\epsilon}\frac{1}{|x-y|}\,|h_{I}(y)\chi_{I}(y)|\,dy \\
&\leq\frac{2}{\epsilon}\int_{|x-y|<\epsilon}|h_{I}(y)\chi_{I}(y)|\,dy \\
&\leq\frac{C}{\epsilon}\int_{|x-y|<\frac{3\epsilon}{2}}|h_{I}(y)|\chi_{I}(y)\,dy.
\end{align*}
这就足以完成证明: 对任何 $\epsilon>0$,
\begin{align*}
&\left|\int_{|x-y|>\epsilon}\frac{1}{x-y}\Bigl(\sum h_{I}\chi_{I}\Bigr)dy\right| \\
&\leq 0+\sum\int_{I}|K(x-y)-K(x-y_{I})|\,|h_{I}(y)|\,dy \\
&\qquad+\sum\frac{2}{\epsilon}\int_{|x-y|<\frac{3\epsilon}{2}}|h_{I}(y)|\chi_{I}(y)\,dy.
\end{align*}
于是, 对 $x\in(B^{\star})^{c}$,
\[
|H^{\star}f(x)|\leq\sum\int_{I}|K(x-y)-K(x-y_{I})|\,|h_{I}(y)|\,dy
+\Bigl[\sum|h_{I}(y)|\chi_{I}\Bigr]^{\star},
\]
这个估计足以让我们完成证明.
\end{proof}

\section{第 2 章评注}
\label{sec:2.6}

\noindent (2.1):\quad Bochner--Riesz 平均由 Riesz 平均 $(1-|\xi|)^{\lambda}$ 演化而来. 两者的乘子性质等价, 因为它们只差一个与 $\mathcal{S}$ 中函数的乘法. 像 Bochner~\cite{bochner} 那样把乘子用 $|\xi|^2$ 加以改动, 使我们能够算出 $T_\lambda$ 的精确公式, 而这些公式在 Bochner 关于局部化失效的证明中起了关键作用.

\noindent (2.4):\quad 关于 Bessel 算子的细节可见 Watson~\cite{watson} 6.1. 这里引用的公式归功于 Sonine; 另见 Stein 与 Weiss~\cite{steinweiss}.

\noindent (2.18):\quad 这个结果的最早版本归功于 Cotlar, 由此才有``Cotlar 引理''的说法. 一个加强的版本由 Stein 得到; 这里的证明取自 Knapp--Stein~\cite{knapp-stein}. 更几何化的结果见 Fefferman~\cite{fefferman-pointwise}. 关于谱半径定理, 见 Rudin~\cite{rudinfa}.

\noindent (2.23):\quad 再说一次, 看这个结果最直观的方式是取 $f$ 为一个 $\delta$ 函数. Stein~\cite{stein-limits} 展示了如何把这类直观严格化.

\noindent (2.26):\quad 推论~\ref{cor:2.26} 有一个逆命题, 它用到 Stein 关于算子序列极限的结果~\cite{stein-limits}. 我们一直在说, 一致估计有助于交换极限; Stein 的结果表明, 推论~\ref{cor:2.26} 中极限的存在性蕴含弱 $(1,1)$ 估计.

\noindent (2.28):\quad 把 $\mathbb{R}$ 分解为好集与坏集的做法, 以~\cite{calderon-zygmund} 命名为 Calder\'on--Zygmund 分解. 若用球或立方体作平均, 这个结果与定理~\ref{thm:2.24} 都可推广到 $\mathbb{R}^n$.

\noindent (2.30):\quad 这个结果归功于 Whitney~\cite{whitney}, 它的建立是为了控制函数在集合边界附近的可微性质. 为什么这样做会有帮助呢?

\noindent (2.32):\quad Hilbert 变换在 $L^{p}$ 上有界的证明归功于 M.~Riesz~\cite{riesz}. Riesz 用的是复变技术; 特别地, Hilbert 变换是把调和函数变成其共轭调和函数的算子 (从而使这一对函数成为解析函数). Kolmogorov~\cite{kolmogorov} 同样用复变方法证明了弱 $(1,1)$ 有界性, 证明了单位圆盘内解析函数的边值的存在性. 这里给出的证明是 Calder\'on 与 Zygmund 证明的一个变体; 它的独特优点是可以推广到多维. 我们采用了 Stein~\cite{stein-singular} 的表述. 核上的条件来自 H\"ormander~\cite{hormander}. 也可以有别的光滑性条件; 见 Stein~\cite{stein-singular}.

\noindent (2.40):\quad 我们陈述这个结果时, 仿佛它适用于 H\"ormander 型的核 $K$; 但事实上我们需要对 $K$ 使用逐点估计, 比如 $|x-y|\leq\epsilon$ 这样的估计. 更一般的处理见 Stein~\cite{stein-singular}.

% 第 3 章 Good Lambda and Weighted Norm Inequalities (原书第 61-74 页)
\chapter{Good \texorpdfstring{$\lambda$}{lambda} 不等式与加权范数不等式}
\label{chap:3}

上一章中, 我们得到了用 $f$ 的大小估计 $Hf$ 大小的整体性结果. 这还不是这个课题的终点; 我们所做的那几类证明提示, 或许存在用 $f(x)$ 或 $f^{\star}(x)$ 估计 $Hf(x)$ 的办法. 我们也许能得到逐点估计, 而不是 $L^p$ 估计.

我们实际上会做一些接近此事、但更弱的事情. 若 $I$ 是一个区间, 我们能指望 $\int_I|Hf|^p\leq C\int_I|f|^p$ 吗? 设想 $I$ 是一个非常小的区间, 这就接近逐点信息了, 但它仍让我们得以使用函数论与 $L^p$ 空间的想法. 为理解这一点, 我们的模型问题是: 若 $\omega>0$, 是否对所有的 $f$ 都有
\[
\int|H f|^{p}\omega\leq C\int|f|^{p}\omega\,?
\]
这称为一个``加权范数不等式''; 我们可以把 $\omega$ 想成把大部分质量集中在一个区间上、而在外面非常小的权重. 我们将用来分析这个问题的技术工具称为 \emph{good-$\lambda$ 不等式} (good-$\lambda$ inequality). 它把 $Hf$ 的大小与 $f^{\star}$ 的大小作比较, 并对``$f^{\star}$ 控制 $Hf$ 的大小''这一直观给出精确的估计.

\section{Good \texorpdfstring{$\lambda$}{lambda} 不等式}
\label{sec:3.1}
\index{good-$\lambda$ 不等式 (good-$\lambda$ inequality)|emph}

我们的直观是: 当 $Hf$ 很大时, $f^{\star}$ 必然也很大. 下面的定理精确地表达了这一点: $Hf$ 很大而 $f^{\star}$ 很小的那个集合, 其测度可以被做得很小.

\begin{theorem}\label{thm:3.1}
存在 $C>0$, 使得对一切 $\gamma<\frac{1}{2}C$\footnote{译者注: 原书此处印作 $\gamma>\frac{1}{2}C$, 应为 $\gamma<\frac{1}{2}C$, 否则与证明及定理~\ref{thm:3.2} 中 $\gamma$ 的取法矛盾.} 与一切 $\lambda$,
\[
|\{x \mid H^{\star}f(x)>2\lambda,\ f^{\star}<\gamma\lambda\}|\leq C\gamma\,|\{x \mid H^{\star}f(x)>\lambda\}|.
\]
\end{theorem}

\begin{proof}
我们先把问题局部化到区间上; 对集合 $B=\{H^{\star}f>\lambda\}$ 应用 Whitney 分解定理, 即定理~\ref{thm:2.30}. 我们得到互不相交的区间 $I$, 满足 $B=\bigcup I$, $|I|\leq d(B^{c},I)\leq 4|I|$. 由于诸区间互不相交, 只需证明
\[
|\{x\in I \mid H^{\star}f(x)>2\lambda,\ f^{\star}<\gamma\lambda\}|\leq C\gamma|I|.
\]
我们用 $f$ 控制 $H^{\star}f$ 的模型是什么? 在定理~\ref{thm:2.40} 的证明中控制 c) 型的项时, $x$ 靠近区间 $J$ 而 $h_j$ 支在 $J$ 中. 这给了我们用 $f^{\star}$ 控制 $H^{\star}f$ 的办法. 这里 $f$ 并不支在 $I$ 中, 对那些不靠近 $I$ 的 $x$ 还会有额外的项. Whitney 分解替我们控制住了这一点: 若 $x$ 不靠近 $I$, 则 $x$ 靠近 $B^{c}$, 而在 $B^{c}$ 上 $H^{\star}f\leq\lambda$.

这一直观提示把 $f$ 分解为靠近 $I$ 的部分与集中于 $\tilde{I}^{c}$ 的部分. 我们选两个参考点来控制 $f$ 的大小: 取 $z_0\in I$ 使 $f^{\star}(z_0)\leq\gamma\lambda$; 如同~\ref*{thm:2.40}c) 中那样, 它将控制 $H^{\star}f$ 的局部部分. 我们再取 $x_0\in B^{c}$, 并免费得到 $H^{\star}f(x_0)\leq\lambda$. 由于 $d(B^{c},I)\leq 4|I|$, 我们可以选取 $x_0$ 使 $d(x_0,I)\leq 5|I|$, 于是这将确定``靠近 $I$''与``远离 $I$''的含义. 但有技术上的困难, 因为 $H^{\star}f(x)$ 与 $H^{\star}f(x_0)$ 不能直接比较; 后者要求我们用在 $x_{0}$ 处居中的区间来度量, 而 $I$ 并不以 $x_0$ 为中心. 我们构造一个比 $I$ 大、但至少以 $x_{0}$ 为中心的区间 $J$. 开始.

取 $z_0\in I$ 使 $f^{\star}(z_0)<\gamma\lambda$; 取 $x_0\in B^{c}$, $d(x_0,I)\leq 5|I|$. 取一个以 $x_0$ 为中心、长度为 $I$ 的 $20$ 倍的区间 $J$. 显然 $I\subset J$, 因为若 $x\in I$, 则 $|x-x_{0}|\leq d(x_{0},I)+|I|\leq 6|I|$.

令 $g=f\chi_J$, $b=f\chi_{J^{c}}$. 我们将证明
\[
\left|\{x\in I \mid H^{\star}g(x)>\frac{\lambda}{2}\}\right|\leq\frac{C}{\lambda}|I|\,f^{\star}(x_{0})\leq C\gamma|I|;
\]
\[
H^{\star}b(x)\leq H^{\star}f(x_{0})+C f^{\star}(x_{0})\leq\lambda+C\gamma\lambda.
\]
有了这些结果,
\begin{align*}
|\{x\in I \mid H^{\star}g(x)>2\lambda,\ f^{\star}<\gamma\lambda\}|
&\leq|\{x\in I \mid H^{\star}g(x)>2\lambda-H^{\star}b,\ f^{\star}<\gamma\lambda\}| \\
&\leq|\{x\in I \mid H^{\star}g(x)>2\lambda-C\gamma\lambda,\ f^{\star}<\gamma\lambda\}| \\
&\leq\left|\{x\in I \mid H^{\star}g(x)>\frac{\lambda}{2}\}\right|\leq C\gamma|I|,
\end{align*}
至少当 $\gamma<\frac{1}{2}C$ 时如此.

对 $H^{\star}g$ 在 $I$ 上的估计并不困难, 尽管这本该是 $H^{\star}$ 最糟糕的部分, 因为我们已经知道 $H^{\star}$ 是弱 $(1,1)$ 的:
\begin{align*}
\left|\left\{x\in I \mid H^{\star}g(x)>\frac{\lambda}{2}\right\}\right|
&\leq\frac{C}{\lambda}\|g\|_{1}=\frac{C}{\lambda}\int_{J}|f| \\
&=\frac{C}{\lambda}|J|\,|J|^{-1}\int_{J}|f|
=20\,\frac{C}{\lambda}|I|\,\frac{1}{|J|}\int_{J}|f|.
\end{align*}
但 $z_0\in I\subset J$, 于是
\[
\frac{1}{|J|}\int_{J}|f|\leq f^{\star}(z_{0})<\gamma\lambda;
\]
整个这一项被 $20C\gamma|I|$ 控制.

为了用 $H^{\star}f(x_{0})$ 控制 $H^{\star}b(x)$, 我们试图把
\[
\int_{|x-y|>\epsilon}\frac{b(y)}{x-y}\,dy
\]
用
\[
\int_{|x_{0}-y|>\epsilon}\frac{b(y)}{x_{0}-y}\,dy
\]
表示出来. 这引入两类误差: 来自 $\frac{1}{x-y}-\frac{1}{x_0-y}$ 的误差, 与来自 $\{|x_0-y|>\epsilon\}\,\Delta\,\{|x-y|>\epsilon\}$ 的误差. 为控制它们, 令
\begin{align*}
S_{1}&=\{y \mid |x_{0}-y|>\epsilon,\ |x-y|<\epsilon\}, \\
S_{2}&=\{y \mid |x_{0}-y|<\epsilon,\ |x-y|<\epsilon\}, \\
E_{j}&=\int_{S_{j}}\frac{|b(y)|}{|x-y|}\,dy.
\end{align*}
则:
\begin{align*}
\left|\int_{|x-y|>\epsilon}\frac{b(y)}{x-y}\,dy\right|
&\leq\left|\int_{|x_{0}-y|>\epsilon}\frac{b(y)}{x-y}\,dy\right|+E_{1}+E_{2} \\
&\leq\left|\int_{|x_{0}-y|>\epsilon}\frac{b(y)}{x_{0}-y}\,dy\right|
+\int_{|x_{0}-y|>\epsilon}\left|\frac{1}{x_{0}-y}-\frac{1}{x-y}\right|dy+E_{1}+E_{2}.
\end{align*}
如所安排, 第一项被 $H^{\star}b(x_0)$ 控制; 不幸的是, 我们在 $x_0$ 处拥有的控制是关于 $H^{\star}f(x_0)$ 而非 $H^{\star}b(x_0)$ 的. 这意味着还要做更多工作, 也显示了我们对 $J$ 的选择何以派上用场. 若取 $y\in J^{c}$, 则 $|y-x_0|>10|I|$. 于是, 若 $y\in\operatorname{supp}b=J^{c}$,
\[
\{y \mid |x_{0}-y|>\epsilon,\ y\in J^{c}\}=\{y \mid |x_{0}-y|>\max\{\epsilon,\ 10|I|\}\},
\]
且
\begin{align*}
\left|\int_{|x_{0}-y|>\epsilon}\frac{b(y)}{x_{0}-y}\,dy\right|
&=\left|\int_{\{y \mid |x_{0}-y|>\max\{\epsilon,\,10|I|\}\}}\frac{f(y)}{x_{0}-y}\,\chi_{J^{c}}\,dy\right| \\
&\leq\left|\sup_{\delta}\int_{|x_{0}-y|>\delta}\frac{f(y)}{x_{0}-y}\,dy\right|=H^{\star}f(x_{0}).
\end{align*}
由卷积核的光滑性控制的那一项, 是我们以前处理过的; 它被
\[
\int_{|x_{0}-y|>\epsilon}|x_{0}-y|^{-1}|x-y|^{-1}|x_{0}-x|\,|f(y)|\,\chi_{J^{c}}\,dy
\]
控制. 我们需要用 $f^{\star}(z_0)$ 控制这一项, 于是需要把 $x_0$ 与 $z_0$ 联系起来, 并用这个关系替换 $|x_0-y|$. 不存在这样的关系, 恰恰意味着我们自己来构造一个. 令
\[
F_{k}=\{y \mid 2^{k+1}\epsilon>|x_{0}-y|>2^{k}\epsilon\}.
\]
我们将证明: 若 $x\in I$, $y\in J^{c}$, $y\in F_k$, 则:
\begin{gather*}
k\geq\log\Bigl(\frac{10|I|}{\epsilon}\Bigr)=C_{\epsilon};\qquad
|x_{0}-y|^{-1}\leq(2^{k}\epsilon)^{-1}; \\
|x-y|^{-1}\leq(2^{k-1}\epsilon)^{-1};\qquad
|z_{0}-x_{0}|<2^{k+1}\epsilon;\qquad
|x-x_{0}|<6|I|.
\end{gather*}
若能证明这些, 则由光滑性控制的那一项被
\[
\sum_{k=C_{\epsilon}}^{\infty}\frac{1}{2^{k}\epsilon}\int_{F_{k}}|f(y)|\,dy\,\frac{|I|}{2^{k-1}\epsilon}
\]
控制, 进而
\begin{align*}
\sum_{k=C_{\epsilon}}^{\infty}&\frac{10|I|}{\epsilon}\,2^{-(k-1)}\,\frac{1}{2^{k}\epsilon}\int_{|y-z_{0}|\leq 2^{k+2}\epsilon}|f(y)|\,dy \\
&\leq\frac{40|I|}{\epsilon}\,f^{\star}(z_{0})\,2^{-C_{\epsilon}}\leq 400\,f^{\star}(z_{0}).
\end{align*}
如今的估计需要的只是耐心; 我们举几个例子. 首先, $y\in F_k$ 蕴含 $2^{k+1}\epsilon>|x_0-y|$, 于是 $y\in F_k\cap J^{c}$ 蕴含 $k\geq C_{\epsilon}$. 由于 $y\in F_k$, 自动地有
\[
|x_{0}-y|^{-1}\leq(2^{k}\epsilon)^{-1};
\]
为估计 $|x-y|$,
\[
|x-y|\geq|x_{0}-y|-|x-x_{0}|>2^{k}\epsilon-6|I|.
\]
但 $2^{k+1}\epsilon>|x_0-y|>10|I|$, 于是 $6|I|<\frac{6}{10}\cdot 2^{k+1}\epsilon$, 从而 $|x-y|\geq 2^{k-1}\epsilon$. 最后,
\[
|z_{0}-x_{0}|\leq d(x_{0},I)+|I|\leq 6|I|<10|I|<2^{k+1}\epsilon,
\]
且
\[
|x_{0}-x|\leq d(x_{0},I)+|I|\leq 6|I|.
\]
为完成证明, 我们必须用 $f^{\star}(z_0)$ 控制误差项 $E_i$; 我们需要把 $|x_0-y|$ 的信息换成 $|z_0-y|$ 的信息. 对 $E_1$ 项:
\[
S_{1}=\{y \mid |x-y|<\epsilon,\ |x_{0}-y|>\epsilon\}.
\]
由于 $z_0\in I$ 且 $x\in I$, $y$ 应当靠近 $x$, 而 $x$ 到 $z_0$ 的距离至多是 $|I|$. 但这必须远小于 $\epsilon$, 因为 $x\in I$ 而 $y\in J^{c}$. 于是:
\[
\epsilon>|x-y|>10|I|,
\]
\[
|z_{0}-y|<|z_{0}-x|+|x-y|<|I|+\epsilon<2\epsilon,
\]
且
\[
S_{1}\subset\{y \mid |y-z_{0}|<2\epsilon\}.
\]
此外,
\[
|x-y|\geq|x_{0}-y|-|x_{0}-x|>\epsilon-6|I|\geq\epsilon-\frac{6}{10}\epsilon=\frac{2}{5}\epsilon,
\]
于是 $|x-y|^{-1}<\frac{5}{2\epsilon}$, 且
\[
E_{1}\leq\frac{5}{2\epsilon}\int_{|y-z_{0}|<2\epsilon}|f(y)|\,dy\leq 5f^{\star}(z_{0}).
\]
对于来自 $S_2=\{y \mid |x_0-y|<\epsilon,\ |x-y|>\epsilon\}$ 的项, 我们注意到 $y$ 必须靠近 $z_0$, 因为 $x_0$ 靠近它; 同样要把 $|I|$ 与 $\epsilon$ 作比较:
\[
\epsilon>|x_{0}-y|>|x-y|-|x-x_{0}|\geq 20|I|-6|I|=14|I|;
\]
\[
|z_{0}-y|\leq|z_{0}-x_{0}|+|x_{0}-y|\leq 6|I|+\epsilon<2\epsilon.
\]
由于 $|x-y|^{-1}<\epsilon^{-1}$, 故 $E_2<2f^{\star}(z_0)$.
\end{proof}

\begin{theorem}\label{thm:3.2}
若 $0<p<\infty$, 则 $\|H^{\star}f\|_p\leq C_p\|f^{\star}\|_p$.
\end{theorem}

\begin{proof}
\begin{align*}
\|H^{\star}f\|_{p}^{p}
&=p\int_{0}^{\infty}|\{H^{\star}f>\lambda\}|\,\lambda^{p-1}\,d\lambda \\
&=2^{p}p\int_{0}^{\infty}|\{H^{\star}f>2\lambda\}|\,\lambda^{p-1}\,d\lambda \\
&\leq 2^{p}p\int_{0}^{\infty}|\{H^{\star}f>2\lambda,\ f^{\star}<\gamma\lambda\}|\,\lambda^{p-1}\,d\lambda
+2^{p}p\int_{0}^{\infty}|\{f^{\star}>\gamma\lambda\}|\,\lambda^{p-1}\,d\lambda \\
&\leq 2^{p}pC\gamma\int_{0}^{\infty}|\{H^{\star}f>\lambda\}|\,\lambda^{p-1}\,d\lambda
+\frac{2^{p}}{\gamma^{p}}\|f^{\star}\|_{p}^{p}.
\end{align*}
选取 $\gamma$ 使 $1-2^{p}C\gamma>0$, 便得到
\[
\left\|H^{\star}f\right\|_{p}\leq 2\gamma\bigl(1-2^{p}C\gamma\bigr)^{-\frac{1}{p}}\left\|f^{\star}\right\|_{p}. \qedhere
\]
\end{proof}

\section{\texorpdfstring{$A_p$}{Ap} 条件}
\label{sec:3.2}

现在我们在区间上控制了 $H^{\star}f$, 并试图用它来回答这个问题: 对哪些 $\omega>0$, 我们能指望
\[
\int|H f|^{p}\omega\leq C\int|f|^{p}\omega\,?
\]
平凡的条件是 $a<\omega<b$; 既然我们能够指望在区间上控制 $Hf$, 我们可以希望有一个条件说: $\omega$ 与它在区间上的平均不能偏离太多; 当然我们希望这种偏离在所有区间上被一致地控制. 表达这一点的方式有很多; 我们可以指望控制 $|\{x\in I \mid \omega(x)>\omega_I\}|$, 也可以指望同时界住 $\int_I\omega$ 与 $\int_I\omega^{-1}$.

\begin{definition}\label{def:3.3}
\index{$A_p$ 权 (condition $A_p$)|emph}
若 $\omega$ 是局部可积的正函数, 我们说 $\omega\in A_p$, 如果:

当 $p=1$ 时,
\[
\omega_{I}\leq C\inf_{x\in I}\omega(x),
\]
其中 $C$ 与 $I$ 无关.

当 $1<p<\infty$ 时,
\[
\left(\frac{1}{|I|}\int_{I}\omega\right)\left(\frac{1}{|I|}\int_{I}\omega^{\frac{1}{p-1}}\right)^{p-1}\leq C,
\]
其中 $C$ 与 $I$ 无关.

当 $p=\infty$ 时,
\[
\frac{\omega(E)}{\omega(I)}\leq\left(\frac{|E|}{|I|}\right)^{\delta}
\]
对某个与 $I$ 无关的 $\delta$ 与一切可测的 $E\subset I$ 成立.
\end{definition}

\begin{remark}\label{rem:3.4}
若 $\omega\in A_\infty$, 令
\[
E=\{x\in I \mid \omega(x)>\beta\omega_{I}\}.
\]
则 $|E|\leq\beta^{\frac{1}{1-\delta}}|I|$, 于是 $A_\infty$ 条件蕴含 $\omega$ 不偏离其平均太多. 反之, 假设
\[
|\{x\in I \mid \omega(x)>\beta\omega_{I}\}|\leq\beta^{\delta}|I|.
\]
则
\begin{align*}
\int_{I}\omega^{-\frac{1}{p-1}}
&=\frac{1}{p-1}\int\left|\left\{\omega^{-1}\geq\lambda\right\}\right|\lambda^{\frac{1}{p-1}-1}\,d\lambda \\
&\leq\frac{1}{p-1}\int_{0}^{\omega_{I}^{-1}}|I|\,\lambda^{\frac{1}{p-1}-1}\,d\lambda
+p\int_{\omega_{I}^{-1}}^{\infty}\left(\lambda^{-1}\omega_{I}^{-1}\right)^{\delta}|I|\,\lambda^{\frac{1}{p-1}-1}\,d\lambda
=C_{\delta}|I|\,\omega_{I}^{-\frac{1}{p-1}}.
\end{align*}
于是
\[
\left(\frac{1}{|I|}\int_{I}\omega\right)\left(\frac{1}{|I|}\int_{I}\omega^{-\frac{1}{p-1}}\right)^{p-1}\leq C_{\delta}^{p-1},
\]
即 $\omega\in A_p$. 对我们而言, 这个计算意味着 $A_p$ 条件与控制函数对其平均的偏离密切相关.
\end{remark}

\begin{lemma}\label{lem:3.5}
若 $\omega\in A_p$, 则
\[
|\{x\in I \mid \omega(x)\leq\beta\omega_{I}\}|\leq C\beta^{\frac{1}{p-1}}|I|.
\]
\end{lemma}

\begin{proof}
把 $A_{p}$ 条件改写,
\begin{align*}
C|I|&\geq(\omega_{I})^{\frac{1}{p-1}}\int_{I}\omega^{-\frac{1}{p-1}}
\geq(\omega_{I})^{\frac{1}{p-1}}\int_{\{\omega<\beta\omega_{I}\}}\omega^{-\frac{1}{p-1}} \\
&\geq(\omega_{I})^{\frac{1}{p-1}}\int_{\{\omega<\beta\omega_{I}\}}(\beta\omega_{I})^{-\frac{1}{p-1}}
=\beta^{-\frac{1}{p-1}}\left|\left\{x\in I \mid \omega<\beta\omega_{I}\right\}\right|. \qedhere
\end{align*}
\end{proof}

\begin{remark}\label{rem:3.6}
对本质上为常数的函数, 不同 $p$ 的 $L^p$ 范数没有差别. ``$A_p$ 函数与其平均相差不多''这一概念使我们能够对 $A_p$ 函数证明基本相同的结果. 我们将证明 $A_p$ 函数既属于 $A_\infty$ 也属于 $A_{p-\epsilon}$. 这两个事实都来自``$L^p$ 范数等价性''的一个精确表述. 记得我们是对 $\omega$ 在区间上作平均, 所以包含关系
\[
L^{p}\subset L^{p+\epsilon}
\]
由 H\"older 不等式得出. 而 $A_p\subset A_{p-\epsilon}$ 这个令人惊讶的结果来自一个``反 H\"older 不等式''.
\end{remark}

\begin{theorem}\label{thm:3.7}
\index{反 H\"older 不等式 (reverse H\"older inequality)}
设 $\omega\in A_p$ 对某个 $p$ 成立, $1<p<\infty$. 则存在 $a$, $q>1$, 使
\[
\left(\frac{1}{|I|}\int_{I}\omega^{q}\right)^{\frac{1}{q}}\leq C\left(\frac{1}{|I|}\int_{I}\omega\right).
\]
\end{theorem}

\begin{proof}
我们先证明: 对某个 $\beta<1$, 有一个 good-$\lambda$ 不等式:
\[
\omega\{x\in I \mid \omega(x)>\lambda\}\leq C\lambda\left|\left\{x\in I \mid \omega(x)>\beta\lambda\right\}\right|.
\]
假设它成立, 定理就容易了:
\[
\int_{I}\omega^{q}=\beta^{-q}C_{q}\int_{0}^{\infty}\lambda^{q-1}\left|\left\{x\in I \mid \omega>\beta\lambda\right\}\right|d\lambda
\]
\begin{align*}
&\geq C'q\int_{\omega_{I}}^{\infty}\lambda^{q-2}\int_{\{x\in I \mid \omega\geq\lambda\}}\omega(x)\,dx\,d\lambda \\
&=C'q\int_{\{x\in I \mid \omega>\omega_{I}\}}\omega(x)\int_{\omega_{I}}^{\omega}\lambda^{q-2}\,d\lambda\,dx \\
&=C'q\int_{\{x\in I \mid \omega>\omega_{I}\}}\omega(x)\left[\frac{\omega^{q-1}(x)}{q-1}-\frac{\omega_{I}^{q-1}}{q-1}\right]dx.
\end{align*}
现在在集合
\[
\{x\in I \mid \omega>\omega_{I}\}
\]
上,
\[
-\omega\,\omega_{I}^{q-1}>-\omega_{I}^{q},
\]
于是
\begin{align*}
C'q\int_{\{x\in I \mid \omega>\omega_{I}\}}\omega(x)\left[\frac{\omega^{q-1}(x)}{q-1}-\frac{\omega_{I}^{q-1}}{q-1}\right]dx
&\geq\frac{C'q}{q-1}\int_{\{x\in I \mid \omega>\omega_{I}\}}\bigl(\omega^{q}-\omega_{I}^{q}\bigr).
\end{align*}
但在 $\{\omega\leq\omega_I\}$ 上, 被积函数是负的, 于是可以把整个这个集合添加进来而不破坏不等式:
\[
\int_{I}\omega^{q}\geq\frac{C'q}{q-1}\int_{I}\bigl(\omega^{q}-\omega_{I}^{q}\bigr),
\]
从而
\[
\left(\frac{C'q}{q-1}-1\right)\frac{1}{|I|}\int_{I}\omega^{q}\leq\left(\frac{1}{|I|}\int_{I}\omega\right)^{q}.
\]
取 $q$ 充分接近 $1$, 证明即告完成.
\end{proof}

还剩下一个 good-$\lambda$ 不等式要证, 它分析 $\omega$ 在区间上的变化. 我们通过一个 Calder\'on--Zygmund 分解来控制它: 构造集合 $B=\{x\in I \mid \omega(x)>\lambda\}$, 并设 $\{J\}$ 是互不相交的区间 $J\subset I$, 满足 $\lambda\leq\omega_J\leq 2\lambda$. 则 $\bigcup J=B$, 且若 $x\in I-B$, 则 $\omega(x)\leq\lambda$. 令 $\alpha=C\beta^{\frac{1}{p-1}}$, 其中 $C$, $\beta$ 如引理~\ref{lem:3.5} 中所述. 若 $\beta<\beta_0$, 则 $\alpha<1$, 且
\[
\omega\{x\in I \mid \omega>\lambda\}=\sum\omega(J)\leq 2\lambda\sum|J|.
\]
但
\[
|J|=|\{x\in J \mid \omega>\beta\omega_{J}\}|+|\{x\in J \mid \omega\leq\beta\omega_{J}\}|.
\]
由于 $\omega_{J}>\lambda$, 它被
\[
|\{x\in J \mid \omega>\beta\omega_{J}\}|+\alpha|J|
\]
控制. 由此得
\[
|J|\leq(1-\alpha)^{-1}\left|\{x\in J \mid \omega>\beta\lambda\}\right|,
\]
最后得
\begin{align*}
\omega\{x\in I \mid \omega>\lambda\}
&\leq C\lambda\sum\left|\left\{x\in J \mid \omega>\beta\lambda\right\}\right| \\
&=C\lambda\left|\left\{x\in B \mid \omega>\beta\lambda\right\}\right|
\leq C\lambda\left|\left\{x\in I \mid \omega>\beta\lambda\right\}\right|.
\end{align*}

\begin{corollary}\label{cor:3.8}
若 $\omega\in A_p$, 则 $\omega\in A_\infty$.
\end{corollary}

\begin{proof}
如定理~\ref{thm:3.7} 那样选取 $q$. 则:
\begin{align*}
\omega(E)&=\int\chi_{E}\,\omega\leq\|\chi_{E}\|_{q'}\,\Bigl(\int\chi_{E}\,\omega^{q}\Bigr)^{\frac{1}{q}}
\leq|E|^{\frac{1}{q'}}\,|I|^{\frac{1}{q}}\left(\frac{1}{|I|}\int_{I}\omega^{q}\right)^{\frac{1}{q}} \\
&\leq C|E|^{\frac{1}{q'}}\,|I|^{\frac{1}{q}}\,\omega_{I}
=C|E|^{\frac{1}{q'}}\,|I|^{\frac{1}{q}-1}\int_{I}\omega
=C\left(\frac{|E|}{|I|}\right)^{\frac{1}{q'}}\omega(I). \qedhere
\end{align*}
\end{proof}

\begin{corollary}\label{cor:3.9}
若 $\omega\in A_p$, 则存在 $\epsilon>0$ 使 $\omega\in A_{p-\epsilon}$.
\end{corollary}

\begin{proof}
改换记号即知 $\omega\in A_p$ 当且仅当
\[
\omega^{-\frac{1}{p-1}}\in A_{p'}.
\]
把定理~\ref{thm:3.7} 应用于 $A_{p'}$, 存在 $q>1$ 使
\[
\left(\frac{1}{|I|}\int_{I}\omega^{-\frac{q}{p-1}}\right)^{\frac{1}{q}}
\leq C\frac{1}{|I|}\int_{I}\omega^{-\frac{1}{p-1}}
\leq C\left(\frac{1}{|I|}\int_{I}\omega\right)^{-\frac{1}{p-1}},
\]
最后一步用了 $\omega\in A_p$. 总之, $\omega\in A_\gamma$, 其中 $\frac{1}{\gamma-1}=\frac{q}{p-1}$; 由于 $q>1$, 故 $\gamma<p$.
\end{proof}

\section{加权范数不等式}
\label{sec:3.3}

\begin{lemma}\label{lem:3.10}
\index{加权范数不等式 (weighted norm inequality)}
若 $\omega\in A_p$, 则
\[
f_{I}\leq C\left(\frac{1}{\omega(I)}\int_{I}|f|^{p}\omega\right)^{\frac{1}{p}}.
\]
若 $\tilde{I}$ 是 $I$ 的倍缩 (即 $\tilde{I}$ 与 $I$ 同心但长度为其三倍), 则 $\omega(\tilde{I})\leq C\omega(I)$; 又若 $J$ 是与 $I$ 同长度、相邻的区间, 则 $\omega(J)\leq c\,\omega(I)$.
\end{lemma}

\begin{proof}
\begin{align*}
f_{I}&=\frac{1}{|I|}\int_{I}f=\frac{1}{|I|}\int_{I}f\,\omega^{\frac{1}{p}}\,\omega^{-\frac{1}{p}} \\
&\leq\frac{1}{|I|}\left(\int_{I}f^{p}\omega\right)^{\frac{1}{p}}\left(\int_{I}\omega^{-\frac{p'}{p}}\right)^{\frac{1}{p'}}.
\end{align*}
但 $\frac{p'}{p}=\frac{1}{p-1}$, 且第二个积分被
\begin{align*}
\left(\int_{I}\omega^{-\frac{p'}{p}}\right)^{\frac{1}{p'}}
&\leq C|I|^{\frac{1}{p'}}\left(\frac{\omega(I)}{|I|}\right)^{-\frac{1}{p(p-1)}} \\
&=C|I|^{\frac{1}{p'}}|I|^{-\frac{1}{p}}\bigl[\omega(I)\bigr]^{-\frac{1}{p}}
\end{align*}
控制.

为得到关于 $\tilde{I}$ 的结果, 令 $f=\chi_{I}$. 则
\[
\frac{1}{2}=f_{\tilde{I}}\leq C\left(\frac{1}{\omega(\tilde{I})}\int_{\tilde{I}}\chi_{I}\,\omega\right)^{\frac{1}{p}},
\]
即
\[
C\leq\frac{\omega(I)}{\omega(\tilde{I})}.
\]
类似的想法可以处理 $J$.
\end{proof}

\begin{lemma}\label{lem:3.11}
\index{倍增条件 (doubling condition)}
设 $f$ 局部可积, 并取 $\omega\in A_p$. 定义
\[
f^{\star\star}(x)=\sup_{x\in I}\frac{1}{\omega(I)}\int_{I}f(y)\,\omega(y)\,dy.
\]
则对 $1<p<\infty$,
\[
\int|f^{\star\star}|^{p}\,\omega\leq C\int|f|^{p}\,\omega.
\]
\end{lemma}

\begin{proof}
回忆定理~\ref{thm:2.24} 的证明; 与那里一样, 我们得到 $E=\{x \mid f^{\star\star}(x)>\lambda\}$ 的一个区间覆盖, 并可以施以选取过程得到互不相交的区间. 我们需要 $E$ 的测度与互不相交区间的测度之和可比, 但 $E\subset\bigcup I^{\star}$, 故 $\omega(E)\leq\sum\omega(I^{\star})$. 而 $I^{\star}$ 的长度是 $I$ 的五倍, 于是引理~\ref{lem:3.10} 告诉我们
\[
\omega(E)\leq\sum\omega(I^{\star})\leq C\sum\omega(I)\leq C\frac{1}{\lambda}\sum\int_{I}f,
\]
等等. 简言之, 只要测度满足``倍增条件''
\[
\omega(I^{\star})\leq C\omega(I),
\]
我们就能对 Lebesgue 测度之外的测度证明极大定理, 即定理~\ref{thm:2.24}.
\end{proof}

\begin{theorem}\label{thm:3.12}
若 $\omega\in A_p$ 对 $1<p<\infty$ 成立, 则
\[
\int|f^{\star}|^{p}\,\omega\leq C\int|f|^{p}\,\omega.
\]
\end{theorem}

\begin{proof}
由于 $\omega\in A_p$, 由推论~\ref{cor:3.9} 存在 $\gamma<p$ 使 $\omega\in A_\gamma$. 则
\[
f_{I}\leq C\left(\frac{1}{\omega(I)}\int_{I}f^{\gamma}\,\omega\right)^{\frac{1}{\gamma}},
\]
从而
\[
f^{\star}(x)\leq C\left[\left(f^{\gamma}\right)^{\star\star}(x)\right]^{\frac{1}{\gamma}}.
\]
于是
\[
\int|f^{\star}|^{p}\,\omega\leq C\int\bigl|\bigl(f^{\gamma}\bigr)^{\star\star}\bigr|^{\frac{p}{\gamma}}\,\omega.
\]
由于 $\frac{p}{\gamma}>1$, 引理~\ref{lem:3.11} 给出最后的积分被
\[
C\int|f^{\gamma}|^{\frac{p}{\gamma}}\,\omega
\]
控制, 这就证明了结果.
\end{proof}

\begin{theorem}\label{thm:3.13}
若 $\omega\in A_p$, 则
\[
\int|H^{\star}f|^{p}\,\omega\leq C\int|f^{\star}|^{p}\,\omega\leq C\int|f|^{p}\,\omega.
\]
\end{theorem}

\begin{proof}
由于 $\omega\in A_p$, 由推论~\ref{cor:3.8} 有 $\omega\in A_\infty$; 按 $A_\infty$ 条件选取 $\delta$. 对 $\{x \mid H^{\star}f(x)>\lambda\}$ 施以 Whitney 分解, 并令
\[
E=\{x\in I \mid H^{\star}f(x)>2\lambda,\ f^{\star}(x)<\gamma\lambda\}.
\]
则由 good-$\lambda$ 不等式, 即定理~\ref{thm:3.1}, 得 $|E|\leq C\gamma|I|$, 因此
\[
\omega(E)\leq\omega(I)\left[\frac{|E|}{|I|}\right]^{\delta}\leq C\gamma^{\delta}\,\omega(I).
\]
于是,
\[
\omega\{x\in I \mid H^{\star}f(x)>2\lambda,\ f^{\star}(x)<\gamma\lambda\}
\leq C\gamma^{\delta}\,\omega\{x\in I \mid H^{\star}f(x)>\lambda\}.
\]
现在 $L^{p}$ 有界性的证明与定理~\ref{thm:3.2} 的证明完全一样.
\end{proof}

\begin{proposition}\label{prop:3.14}
若 $1<p<\infty$, 且
\[
\int|H f|^{p}\,\omega\leq C\int|f|^{p}\,\omega
\]
对所有的 $f\in L^{p}$ 成立, 则 $\omega\in A_{p}$.
\end{proposition}

\begin{proof}
设 $I$, $J$ 如引理~\ref{lem:3.10} 中那样. 则若 $x\in J$, $y\in I$, 有 $|x-y|\leq 2|I|$, 于是若 $f$ 支在 $I$ 中,
\[
|H f(x)|\,\chi_{J}(x)=\left|\int\frac{f(y)}{x-y}\,dy\right|\chi_{J}(x)\geq(2|I|)^{-1}\int_{I}|f|\,\chi_{J}(x).
\]
则
\begin{align*}
\left(\int_{J}\omega\right)|f_{I}|^{p}
&\leq\int|H f|^{p}\,\omega\,\chi_{J}\leq\int|H f|^{p}\,\omega \\
&\leq C\int|f|^{p}\,\omega=C\int_{I}|f|^{p}\,\omega.
\end{align*}
在特殊情形 $f=\chi_{I}$, $f_{I}=1$, 有
\[
\int_{J}\omega\leq C\int_{I}\omega,
\]
于是利用 $I$ 与 $J$ 的对称性,
\[
\omega(I)\leq C\,\omega(J).
\]
于是对支在 $I$ 中的任意 $f$,
\[
\left(\int_{I}\omega\right)|I|^{-p}\left(\int_{I}\omega^{-\frac{1}{p-1}}\right)^{p}
\leq C\int_{I}\omega^{-\frac{p}{p-1}+1},
\]
等等, 等等.
\end{proof}

\begin{lemma}\label{lem:3.15}
若 $\omega\in A_1$, 则 $\omega\in A_p$.
\end{lemma}

\begin{proof}
\begin{align*}
\left(\frac{1}{|I|}\int_{I}\omega\right)\left(\frac{1}{|I|}\int_{I}\omega^{-\frac{1}{p-1}}\right)^{p-1}
&\leq\omega_{I}\left(\frac{1}{|I|}\int_{I}1\right)\left(\inf_{I}\omega\right)^{-1}
=\frac{\omega_{I}}{\inf_{I}\omega}\leq C. \qedhere
\end{align*}
\end{proof}

\begin{corollary}\label{cor:3.16}
\index{$M_s(g)$ 极大函数 (maximal function $M_s(g)$)}
取 $s>1$, 设 $g\geq 0$ 局部可积; 定义 $M_s(g)=[(g^{s})^{\star}]^{\frac{1}{s}}$. 则 $M_s(g)\in A_1$.
\end{corollary}

\begin{proof}
取定 $I$ 与 $\tilde{I}$; 令 $\omega_1=g^{s}\chi_{\tilde{I}}$ 与 $\omega_2=\omega-\omega_1$. 由于 $\frac{1}{s}<1$, 我们可以用推论~\ref{cor:2.27} 控制 $\omega_1^{\star}$:
\begin{align*}
\frac{1}{|I|}\int\bigl(\omega_{1}^{\star}\bigr)^{\frac{1}{s}}
&\leq C\frac{1}{|I|}\int\bigl(\omega_{1}\bigr)^{\frac{1}{s}} \\
&\leq C\frac{1}{|I|}\int_{I}\bigl(\omega_{1}\bigr)^{\frac{1}{s}}
\leq C\bigl(\omega^{\star}(y)\bigr)^{\frac{1}{s}},
\end{align*}
对任意 $y\in I$ 成立. 对所有 $y\in I$ 取下确界, 得
\[
\left[(\omega_{1}^{\star})^{\frac{1}{s}}\right]_{I}\leq C\left[\inf_{y\in I}\omega^{\star}(y)\right]^{\frac{1}{s}}
=C\inf_{y\in I}\bigl[\omega^{\star}\bigr]^{\frac{1}{s}}.
\]
于是 $\omega_{1}$ 被 $A_{1}$ 条件控制.

为控制 $\omega_2$ 这一部分, 取 $x$, $y\in I$ 并计算
\[
\omega_{2}^{\star}(x)=\sup_{J}\frac{1}{|J|}\int_{J}\omega_{2}.
\]
若它非零, 则区间 $J$ 必须与 $\omega_2$ 的支集相交. 另一方面, $J$ 包含 $x$, 所以 $J$ 的长度至少是 $I$ 的长度. 再者, 由于 $x\in J$, $\tilde{J}$ 包含 $y$, 因此
\begin{align*}
\frac{1}{|J|}\int_{J}\omega_{2}
&\leq 2\,\frac{1}{|\tilde{J}|}\int_{\tilde{J}}\omega_{2}
\leq 2\sup_{J}\frac{1}{|J|}\int_{J}\omega=2\omega^{\star}(y).
\end{align*}
现在对所有包含 $x$ 的 $J$ 取上确界, 得 $\omega_2^{\star}(x)\leq 2\omega^{\star}(y)$. 对所有包含 $y$ 的 $I$ 取下确界, 得
\[
\omega_{2}^{\star}(x)\leq 2\inf_{I}\omega^{\star}(y).
\]
于是
\[
\frac{1}{\left|I\right|}\int_{I}\left(\omega_{2}^{\star}\right)^{\frac{1}{s}}
\leq 2\inf_{I}\left(\omega^{\star}\right)^{\frac{1}{s}}(y). \qedhere
\]
\end{proof}

\section{第 3 章评注}
\label{sec:3.4}

\noindent (3.1):\quad 这个结果归功于 R.~Coifman. 这里的证明取自 Coifman 与 C.~Fefferman 的文章~\cite{coifman-fefferman}; 参见该文的致谢. 这类结果同样可以以类似的证明推广到更高维; 如同极大 Hilbert 变换的情形, 我们使用了 $K$ 的逐点估计, 这些证明并不适用于 H\"ormander 型的核.

\noindent (3.2):\quad 利用 good-$\lambda$ 不等式对 $0<p<\infty$ 得到 $L^{p}$ 估计, 是由 D.~Burkholder 从概率论引入的一种技术, 他对与此相关的结果给出了概率论的证明. 见综述文章~\cite{burkholder}.

\noindent (3.3):\quad $A_p$ 条件由 B.~Muckenhoupt~\cite{muckenhoupt-maximal} 在对极大函数与 Hilbert 变换在加权 $L^p$ 空间上行为的一般研究中发现. 诸如 $|x|^\alpha$ 的特殊情形早已为人所知; 见 Hardy--Littlewood--P\'olya~\cite{hardy-littlewood-polya} 第~IX 章的讨论. 我们的方法与 Muckenhoupt 的不同, 同样取自~\cite{coifman-fefferman}. 关于加权范数不等式的更多问题, 见 Muckenhoupt 的综述文章~\cite{muckenhoupt-classical}, 尽管那里提到的问题已经取得了大量进展. $A_2$ 权的理论在调和函数边值的分析中扮演重要角色; 讨论见 P.~Jones~\cite{jones} 的文章.

\noindent (3.16):\quad 这个结果最初由 A.~C\'ordoba 与 C.~Fefferman~\cite{cordoba-feffermanC-weighted} 证明; 这里的证明取自 R.~Coifman 与 R.~Rochberg~\cite{coifman-rochberg}.

% 第 4 章 Multipliers with Singularities (原书第 75-88 页)
\chapter{带奇性的乘子}
\label{chap:4}

在本章中, 我们将推广第~\ref{chap:2} 章的标准模型. 该模型告诉我们: 若 $\mu=\sum\mu_j\phi_j$, 其中 $\sum|\mu_j|<\infty$, 则 $\mu$ 是所有 $L^p$ 的乘子. 在 Hilbert 变换的情形, $\sum|\mu_j|=\infty$; 然而 Hilbert 变换却是 $1<p<\infty$ 时 $L^p$ 的乘子. 本章的目的是把我们对 Hilbert 变换的理解推广到更一般的乘子.

\section{H\"ormander--Mikhlin 定理}
\label{sec:4.1}
\index{H\"ormander--Mikhlin 定理 (H\"ormander--Mikhlin theorem)|emph}

\begin{proposition}\label{prop:4.1}
设 $\mu\in L^\infty$ 且 $\mu'\in L^1$. 则 $\mu$ 是 $L^p(\mathbb{R}^1)$ ($1<p<\infty$) 的乘子.
\end{proposition}

\begin{proof}
设 $\psi$ 为光滑函数, 在 $[-N,N]$ 上 $\psi=1$, 在 $[-(N+1),(N+1)]$ 之外 $\psi=0$. 令 $\mu_0=\mu\psi$; 我们将证明 ${}_p\|T_0\|\leq C$, 其中 $C$ 与 $N$ 无关. 由一个标准的极限论证, 这蕴含 $\mu$ 是乘子. 照例, 我们使用对偶技巧: 设 $f,g\in\mathcal{S}$, $\|f\|_p=\|g\|_{p'}=1$. 则
\begin{align*}
\int T_{0}fg&=\int\mu_{0}(\xi)\hat{f}(\xi)\hat{g}(\xi)\,d\xi
=\int_{-\infty}^{\infty}\int_{-\infty}^{\xi}\mu_{0}'(\eta)\,d\eta\,\hat{f}(\xi)\hat{g}(\xi)\,d\xi \\
&=\int_{-\infty}^{\infty}\int_{-\infty}^{\infty}\mu_{0}'(\eta)\,\chi_{(-\infty,\xi)}(\eta)\,\hat{f}(\xi)\hat{g}(\xi)\,d\xi\,d\eta \\
&=\int_{-\infty}^{\infty}\int_{-\infty}^{\infty}\chi_{(\eta,\infty)}(\xi)\,\hat{f}(\xi)\hat{g}(\xi)\,d\xi\,\mu_{0}'(\eta)\,d\eta.
\end{align*}
分部积分中不出现边界项, 因为 $\mu_0$ 在 $\infty$ 处消失; 由于 $\mu_0,\mu_0'$ 属于 $L^1$ 且 $\hat{f},\hat{g}$ 属于 $\mathcal{S}$, 交换积分的技术细节是容易的. 现在设 $T_\eta$ 是以
\[
\chi_{(\eta,\infty)}(\xi)=\frac{1}{2}\bigl(1-i\,\mathrm{sign}(\xi-\eta)\bigr)
\]
为乘子的算子, 于是 $T_{\eta}=\frac{1}{2}\bigl(I-iH_{\eta}\bigr)$, 其中 $H_{\eta}$ 是 Hilbert 变换的一个平移. 由此诸 $T_{\eta}$ 在 $L^{p}$ 上一致有界, 于是我们可以这样完成估计:
\begin{align*}
\int T_{0}fg&=\int_{-\infty}^{\infty}T_{\eta}f(x)\,g(x)\,dx\,\mu_{0}'(\eta)\,d\eta \\
&\leq\int\|T_{\eta}\|_{p}\,\|g\|_{p'}\,|\mu_{0}'(\eta)|\,d\eta
\leq C\|f\|_{p}\,\|g\|_{p'}\int|\mu_{0}'(\eta)|\,d\eta.
\end{align*}
为完成证明, 注意
\[
\|\mu_{0}'\|_{1}\leq\|\mu'\psi\|_{1}+\|\mu\psi'\|_{1}\leq C\|\mu'\|_{1}+\|\mu\|_{\infty}\|\psi'\|_{1}.
\]
我们只须选取 $\psi$ 使 $\|\psi'\|_{1}\leq C$, 其中 $C$ 与 $N$ 无关. 这是可能的, 因为 $\psi'$ 仅在 $(-(N+1),-N)$ 与 $(N,N+1)$ 上非零.
\end{proof}

\begin{remark}\label{rem:4.2}
我们稍稍改变了视角; 我们不再考察 $\sum\mu_j\psi_j$ (其中诸 $\psi_j$ 紧支集), 而是考察 Hilbert 变换的加权平均. 这两种看法是一样的, 正如积分是和的连续版本.

牢记这一点, 条件 $\mu'\in L^1$ 便类似于条件 $\sum|\mu_j|<\infty$. 简言之, 除了记号之外, 该命题并没有真正使我们跳出标准模型.

我们真正想知道的, 是条件``$\mu=\sum\psi_j$ 而前面没有收敛因子''的连续类比. 这不难看出; 直观上, $\psi_j'=2^j$ 于区间 $[2^{-(j+1)},2^{-j}]$ 上. 用 $\mu$ 的语言, $\mu'(\xi)=\frac{1}{\xi}$. 正确的乘子条件是 $\xi\mu'(\xi)\in L^{\infty}$.
\end{remark}

\begin{theorem}\label{thm:4.3}
设 $\mu\in L^{\infty}$ 且 $\mu$ 在 $0$ 附近之外是 $C^{1}$ 的. 若 $\xi\mu'(\xi)\in L^{\infty}$, 则 $\mu$ 是 $L^{p}$ ($1<p<\infty$) 的乘子.
\end{theorem}

\begin{proof}
我们打算应用 Hilbert 变换理论; 我们需要把二进分解 $\mu=\sum\psi_j$ 转移到对卷积核的估计上. 为此, 对 $\mu$ 的分解必须光滑地进行, 使得逆 Fourier 变换把 $\mathcal{S}$ 函数映到 $\mathcal{S}$ 函数. 技术细节: 设 $\varphi\geq 0$ 光滑, 且满足
\[
\varphi(\xi)=1\quad\text{若}\ \frac{3}{4}<\xi<\frac{3}{2},
\]
\[
\varphi(\xi)=0\quad\text{若}\ \xi<\frac{1}{2}\ \text{或}\ \xi>2.
\]
令
\[
\psi_{j}(\xi)=\frac{\varphi\bigl(2^{-j}|\xi|\bigr)}{\sum_{k=-\infty}^{\infty}\varphi\bigl(2^{-k}|\xi|\bigr)}.
\]
再令
\[
\psi_{j}(\xi)=\psi_{0}\bigl(2^{-j}\xi\bigr);
\]
设 $\mu_j=\psi_j\mu$, 并定义乘子三元组 $(T,K,\mu)$ 与 $(T_j,K_j,\mu_j)$. 我们的希望是应用 Hilbert 变换的思想, 于是我们希望
\[
\int_{|x|>2|y|}|K_{j}(x-y)-K_{j}(x)|\,dx\leq C2^{-j}.
\]
不确定性原理阻止了这一点; 由于 $\mu_{j}$ 支在长度为 $2^{j}$ 的区间上, $K_{j}$ 直观上集中在一个长度为 $2^{-j}$ 的区间上. 因此, 若 $|y|>2^{-j}$, 则 $K_{j}(x-y)$ 绝不会看起来像 $K_{j}(x)$, 从而我们无法利用光滑性带来的控制. 我们唯一的希望是 $K_{j}$ 满足估计:
\begin{align*}
\int_{|x|>1}|K_{j}(x)|\,dx
&\leq C\left(\int_{|x|>1}|x|^{2}|K_{j}(x)|\,dx\right)^{\frac{1}{2}}
\leq C\|xK_{j}\|_{2}=C\Bigl\|\frac{d}{d\xi}\hat{K}_{j}\Bigr\|_{2}=C\|\mu_{j}'\|_{2}.
\end{align*}
但是
\begin{align*}
\|\left(\psi_{j}\mu\right)^{\prime}\|_{2}&\leq\|\psi_{j}\mu^{\prime}\|_{2}+\|\psi_{j}^{\prime}\mu\|_{2}, \\
\int\bigl|D^{1-k}\psi_{j}\,D^{k}\mu\bigr|^{2}
&=\int_{\operatorname{supp}\psi_{j}}\left|2^{-j(1-k)}\psi_{0}^{(1-k)}\bigl(2^{-j}|\xi|\bigr)\,D^{k}\mu\right|^{2} \\
&\leq C\,2^{-2(1-k)j}\int_{\operatorname{supp}\psi_{j}}\bigl|D^{k}\mu\bigr|^{2}
\leq C\,2^{-2(1-k)j}\int_{2^{j-1}}^{2^{j+1}}|\xi|^{-2k}\,d\xi
=C\bigl(2^{-\frac{j}{2}}\bigr)^{2},
\end{align*}
因此 $\|xK_j\|_2\leq C2^{-\frac{j}{2}}$. 当无法使用光滑性时, 我们将应用这类估计; 于是对 $2^j|y|\geq 1$,
\[
\int_{|x|>2|y|}|K_{j}(x-y)-K_{j}(x)|\,dx
\leq\int_{|x|>2|y|}|K_{j}(x-y)|\,dx+\int_{|x|>2|y|}|K_{j}(x)|\,dx.
\]
在第一个积分中, 变量替换把我们带到集合
\[
\{x \mid |x+y|>2|y|\}\subset\{|x|+|y|>2|y|\}=\{|x|>|y|\},
\]
而两个积分都被
\begin{align*}
\int_{|x|>|y|}|K_{j}(x)|\,dx
&\leq\left(\int_{|x|>|y|}|x|^{-2}\,dx\right)^{\frac{1}{2}}\|xK_{j}\|_{2} \\
&\leq C|y|^{-\frac{1}{2}}\,2^{-\frac{j}{2}}=C\bigl(2^{j}|y|\bigr)^{-\frac{1}{2}}
\end{align*}
控制. 接着, 在区域 $|y|<2^{-j}$ 中, 我们预计可以使用光滑性. 若 $|y|<2^{-j}$ 且 $|x|>2|y|$, 则
\begin{align*}
\int_{\{|x|>2|y|\}}|K_{j}(x-y)-K_{j}(x)|\,dx
&\leq\left(\int_{|x|\leq 2^{-j}}1\,dx\right)^{\frac{1}{2}}\|K_{j}(x-y)-K_{j}(x)\|_{2} \\
&=2^{-\frac{j}{2}}\bigl\|\bigl(e^{2\pi iy\xi}-1\bigr)\psi_{j}(\xi)\mu(\xi)\bigr\|_{2}
\leq C2^{-\frac{j}{2}}\|\xi\psi_{j}\mu\|_{2}\,|y|.
\end{align*}
在 $\psi_j$ 的支集上, $|\xi|\leq 2^j$ 且 $|\psi_j\mu|\leq C$, 所以这一项被
\[
C2^{-\frac{j}{2}}\,2^{j}|y|\,|\operatorname{supp}\psi_{j}|^{\frac{1}{2}}=C\bigl(2^{j}|y|\bigr)
\]
控制. 最后, 若 $|x|>2^{-j}$,
\begin{align*}
\int_{|x|>2^{-j}}\left|K_{j}(x-y)-K_{j}(x)\right|dx
&\leq\left(\int_{|x|>2^{-j}}|x|^{-2}\,dx\right)^{\frac{1}{2}}\bigl\|x\bigl[K_{j}(x-y)-K_{j}(x)\bigr]\bigr\|_{2} \\
&=\;C2^{\frac{j}{2}}\bigl\|\bigl(e^{2\pi iy\xi}-1\bigr)\bigl(\psi_{j}(\xi)\mu(\xi)\bigr)^{\prime}\bigr\|_{2} \\
&\qquad+2^{-\frac{j}{2}}\bigl\|\bigl(e^{2\pi iy\xi}-1\bigr)\psi_{j}(\xi)\mu(\xi)\bigr\|_{2} \\
&\leq C2^{\frac{j}{2}}|y|\,\bigl\|\xi(\psi_{j}\mu)^{\prime}\bigr\|_{2}+C2^{\frac{j}{2}}|y|\,\|\psi_{j}\mu\|_{2}.
\end{align*}
在第一项中, 在 $\psi_j$ 的支集上 $|\xi|\leq 2^j$, 于是
\[
\bigl\|\bigl(\psi_{j}\mu\bigr)^{\prime}\bigr\|_{2}\leq 2^{-\frac{j}{2}}.
\]
在第二项中,
\[
\|\psi_{j}\mu\|_{2}\leq\|\mu\|_{\infty}\|\psi_{j}\|_{2}=C2^{\frac{j}{2}},
\]
而且两项都被 $C2^{j}|y|$ 控制.

为完成证明, 设 $\mu_N=\sum_{j=-N}^{N}\psi_j\mu$, 并设 $(T_N,K_N,\mu_N)$ 为相应的乘子三元组. 由于
\[
\sum_{j=-N}^{N}\psi_{j}\leq 1,\qquad\|\mu_{N}\|_{\infty}\leq\|\mu\|_{\infty},
\]
而且我们有主要估计
\[
\int_{|x|>2|y|}|K_{N}(x-y)-K_{N}(x)|\,dx
\leq C\sum_{2^{j}|y|\leq 1}2^{j}|y|+C\sum_{2^{j}|y|>1}\bigl(2^{j}|y|\bigr)^{-\frac{1}{2}}.
\]
这两个和关于 $y$ 一致有界; 例如若 $|y|>1$, 则第一个和中的指数 $j$ 全是负的, 其贡献为
\[
C|y|\sum_{j\leq-\log|y|}2^{j}=C|y|\,|y|^{-1},
\]
而在第二项中,
\[
|y|^{-\frac{1}{2}}\sum_{j\geq-\log|y|}2^{-\frac{j}{2}}
\leq C\bigl(1+|y|^{-\frac{1}{2}}\bigr).
\]
但 $|y|^{-\frac{1}{2}}\leq 1$. 对 $|y|\leq 1$ 有类似的估计.

现在只需对 $(T_N,K_N,\mu_N)$ 应用定理~\ref{thm:2.32}, 得到它们在 $L^p$ 上有界, 且界与 $N$ 无关. 一个快速的极限论证告诉我们如何把算子 $T$ 扩张到整个 $L^p$: 若 $f\in\mathcal{S}$, 则 $\sum\psi_j\leq 1$ 且 $\mu\in L^\infty$, 于是 $\mu_N\hat{f}$ 控制收敛且在 $L^2$ 中收敛到 $\mu\hat{f}$. 算子 $T_N$ 一致有界且在一个稠密集上收敛到 $T$, 故 $T$ 也有界. 在 $\mathcal{S}$ 上, $T$ 由 $\mu$ 给出.
\end{proof}

\begin{remark}\label{rem:4.4}
\begin{enumerate}[label=\alph*)]
\item Hilbert 变换及其在命题~\ref{prop:4.1} 中的简单变体, 证明了具有简单跳跃间断的乘子的有界性. 我们新的、改进了的结果涵盖了一些振荡型的间断; 例如 $\mu(\xi)=|\xi|^{it}$ 满足定理的假设.
\item 有时, 检验较弱的条件更容易:
\[
\sup_{R>0}\frac{1}{R^{1+2j}}\int_{R<|\xi|<2R}\bigl|D^{j}\mu\bigr|^{2}\leq C,\qquad j=0,\ 1.
\]
在定理~\ref{thm:4.3} 的证明中, 这实际上就是我们真正需要的全部.
\end{enumerate}
\end{remark}

\essayhead{ess:4.5}{关于标准模型的注记}

我们刚完成的证明只在 $f_j$ 的支集上估计 $D^j\mu$, 而我们取后者为一个二进区间. 在这个意义上, 我们把 $\mu_j=\mu\psi_j$ 视为彼此独立地行动的乘子. 在本章的其余部分, 我们将把这一想法严格化. 如同上面的证明, 我们将给出三个版本. 首先考察没有光滑性的乘子, 比如 Hilbert 变换:
\[
\widehat{\left(P_{j}f\right)}(\xi)=\chi_{(2^{j},2^{j+1})}(\xi)\,\hat{f}(\xi).
\]
接下来我们证明光滑版本, 其中 $\varphi\in C^\infty$ 且 $\operatorname{supp}\varphi\subset[1,2]$; 此时
\[
\widehat{\left(T_{j}f\right)}(\xi)=\varphi(2^{-j}\xi)\,\hat{f}(\xi).
\]
最后是乘子版本:
\[
\widehat{\left(T_{j}f\right)}(\xi)=\mu(\xi)\,\varphi(2^{-j}\xi)\,\hat{f}(\xi),
\]
其中 $\mu$ 满足与~\ref*{thm:4.3} 中类似的估计. 与之前一样, 我们将用不连续的 $P_j$ 作为证明乘子定理的技术工具; 这最终将消除对 $K$ 的核估计, 使我们能完全集中在 $\mu$ 上. 最终的好处是我们可以完全用标准模型中的想法工作, 而不必在 $T$ 与 $\mu$ 的估计之间来回切换.

这个纲领基于对``诸 $\mu_j$ 算子彼此独立地行动''的含义的理解. 若转而考察以 $\chi_{(2^j,2^{j+1})}$ 为乘子的算子 $P_j$, 独立性表现为这些算子是正交投影这一事实:
\[
\|f\|_{2}^{2}=\|\hat{f}\|_{2}^{2}=\sum_{j=-\infty}^{\infty}\int_{2^{j}}^{2^{j+1}}|\hat{f}|^{2}
=\sum\int|P_{j}f|^{2}=\int\Bigl(\sum|P_{j}f|^{2}\Bigr).
\]
用 $L^{2}$ 估计的语言, 独立性的正确表达是:
\[
\|f\|_{2}=\Bigl\|\Bigl(\sum|P_{j}f|^{2}\Bigr)^{\frac{1}{2}}\Bigr\|_{2}.
\]
现在我们必须猜测``独立性''或``正交性''在 $L^p$ 上的类比; 我们的出发点是量 $\bigl(\sum|P_jf|^2\bigr)^{\frac{1}{2}}$, 只要并非所有的 $|P_jf|$ 同时很大, 它就像 $|P_jf|$ 一样. 定理将给出: 这就是正交性的正确观念:
\[
\Bigl\|\Bigl(\sum|P_{j}f|^{2}\Bigr)^{\frac{1}{2}}\Bigr\|_{p}\leq C_{p}\|f\|_{p}.
\]

\begin{lemma}\label{lem:4.6}
设
\[
\Bigl\|\Bigl(\sum|P_{j}f|^{2}\Bigr)^{\frac{1}{2}}\Bigr\|_{p}\leq C_{p}\|f\|_{p}
\]
对所有的 $f\in L^{p}$ 成立. 则对所有的 $g\in L^{p'}$,
\[
\|g\|_{p'}\leq C_{p}\Bigl\|\Bigl(\sum|P_{j}g|^{2}\Bigr)^{\frac{1}{2}}\Bigr\|_{p'}.
\]
\end{lemma}

\begin{proof}
\begin{align*}
\int fg&=\int\hat{f}\,\check{g}=\sum\int\widehat{P_{j}f}\,\widehat{P_{j}g} \\
&=\int\sum\bigl(P_{j} f\bigr)\bigl(P_{j} g\bigr) \\
&\leq\int\Bigl(\sum|P_{j} f|^{2}\Bigr)^{\frac{1}{2}}\Bigl(\sum|P_{j} g|^{2}\Bigr)^{\frac{1}{2}} \\
&\leq\Bigl\|\Bigl(\sum|P_{j} f|^{2}\Bigr)^{\frac{1}{2}}\Bigr\|_{p}\Bigl\|\Bigl(\sum|P_{j} g|^{2}\Bigr)^{\frac{1}{2}}\Bigr\|_{p'} \\
&\leq C_{p}\|f\|_{p}\Bigl\|\Bigl(\sum|P_{j} g|^{2}\Bigr)^{\frac{1}{2}}\Bigr\|_{p'}.
\end{align*}
现在对所有满足 $\|f\|_p=1$ 的 $f$ 取上确界, 即得结论.
\end{proof}

\begin{remark}\label{rem:4.7}
\index{独立性 (independence)}
这个结果的全部要点在于下述推论: 当我们对 $1<p<\infty$ 证明了
\[
\Bigl\|\Bigl(\sum|P_{j}f|^{2}\Bigr)^{\frac{1}{2}}\Bigr\|_{p}\leq C_{p}\|f\|_{p}
\]
时, 我们自动得到
\[
\|f\|_{p}\leq C_{p}\Bigl\|\Bigl(\sum|P_{j}f|^{2}\Bigr)^{\frac{1}{2}}\Bigr\|_{p}\leq C_{p}\|f\|_{p}.
\]
为了比较 $f$ 与 $\bigl(\sum|P_jf|^2\bigr)^{1/2}$ 的 $L^p$ 范数, 我们需要几个新想法; 其复杂性来自这样一个事实: 我们是在为 $p\neq 2$ 的 Banach 空间 $L^p$ 定义一种``正交性''的观念, 而在那里正交性根本没有任何意义. 正确的想法来自概率论, 称为独立性 (independence). 若 $f$, $g$ 是可测函数, 我们说它们独立, 如果对一切可测集 $E$ 有 $\int_Efg=\int_Ef\int_Eg$; 其想法是: $f$ 的大小必须在逐点意义上影响 $g$ 的大小, 但在小集合上取平均时, 两个函数互不影响. 若用二项式定理把一个和展开到非整数次幂, 你就必须估计的正是这类对象.
\end{remark}

\section{Littlewood--Paley 与 Marcinkiewicz 定理}
\label{sec:4.2}

\begin{definition}\label{def:4.8}
\index{Rademacher 函数 (Rademacher functions)}
Rademacher 函数 $r_{m}(t)$ 由以下要求定义: 它们以 $2$ 为周期, 且
\begin{gather*}
r_{0}(t)=r_{0}(t+1); \\
r_{0}(t)=1\quad\text{若}\ 0\leq t\leq\frac{1}{2}; \\
r_{0}(t)=-1\quad\text{若}\ \frac{1}{2}<t\leq 1; \\
r_{m}(t)=r_{0}(2^{m}t).
\end{gather*}
\end{definition}

\begin{theorem}\label{thm:4.9}
Rademacher 函数在 $[0,1]$ 上独立; 也就是说, 对 $m_1,m_2,\ldots,m_k$ 与 $\alpha_1,\allowbreak\alpha_2,\ldots,\alpha_k$,
\[
|\{t \mid r_{m_{i}}(t)\leq\alpha_{i}\ \text{对所有的}\ i\}|=\prod_{i=1}^{k}|\{t \mid r_{m_{i}}(t)\leq\alpha_{i}\}|,
\]
而且对可测函数 $f_i$ 与任意可测集 $E$,
\[
\int_{E}\prod_{i}f_{i}\bigl(r_{m_{i}}(t)\bigr)\,dt=\prod_{i}\int_{E}f_{i}\bigl(r_{m_{i}}(t)\bigr)\,dt.
\]
\end{theorem}

\begin{proof}
关于积分的结果源于关于阶梯函数的类似结果, 后者源于关于 Borel 集的结果, 再源于关于区间的结果, 而后者又源于关于半无穷区间的结果, 我们就来证明它. 我们只处理 $(-\infty,\alpha_i]$ 的情形.

若对某个 $i$ 有 $\alpha_i<-1$, 则两边都有一些集合是空的, 乘积为零. 若对某个 $i$ 有 $\alpha_i>1$, 则条件 $r_{m_i}(t)\leq\alpha_i$ 总被满足, 于是它在左端没有贡献, 而在右端贡献因子 $1$. 类似地, 条件 $r_m(t)\leq\alpha_m<1$ 等价于 $r_m(t)\leq-1$, 即 $r_m(t)=-1$. 于是, 定理由等式
\[
|\{t \mid r_{m_{i}}(t)=-1\ \text{对所有的}\ i\}|
=\prod_{i=1}^{k}|\{t \mid r_{m_{i}}(t)=-1\}|=\left(\frac{1}{2}\right)^{k}
\]
得出.

这样的计算是一个归纳法, 其中除最后一步外都被跳过了. 我们重新排序使 $m_1<m_2<\cdots<m_k$. 归纳假设是: $\{t \mid r_{m_i}(t)=-1,\ 1\leq i\leq k-1\}$ 是二进区间的无交并, 每个的长度为 $2^{-m_{k-1}}$, 而该集合的总测度为 $\bigl(\frac{1}{2}\bigr)^{k-1}$. 现在再加上条件 $r_{m_k}(t)=-1$. 由于 $m_k>m_{k-1}$, $r_{m_k}$ 在这些二进区间的每一个上恰在一半的时间里取负值 (为什么? 你必须用到 $r_{m_k}$ 是 $r_{m_{k-1}}$ 的伸缩这一事实).
\end{proof}

\begin{theorem}\label{thm:4.10}
设 $F(t)=\sum_{k=0}^{N}a_kr_k(t)$. 若 $\sum a_k^2=1$, 则
\[
|\{t \mid |F(t)|>\lambda\}|\leq 2e^{-\frac{\lambda^{2}}{4}}.
\]
\end{theorem}

\begin{proof}
\begin{align*}
e^{\frac{\lambda^{2}}{2}}\left|\left\{t \mid |F(t)|>\lambda\right\}\right|
&\leq\int_{|F|>\lambda}e^{\frac{\lambda}{2}|F(t)|}\,dt
\leq\int_{0}^{1}e^{\frac{\lambda}{2}|F|}\,dt \\
&\leq\int_{0}^{1}\left(e^{\frac{\lambda}{2}F}+e^{-\frac{\lambda}{2}F}\right)dt.
\end{align*}
但是
\begin{align*}
\int_{0}^{1}e^{-\frac{\lambda}{2}F}\,dt
&=\int_{0}^{1}\prod_{k}e^{-\frac{\lambda}{2}a_{k}r_{k}}\,dt
=\prod_{k}\int_{0}^{1}e^{-\frac{\lambda}{2}a_{k}r_{k}}\,dt
=\prod_{k}2^{-k}\int_{0}^{2^{k}}e^{-\frac{\lambda}{2}a_{k}r_{0}}\,dt \\
&=\prod_{k=1}^{N}\left[\sum_{n=0}^{2^{k-1}}\left(\int_{n}^{n+\frac{1}{2}}e^{-\frac{\lambda}{2}a_{k}}\,dt+\int_{n+\frac{1}{2}}^{n+1}e^{\frac{\lambda}{2}a_{k}}\,dt\right)\right] \\
&=\prod_{k=1}^{N}2^{-k}\left(\frac{e^{\frac{\lambda}{2}a_{k}}+e^{-\frac{\lambda}{2}a_{k}}}{2}\right)\sum_{n=0}^{2^{k-1}}1 \\
&=\prod_{k=1}^{N}\cosh\Bigl(\frac{\lambda}{2}a_{k}\Bigr)
\leq\prod_{k=1}^{N}e^{\frac{\lambda^{2}}{4}a_{k}^{2}}
=e^{\frac{\lambda^{2}}{4}\sum a_{k}^{2}}=e^{\frac{\lambda^{2}}{4}}.
\end{align*}
对另一项有类似的估计, 这带来了因子 $2$.
\end{proof}

\begin{corollary}\label{cor:4.11}
设 $F(t)=\sum a_kr_k(t)$. 若 $1<p<\infty$, 则
\[
\|F\|_{p}\leq C_{p}\|F\|_{2}\leq C_{p}'\|F\|_{p}.
\]
\end{corollary}

\begin{proof}
令
\[
G(t)=\|F\|_{2}^{-1}F(t)=\left(\sum a_{k}^{2}\right)^{-\frac{1}{2}}F(t).
\]
则 $G=\sum b_kr_k$, 其中
\[
\|G\|_{2}=\Bigl(\sum b_{k}^{2}\Bigr)^{\frac{1}{2}}=1.
\]
因此,
\begin{align*}
\|G\|_{p}^{p}&=p\int_{0}^{\infty}\lambda^{p-1}|\{|G|>\lambda\}|\,d\lambda
\leq 2p\int_{0}^{\infty}\lambda^{p-1}e^{-\frac{\lambda^{2}}{4}}\,d\lambda=C_{p},
\end{align*}
于是对 $1\leq p<\infty$, $\|F\|_p\leq C_p\|F\|_2$. 重复引理~\ref{lem:4.6} 的证明, 我们得到 $\|F\|_2^2\leq\|F\|_p\|F\|_r$, 其中 $\frac{1}{2}=\frac{1}{2r}+\frac{1}{2p}$. 若 $p>1$, 则 $r<\infty$, 于是 $\|F\|_r\leq C_p\|F\|_2$.
\end{proof}

\noindent\textbf{注:} \index{向量值不等式 (vector-valued inequalities)}这就是独立性的要点所在: $L^p$ 范数与 $L^2$ 范数等价. $L^2$ 是 Hilbert 空间, 在那里我们可以使用正交性; 该推论使我们能把它转移到 $L^p$. 原始的证明是用二项式定理把和展开成幂, 再利用独立性把乘积的积分化成积分的乘积. 通常这只能在 $L^2$ 上借助那里的正交性完成; 独立性提供了一个替代品.

\begin{proposition}\label{prop:4.12}
设 $\{I_{j}\}$ 是区间的任意集合族, 并设
\[
\widehat{\left(S_{j}f\right)}(\xi)=\chi_{I_{j}}(\xi)\,\hat{f}(\xi).
\]
则对 $1<p<\infty$,
\[
\Bigl\|\Bigl(\sum|S_{j}f_{j}|^{2}\Bigr)^{\frac{1}{2}}\Bigr\|_{p}\leq C_{p}\Bigl\|\Bigl(\sum|f_{j}|^{2}\Bigr)^{\frac{1}{2}}\Bigr\|_{p}.
\]
\end{proposition}

\begin{proof}
由于区间的特征函数可以由 Hilbert 变换构造出来, 我们先对它们证明一个类似的结果. 若 $p>2$, 令 $r=\bigl(\frac{p}{2}\bigr)'$, 并任取 $g\in L^r$. 则对几乎所有的 $x$,
\begin{align*}
|g(x)|&=\lim_{|I|\to0}\left(\frac{1}{|I|}\int_{I}|g(y)|^{s}\,dy\right)^{\frac{1}{s}}
\leq\sup_{I}\left(\frac{1}{|I|}\int_{I}|g(y)|^{s}\,dy\right)^{\frac{1}{s}}=M_{s}(g).
\end{align*}
但定理~\ref{cor:3.16} 告诉我们 $M_{s}(g)$ 属于 $A_{1}$, 而引理~\ref{lem:3.15} 告诉我们 $M_{s}(g)$ 属于 $A_{2}$, 于是定理~\ref{thm:3.13} 与命题~\ref{prop:3.14} 表明
\[
\int|H f|^{2}\,M_{s}(g)\leq C\int|f|^{2}\,M_{s}(g).
\]
因此,
\begin{align*}
\int\Bigl(\sum|H f_{j}|^{2}\Bigr)g
&\leq\int\Bigl(\sum|H f_{j}|^{2}\Bigr)M_{s}(g)
\leq C\int\Bigl(\sum|f_{j}|^{2}\Bigr)M_{s}(g) \\
&\leq C\Bigl\|\sum|f_{j}|^{2}\Bigr\|_{\frac{p}{2}}\bigl\|M_{s}(g)\bigr\|_{(\frac{p}{2})'}
=C\Bigl\|\Bigl(\sum|f_{j}|^{2}\Bigr)^{\frac{1}{2}}\Bigr\|_{p}^{2}\,\|M_{s}(g)\|_{r}.
\end{align*}
但 $r>1$, 于是我们可以找到 $s>1$ 使 $\frac{r}{s}>1$, 而推论~\ref{cor:2.25} 给出估计
\begin{align*}
\|M_{s}(g)\|_{r}^{r}&=\int\bigl[(|g|^{s})^{\star}\bigr]^{\frac{r}{s}}
\leq C\int\bigl[|g|^{s}\bigr]^{\frac{r}{s}}=C\int|g|^{r}=C\|g\|_{r}^{r}.
\end{align*}
总之,
\[
\int\Bigl(\sum|H f_{j}|^{2}\Bigr)g\leq C\Bigl\|\Bigl(\sum|f_{j}|^{2}\Bigr)^{\frac{1}{2}}\Bigr\|_{p}^{2}\|g\|_{r}.
\]
若对所有满足 $\|g\|_r=1$ 的 $g$ 取上确界, 我们便在左端得到 $r'=\frac{p}{2}$ 范数, 而这正是
\[
\left(\int\Bigl(\sum|H f_{j}|^{2}\Bigr)^{\frac{p}{2}}\right)^{\frac{2}{p}}
=\Bigl\|\Bigl(\sum|H f_{j}|^{2}\Bigr)^{\frac{1}{2}}\Bigr\|_{p}^{2}.
\]
在 $p=2$ 的情形,
\begin{align*}
\Bigl\|\Bigl(\sum|H f_{j}|^{2}\Bigr)^{\frac{1}{2}}\Bigr\|_{2}^{2}
&=\int\sum|H f_{j}|^{2}=\int\sum|f_{j}|^{2}
=\Bigl\|\Bigl(\sum|f_{j}|^{2}\Bigr)^{\frac{1}{2}}\Bigr\|_{2}^{2},
\end{align*}
因为 $H$ 是 $L^{2}$ 上的等距算子. $1<p<2$ 的情形用对偶性完成, 正如引理~\ref{lem:4.6} 的证明那样.

接下来我们需要从 $H$ 过渡到区间. 设 $I_{j}=[a_{j},b_{j}]$; 则
\[
S_{j}=\frac{1}{4}\bigl(I+iH_{a_{j}}\bigr)\bigl(I-iH_{-b_{j}}\bigr),
\]
正如引理~\ref{lem:2.12} 中那样. 使用定理~\ref{thm:1.23}d) 的记号, $H_{c}f=J(Hf)J^{-1}$, 其中 $Jf=e^{2\pi icx}f$. 由此, $\bigl(I+iH_{a_{j}}\bigr)\bigl(I-iH_{-b_{j}}\bigr)$ 可以由四个算子控制, 其中每一个都由 $H$ 或 $I$ 控制.
\end{proof}

\begin{theorem}\label{thm:4.13}
\index{Littlewood--Paley 分解 (Littlewood--Paley decomposition)}
设 $\Delta_j=(2^j,2^{j+1})\cup(-2^{j+1},-2^j)$, $-\infty<j<\infty$. 设 $P_j$ 是以 $\chi_{\Delta_j}$ 为乘子的乘子算子. 则对 $1<p<\infty$,
\[
\|f\|_{p}\leq C_{p}\Bigl\|\Bigl(\sum|P_{j}f|^{2}\Bigr)^{\frac{1}{2}}\Bigr\|_{p}\leq C_{p}'\|f\|_{p}.
\]
\end{theorem}

\begin{proof}
在 $p=2$ 处成立的是精确等式, 于是引理~\ref{lem:4.6} 告诉我们, 只需证明右端的不等式. 我们将使用给出有界乘子的 $\chi_{(-\infty,0)}$, 分别构造 $(2^j,2^{j+1})$ 及其负部. 我们的策略是用命题~\ref{prop:4.12} 比较
\[
\Bigl\|\Bigl(\sum|P_{j}f|^{2}\Bigr)^{\frac{1}{2}}\Bigr\|_{p}
\]
与
\[
\Bigl\|\Bigl(\sum|T_{j}f|^{2}\Bigr)^{\frac{1}{2}}\Bigr\|_{p},
\]
其中诸 $T_{j}$ 是光滑的. 然后我们将用推论~\ref{cor:4.11} 比较
\[
\Bigl\|\Bigl(\sum|T_{j}f|^{2}\Bigr)^{\frac{1}{2}}\Bigr\|_{p}
\]
与 $\bigl\|\sum T_jf\bigr\|_p$. 最后, 我们将利用 Rademacher 函数把两者拼合起来. 设
\begin{gather*}
\phi(\xi)=1\quad\text{对}\ 1\leq\xi\leq 2; \\
\phi(\xi)=0\quad\text{对}\ \xi\notin\Bigl[\frac{1}{2},\frac{5}{2}\Bigr]; \\
\phi_{j}(\xi)=\phi(2^{j}\xi).
\end{gather*}
则 $\chi_{\Delta_j}=\phi_j\chi_{\Delta_j}$, 且若 $T_j$ 表示以 $\phi_j$ 为乘子的算子, 则 $P_jf=P_jT_jf$, 于是
\begin{align*}
\Bigl\|\Bigl(\sum|P_{j} f|^{2}\Bigr)^{\frac{1}{2}}\Bigr\|_{p}
&=\Bigl\|\Bigl(\sum|P_{j} T_{j} f|^{2}\Bigr)^{\frac{1}{2}}\Bigr\|_{p}
\leq C\Bigl\|\Bigl(\sum|T_{j} f|^{2}\Bigr)^{\frac{1}{2}}\Bigr\|_{p}.
\end{align*}
现在令
\[
\mu_{o}(\xi,t)=\sum_{j\ \text{奇}}\phi_{j}(\xi)\,r_{j}(t);
\]
对 $\mu_{e}$ (偶数指标) 类似. 我们将证明 $\mu_{o},\mu_{e}$ 关于 $t$ 一致地满足定理~\ref{thm:4.3} 的假设. 设 $T_{o}$, $T_{e}$ 为相应的算子. 则由推论~\ref{cor:4.11},
\[
\Bigl(\sum|T_{j}f|^{2}\Bigr)^{\frac{1}{2}}\leq C\Bigl\|\sum T_{j}f(x)\,r_{j}(t)\Bigr\|_{p,t}.
\]
于是
\begin{align*}
\int\Bigl(\sum|T_{j}f|^{2}\Bigr)^{\frac{p}{2}}dx
&\leq C\int_{0}^{1}\int\Bigl|\sum T_{j}f(x)\,r_{j}(t)\Bigr|^{p}\,dx\,dt
=C\int_{0}^{1}\int|T_{o}f(x,t)|^{p}\,dx\,dt \\
&\leq C\int_{0}^{1}\int|f(x)|^{p}\,dx\,dt=C\|f\|_{p}^{p}.
\end{align*}
这就证明了定理; 为得到关于 $\mu_{o}$ 的估计, 注意奇数的诸 $\phi_{j}$ 具有互不相交的支集, 于是
\[
|\mu_{o}(\xi,t)|\leq\sum|\phi_{j}(\xi)|\leq 1.
\]
类似地, 若 $\xi\in\operatorname{support}\phi_{j}$, 则
\begin{align*}
|\xi\,\mu_{0}'(\xi,t)|
&\leq|\xi|\,2^{-j}\bigl|\phi_{0}'\bigl(2^{-j}\xi\bigr)\bigr|
\leq 2^{j+1}\,2^{-j}\,\|\phi_{0}'\|_{\infty}\leq C. \qedhere
\end{align*}
\end{proof}

\begin{theorem}\label{thm:4.14}
\index{Marcinkiewicz 乘子定理 (Marcinkiewicz multiplier theorem)}
设 $\mu\in L^{\infty}$, 且对每个二进区间 $I$,
\[
\int_{I}|\mu^{\prime}|\leq C
\]
关于 $I$ 一致成立. 则 $\mu$ 给出 $L^{p}$ ($1<p<\infty$) 的一个乘子.
\end{theorem}

\begin{proof}
首先注意
\[
\left\|T f\right\|_{p}\leq C\Bigl\|\Bigl(\sum|P_{j}T f|^{2}\Bigr)^{\frac{1}{2}}\Bigr\|_{p}.
\]
若 $I_{j}=(2^{j},2^{j+1})$, 则
\begin{align*}
P_{j}Tf(x)&=\int_{I_{j}}\mu(\xi)\hat{f}(\xi)e^{2\pi ix\xi}\,d\xi \\
&=\int_{I_{j}}\int_{2^{j}}^{\xi}\mu'(\eta)\,d\eta\,\hat{f}(\xi)\chi_{I_{j}}(\xi)\,e^{2\pi ix\xi}\,d\xi
+\mu(2^{j})\int\hat{f}(\xi)\chi_{I_{j}}(\xi)\,e^{2\pi ix\xi}\,d\xi \\
&=\int_{I_{j}}\mu'(\eta)\int_{\eta}^{2^{j+1}}\hat{f}(\xi)\chi_{I_{j}}(\xi)\,e^{2\pi ix\xi}\,d\xi\,d\eta
+\mu(2^{j})P_{j}f(x) \\
&=\int_{I_{j}}\mu'(\eta)\,\frac{(I+iH_{\eta})}{2}\,P_{j}f(x)\,d\eta+\mu(2^{j})P_{j}f(x).
\end{align*}
由于 $|\mu(2^j)|\leq\|\mu\|_\infty$, 我们可以把第二项分离出来, 然后使用估计
\[
\Bigl\|\Bigl(\sum|P_{j}f|^{2}\Bigr)^{\frac{1}{2}}\Bigr\|_{p}\leq C\|f\|_{p}.
\]
由于
\[
\int_{I}\mu'(\eta)\,\frac{1}{2}I\bigl(P_{j}f(x)\bigr)\,d\eta
=\frac{1}{2}\left(\int_{I}\mu'\right)P_{j}f(x)\leq C\,P_{j}f(x),
\]
这一项也可以被控制. ``主要的''项总是你最后才能对付的那一项; 在此情形, $|P_jTf|^2$ 中的主要项是
\begin{align*}
\left|\int_{I_{j}}\mu'(\eta)\,H_\eta P_{j}f(x)\,d\eta\right|^{2}
&\leq\left(\int_{I_{j}}|\mu'(\eta)|\right)\int_{I_{j}}|\mu'|\,|H_\eta P_{j}f(x)|^{2}\,d\eta.
\end{align*}
如同命题~\ref{prop:4.12} 的证明, 我们取 $g\in L^r$, $r=\bigl(\frac{p}{2}\bigr)'$, 并注意到 $g\leq M_s(g)$. 则
\begin{align*}
\int\sum|P_{j}T f|^{2}\,g(x)\,dx
&\leq 2C\int\sum|P_{j}f|^{2}\,g
+C\sum\int_{I_{j}}|\mu'(\eta)|\int|H_\eta P_{j}f|^{2}\,g(x)\,dx\,d\eta \\
&\leq C\Bigl\|\Bigl(\sum|P_{j}f|^{2}\Bigr)^{\frac{1}{2}}\Bigr\|_{p}^{2}\,\|g\|_{r} \\
&\qquad+C\sum\int_{I_{j}}|\mu'(\eta)|\int|H_\eta P_{j}f|^{2}\,M_{s}(g)\,dx\,d\eta.
\end{align*}
$A_{p}$ 理论适用于诸 $H_{\eta}$ 算子; 界与 $\eta$ 无关:
\[
\int\sum|P_{j}T f|^{2}\,g(x)\,dx
\leq C\|f\|_{p}^{2}\|g\|_{r}+C\Bigl\|\Bigl(\sum|P_{j}f|^{2}\Bigr)^{\frac{1}{2}}\Bigr\|_{p}^{2}\,\|M_{s}g\|_{r}.
\]
既然我们在证明一个乘子定理, 可以假设 $p>2$; 于是存在 $s>1$ 使 $\|M_s(g)\|_r\leq C\|g\|_r$. 总之,
\[
\|T f\|_{p}\leq C\Bigl\|\Bigl(\sum|P_{j}f|^{2}\Bigr)^{\frac{1}{2}}\Bigr\|_{p}\leq C\|f\|_{p}. \qedhere
\]
\end{proof}

\section{评注与推广}
\label{sec:4.3}

\noindent (4.3):\quad 这里陈述的结果与给出的证明取自 H\"ormander~\cite{hormander}; Mikhlin~\cite{mikhlin} 有一个类似的结果. 尽管 $|x|^{-1}$ 在无穷远处不可积, 这些结果在一般维数仍然成立. 这意味着一般的定理需要更多的导数.

\noindent (4.9):\quad 这个结果蕴含 Rademacher 函数是规范正交的, 但它们不构成完备集.

\noindent (4.11):\quad 这里的叙述取自 Stein~\cite{stein-singular}; 原始结果见 Paley~\cite{paley}.

\noindent (4.12):\quad 这里给出的证明仿照了 A.~C\'ordoba 与 C.~Fefferman~\cite{cordoba-feffermanC-weighted} 中的结果, $M_{s}(g)$ 就是在那里引进的. 也有一个使用向量值 Hilbert 变换理论的替代证明; 见 Stein~\cite{stein-singular}.

\noindent (4.13):\quad 这称为 Littlewood--Paley 分解, 得名于圆周上的类似结果. 见 Stein~\cite{stein-square} 的文章.

\noindent (4.14):\quad 这归功于 Marcinkiewicz~\cite{marcinkiewicz-interpolation}; 它在一般维数同样成立. 需要指出的是, 在 $\xi=0$ 处没有任何光滑性条件. 高维的结果在坐标轴上也没有光滑性条件. 这里的证明是 Jodeit 向我们展示的 (未发表).

% 第 5 章 Singularities along Curves (原书第 89-100 页)
\chapter{沿曲线的奇性}
\label{chap:5}

\setcounter{mtthm}{-1}

\section{Bessel 函数的渐近性}
\label{sec:5.1}
\index{Bessel 函数!的渐近性 (asymptotics of Bessel functions)}

在本节中, 我们将计算乘子 $\mu_{\lambda}(\xi)=\bigl(1-|\xi|^2\bigr)_+^{\lambda}$ 的 Fourier 变换 $K_{\lambda}$; 我们的目标是估计
\[
K_{\lambda}(x)=K(x)+E(x);
\]
这里
\begin{align*}
K(x)&=C_{\lambda}|x|^{-\frac{n+1+2\lambda}{2}}\cos\Bigl(2\pi|x|-(n+1)\pi-\frac{\lambda\pi}{2}\Bigr), \\
E(x)&=O\Bigl(|x|^{-\frac{n+8+2\lambda}{2}}\Bigr).
\end{align*}
我们遵循 Stein 与 Weiss 的处理方式, 从 Bessel 函数的性质出发详细推导这些估计.

\begin{fact}\label{fact:5.0}
\[
J_{\lambda}(R)=C_{\lambda}R^{\lambda}\int_{-1}^{1}e^{iRt}\bigl(1-t^{2}\bigr)^{\lambda-\frac{1}{2}}\,dt.
\]
\end{fact}

\begin{proof}
这在命题~\ref{prop:2.3} 中已经讨论过.
\end{proof}

\begin{lemma}\label{lem:5.1}
若 $d\theta$ 表示单位球面 $S^{n-1}$ 上的曲面测度, 则
\[
d\check{\theta}(x)=\int_{S^{n-1}}e^{2\pi ix\xi}\,d\theta(\xi)
=c_{n}\frac{J_{\frac{n-2}{2}}\bigl(2\pi|x|\bigr)}{|x|^{\frac{n-2}{2}}}.
\]
\end{lemma}

\begin{proof}
设 $x=|x|x'$, 其中 $|x'|=1$. 我们用低一维的球面来对球面作参数化; 在 $\mathbb{R}^3$ 中这便是平行于赤道而半径各异的圆. 若取北极固定为 $\bar{n}$, 则顶点在单位球面上、且沿这些圆之一的向量 $\bar{\xi}'$ 由条件 $\bar{n}\cdot\bar{\xi}'=\text{常数}$ 刻画. 极点处当然有 $\bar{n}\cdot\bar{\xi}'=\pm 1$; 赤道处 $\bar{n}\cdot\bar{\xi}'=0$.

北极可以任意选取; 为了简单起见, 我们把 $x\in\mathbb{R}^n$ 视为固定的, 并把北极安排为 $\bar{x}'$. 于是对顶点在球面上的 $\bar{\xi}'$, $\bar{x}'$ 与 $\xi'$ 之间的夹角 $\theta$ 属于 $(-\pi,\pi]$. 若令 $L_\theta=\{\bar{\xi}'\in S^{n-1} \mid \bar{x}'\cdot\bar{\xi}'=\cos\theta\}$, 便得到 $S^{n-1}$ 按 $S^{n-2}$ 球面的一个分解. 这是以 $\bar{x}'$ 为北极的球极坐标. 球面 $L_\theta$ 的半径为 $|\sin\theta|$, 故 $L_\theta$ 的体积为 $c_n|\sin\theta|^{n-2}$. 最后我们可以计算:
\begin{align*}
d\check{\theta}(x)&=\int_{-\pi}^{\pi}\int_{L_{\theta}}e^{2\pi i|x|\,\bar{x}'\cdot\bar{\xi}'}\,d\eta\,d\theta \\
&=\int_{-\pi}^{\pi}e^{2\pi i|x|\cos\theta}\int_{L_{\theta}}1\,d\eta\,d\theta
=\int_{-\pi}^{\pi}e^{2\pi i|x|\cos\theta}|L_{\theta}|\,d\theta
=c_{n}\int_{-\pi}^{\pi}e^{2\pi i|x|\cos\theta}\,|\sin\theta|^{n-2}\,d\theta \\
&=c_{n}\int_{-1}^{1}e^{2\pi i|x|y}\bigl(1-y^{2}\bigr)^{\frac{n-3}{2}}\,dy
=c_{n}\frac{J_{\frac{n-2}{2}}(2\pi|x|)}{|x|^{\frac{n-2}{2}}}. \qedhere
\end{align*}
\end{proof}

\begin{proposition}\label{prop:5.2}
\[
\check{\mu}_{\lambda}(x)=K_{\lambda}(x)
=c_{\lambda}\frac{J_{\frac{n}{2}+\lambda}(2\pi|x|)}{|x|^{\frac{n}{2}+\lambda}}.
\]
\end{proposition}

\begin{proof}
对 Bessel 微分方程应用幂级数技巧, 得
\[
J_{\mu}(t)=\sum_{j=0}^{\infty}\frac{(-1)^{j}\bigl(\frac{t}{2}\bigr)^{\mu+j}}{\Gamma(j+\mu+1)},
\]
从而
\begin{align*}
\int_{0}^{1}J_{\mu}(st)\,s^{\mu+1}\bigl(1-s^{2}\bigr)^{\nu}\,ds
&=\sum_{j=0}^{\infty}\frac{(-1)^{j}\bigl(\frac{t}{2}\bigr)^{\mu+j}}{\Gamma(j+\mu+1)}
\int_{0}^{1}s^{\mu+2j}\bigl(1-s^{2}\bigr)^{\nu}s\,ds \\
&=\sum_{j=0}^{\infty}\frac{(-1)^{j}\bigl(\frac{t}{2}\bigr)^{\mu+j}}{\Gamma(j+\mu+1)}\cdot\frac{1}{2}\int_{0}^{1}r^{\mu+j}(1-r)^{\nu}\,dr.
\end{align*}
回忆 Beta 积分, 这就是
\begin{align*}
\frac{1}{2}\sum_{j=0}^{\infty}&\frac{(-1)^{j}\bigl(\frac{t}{2}\bigr)^{\mu+j}}{j!}\,
\frac{\Gamma(\nu+1)\,\Gamma(\mu+j+1)}{\Gamma(\mu+\nu+j+2)\,\Gamma(\mu+j+1)} \\
&=\frac{\Gamma(\nu+1)\,2^{\nu+1}}{t^{\nu+1}}\sum_{j=0}^{\infty}\frac{(-1)^{j}\bigl(\frac{t}{2}\bigr)^{\mu+\nu+1+j}}{\Gamma(j+\nu+\mu+2)}
=c_{\nu}\frac{J_{\mu+\nu+1}(t)}{t^{\nu+1}}.
\end{align*}
在我们的情形, 要计算的是 $K_{\lambda}(x)$, 它是一个 Fourier 变换. 利用极坐标与 $d\check{\theta}(x)$ 的 Bessel 函数表达式, 我们得到
\begin{align*}
K_{\lambda}(x)&=c_{\lambda}\int_{0}^{1}(1-r^{2})^{\lambda}\,r^{n-1}\,
\frac{J_{\frac{n-2}{2}}(2\pi r|x|)}{(r|x|)^{\frac{n-2}{2}}}\,dr \\
&=c_{\lambda}|x|^{\frac{n-2}{2}}(2\pi|x|)^{-(\lambda+1)}\,J_{\frac{n-2}{2}+\lambda+1}(2\pi|x|).
\end{align*}
\end{proof}

\begin{proposition}\label{prop:5.3}
\[
J_{\lambda}(r)=c_{\lambda}r^{-\frac{1}{2}}\cos\Bigl(r-\frac{\pi\lambda}{2}-\frac{\pi}{4}\Bigr)+O\Bigl(r^{-\frac{3}{2}}\Bigr).
\]
\end{proposition}

\begin{proof}
这个课题中的标准技巧, 是利用表达式
\[
r^{-\lambda}J_{\lambda}(r)=c\int_{-1}^{1}e^{irs}\bigl(1-s^{2}\bigr)^{\lambda-\frac{1}{2}}\,ds
\]
中被积函数的解析性来改变围道; 若做得巧妙, 项 $(1-s^2)^{\lambda-\frac{1}{2}}$ 就会变得像 $s^{\lambda-\frac{1}{2}}$ 一样. 这要好得多, 因为我们可以利用齐次性来求出对 $r$ 的依赖关系. 当然, 指数的相位会随围道而改变, 这正是展开式中那些 $\cos$ 项的来历.

那就开始吧. 我们写
\[
\int_{-1}^{1}e^{irs}\bigl(1-s^{2}\bigr)^{\lambda-\frac{1}{2}}\,ds
=\lim_{\epsilon\to0}\int_{-1+\epsilon}^{1-\epsilon}e^{irs}\bigl(1-s^{2}\bigr)^{\lambda-\frac{1}{2}}\,ds,
\]
这使我们避开了非解析点. 于是 $e^{irz}(1-z^2)^{\lambda-\frac{1}{2}}$ 在单连通区域 $\mathbb{C}-\{(-\infty,-1+\epsilon]\cup[1-\epsilon,\infty)\}$ 上有一个解析分支. 下面我们会稍微马虎地对待这些 $\epsilon$.

我们取一个高为 $a$、底为 $[-1,1]$ 的矩形, 并沿逆时针方向走过其边界. 边界曲线 $\gamma$ 有四段:
\begin{gather*}
\gamma_{1}=\{s+i0 \mid -1\leq s\leq 1\};\qquad \gamma_{2}=\{1+iy \mid 0\leq y\leq a\}; \\
\gamma_{3}=\{s+ia \mid -1\leq s\leq 1\};\qquad \gamma_{4}=\{-1+i(a-y) \mid 0\leq y\leq a\}.
\end{gather*}
则
\[
\int_{\gamma_{1}}=-\int_{\gamma_{2}}-\int_{\gamma_{3}}-\int_{\gamma_{4}},
\]
我们取 $a\to\infty$ 的极限. 此时
\[
\int_{\gamma_{3}}=-\int_{-1}^{1}e^{-ar}e^{irs}\bigl(1-(s+ia)^{2}\bigr)^{\lambda-\frac{1}{2}}\,ds,
\]
指数因子保证了它收敛于零. 对 $\gamma_{2}$,
\begin{align*}
\bigl(1-(1+iy)^{2}\bigr)^{\lambda-\frac{1}{2}}
&=\bigl(y^{2}-2iy\bigr)^{\lambda-\frac{1}{2}}
=y^{\lambda-\frac{1}{2}}(y-2i)^{\lambda-\frac{1}{2}},
\end{align*}
对 $\gamma_{4}$ 类似. 这两个积分合起来给出
\[
\int_{0}^{\infty}e^{ir(1+y)}\,y^{\lambda-\frac{1}{2}}(y-2i)^{\lambda-\frac{1}{2}}\,dy
-\int_{0}^{\infty}e^{ir(-1+y)}\,y^{\lambda-\frac{1}{2}}(y+2i)^{\lambda-\frac{1}{2}}\,dy.
\]
现在我们把 $(y+2i)^{\lambda-\frac{1}{2}}$ 展开为幂级数, 其首项为 $(2i)^{\lambda-\frac{1}{2}}$; 在每个被积函数中它贡献一项形如 $2^{\lambda-\frac{1}{2}}e^{(\lambda-\frac{1}{2})i\pi/2}$ 的量; 两个积分合起来产生
\[
\cos\Bigl(r-\frac{\pi\lambda}{2}-\frac{\pi}{4}\Bigr)\int_{0}^{\infty}e^{-ry}\,y^{\frac{1}{2}+\lambda}\,dy
=\Gamma\bigl(\lambda+\tfrac{1}{2}\bigr)\,r^{-(\frac{1}{2}+\lambda)}\cos\Bigl(r-\frac{\pi\lambda}{2}-\frac{\pi}{4}\Bigr).
\]
如果我们能控制由误差 $(y+2i)^{\lambda-\frac{1}{2}}-(2i)^{\lambda-\frac{1}{2}}$ 产生的项, 证明即告完成. 若 $y\leq 1$, 我们利用光滑性, 均值定理给出这个差被 $Cy$ 控制. 若 $y>1$, 我们提出因子 $y$ 并再次使用均值定理. 于是
\begin{align*}
\int_{0}^{\infty}&E(y)\,e^{-ry}\,y^{\lambda-\frac{1}{2}}\,dy
\leq\int_{0}^{1}e^{-ry}y^{\lambda+\frac{1}{2}}\,dy+\int_{1}^{\infty}e^{-ry}y^{2\lambda-1}\,dy \\
&=r^{-(\lambda+\frac{3}{2})}\int_{0}^{r}e^{-y}y^{\lambda+\frac{1}{2}}\,dy
+r^{-2\lambda}\int_{r}^{\infty}e^{-y}y^{2\lambda+1}\,dy \\
&=r^{-(\lambda+\frac{3}{2})}\,\Gamma\bigl(\lambda+\tfrac{3}{2}\bigr)
+C_{k}\,r^{-2\lambda}\int_{r}^{\infty}y^{-k}y^{2\lambda+1}\,dy
=O\bigl(r^{-(\lambda+\frac{3}{2})}\bigr). \qedhere
\end{align*}
\end{proof}

\begin{corollary}\label{cor:5.4}
当 $\lambda>\frac{n-1}{2}$ 时, $T_\lambda$ 在 $L^p(\mathbb{R}^n)$ ($1\leq p\leq\infty$) 上有界; 而当 $\lambda<\frac{n-1}{2}$ 且 $p$ 不在范围
\[
\frac{2n}{n+1+2\lambda}<p<\frac{2n}{n-1-2\lambda}
\]
之内时, $T_\lambda$ 在 $L^p(\mathbb{R}^n)$ 上无界 (对照图~\ref{fig:2}).
\end{corollary}

\begin{proof}
若 $\lambda>\frac{n-1}{2}$, 则 $K_\lambda\in L^1$. 为得到无界性结果, 我们想取 $f$ 为一个 $\delta$ 函数; 渐近性将会表明
\[
K_{\lambda}\star f\approx|x|^{-\frac{n+1+2\lambda}{2}},
\]
而除非 $\frac{2n}{n+1+2\lambda}<p$, 这不可能属于 $L^p$. 由于 $\delta$ 函数不在 $L^p$ 中, 我们必须更巧妙一些; 关键在于取一个支集非常小的 $f$, 使得卷积核中的振荡可以被抑制. 若把 $f$ 想成紧支集的, 卷积就是把 $f$ 对 $K_\lambda$ 的平移作平均, 于是有时这个平均几乎为零. 我们必须分析 $f$ 与 $K_\lambda$ 的峰相遇的那些位置.

设 $D$ 是半径为 $\frac{1}{100}$ 的球, 并设 $A_{k}$ 表示环形区域
\[
\{x \mid \bigl|2\pi|x|-\lambda\pi-\frac{\pi}{2}-2k\pi\bigr|\leq 10^{-2}\}.
\]
注意 $A_{k}$ 的测度约为 $Ck^{n-1}$, 且若 $x\in A_{k}$, $y\in D$, 则
\[
\cos\Bigl(2\pi|x-y|-(n+1)\pi-\frac{1}{2}\bigl(\frac{n}{2}+\lambda\bigr)\pi\Bigr)\geq\frac{1}{2}.
\]
于是, 若取 $f=\chi_{D}$,
\[
\|T_{\lambda}f\|_{p}\geq\bigl\|(K_{\lambda}\star f)\,\chi_{\bigcup A_{k}}\bigr\|_{p}
-\bigl\|(E\star f)\,\chi_{\bigcup A_{k}}\bigr\|_{p}.
\]
但
\begin{align*}
(K_{\lambda}\star f)(x)\,\chi_{\bigcup A_{k}}(x)
&=\int_{D}\frac{\cos\bigl(2\pi|x-y|+\mathrm{mess}\bigr)}{|x-y|^{\frac{n+1+2\lambda}{2}}}\,dy
\geq C\,k^{-(\frac{n+1+2\lambda}{2})};
\end{align*}
类似地,
\[
|E\star f|\,\chi_{\bigcup A_{k}}\leq C\,k^{-(\frac{n+3+2\lambda}{2})}.
\]
最后,
\[
\|T_{\lambda}f\|_{p}\geq C\left(\sum_{k}k^{-\frac{n+1+2\lambda}{2}}|A_{k}|\right)^{\frac{1}{p}}.
\]
这个级数有限当且仅当 $(n-1)-(n+1+2\lambda)\frac{p}{2}<-1$.
\end{proof}

\begin{remark}\label{rem:5.5}
我们已经看到, 沿曲线的奇性比沿直线的奇性糟糕得多. 单位正方形的特征函数的卷积核是
\[
\frac{\sin x_{1}}{x_{1}}\cdot\frac{\sin x_{2}}{x_{2}},
\]
它对所有 $p>1$ 属于 $L^{p}$. 与之相对, 单位圆盘的特征函数的变换是
\[
\frac{J_{1}(|x|)}{|x|},
\]
它仅在 $p>\frac{4}{3}$ 时属于 $L^p$. 这两个结果的详细证明我们都已看到; 但隐藏在这一区别之下的几何却仍让我们困惑.

进入 $L^p$ 的标准途径是从 $L^1$ 与 $L^2$ 内插; 由于对小的 $\lambda$, $T_\lambda$ 在 $L^1$ 上无界, 这条路走不通. 我们求助于标准模型: 把光滑紧支集函数 $\psi_j$ 扩张到 $\mathbb{R}^n$, 办法是从定义在 $\mathbb{R}^1$ 上的 $\psi_0$ 出发, 令 $\psi_j(\xi)=\psi_0\bigl(2^j(|\xi|-1)\bigr)$. 这些 $\psi_j$ 的支集在一个厚度约 $2^{-j}$、以 $|\xi|=1$ 为中心的环形区域中. (这与一维的图景颇为不同, 但请记住, 我们希望解决的奇性正在 $|\xi|=1$ 处.) 我们预计 $\mu_\lambda\approx 2^{-j\lambda}$, 且 ${}_p\|T_\lambda\|\leq\sum 2^{-j\lambda}\,{}_p\|T_j\|$. 然而与一维情形不同, 这些 $\psi_j$ 并不是同一个函数的伸缩, 因此诸 $T_j$ 的算子范数并不相等. 为了控制这些范数, 我们考察 $L^1$ 与 $L^2$ 上的范数并作内插.

在 $L^2$ 上, ${}_2\|T_j\|=\|\psi_j\|_\infty=1$. 在 $L^1$ 上, ${}_1\|T_j\|=\|\check{\psi}_j\|_1$. 直观上, $\check{\psi}_j$ 恰是圆盘特征函数的 Fourier 变换限制在环形区域 $\{2^j<|x|<2^{j+1}\}$ 上的部分. 一个快速的计算表明我们预计 $\|\check{\psi}_j\|_1=C2^{\frac{n-1}{2}j}$. 利用渐近性可以把这个算得更精确:
\[
\check{\psi}_{j}(x)=\int\frac{J_{\frac{n-2}{2}}(2\pi|x|r)}{(|x|r)^{\frac{n-2}{2}}}\,\psi\bigl(2^{j}(|r|-1)\bigr)\,r^{n-1}\,dr.
\]
在区域 $|x|<1$ 中, $|\check{\psi}_j|\leq\|\psi_j\|_1\leq 1$; 在区域 $|x|>1$ 中, 对那些 $r\in\operatorname{support}\psi_j$, $r$ 接近 $1$, 于是 $r|x|>1$, 即我们处于 Bessel 函数渐近性成立的区域. 主项贡献了
\[
C|x|^{-\frac{n-2}{2}}\int e^{2\pi ir|x|}\,\psi\bigl(2^{j}(|r|-1)\bigr)\,r^{\frac{n-1}{2}}\,dr.
\]
令 $s=2^{j}(|r|-1)$; 我们得到
\[
C|x|^{-\frac{n-2}{2}}e^{2\pi i|x|}2^{-j}\int e^{2\pi is2^{-j}|x|}\,\psi(s)\,(1+s2^{-j})^{\frac{n-1}{2}}\,ds.
\]
若把 $(1+s2^{-j})^{\frac{n-1}{2}}$ 在 $s=0$ 附近展开, 主项是
\begin{align*}
C|x|^{-\frac{n-2}{2}}e^{2\pi i|x|}2^{-j}\int e^{2\pi is2^{-j}|x|}\,\psi(s)\,ds
&=C|x|^{-\frac{n-2}{2}}e^{2\pi i|x|}2^{-j}\,\check{\psi}_{0}(2^{-j}|x|).
\end{align*}
注意到 $\psi_0\in\mathcal{S}$; 总之, 若令
\[
\varphi(x)=|x|^{-\frac{n-1}{2}}\check{\psi}_{0}(|x|),
\]
则 $\varphi\in L^1$, 且
\[
|\check{\psi}_{j}(x)|\leq C\,2^{\frac{n-1}{2}j}\,2^{-jn}\,\varphi(2^{-j}x).
\]
于是 $\|\check{\psi}_j\|_1\leq C2^{\frac{n-1}{2}j}\|\varphi\|_1=C2^{\frac{n-1}{2}j}$. 这印证了我们的直观, 但应当提醒读者: 在上面的计算中, 我们对误差项的控制是非常马虎的.

现在可以内插了: 若 $\frac{1}{p}=\frac{t}{1}+\frac{1-t}{2}$, 即 $t=\frac{2}{p}-1$, 则
\begin{align*}
{}_p\|T_{j}\|&\leq{}_{1}\|T_{j}\|^{t}\,{}_{2}\|T_{j}\|^{1-t}
\leq 2^{j\left(\frac{n-1}{2}\right)t}\cdot 1^{t}\leq 2^{j(n-1)\left(\frac{1}{p}-\frac{1}{2}\right)}.
\end{align*}
于是, 当 $p\leq 2$ 时, 只要
\[
\lambda>(n-1)\left(\frac{1}{p}-\frac{1}{2}\right),
\]
我们就能使几何级数求和.
\end{remark}

\begin{figure}[htbp]
\centering
\begin{tikzpicture}[scale=0.9]
  % 阴影区域 (竖线填充)
  \fill[pattern=vertical lines, pattern color=black!70]
    (0,3.4) -- (2.5,0) -- (5,3.4) -- (5,4.2) -- (0,4.2) -- cycle;
  % 边框
  \draw (0,0) rectangle (5,4.2);
  % 虚线 V 形边界
  \draw[dashed] (0,3.4) -- (2.5,0);
  \draw[dashed] (2.5,0) -- (5,3.4);
  % 标记点
  \draw (0,3.4) circle (2.4pt);
  \draw (5,3.4) circle (2.4pt);
  \fill (2.5,0) circle (2.4pt);
  % 标注
  \node[left] at (-0.1,3.4) {$\alpha=\frac{n-1}{2}$};
\end{tikzpicture}
\caption{Bessel 函数的振荡与截断}
\label{fig:3}
\end{figure}

\section{de Leeuw 与 Stein 的一个定理}
\label{sec:5.2}

\begin{theorem}\label{thm:5.6}
若 $0<\lambda<\frac{n-1}{2}$, 则当 $\lambda>(n-1)\bigl|\frac{1}{p}-\frac{1}{2}\bigr|$ 时, $T_\lambda$ 在 $L^p(\mathbb{R}^n)$ 上有界 (对照图~\ref{fig:3}).
\end{theorem}

\begin{proof}
当然, 我们可以沿标准模型的轮廓给出一个证明; 不过为了展示技术的多样性, 我们将给出一个对 $K_{\lambda}$ (而非 $\mu_{\lambda}$) 作分解的证明.

标准分解提示我们把 $K_\lambda$ 分解为支在二进环形区域上的小块. 为此, 设 $K\in\mathcal{S}$ 满足 $0\leq K\leq 1$, 当 $|x|\leq 100$ 时 $K(x)=1$, 当 $|x|>200$ 时 $K(x)=0$. 令 $\phi_j(x)=K\bigl(\frac{|x|}{2^j}\bigr)-K\bigl(\frac{|x|}{2^{j+1}}\bigr)$. 注意 $\phi_j$ 支在环形区域 $100\cdot 2^j<|x|<200\cdot 2^j$ 上, 于是 $\phi_j$ 的作用如同某个环形区域特征函数的光滑化版本.

现在令 $K_j=\phi_jK_\lambda$. 则 $\sum_j\phi_j(x)=K(x)-K(0)$, 于是 $K(x)-\sum_j\phi_j(x)=K(0)=1$; 因此 $K_\lambda K(x)=\sum_jK_j+K_\lambda$, 从而 $K_\lambda$ 与 $\sum_jK_j$ 只相差 $K_\lambda\bigl(1-K(x)\bigr)$. 由于 $K_\lambda\bigl(1-K(x)\bigr)$ 有界且紧支集, 它给出一个 $L^1$ 误差项, 在所有 $L^p$ 上是有界算子.

在这幅图景中, $\|K_j\|_1$ 容易计算, 而 $\|\hat{K}_j\|_\infty$ 却难. 前者给出:
\[
\int\left|K_{\lambda}(x)\right|\,\left|\phi_{j}(x)\right|dx
\leq C\int_{100\cdot 2^{j}}^{200\cdot 2^{j}}r^{-\frac{n+1+2\lambda}{2}}\,r^{n-1}\,dr
=C\,2^{-\lambda j}\,2^{j\frac{n-1}{2}},
\]
与之前一样. 为了对 $\|\hat{K}_j\|_\infty$ 可能有多大有一个直观感受, 注意 $\hat{K}_j=\hat{K}_\lambda\star\hat{\phi}_j$, 而
\[
\hat{\phi}_{j}(\xi)=2^{nj}\hat{K}(2^{j}\xi)-2^{n(j-1)}\hat{K}(2^{j-1}\xi).
\]
若令 $\delta=2^{-j}$, $\psi(\xi)=2\hat{K}(\xi)-\hat{K}\bigl(\frac{\xi}{2}\bigr)$, 则可以把这些改写为
\[
\psi_{\delta}(\xi)=\delta^{-n}\psi(\delta^{-1}\xi)=\hat{\phi}_{j}(\xi).
\]
现在 $\psi_\delta$ 是一个标准的恒等逼近, 只是 $\int\psi_\delta=0$. 当 $\delta$ 变小, 即 $j\to\infty$ 时, 我们预计 $\psi_\delta\star\mu_\lambda$ 将收敛到 $\bigl(\int\psi_\delta\bigr)\mu_\lambda=0$, 并且预计收敛速度由 $\mu_\lambda$ 的光滑性控制. 于是我们预计有估计:
\[
|\psi_{\delta}\star\mu_{\lambda}-0|\leq C\delta^{\lambda}=C2^{-j\lambda}.
\]
有一种更直观的思考方式: 把 $\psi$ 想成一个圆盘的特征函数; 那么 $\psi_{\delta}\star f$ 就是在一个更小的圆盘上对 $f$ 作平均. 在 $|\xi|=1$ 附近, $\mu_{\lambda}$ 的最大值是 $C\delta^{\lambda}$, 它支配着这个平均. 远离 $|\xi|=1$ 处, $\mu_{\lambda}$ 光滑, 于是这个平均如同 $\mu_{\lambda}\int\psi_{\delta}=0$.

直观分析至此完成; 如同本节开头的论证, 估计
\[
\|K_{j}\|_{1}\leq C\,2^{j\frac{n-1}{2}}\,2^{-j\lambda};\qquad
\|\hat{K}_{j}\|_{\infty}\leq C2^{-j\lambda}
\]
足以通过内插完成定理的证明. 我们现在的义务, 是给这些估计背后的直观一个严格 (尽管乏味) 的解读. 不过这个证明会引出几个重要的技巧, 所以它是一定要看的.

在区域 $|\xi|>2$ 中, $\psi_\delta$ 与 $\mu_\lambda$ 的``支集''不相交, 于是 $\psi_\delta\star\mu_\lambda$ 应当为零. 由于 $\psi_\delta$ 实际上并不紧支集, 我们需要另一个窍门. $\psi\in\mathcal{S}$ 的急速下降应当能解决问题. 于是,
\[
|\psi_{\delta}(\eta)|=|\delta^{-n}\psi(\delta^{-1}\eta)|
\leq C\delta^{-n}\,|\delta^{-1}\eta|^{-(n+i)}
=C\delta^{i}\,|\eta|^{-(n+i)},
\]
\[
|K_{j}(\xi)|\leq\int|\mu_{\lambda}(\xi-\eta)|\,|\psi_{\delta}(\eta)|\,d\eta
\leq C\|\mu_{\lambda}\|_{\infty}\,\delta^{i}\int\chi_{D}(\xi-\eta)\,|\eta|^{-(n+i)}\,d\eta;
\]
这里 $D=\operatorname{support}\mu_\lambda$. 由于必须有 $|\xi-\eta|\leq 1$, 有 $|\xi|-|\eta|<1$, 即 $|\eta|>|\xi|-1>1$, 于是这个积分被 $C$ 控制, 取 $i$ 充分大便可使其任意小.

在区域 $|\xi|<2$ 中, 我们预计可以用 $\psi_\delta\star\mu_\lambda$ 收敛到零的速度来控制其大小. 典型地, 我们会估计
\begin{align*}
\left|\psi_{\delta}\star\mu_{\lambda}(\xi)-\Bigl(\int\psi_{\delta}\Bigr)\mu_{\lambda}(\xi)\right|
&\leq\int|\mu_{\lambda}(\xi-\delta\eta)-\mu_{\lambda}(\xi)|\,|\psi(\eta)|\,d\eta.
\end{align*}
现在我们想用估计
\[
|\mu_{\lambda}(\xi-\delta\eta)-\mu_{\lambda}(\xi)|\leq C|\delta\eta|^{\lambda}
\]
来完成证明. 不幸的是, 当 $\lambda>1$ 时这个估计是错的; 不过, 若 $\lambda\leq 1$, 我们可以完成证明:
\begin{align*}
\int|\mu_{\lambda}(\xi-\delta\eta)-\mu_{\lambda}(\xi)|\,|\psi(\eta)|\,d\eta
&\leq C\int|\delta\eta|^{\lambda}\,|\psi(\eta)|\,d\eta\leq C\delta^{\lambda}.
\end{align*}
若 $\lambda>1$, 我们需要分部积分 $k=[\lambda]$ 次; 这将把 $\mu_\lambda$ 微分成一团乱麻乘以 $(1-|\xi|^2)^{\lambda-k}$, 而此时 $\lambda-k\leq 1$. 为使奏效, $\psi$ 的 $k$ 次积分 (称之为 $\phi$) 必须具有与 $\psi$ 相同的结构, 即 $\phi\in\mathcal{S}$ 且 $\int\phi=0$. 当然, 对任意的 $\psi$, 它的不定积分一般不满足这一点; 我们必须小心. 如下定义 $\phi$:
\[
\check{\phi}(x)=\frac{\check{\psi}(x)}{(2\pi ix_{1})\cdots(2\pi ix_{k})}.
\]
现在 $\tilde{\psi}=K_1$ 在原点的一个邻域中消失, 于是 $\tilde{\phi}\in\mathcal{S}$, $\phi\in\mathcal{S}$; 且 $\int\phi=\check{\phi}(0)=0$. 此外, $D_{\xi_1}\cdots D_{\xi_k}\phi=\psi$, 并且
\[
\int\eta_{1}\cdots\eta_{k}\,\phi(\eta)\,d\eta
=\bigl(D_{x_{1}}\cdots D_{x_{k}}\check{\phi}\bigr)(0)=0.
\]
再者,
\[
D_{\eta_{1}}\cdots D_{\eta_{k}}\,\mu_{\lambda}(\xi-\delta\eta)
=C\,(\delta\eta_{1})\cdots(\delta\eta_{k})\,\mu_{\lambda-k}(\xi-\delta\eta).
\]
于是我们可以分部积分 $k$ 次而不产生任何边界项, 因为 $\mu_{\lambda-k}$ 与 $\phi$ 都在无穷远处消失. 我们要指出, 若诸 $\eta$ 并非互不相同, 这些公式并不准确; 另一方面, 我们只需处理 $\lambda<\frac{n-1}{2}<n$ 的情形, 所以我们有用之不竭的余地来选择互不相同的 $\eta$. 为完成证明:
\begin{align*}
\bigl|\hat{K}_{j}(\xi)\bigr|&=\left|\int\mu_{\lambda}(\xi-\delta\eta)\,\psi(\eta)\,d\eta\right| \\
&=C\delta^{k}\left|\int\mu_{\lambda-k}(\xi-\delta\eta)\,\eta_{1}\cdots\eta_{k}\,\phi(\eta)\,d\eta
-\mu_{\lambda-k}\int\eta_{1}\cdots\eta_{k}\,\phi(\eta)\,d\eta\right| \\
&\leq C\delta^{k}\int|\mu_{\lambda-k}(\eta-\delta\eta)-\mu_{\lambda-k}(\eta)|\,\bigl|\eta_{1}\cdots\eta_{k}\bigr|\,|\phi(\eta)|\,d\eta \\
&\leq C\delta^{k}\,\delta^{\lambda-k}\int|\eta|^{\lambda-k}\,|\eta|^{k}\,|\phi(\eta)|\,d\eta
\leq C\delta^{\lambda}. \qedhere
\end{align*}
\end{proof}

看一眼图~\ref{fig:3} 中的直线段便可知发生过某种内插. 但那是哪种内插呢? 毕竟, 图~\ref{fig:3} 画的并不是某个算子的 $L^p$--$L^q$ 有界性, 而是一族算子的 $L^p$--$L^p$ 有界性. 要作内插, 需要 ${}_p\|T_\lambda\|$ 作为 $\lambda$ 的函数具有某种凸性. 如果研究一下 Riesz--Thorin 定理的证明, 其要求是 $F(z)=\int T_zfg$ 在带形区域上解析、在边界上连续、且增长不快于 $Ce^{-b|y|}$. 我们把这样的算子族称为\emph{容许} (admissible) 的.

\begin{theorem}\label{thm:5.7}
\index{内插!复内插 (complex interpolation)}
\index{容许算子族 (admissible family)}
设 $T_{z}$ 是区域 $0\leq\operatorname{Re}z\leq 1$ 中的一个容许算子族. 设 $T_{0+iy}$ 在 $L^{p_0}$ 上有界, $T_{1+iy}$ 在 $L^{p_1}$ 上有界. 若
\[
{}_{p_{j}}\|T_{j+iy}\|\leq e^{C\,e^{-b|y|}}
\]
对某个 $b<\pi$ 成立, 则对 $\frac{1}{p}=\frac{t}{p_0}+\frac{1-t}{p_1}$, $T_t$ 在 $L^p$ 上有界.
\end{theorem}

\begin{proof}
这个称为复内插的结果归功于 Stein. 尽管这是一个重要的结果, 详细的证明会使我们离题太远; 参见章末的文献.

为证明该定理, Stein 首先把三线定理推广, 得到估计
\[
\log|F(z)|\leq\frac{1}{2}\sin\pi x\int_{-\infty}^{\infty}
\left[\frac{\log|F(iy)|}{\cosh\pi y-\cos\pi x}
+\frac{\log|F(1+iy)|}{\cosh\pi y+\cos\pi x}\right]dy.
\]
Stein 现在取 $f=\sum a_j\chi_{E_j}$, $g=\sum b_k\chi_{F_k}$ 为满足 $\|f\|_p=\|g\|_{p'}=1$ 的简单函数. 使用定理~\ref{thm:1.13} 的记号, 令
\[
F(z)=\int T_{z}\left(\sum|a_{j}|^{p\alpha}e^{i\theta_{j}}\chi_{E_{j}}\right)\left(\sum|b_{k}|^{p'(1-\alpha)}e^{i\psi_{k}}\chi_{F_{k}}\right),
\]
其中 $\alpha(z)=\frac{1-z}{p_0}+\frac{z}{p_1}$. 则与定理~\ref{thm:1.13} 的证明一样,
\begin{align*}
|F(1+iy)|&\leq{}_{p_{1}}\|T_{1+iy}\|\leq M_{1}(y), \\
|F(0+iy)|&\leq{}_{p_{0}}\|T_{0+iy}\|\leq M_{0}(y),
\end{align*}
这里诸 $M_j(y)$ 的增长不快于 $e^{Ce^{-b|y|}}$. 于是修改后的三线定理给出
\begin{align*}
|F(z)|\leq\exp\frac{1}{2}\sin\pi x\int_{-\infty}^{\infty}
&\frac{\log\bigl|M_{0}(y)\bigr|}{\cosh\pi y-\cos\pi x} \\
&+\frac{\log\bigl|M_{1}(y)\bigr|}{\cosh\pi y+\cos\pi x}\,dy,
\end{align*}
而 $M_j(y)$ 的增长保证了积分的收敛性. 于是 $|F(z)|\leq C_t$, 正如所需.
\end{proof}

\begin{corollary}\label{cor:5.8}
设 $0<\lambda<\frac{n-1}{2}$. 则当 $\lambda>(n-1)\bigl|\frac{1}{p}-\frac{1}{2}\bigr|$ 时, $T_\lambda$ 在 $L^p$ 上有界.
\end{corollary}

\begin{proof}
令
\[
K_{z}(x)=c_{z}\frac{J_{\frac{n}{2}+z}(2\pi|x|)}{|x|^{\frac{n}{2}+z}},
\]
其中 $c_{z}$ 待定. 为检验容许性,
\[
\int T_{z}(f)\,g\leq\|T_{z}f\|_{2}\,\|g\|_{2}\leq\|f\|_{2}\,\|g\|_{2}.
\]
类似地,
\[
{}_{2}\|T_{0+iy}\|=\|\mu_{0+iy}\|_{\infty}\leq 1.
\]
但
\[
{}_{1}\|T_{1+iy}\|\leq\|K_{\lambda+iy}\|_{1}.
\]
现在归一化常数变得重要了! 由引理~\ref{lem:5.1} 与命题~\ref{prop:5.2}, Bessel 函数的标准归一化涉及 Gamma 函数的商. 若考察 Gamma 函数作为积分的定义, 就会看到它在复平面上呈指数增长. Bessel 函数本身也有类似的增长; 只需看看它的幂级数展开. 因此平衡这两种增长十分重要; 实际的工作见 Stein~\cite{stein-interpolation}.
\end{proof}

\section{评注与推广}
\label{sec:5.3}

\noindent (5.4):\quad 这里的简单反例之所以有效, 只是因为核的振荡
\[
\cos(2\pi|x|)
\]
是缓慢的; 它对表面相似而振荡为 $\cos(2\pi|x|^2)$ 的核就不会奏效. 这类核的 $L^p$ 行为非常不同; 关于 $\frac{e^{i|x|^\alpha}}{|x|^\beta}$ 的处理, 见 C.~Fefferman~\cite{fefferman-stronglysingular} 与 S.~Wainger~\cite{wainger-special}.

\noindent (5.5):\quad 我们将在第~\ref{chap:6} 章第~1 节中重新着手 $K_{\lambda}$ 的局部计算.

\noindent (5.6):\quad 这里给出的证明归功于 de Leeuw~\cite{deleeuw}.

\noindent (5.7):\quad 这个结果归功于 Stein~\cite{stein-interpolation}; 关于其他内插技术, 见 Stein 与 Weiss~\cite{steinweiss} 第~209 页的讨论.

\noindent (5.8):\quad 这个结果归功于 Stein~\cite{stein-interpolation}.

% 第 6 章 Restriction Theorems (原书第 101-113 页)
\chapter{限制性定理}
\label{chap:6}

本章的目的是证明 Bochner--Riesz 算子 $T_\lambda$ 在最优范围 $\frac{2n}{n+1+2\lambda}\leq p\leq\frac{2n}{n-1-2\lambda}$ 内的 $L^p$ 有界性, 但仅当 $\lambda$ 相当大时如此. 问题在于控制振荡: 在一维情形, 我们曾把卷积核 $\frac{\sin 2\pi x}{x}$ 换成 $\frac{1}{x}$, 但在更高维数, 指数 $e^{i|x|}$ 无法简化. 我们将叙述 C.~Fefferman 与 E.~M.~Stein 的理论, 它分析的正是这类振荡核.

\section{限制}
\label{sec:6.1}
\index{限制性定理 (restriction theorem)|emph}

为了导出振荡效应的一个简单模型, 我们分析 $\mathbb{R}^{2}$ 中的一个简单的核:
\[
K(x)=\frac{e^{i|x|}}{|x|^{\frac{3}{2}+\lambda}}.
\]
设 $f$ 紧支集在 $D=\{y \mid |y|<1\}$ 中, 并设 $x$ 很大. 控制 $|x-y|$ 的几何如图~\ref{fig:4} 所示. 我们可以取 $|x-y|$ 近似等于 $|x|$ 减去 $y$ 在 $x$ 方向上的投影长度. 于是
\[
|x-y|\approx|x|-y\cdot\frac{x}{|x|},
\qquad
e^{i|x-y|}\approx e^{i|x|}\,e^{-iy\cdot\frac{x}{|x|}}.
\]
我们还可以估计
\[
|x-y|^{\frac{3}{2}+\lambda}\approx|x|^{\frac{3}{2}+\lambda}+\text{误差},
\]
于是
\begin{align*}
K\star f(x)=\int\frac{e^{i|x-y|}}{|x-y|^{\frac{3}{2}+\lambda}}\,f(y)\,dy
&\approx\frac{e^{i|x|}}{|x|^{\frac{3}{2}+\lambda}}\int e^{-iy\cdot\frac{x}{|x|}}\,f(y)\,dy \\
&=\frac{e^{i|x|}}{|x|^{\frac{3}{2}+\lambda}}\,\hat{f}\Bigl(\frac{x}{|x|}\Bigr).
\end{align*}
推论~\ref{cor:5.4} 的反例使我们预计 $T_{\lambda}f$ 会看起来像 $|x|^{-(\frac{3}{2}+\lambda)}$, 但这个新估计还要更好: 它把 $K$ 的作用分离成两个部分: 球面变量方向上的一个 Fourier 变换, 与无穷远处一个纯径向的衰减. 在极坐标下, 我们于是预计
\begin{align*}
\|K\star f\|_{p}^{p}&\geq\int_{2}^{\infty}\int_{0}^{2\pi}|K\star f(r,\theta)|^{p}\,d\theta\,r\,dr \\
&=\int_{2}^{\infty}r\,r^{-(\frac{3}{2}+\lambda)p}\,dr\cdot\int_{0}^{2\pi}\bigl|\hat{f}(1,\theta)\bigr|^{p}\,d\theta.
\end{align*}
由此可知, 除非 $p>\frac{4}{3+2\lambda}$, 否则 $K\star f$ 不属于 $L^p$; 若这一点满足, 则 $K$ 的 $L^p$ 有界性将蕴含
\[
\int_{0}^{2\pi}\bigl|\hat{f}(1,\theta)\bigr|^{p}\,d\theta\leq C\|f\|_{p}^{p}.
\]
这种类型的不等式称为\emph{限制性定理} (restriction theorem), 因为它支配着 $f$ 的 Fourier 变换在集合 $\{\xi \mid |\xi|=1\}$ 上的限制, 而这毕竟是一个零测集. 这是了不起的, 因为任意的可测函数是无法限制到零测集上而不出现无穷的. 另一方面, 若 $f\in L^p$, $1<p\leq 2$, 则 $\hat{f}$ 是由 $L^{p'}$ 范数中的极限过程计算的: 若 $f_n\in\mathcal{S}$ 且 $f_n\to f$, 则 Hausdorff--Young 不等式
\[
\|\hat{f}\|_{p'}\leq\|f\|_{p}
\]
保证诸 $\hat{f}_n$ 按 $L^{p'}$ 范数收敛到某个函数 $\hat{f}$. 一个子列几乎处处收敛, 于是 $\hat{f}$ 几乎处处有定义. 但在一个零测集上定义 $\hat{f}$ 看起来是荒谬的.

\begin{figure}[htbp]
\centering
\begin{tikzpicture}[scale=0.85]
  % 小球 D 与向量 y
  \fill (0,0) circle (1.2pt);
  \draw[->] (0,0) -- (1.2,1.5) node[midway, above left] {$y$};
  \draw[->] (0,0) -- (6,0) node[midway, below] {$x$};
  \draw[->] (1.2,1.5) -- (6,0) node[midway, above right] {$x-y$};
  % 投影虚线
  \draw[dashed] (1.2,1.5) -- (1.2,0);
  \fill (1.2,0) circle (1.2pt);
  \node[below] at (0.9,-0.15) {$\bigl(y\cdot\frac{x}{|x|}\bigr)\frac{x}{|x|}$};
\end{tikzpicture}
\caption{球面上两点间距离与法向几何}
\label{fig:4}
\end{figure}

由于限制性定理令人惊讶到使人怀疑其正确性, 并让人对 $T_{\lambda}$ 的有界性起疑, 我们就从限制性定理讲起.

\begin{proposition}\label{prop:6.1}
设 $f$ 是 $\mathbb{R}^n$ 上的径向函数. 则
\[
\bigl|\hat{f}(1,\theta)\bigr|\leq C\|f\|_{p}
\]
对 $L^p$ 中所有的径向 $f$ 成立, 当且仅当 $1\leq p<\frac{2n}{n+1}$.
\end{proposition}

\begin{proof}
由于 $f$ 径向, $\hat{f}$ 也径向, 于是 $\hat{f}(1,\theta)$ 与 $\theta$ 无关. 那么线性泛函 $f\mapsto\hat{f}(1,\theta)$ 在 $L^p$ 上有界当且仅当它由对一个 $L^{p'}$ 函数的积分给出. 我们把 $\hat{f}(1,\theta)$ 缩写为 $\hat{f}(1)$; 则
\begin{align*}
\hat{f}(1)&=\int_{0}^{\infty}f(r)\int_{S^{n-1}}e^{2\pi i\sigma r\xi}\,d\sigma\,r^{n-1}\,dr \\
&=\int_{0}^{\infty}f(r)\,\frac{J_{\frac{n-2}{2}}(2\pi r|\xi|)}{(r|\xi|)^{\frac{n-2}{2}}}\,r^{n-1}\,dr.
\end{align*}
但
\[
\frac{J_{\frac{n-2}{2}}(2\pi r)}{r^{\frac{n-2}{2}}}
\]
属于 $L^{p'}$ 当且仅当 $p'>\frac{2n}{n-1}$. 结论得证.
\end{proof}

这只是通往定理的一个初步近似, 因为它只处理径向函数. 尽管如此, 它表明即便是这些函数的 Fourier 变换也具有令人惊讶的光滑性 (在一定程度上). 下一个结果试图对限制中出现的临界指标 $\frac{2n}{n+1}$ 给出一个几何层面的理解; 我们想限制从 $L^p$ 到 $L^q$ 的限制性可能成立的 $(p,q)$ 的范围. 我们关于 Fourier 变换最初得到的估计是在 Hausdorff--Young 定理中; 我们利用 Fourier 变换的伸缩结构证明了这些估计是最优的: $f(\delta x)$ 被变换为 $\delta^{-n}\hat{f}(\delta^{-1}\xi)$. 但圆对伸缩并不是不变的, 我们必须更狡猾一些.

直观上, 我们想让 $\hat{f}|_{S^1}$ 尽可能大. 利用伸缩, 我们设想 $\hat{f}$ 支在一个小矩形上, 并把它定向为与 $S^1$ 相交尽可能多. 见图~\ref{fig:5}. 支集 $\hat{f}$ 的什么样的尺寸能在最小化 $\|f\|_p$ 的同时最大化 $\hat{f}|_{S^1}$? 取矩形的大小为 $\delta\times\delta^{\alpha}$, 其中 $\alpha<1$, 从而这是长边. 我们选取 $\alpha$ 使得矩形的角恰好与圆相交; 这使与圆的相交达到最大. 由于 $x=\sqrt{1-y^2}=\sqrt{1-(1-\delta)^2}\approx\sqrt{\delta}$, 我们取 $\alpha=\frac{1}{2}$.

\begin{figure}[htbp]
\centering
\begin{tikzpicture}[scale=0.85]
  % 圆弧 (大半径圆的一段)
  \draw (-3.4,-1.2) arc (250:290:9);
  % 矩形
  \draw (0,0) rectangle (4.5,-1.2);
  \node[above] at (2.25,0.05) {$\delta^{\alpha}$};
  \node[right] at (4.55,-0.6) {$\delta$};
\end{tikzpicture}
\caption{Knapp 反例函数的支撑集几何}
\label{fig:5}
\end{figure}

% 原书编号由 6.1 直接跳至 6.3, 此处跳过 6.2 以与原书一致.
\addtocounter{mtthm}{1}

\begin{proposition}\label{prop:6.3}
不等式
\[
\|\hat{f}|_{S^{n-1}}\|_{q}\leq C\|f\|_{p}
\]
在 $q>\frac{n-1}{n+1}p'$ 时不成立.
\end{proposition}

\begin{proof}
设 $\psi$ 是光滑的紧支集函数, 满足当 $|\xi|<\sqrt{n}$ 时 $\psi(\xi)=1$. 定义
\[
\hat{f}_{\delta}(\xi_{1},\ldots,\xi_{n})=\psi\bigl(\delta^{-2}(\xi_{1}-1),\,\delta^{-1}\xi_{2},\,\ldots,\,\delta^{-1}\xi_{n}\bigr).
\]
若 $\delta<1$, $\xi\in S^{n-1}$, $|\xi_j|\leq\frac{\delta}{\sqrt{n}}$, 我们断言 $\hat{f}_\delta=1$. 只需证 $\delta^{-2}|\xi_1-1|<1$, 即 $1-\xi_1<\delta^2$. 但
\begin{align*}
1-\xi_{1}&=1-\sqrt{1-(\xi_{2}^{2}+\cdots+\xi_{n}^{2})}
\leq 1-\bigl(1-(\xi_{2}^{2}+\cdots+\xi_{n}^{2})\bigr) \\
&=\xi_{2}^{2}+\cdots+\xi_{n}^{2}\leq(n-1)\left(\frac{\delta}{\sqrt{n}}\right)^{2}\leq\delta^{2}.
\end{align*}
因此,
\begin{align*}
\|\hat{f}|_{S^{n-1}}\|_{q}
&\geq\left|\{(\xi_{1},\ldots,\xi_{n}) \mid |\xi_{j}|\leq\frac{\delta}{\sqrt{n}}\}\right|^{\frac{1}{q}}
\geq C\left(\frac{\delta}{\sqrt{n}}\right)^{\frac{n-1}{q}}
=C\delta^{\frac{n-1}{q}}.
\end{align*}
另一方面,
\[
f_{\delta}(x_{1},\ldots,x_{n})=\delta^{n+1}e^{2\pi ix_{1}}\,\check{\psi}(\delta^{2}x_{1},\,\delta x_{2},\,\ldots,\,\delta x_{n}),
\]
于是
\[
\|f_{\delta}\|_{p}=\delta^{\frac{n+1}{p'}}\|\check{\psi}\|_{p}=C\delta^{\frac{n+1}{p'}}.
\]
这样, 对所有小的 $\delta$, 有 $\delta^{\frac{n+1}{p'}}\geq C\delta^{\frac{n-1}{q}}$, 定理随之得证.
\end{proof}

\begin{remark}\label{rem:6.4}
在图~\ref{fig:6} 中, 我们画出了已知限制性失效的区域. 虚线表示上述结果中区域的边界; 竖直线来自要求 $1\leq p<\frac{2n}{n+1}$, 即便对径向函数这也是必要的. 两条线相交于 $p=q=\frac{2n}{n+1}$.
\end{remark}

\begin{figure}[htbp]
\centering
\begin{minipage}[b]{0.45\textwidth}
\centering
\begin{tikzpicture}[scale=0.75]
  % 图 6: 失效区域 (斜线阴影在虚线上方/左侧)
  \fill[pattern=north west lines, pattern color=black!70]
    (0,0) -- (0,4.2) -- (3.1,4.2) -- (3.1,3.1) -- (4.2,0) -- cycle;
  \draw (0,0) rectangle (4.2,4.2);
  \draw[dashed] (3.1,3.1) -- (4.2,0);
  \draw (3.1,0) -- (3.1,4.2);
  \fill (3.1,3.1) circle (2.2pt);
  \fill (4.2,0) circle (2.2pt);
  \node[right] at (4.3,3.1) {$\frac{n+1}{2n}$};
  \node[below] at (3.1,-0.2) {$\frac{n+1}{2n}$};
  \node[below=0.5cm] at (2.1,-0.3) {};
\end{tikzpicture}
\caption{限制性估计失效的指标区域}
\label{fig:6}
\end{minipage}\hfill
\begin{minipage}[b]{0.45\textwidth}
\centering
\begin{tikzpicture}[scale=0.75]
  % 图 7: 有界区域 (斜线阴影在右上方窄条)
  \fill[pattern=north west lines, pattern color=black!70]
    (3.1,4.2) -- (3.1,2.3) -- (4.2,0) -- (4.2,4.2) -- cycle;
  \draw (0,0) rectangle (4.2,4.2);
  \draw[dashed] (3.1,4.2) -- (3.1,2.3);
  \draw (3.1,2.3) -- (4.2,0);
  \draw (3.1,2.3) circle (2.2pt);
  \fill (3.5,1.15) circle (2.2pt);
\end{tikzpicture}
\caption{Tomas--Stein 限制性有界区域}
\label{fig:7}
\end{minipage}
\end{figure}

\begin{proposition}\label{prop:6.5}
\index{球面测度 (surface measure $d\sigma$)}
若 $1\leq p\leq\frac{2(n+1)}{n+3}$, 则
\[
\|\hat{f}|_{S^{n-1}}\|_{2}\leq C\|f\|_{p}.
\]
\end{proposition}

\begin{proof}
我们可以假设 $f\in\mathcal{S}$, 这为下面的形式运算提供了依据. 则
\begin{align*}
\|\hat{f}|_{S^{n-1}}\|_{2}^{2}
&=\int\hat{f}(\sigma)\,\overline{\hat{f}(\sigma)}\,d\sigma
=\int(f\star\tilde{f})\,\widehat{d\sigma}
=\int f\star\tilde{f}(x)\,\widehat{d\sigma}(x)\,dx \\
&=\int\check{d\sigma}\star f(x)\,f(x)\,dx
\leq\|f\|_{p}\,\|\check{d\sigma}\star f\|_{p'}.
\end{align*}
只需证明 $Rf=\check{d\sigma}\star f$ 是从 $L^p$ 到 $L^{p'}$ 的有界映射. 我们有两种技术: 一是把 $\check{d\sigma}(x)$ 切成碎块; 二是使用复内插. 这里我们做内插; 在定理之后的讨论中, 我们会看到标准模型能告诉我们什么.

我们需要把算子 $R$ 嵌入一个解析的算子族; 令 $R_{z}f=K_{z}\star f$; 照例, 我们只能控制在 $p=1$ 与 $p=2$ 处发生的事情. 于是在 $p'=\infty$ 与 $p'=2$ 处, 我们用
\begin{align*}
\|K_{z}\star f\|_{\infty}&\leq\|K_{z}\|_{\infty}\,\|f\|_{1}, \\
\|K_{z}\star f\|_{2}&\leq\|\hat{K}_{z}\|_{\infty}\,\|f\|_{2}.
\end{align*}
啊, 又是我们的老对手——对 $K$ 与 $\hat{K}$ 的同时控制! 现在,
\[
K_{0}=\widehat{d\sigma}(x)=\frac{J_{\frac{n-2}{2}}(2\pi|x|)}{|x|^{\frac{n-2}{2}}},
\]
它不仅属于 $L^\infty$, 而且 $|x|^{\frac{n-1}{2}}K_0\in L^\infty$. 所以, 我们可以容许把 $K_0$ 弄得更糟一些而仍保持第一个不等式. 另一方面, $\hat{K}_0=d\sigma$, 它根本不是一个 $L^\infty$ 函数, 所以必须对它作相当大的改进. 在改进与变糟之间将落下 $K_z$, 而我们将得到一个使 $K_0$ 给出有界算子的 $p$ 的范围. 令
\[
K_{z}(x)=c_{z}\frac{J_{\frac{n-2}{2}+z}\bigl(2\pi|x|\bigr)}{|x|^{\frac{n-2}{2}+z}}.
\]
若 $z=-\frac{n-1}{2}$, 则 $\frac{n-2}{2}+z=-\frac{1}{2}$, 并且由 Watson 的公式 3.4(6) 我们注意到
\[
J_{-\frac{1}{2}}(2\pi|x|)=\frac{\cos 2\pi|x|}{(\pi|x|)^{\frac{1}{2}}}.
\]
于是, 在 $z=-\frac{n-1}{2}$ 处, $K_z\in L^\infty$ (事实上, 这正是我们从 Bessel 函数的渐近性所预期的).

另一方面, 若 $z>0$, 则 $\hat{K}_z(\xi)=(1-|\xi|^2)_+^{z-1}$, 它首次属于 $L^\infty$ 是在 $z-1=0$ 时. 复内插给出 $z=\frac{n+3}{2}-\frac{n+1}{p}$. 在 $z=0$ 处, 这就给出定理.
\end{proof}

\begin{remark}\label{rem:6.6}
\begin{enumerate}[label=\alph*)]
\item 在图~\ref{fig:7} 中, 我们画出了限制性已知有界的区域. 空心圆与虚线表示已知无界的区域.
\item 若把 $f$ 换成 $\delta^{-n}f(\delta^{-1}x)$, 则 $\hat{f}$ 被换成 $\hat{f}(\delta\xi)$, 而 $\hat{f}_\delta|_{S^{n-1}}$ 与限制在半径为 $\delta^{-1}$ 的球面上的 $\hat{f}$ 相同. 把变量替换进行到底, 得到
\[
\left(\int_{S^{n-1}}\bigl|\hat{f}(R,\sigma)\bigr|^{2}\,d\sigma\right)^{\frac{1}{2}}
\leq C_{p}\,R^{-\frac{n}{p'}}\|f\|_{p}. \qedhere
\]
\end{enumerate}
\end{remark}

\essayhead{ess:6.7}{限制性再论}

利用标准模型可以证明一个较弱版本的限制性定理; 如同定理~\ref{thm:5.6} 的证明, 我们写 $\widehat{d\sigma}=\sum K_j$. 由于
\[
|\widehat{d\sigma}(x)|\leq C|x|^{-\frac{n-1}{2}},
\]
故
\[
\|K_{j}\|_{\infty}\leq C2^{-j\frac{n-1}{2}}.
\]
为估计 $\|\hat{K}_j\|_\infty$, 我们依靠~\ref*{prop:5.2} 中的想法: $\hat{K}_j=\psi_\delta\star d\sigma$, 其中 $\delta=2^{-j}$. 直观上, $\psi_\delta$ 紧支集在半径为 $\delta$ 的球中, 于是 $\psi_\delta\star d\sigma$ 在环形区域 $1-\delta<|\xi|<1+\delta$ 之外为零. 为使 $\psi_\delta\star d\sigma$ 达到最大, 我们看 $|\xi|=1$; 此时 $\psi_\delta\star d\sigma$ 应当被看作 $\psi_\delta$ 对 $d\sigma$ 的平均. 也就是说, 它是半径为 $\delta$ 的球在球面上的平均. 该球与球面相交处的曲面测度为 $\delta^{n-1}$; 而球本身的测度为 $\delta^n$, 于是这个平均是 $\delta^{-1}=2^j$. 因此我们预计:
\[
\|K_{j}\|_{\infty}\leq C2^{-j\frac{n-1}{2}},
\qquad
\|\hat{K}_{j}\|_{\infty}\leq C2^{j}.
\]
于是, 与 $K_j$ 作卷积的算子的范数是 $2^{\epsilon j}$, 其中
$\epsilon=2\bigl(1-\frac{1}{p}\bigr)-\frac{n-1}{2}\Bigl(1-2\bigl(1-\frac{1}{p}\bigr)\Bigr)$.
当 $p<2\frac{n+1}{n+3}$ 时它是负的.

用我们的标准模型解释限制性, 仍然没有告诉我们这一切究竟意味着什么. 我们将编织一个相当复杂的论证, 以展示使限制性得以成立的一些想法. 尽管如此, 时至今日仍无人能给出一个几何判别法, 来判定哪些零测集具有限制性而哪些没有, 所以我们很难说已经理解了限制性.

我们将首先考察曲率的影响. 若限制性对某条曲线 $S\subset\mathbb{R}^2$ 成立, 而 $S$ 包含一段线段, 则通过平移、旋转、伸缩与截断, 我们可以假设 $S=[0,1]\times\{0\}$. 现在令 $f(x_1,x_2)=g(x_1)h(x_2)$, 并注意 $\hat{f}(\xi_1,\xi_2)|_S=\hat{g}(\xi_1)|_I\,\hat{h}(0)$; 于是若取固定的 $g\in\mathcal{S}$ 使 $\hat{g}|_I=1$, 则限制性不等式将蕴含 $|\hat{h}(0)|\leq C\|h\|_{p}$. 于是映射 $h\mapsto\hat{h}(0)$ 是 $L^p$ 上的有界线性泛函, 这蕴含 $e^{2\pi i0x}\in L^{p'}$. 这只有在 $p=1$ 时才成立. 寓意是: 对任何包含一片超平面的曲面, 限制性都不可能成立.

我们的第二个观察把 $S$ 上的限制性与支在 $S$ 上的测度的 Fourier 变换联系起来. 设 $S$ 是光滑紧曲面, $d\sigma$ 为曲面测度; 我们断言: 不等式 $\|\hat{f}|_{S}\|_{q}\leq C\|f\|_{p}$ 对某个 $p>1$ 成立, 当且仅当 $\widehat{d\sigma}\in L^{p_0}$ 对某个 $p_0<\infty$ 成立.

先设限制性成立. 定义 $R:L^p(\mathbb{R}^n)\to L^q(S)$ 为 $Rf=\hat{f}|_S$. 则伴随算子 $R^{\star}:L^{q'}(S)\to L^{p'}(\mathbb{R}^n)$ 也有界, 并且我们很快可以算出
\begin{align*}
\int\bigl(R^{\star}g\bigr)(x)\,f(x)\,dx
&=\int g(\sigma)\,Rf(\sigma)\,d\sigma
=\int g(\sigma)\,\hat{f}(\sigma)\,d\sigma
=\int\widehat{g\,d\sigma}(x)\,f(x)\,dx.
\end{align*}
于是 $R^{\star}g=\widehat{g\,d\sigma}$; 这必须在分布的意义下理解, 而且就我们的目的而言, 不妨设 $g\in C^\infty(S)$. 取 $g\equiv 1$,
\[
\|\widehat{d\sigma}\|_{p'}\leq C|S|^{\frac{1}{q'}}.
\]
若 $p>1$, 则 $p'<\infty$.

反之, 设 $\widehat{d\sigma}\in L^{p_0}$ 对某个 $p_0<\infty$ 成立. 则
\[
\int|\hat{f}|^{2}\,d\sigma=\int f\star\widehat{f\,d\sigma}
\leq\|\widehat{d\sigma}\|_{p_{0}}\|f\star f\|_{p_{0}'}
\leq C\|f\|_{p}^{2};
\]
这里我们用了 Young 卷积不等式
\[
\|h\star g\|_{r}\leq\|h\|_{p}\,\|g\|_{q},
\qquad
1+\frac{1}{r}=\frac{1}{p}+\frac{1}{q}.
\]
由于 $r=p_{0}'$, $p=q$, 故 $p=\frac{2p_{0}}{2p_{0}-1}>1$.

既然 $d\sigma$ 的 Fourier 变换的 $L^p$ 行为刻画了 $S$ 的可限制性, 对限制性的几何理解就必须从``$S$ 的曲率如何影响 $\widehat{d\sigma}$ 的大小''开始. 正确的结果是: 若 $S$ 的曲率在 $S$ 的任何一点处都不无穷阶消失, 则 $\widehat{d\sigma}\in L^{p_0}$. 与其在这里试图证明它, 不如看看标准模型如何处理直线与圆的差别. 设 $\phi$ 是光滑函数, 当 $|x|<1$ 时 $\phi(x)=1$, 当 $|x|>2$ 时 $\phi(x)=0$. 则当 $\delta\to 0$ 时, $\delta^{-1}\phi(\delta^{-1}x_2)\phi(x_1)$ 按分布收敛到支在 $[-1,1]\times\{0\}$ 上的测度. 其 Fourier 变换的行为如同
\[
\hat{\phi}(\xi_{1})\,\check{\phi}(\delta\xi_{2})\to\hat{\phi}(\xi_{1})\check{\phi}(0)=\check{\phi}(\xi_{1}).
\]
由于在 $\xi_{2}$ 方向上完全没有衰减, 它不属于任何 $p<\infty$ 的 $L^{p}$.

那么单位圆有何不同? 我们有
\[
\lim_{\delta\to0}\delta^{-1}\psi\bigl(\delta^{-1}(|\xi|-1)\bigr),
\]
而且有估计
\[
\bigl|\bigl[\psi\bigl(\delta^{-1}(|\xi|-1)\bigr)\bigr]\check{}\,(x)\bigr|\leq C\delta\,|x|^{-\frac{n-1}{2}}\,\check{\psi}(\delta|x|).
\]
于是我们预计在极限中
\[
\bigl|\widehat{d\sigma}(x)\bigr|\leq C|x|^{-\frac{n-1}{2}},
\]
而这确实成立. 我们所缺乏的, 是对 $\check{\psi}$ 的估计的几何理解. 命题~\ref{prop:5.2} 中的计算始于 Bessel 函数的渐近性, 而当我们利用球面对称计算 $\widehat{d\sigma}$ 时, Bessel 函数又出现了. 这类公式唯一的问题在于它们平均得太多; 我们想要局部化. 我们从对 $\mathbb{R}^2$ 中一个环形区域的 Fourier 变换的几何理解开始.

\essayhead{ess:6.8}{重访``凸起''}
\index{曲率 (curvature)}

我们的问题在于: Fourier 变换非常适合把 $\mathbb{R}^2$ 看作一维空间的乘积; 计算矩形的变换很容易, 球则难一些. 因此, 若把一个环形区域分解成矩形, 计算会容易得多. 我们用单位分解把它分解成环状片段 $R_j$, 每个的大小为 $\sqrt{\delta}\times\delta$; 这给出 $\frac{2\pi}{\sqrt{\delta}}$ 块. 若 $v_j$ 表示顶点在 $R_j$ 中心的单位向量, 则 $v_j=(\cos\theta_j,\sin\theta_j)$, 其中 $\theta_j\approx j\sqrt{\delta}$. 我们可以把这些矩形视为以原点为中心的矩形 $R_j^0$ 的平移:
\[
\chi_{R_{j}}(\xi)=\chi_{R_{j}^{0}}(\xi-v_{j});
\qquad
\check{\chi}_{R_{j}}(x)=e^{2\pi iv_{j}x}\,\check{\chi}_{R_{j}^{0}}(x).
\]
为理解 $\check{\chi}_{R_{j}^{0}}(x)$, 我们把积分
\[
\int_{R_{j}^{0}}e^{2\pi ix\xi}\,d\xi
\]
看作让某些波对矩形作平均的过程. 在二维, 这些是沿方向 $\frac{x}{|x|}$ 传播、频率为 $|x|$ 的平面波. 哪些方向与哪些频率对 $\chi_{R_j^0}$ 的 Fourier 变换有贡献?

若方向平行于 $v_j$, 即 $x=v_j|x|$, 则波对一个厚度为 $\delta$ 的区间作平均, 于是使这种平均非零的适当频率是 $|x|=\delta^{-1}$. 此时平均给出矩形的大小, 即 $\delta^{\frac{3}{2}}$, 因而这个频率与方向贡献 $e^{2\pi ixv_j}\delta^{\frac{3}{2}}$. 但 $x=v_j|x|$, $|x|=\delta^{-1}$, 于是我们得到大小为
\[
\frac{e^{2\pi i|x|}}{|x|^{\frac{3}{2}}}
\]
的贡献.

那么还有哪些方向有贡献? 在 $\pi$ 方向, 我们得到与 $0$ 方向同样大小的贡献, 但相位不同, 于是求和给出一个余弦. 在垂直于 $v_j$ 的方向, 进来的波``看到''的厚度是 $\sqrt{\delta}$, 于是它挑选频率 $|x|=\delta^{-\frac{1}{2}}$, 其中 $x=|x|w_j$, $v_j\cdot w_j=0$. 这样的波对 Fourier 变换的贡献是 $\delta^{\frac{3}{2}}=|x|^{-\frac{3}{2}}$. 由于这小得多, 我们采取这样的直观: 凡是不平行于 $v_{j}$ 的方向都可以忽略, 尽管这当然是非常简化的. 总而言之,
\[
\check{\chi}_{R_{j}}(x)=\frac{e^{2\pi i|x|}}{|x|^{\frac{3}{2}}}.
\]
现在为计算 $\psi_\delta$, 我们需要把来自诸 $R_j$ 的贡献加起来. 它们会在 $x$ 一侧相互作用, 我们需要控制这一点. 直观上, 我们把 $R_j^0$ 想成在 $v_j$ 与 $w_j$ 方向上的矩形, 于是其 Fourier 变换又是一个矩形 $S_j$, 在 $w_j$ 方向上大小为 $\delta^{-\frac{1}{2}}$, 在 $v_j$ 方向上大小为 $\delta^{-1}$. 我们断言这些矩形本质上互不重叠. 诸 $S_j$ 与原点的距离为 $|x|=\delta^{-1}$, 而相邻矩形之间的夹角为 $C\delta^{\frac{1}{2}}$. 于是, 中心之间的距离为
\[
\delta^{-1}\,\delta^{\frac{1}{2}}=\delta^{-\frac{1}{2}}
\]
(在 $w_j$ 方向). 这样 $S_j$ 与 $S_{j+1}$ 本质上不相交, 而且随着 $j$ 的变动, 整个环形区域被诸 $S_j$ 覆盖. 我们得出结论:
\[
\check{\psi}_{\delta}(x)=\frac{\cos 2\pi|x|}{|x|^{\frac{3}{2}}},
\qquad |x|\approx\delta^{-1}.
\]
把诸 $\psi_{2^{-j}}$ 加起来便得到 $\chi_{D}$, 或者取极限便得到 $\check{d\sigma}$.

\section{Bochner--Riesz 平均}
\label{sec:6.2}

\begin{theorem}\label{thm:6.9}
若 $\lambda>\frac{n-1}{2(n+1)}$, 则对最优范围
\[
\frac{2n}{n+1+2\lambda}<p<\frac{2n}{n-1-2\lambda},
\]
$T_\lambda$ 在 $L^p(\mathbb{R}^n)$ 上有界.
\end{theorem}

\begin{proof}
我们像之前那样分解 $K_{\lambda}=\sum K_j$, 但不做内插, 而是采用这样的哲学: $\hat{K}_j=\psi_\delta\star\mu_\lambda$ 的行为如同集中在厚度为 $\delta$ 的环形区域上. 于是~\ref*{rem:6.6}b) 控制了到这类环形区域上的限制:
\begin{align*}
\|K_{j}\star f\|_{2}^{2}&=\|\hat{K}_{j}\hat{f}\|_{2}^{2}
=\int_{1-\delta<|\xi|<1+\delta}|\mu_{\lambda}(\xi)|^{2}\,\bigl|\hat{f}(\xi)\bigr|^{2}\,d\xi \\
&\leq\delta^{\frac{1}{2}+\lambda}\left(\frac{1}{\delta}\int_{1-\delta}^{1+\delta}\int_{S^{n-1}}\bigl|\hat{f}(\sigma)\bigr|^{2}\,d\sigma\,dr\right) \\
&\leq C\delta^{\frac{1}{2}+\lambda}\bigl\|\hat{f}|_{S^{n-1}}\bigr\|_{2}^{2}
\leq C\delta^{\frac{1}{2}+\lambda}\|f\|_{p}^{2}.
\end{align*}
这里有一个小问题: 我们需要控制的不是 $\|K_j\star f\|_2$ 而是 $\|K_j\star f\|_p$; 另一方面, 除非从 $L^2$ 出发, 我们无法到 Fourier 变换一侧去使用限制性. 我们用一个绝妙的新技巧来解决这个问题.

把 $\mathbb{R}^n$ 分解成对角线长度为 $200\cdot 2^j$ 的互不相交的立方体 $Q$, 并令 $\tilde{Q}$ 是与 $Q$ 同心但对角线长度两倍的立方体. 现在
\[
K_{j}\star f=\sum_{Q}K_{j}\star\bigl(f\chi_{Q}\bigr)
=\sum_{Q}\int_{Q}K_{j}(x-y)\,f(y)\,dy.
\]
除非存在 $y\in Q$ 使 $x-y\in\operatorname{support}K_j$, 即 $|x-y|<200\cdot 2^j=\operatorname{diameter}Q$, 被积函数为零. 于是 $|x-\operatorname{center}Q|\leq 2\,\operatorname{diameter}Q$, 且
\[
\operatorname{support}K_{j}\star(f\chi_{Q})\subset\tilde{Q},
\]
从而
\[
K_{j}\star(f\chi_{Q})=\sum\bigl[K_{j}\star(f\chi_{Q})\bigr]\chi_{\tilde{Q}}.
\]
现在对每个 $x\in\mathbb{R}^n$, 至多有 $2^n$ 个立方体 $\tilde{Q}$ 使 $x\in\tilde{Q}$, 于是
\begin{align*}
\left|\sum\bigl[K_{j}\star(f\chi_{Q})\bigr]\chi_{\tilde{Q}}(x)\right|^{p}
&=\left|\sum_{j=1}^{2^{n+1}}a_{j}\right|^{p}\leq C_{p,n}\sum a_{j}^{p} \\
&=\sum\left|\sum\bigl[K_{j}\star(f\chi_{Q})\bigr]\chi_{\tilde{Q}}\right|^{p},
\end{align*}
从而
\[
\int|K_{j}\star f|^{p}\leq C_{p,n}\sum_{Q}\int_{\tilde{Q}}|K_{j}\star(f\chi_{Q})|^{p}.
\]
我们将证明
\[
\bigl\|K_{j}\star(f\chi_{Q})\,\chi_{\tilde{Q}}\bigr\|_{p}^{p}
\leq C\,2^{-j(\frac{1}{2}+\lambda)}\,2^{jn(\frac{1}{p}-\frac{1}{2})p}\,\|f\chi_{Q}\|_{p}^{p},
\]
只要 $1\leq p\leq\frac{2(n+1)}{n+3}$. 于是, 对这个 $p$ 的范围,
\begin{align*}
\int|K_{j}\star f|^{p}
&\leq C\sum_{Q}\int_{\tilde{Q}}|K_{j}\star f\chi_{Q}|^{p}
\leq C2^{-\epsilon j}\sum_{Q}\int_{Q}|f|^{p}
=C2^{-\epsilon j}\|f\|_{p}^{p},
\end{align*}
从而当
\[
\sum_{j}2^{-j\left[\left(\frac{1}{2}+\lambda\right)-n\left(1-\frac{p}{2}\right)\right]}<\infty
\]
时, $\sum_j{}_p\|K_j\|<\infty$.

为得到这个估计, 注意
\[
\int_{\tilde{Q}}|K_{j}\star f|^{p}\leq|\tilde{Q}|^{\frac{1}{r}}\left(\int|K_{j}\star f|^{s}\right)^{\frac{1}{s}}.
\]
若选取 $ps=2$, $\frac{1}{r}=1-\frac{p}{2}$, 则 $|\tilde{Q}|^{\frac{1}{r}}=2^{j(1-\frac{p}{2})n}$. 定理将由估计
\[
\|K_{j}\star f\|_{2}\leq C2^{-j(\frac{1}{2}+\lambda)}
\]
得出.

当然, 开头 我们的直观本应助一臂之力; 不幸的是, $\psi_\delta$ 并非真的紧支集, $\hat{K}_j$ 也并非真的生活在环形区域上. 于是会有来自 $0$ 与 $\infty$ 的贡献. 避开在零处的问题的狡猾办法是干脆把它定义掉. 设 $\chi\in\mathcal{S}$ 满足当 $|r|>\frac{3}{4}$ 时 $\chi(r)=0$, 当 $|r|<\frac{1}{2}$ 时 $\chi(r)=1$. 令
\begin{gather*}
p_{\lambda}(\xi)=\mu_{\lambda}(\xi)\,\chi(|\xi|),\qquad
q_{\lambda}(\xi)=\mu_{\lambda}(\xi)\bigl(1-\chi(|\xi|)\bigr), \\
k_{\lambda}=\check{\mu}_{\lambda}=\check{p}_{\lambda}+\check{q}_{\lambda}\equiv g_{\lambda}+b_{\lambda}.
\end{gather*}
由于 $\mu_{\lambda}$ 在 $\chi$ 的支集上光滑, 故 $p_{\lambda}\in\mathcal{S}$, $g_{\lambda}\in\mathcal{S}$. 于是若 $K_{j}=g_{j}+b_{j}$, 则对每个 $k$, $j$ 有 $\|g_{j}\|_{1}\leq c_{k}2^{-kj}$. 这样, $g$ 部分贡献的是我们能够处理的误差项. 只需对 $b_{j}$ 从头再来. 快速跳过第一部分, 只需证明估计:
\[
\bigl\|b_{j}\star f\bigr\|_{2}\leq C2^{-j(\frac{1}{2}+\lambda)}\,\|f\|_{p}.
\]
我们断言: 若 $|\xi|<\frac{1}{4}$, 则对所有的 $k$ 有 $|\hat{b}_j(\xi)|\leq C_k2^{-kj}$. 注意
\[
\hat{b}_{j}(\xi)=\int\psi_{\delta}(\xi-\eta)\,\mu_{\lambda}(\eta)\bigl[1-\chi(|\eta|)\bigr]\,d\eta,
\]
但 $\mu_{\lambda}\in L^{\infty}$ 且 $\psi_{\delta}\in\mathcal{S}$, 于是
\[
|\psi_{\delta}(x)|\leq C\delta^{-n}\,|\delta^{-1}x|^{-(n+k)},
\]
从而
\[
\bigl|\hat{b}_{j}(\xi)\bigr|\leq C\delta^{k}\int\frac{|1-\chi(|\eta|)|}{|\xi-\eta|^{n+k}}\,d\eta.
\]
由于除非 $|\eta|>\frac{1}{2}$ 否则 $1-\chi=0$, 而 $|\eta-\xi|\geq|\eta|-|\xi|\geq\frac{1}{2}-\frac{1}{4}$, 积分收敛, 断言成立.

对靠近无穷的部分, 我们用估计
\[
\int\bigl|\hat{f}(r,\sigma)\bigr|^{2}\,d\sigma\leq C\,r^{-\frac{n}{p'}}\|f\|_{p}^{2}.
\]
则
\begin{align*}
\|b_{j}\star f\|_{2}^{2}
&=\int_{0}^{\infty}\int_{S^{n-1}}\bigl|\hat{b}_{j}(r)\bigr|^{2}\bigl|\hat{f}(r,\sigma)\bigr|^{2}\,r^{n-1}\,dr\,d\sigma \\
&=\int_{0}^{\infty}\bigl|\hat{b}_{j}(r)\bigr|^{2}\int_{S^{n-1}}\bigl|\hat{f}(r,\sigma)\bigr|^{2}\,d\sigma\,r^{n-1}\,dr \\
&\leq C\|f\|_{p}^{2}\int_{0}^{\infty}\bigl|\hat{b}_{j}(r)\bigr|^{2}\,r^{n-1}\,r^{-\frac{n}{p'}}\,dr \\
&\leq C\|f\|_{p}^{2}\left[C2^{-kj}\int_{0}^{\frac{1}{4}}r^{\frac{n}{p}-1}\,dr+\int_{\frac{1}{4}}^{\infty}|\hat{b}_{j}(r)|^{2}\,r^{n-1}\,dr\right] \\
&\leq C\|f\|_{p}^{2}\bigl[C2^{-kj}+\|\hat{b}_{j}\|_{2}^{2}\bigr].
\end{align*}
但
\begin{align*}
\|\hat{b}_{j}\|_{2}^{2}=\|b_{j}\|_{2}^{2}
&\leq 2\|K_{j}\|_{2}^{2}+\|g_{j}\|_{2}^{2} \\
&\leq 2\int_{100\cdot 2^{j}}^{200\cdot 2^{j}}r^{-\frac{n+1+2\lambda}{2}}\,r^{n-1}\,dr+C2^{-kj}
\leq C2^{-j(1+2\lambda)}. \qedhere
\end{align*}
\end{proof}

\section{推广与评注}
\label{sec:6.3}

\noindent (引言):\quad 关于 Bochner--Riesz 平均与限制性之间关系的分析取自 C.~Fefferman~\cite{fefferman-spherical}. 注意, 在使用这一关系时必须相当小心, 因为事实上 Bochner--Riesz 平均在 $p>\frac{4}{3}$ 的范围内有界, 而限制性恰恰在这个范围内失效.

\noindent (命题~\ref{prop:6.1}):\quad 这些估计与想法起源于谱合成理论; 径向的限制性似乎归功于 Stein, 未发表.

既然对 Bessel 函数的估计会产生限制性, 人们可以用各种方式推广径向限制定理, 例如证明: 若 $p$ 接近 $1$, 则 $f\in L^p$ 必有径向导数; 或者证明: $L^p$ 的径向乘子必具有 H\"older 连续性.

\noindent (命题~\ref{prop:6.3}):\quad 这最初由 C.~Fefferman 与 E.~M.~Stein 用球面调和的论证证明. 这里给出的证明归功于 A.~W.~Knapp, 未发表.

\noindent (命题~\ref{prop:6.5}):\quad $p<\frac{2(n+1)}{n+3}$ 的结果来自 Tomas~\cite{tomas1}, 证明见定理之后的注记. 强结果及其证明归功于 Stein, 未发表.

\noindent (短论~\ref{ess:6.7}):\quad ``$d\sigma$ 的衰减等价于限制性''这一结果归功于 Stein; 见 C.~Fefferman~\cite{fefferman-stronglysingular}.

\noindent (定理~\ref{thm:6.9}):\quad 这个结果归功于 C.~Fefferman~\cite{fefferman-stronglysingular}; 这里给出的证明归功于 Stein; 见~\cite{fefferman-spherical}.

\noindent (进一步的评注):\quad $\mathbb{R}^n$ 中的限制性定理, 应当被视为把 $\mathbb{R}^n$ 中的函数限制到余维数为正的集合上. 例如, 可以限制到 $\mathbb{R}^3$ 中的曲线上而非曲面上; 这最早由 Sj\"olin 对一般曲线作了分析. 关于这一主题有大量文献, 我们无法指望公平地处理. 截至写作之时, 最好的综述与最好的结果归功于 Christ~\cite{christ}.

% 第 7 章 The Multiplier Theorem for the Disc (原书第 114-125 页)
\chapter{圆盘的乘子定理}
\label{chap:7}

\section{Meyer 引理}
\label{sec:7.1}

我们如何猜出一个反例? 让我们从 Fourier 变换一侧出发, 想象一个反例函数 $f$, 其 $\hat{f}$ 支在一个小矩形上. 若圆盘切过这个矩形, 那么在小尺度上, 圆盘的作用就如同沿圆盘法线方向的半平面的特征函数. 由于半平面的特征函数给出有界乘子, 这使我们一无所获.

幸运的是, 我们已经理解圆盘与直线有何不同; 关键全在于不同的矩形如何沿不同方向堆积. 我们应该做的, 是考察大量的矩形 $\hat{f}_j$, 其中诸 $\hat{f}_j$ 支在互不相交的矩形 $R_j$ 上, 而这些矩形覆盖圆盘的边缘. 在每个 $R_j$ 上, 圆盘的作用如同半平面, 我们要做的只是分析各个 $f_j$ 之间的重叠. 问题在于仔细分析 $\bigl\|\sum\chi_{R_j}\bigr\|_p$; 若能取它们互不相交就好了. 我们转而使用第~\ref{chap:4} 章的技术.

\begin{theorem}\label{thm:7.1}
若 $T_0$ 在某个 $L^p$ 上有界, 则对每一个函数序列 $\{f_j\}$,
\[
\Bigl\|\Bigl(\sum|T_{0}f_{j}|^{2}\Bigr)^{\frac{1}{2}}\Bigr\|_{p}\leq C\Bigl\|\Bigl(\sum|f_{j}|^{2}\Bigr)^{\frac{1}{2}}\Bigr\|_{p}.
\]
\end{theorem}

\begin{proof}
设 $\bar{f}=(f_1,\ldots,f_k)$. 注意 $T_0$ 具有实值的卷积核, 因此它保持函数的实部与虚部. 于是我们可以假设诸 $f_j$ 是实值的. 设 $\bar{v}=(v_1,\ldots,v_k)$ 是单位向量, 并注意
\[
\bar{v}\cdot\overline{T_{0}f}=\sum v_{j}T_{0}f_{j}=T_{0}(\bar{v}\cdot\bar{f}).
\]
于是
\[
\int_{S^{k-1}}\int_{\mathbb{R}^{n}}\bigl|\bar{v}\cdot\overline{T_{0}f}(x)\bigr|^{p}\,dx\,dv
\leq{}_{p}\|T_{0}\|^{p}\iint\bigl|\bar{v}\cdot\bar{f}\bigr|^{p}\,dx\,dv.
\]
另一方面,
\[
\bar{v}\cdot\bar{f}=\|\bar{v}\|\,\|\bar{f}\|\cos\theta=\|\bar{f}\|\cos\theta,
\]
于是, 交换积分次序后,
\[
\int_{\mathbb{R}^{n}}\int_{S^{k-1}}|\cos\theta|^{p}\,dv\;\bigl\|\overline{T_{0}f}(x)\bigr\|^{p}\,dx
\leq{}_{p}\|T_{0}\|\iint|\cos\theta|^{p}\,dv\;\|\bar{f}\|^{p}\,dx.
\]
现在, 对固定的 $(f_1,\ldots,f_k)$, 向量 $\bar{f}$ 为 $S^{k-1}$ 定义了一个北极, 于是 $|\cos\theta|$ 的积分与这个向量无关. 特别地, 它与 $x$ 无关, 从而作为常数从不等式两边提出来. 于是
\[
\int\bigl\|\overline{T_{0}f}\bigr\|^{p}\,dx\leq{}_{p}\|T_{0}\|^{p}\int\|\bar{f}\|^{p}\,dx,
\]
即
\[
\int\left[\Bigl(\sum|T_{0}f_{j}|^{2}\Bigr)^{\frac{1}{2}}\right]^{p}dx
\leq{}_{p}\|T_{0}\|^{p}\int\left[\Bigl(\sum|f_{j}|^{2}\Bigr)^{\frac{1}{2}}\right]^{p}dx. \qedhere
\]
\end{proof}

为了完成我们的直观论证, 我们把``圆盘的作用如同半平面的特征函数''这一想法严格化.

\begin{theorem}\label{thm:7.2}
\index{Meyer 引理 (Meyer's lemma)}
设 $v_j$ 是 $\mathbb{R}^n$ 中的单位向量; 设 $P_j=\{\xi \mid \xi\cdot v_j\geq 0\}$, 并设 $H_j$ 是以 $\chi_{P_j}$ 为乘子的算子. 则
\[
\Bigl\|\Bigl(\sum|H_{j}f_{j}|^{2}\Bigr)^{\frac{1}{2}}\Bigr\|_{p}
\leq{}_{p}\|T_{0}\|\,\Bigl\|\Bigl(\sum|f_{j}|^{2}\Bigr)^{\frac{1}{2}}\Bigr\|_{p}.
\]
\end{theorem}

\begin{proof}
设 $T_{j}^{R}$ 是对应于乘子
\[
\chi_{D}\Bigl(\frac{\xi-Rv_{j}}{R}\Bigr)
\]
的算子; 设 $T_0^R$ 对应于 $\chi_D\bigl(\frac{\xi}{R}\bigr)$. 则 $T_j^R=JT_0^RJ^{-1}$, 其中 $J$ 是乘以 $e^{2\pi ixv_j}$ 的等距算子. 于是
\begin{align*}
\Bigl\|\Bigl(\sum|T_{j}^{R} f_{j}|^{2}\Bigr)^{\frac{1}{2}}\Bigr\|_{p}
&=\Bigl\|\Bigl(\sum|JT_{0}^{R}J^{-1}f_{j}|^{2}\Bigr)^{\frac{1}{2}}\Bigr\|_{p} \\
&\leq{}_{p}\|T_{0}^{R}\|\,\Bigl\|\Bigl(\sum|J^{-1}f_{j}|^{2}\Bigr)^{\frac{1}{2}}\Bigr\|_{p} \\
&\leq{}_{p}\|T_{0}\|\,\Bigl\|\Bigl(\sum|f_{j}|^{2}\Bigr)^{\frac{1}{2}}\Bigr\|_{p}.
\end{align*}
现在取极限. 断言是
\[
\lim_{R\to\infty}\chi_{D}\Bigl(\frac{\xi-Rv_{j}}{R}\Bigr)=\chi_{P_{j}}(\xi).
\]
几何上这非常清楚; 我们只是在说: 取极限后, 一个平移、伸缩过的球面变成一张平面 (曲率 $0$). 分析上, 若 $\bigl|\frac{\xi-Rv_j}{R}\bigr|<1$ 则极限为 $1$, 否则为零. 于是, 极限为 $1$ 当且仅当
\[
\lim\frac{\xi-Rv_{j}}{R}\cdot\frac{\xi-Rv_{j}}{R}<1,
\]
即
\[
\lim\left(\frac{|\xi|^{2}}{R^{2}}-\frac{2v_{j}\cdot\xi}{R}+\frac{R^{2}|v_{j}|^{2}}{R^{2}}\right)<1,
\]
亦即 $-2v_j\cdot\xi+1<1$.

现在若诸 $f_j$ 属于 $\mathcal{S}$, 则 $\chi_D\bigl(\frac{\xi-Rv_j}{R}\bigr)\hat{f}_j$ 控制收敛到 $\chi_{P_j}$, 于是 $\lim T_j^Rf_j=H_jf_j$ 在 $L^2$ 中收敛; 一个子列几乎处处收敛. 则
\begin{align*}
\Bigl\|\Bigl(\sum|H_{j}f_{j}|^{2}\Bigr)^{\frac{1}{2}}\Bigr\|_{p}
&=\Bigl\|\liminf_{k}\Bigl(\sum|T_{j}^{R_{k}}f_{j}|^{2}\Bigr)^{\frac{1}{2}}\Bigr\|_{p} \\
&\leq\liminf_{k}\Bigl\|\Bigl(\sum|T_{j}^{R_{k}}f_{j}|^{2}\Bigr)^{\frac{1}{2}}\Bigr\|_{p} \\
&\leq{}_{p}\|T_{0}\|\,\Bigl\|\Bigl(\sum|f_{j}|^{2}\Bigr)^{\frac{1}{2}}\Bigr\|_{p}. \qedhere
\end{align*}
\end{proof}

\essayhead{ess:7.3}{猜测反例, 第二部分}

在选取构成反例的 $f_j$ 时, 我们预计它们的 Fourier 变换会是矩形的特征函数. 与圆盘最匹配的矩形大小为 $\delta^{\frac{1}{2}}\times\delta$, 最长方向与圆盘在 $v_j$ 处相切. 于是直观上, $f_j$ 是大小为 $N^{\frac{1}{2}}\times N$、最长方向平行于 $v_j$ 的矩形的特征函数. 由于上一结果的变换 $H_j$ 沿方向 $v_j$, 我们可以预计在 $\tilde{R}_j$ 上 $H_jf_j\approx\frac{1}{2}$, 于是
\[
\Bigl\|\Bigl(\sum\chi_{\tilde{R}_{j}}\Bigr)^{\frac{1}{2}}\Bigr\|_{p}\leq C\Bigl\|\Bigl(\sum\chi_{R_{j}}\Bigr)^{\frac{1}{2}}\Bigr\|_{p}.
\]
诸 $R_j$ 互不相交, 因此 $\bigl\|\bigl(\sum\chi_{R_j}\bigr)^{1/2}\bigr\|_p$ 容易计算; 但诸 $\tilde{R}_j$ 相互重叠, $\sum\chi_{\tilde{R}_j}$ 的范数相当复杂. 恰好有一个容易的情形: $p=4$. 别管我们在这里已经知道答案; 也许我们能学到些东西.
\begin{align*}
\Bigl\|\Bigl(\sum\chi_{\tilde{R}_{j}}\Bigr)^{\frac{1}{2}}\Bigr\|_{4}^{4}
&=\int\Bigl(\sum\chi_{\tilde{R}_{j}}\Bigr)^{2}
=\sum_{j,k}\bigl|\tilde{R}_{j}\cap\tilde{R}_{k}\bigr|.
\end{align*}
现在看图~\ref{fig:8}. 设 $\theta_j$ 是 $v_j$ 与 $x$ 轴之间的夹角. 回忆~\ref*{ess:6.7} 中的讨论, $\theta_j\approx j\delta^{\frac{1}{2}}=jN^{-\frac{1}{2}}$. 于是 $\tilde{R}_j\cap\tilde{R}_k$ 的面积是一个平行四边形的面积, 即底乘高. 图中的虚线长度为 $N^{\frac{1}{2}}$, 而 $h=N^{\frac{1}{2}}$,
\[
b=\frac{N^{\frac{1}{2}}}{\sin(\theta_{j}-\theta_{k})}\approx\frac{N}{|j-k|}.
\]
若 $j=k$, 估计 $\bigl|\tilde{R}_j\cap\tilde{R}_k\bigr|=N^{\frac{3}{2}}$ 更好; 总之,
\[
\bigl|\tilde{R}_{j}\cap\tilde{R}_{k}\bigr|\geq\frac{N^{\frac{3}{2}}}{|j-k|+1}.
\]
于是
\begin{align*}
\Bigl\|\Bigl(\sum\chi_{\tilde{R}_{j}}\Bigr)^{\frac{1}{2}}\Bigr\|_{4}^{4}
&\geq\sum_{j,k}\frac{N^{\frac{3}{2}}}{|j-k|+1}
=C\bigl(N^{\frac{3}{2}}N^{\frac{1}{2}}\log N\bigr)
=C(\log N)\Bigl(\sum|R_{j}|\Bigr).
\end{align*}
总之,
\[
{}_{4}\|T_{0}\|\geq C(\log N)^{\frac{1}{4}}.
\]
当 $\delta\to 0$ 时, $N\to\infty$, 于是 $T_0$ 不可能是 $L^4$ 有界的.

在某种意义上, 这个窍门不过是计算了一个环形区域的 Fourier 变换, 并表明它不能一致地属于 $L^p$. 若 $p<4$, 我们将需要一个全新的想法, 因为环形区域在这个范围内表现非常好. 我们要做的是: 设 $B$ 为所有 $\tilde{R}_j$ 彼此重叠的集合; 其实我们只需要在 $B$ 上估计 $\bigl(\sum\chi_{\tilde{R}_j}\bigr)^{1/2}$, 因为在 $B$ 之外, 诸 $\tilde{R}_j$ 互不相交, 从而
\[
\Bigl(\sum\chi_{\tilde{R}_{j}}\Bigr)^{\frac{1}{2}}\approx\sum\bigl(\chi_{\tilde{R}_{j}}\bigr)^{\frac{1}{2}},
\]
于是在 $B$ 之外,
\begin{align*}
\Bigl\|\Bigl(\sum\chi_{\tilde{R}_{j}}\Bigr)^{\frac{1}{2}}\Bigr\|_{p}
&\approx\Bigl\|\Bigl(\sum\chi_{R_{j}}\Bigr)^{\frac{1}{2}}\Bigr\|_{p}
\approx\Bigl(\sum|R_{j}|\Bigr)^{\frac{1}{p}}.
\end{align*}
在 $B$ 上, 我们需要正视 $\bigl(\sum\chi_{\tilde{R}_j}\bigr)^{\frac{n}{2}}$. 设 $|B\cap\tilde{R}_j|=c|R_j|$, 其中 (可以想见) $c$ 与 $j$ 无关. 则
\begin{align*}
c\sum|R_{j}|&=c\int_{B}\sum\chi_{\tilde{R}_{j}}
\leq c\int_{B}\sum|H_{j}f_{j}|\,\chi_{\tilde{R}_{j}} \\
&\leq|B|^{\frac{1}{r}}\Bigl\|\Bigl(\sum|H_{j}f_{j}|^{2}\Bigr)^{\frac{1}{2}}\Bigr\|_{p}^{2}
\leq|B|^{\frac{1}{r}}\,{}_p\|T_{0}\|\Bigl(\sum|R_{j}|\Bigr)^{\frac{2}{p}}.
\end{align*}
由于 $r=\bigl(\frac{p}{2}\bigr)'$, $\frac{1}{r}=1-\frac{2}{p}$, 我们得到下界
\[
{}_{p}\|T_{0}\|\geq C\left(\frac{\sum|R_{j}|}{|B|}\right)^{1-\frac{2}{p}}.
\]
若 $p<2$, 这行不通; 但若 $p>2$, 只要 $\frac{\sum|R_j|}{|B|}$ 很大, 我们就能从中获益. 这意味着 $B$ 的测度要比 $\sum|R_j|$ 小得多; 另一方面, $|\tilde{R}_j\cap B|$ 又必须很大——与 $|R_j|$ 可比. 于是诸 $R_j$ 必须精心选取, 使诸 $\tilde{R}_j$ 相交的集合尽可能小.

\begin{figure}[htbp]
\centering
\begin{tikzpicture}[scale=0.8]
  % 两个交叉的长矩形
  \draw[rotate=8] (-4,-0.3) rectangle (4,0.3);
  \draw[rotate=-8] (-4,-0.3) rectangle (4,0.3);
  % 平行四边形 (交叠区域) 与虚线
  \draw[dashed] (-1.3,-0.55) -- (-0.5,0.65);
  \draw[dashed] (0.5,0.65) -- (1.3,-0.55);
  \node at (0,0.95) {$h$};
  \node at (2.2,-0.8) {$b$};
  \node at (-2.5,0.9) {$\theta_j$};
  \node at (2.5,-1.0) {$\theta_k$};
\end{tikzpicture}
\caption{相交矩形的几何与方向夹角}
\label{fig:8}
\end{figure}

\section{Kakeya 集}
\label{sec:7.2}

\begin{theorem}\label{thm:7.4}
\index{Kakeya 集 (Kakeya set)|emph}
\index{发芽 (sprouting)}
对每个 $k>k_0$, 存在集合 $E\subset\mathbb{R}^2$ 与互不相交的矩形族 $\{R_j\}$, 使得若 $\tilde{R}_j$ 是与 $R_j$ 全等、但沿 $R_j$ 最长边的方向平移所得的矩形对, 则
\[
\bigl|\tilde{R}_{j}\cap E\bigr|\geq\frac{1}{4}|R_{j}|,
\qquad
\frac{\sum|R_{j}|}{|E|}\geq C\log k.
\]
\end{theorem}

\begin{proof}
我们通过迭代一个称为``发芽'' (sprouting) 的过程得到集合 $E$. 这里的愿望是构造朝各个方向延伸、而公共交集很小的三角形. 我们先看看如何对两个三角形做到这一点. 设 $ABC$ 是图~\ref{fig:9} 中的三角形. 如图所示, 我们构造新三角形 $AA'M$ 与 $BB'M$: 把 $AC$ 延长到 $AA'$, 把 $BC$ 延长到 $BB'$. 设 $M$ 表示 $AB$ 的中点, 并作出三角形 $AA'M$ 与 $BB'M$.

这一过程构造了两个朝不同方向延伸的新三角形; 我们先证明没有增加多少新的面积. 设 $b$ 为 $AB$ 的底长, $h_{0}$ 为 $ABC$ 的高, $h_{1}$ 为 $AA'M$ 与 $BB'M$ 的高. 我们断言: 两个新三角形的并的面积至多为
\[
\frac{bh_{0}}{2}+2\,\frac{bh_{1}^{2}}{h_{0}}.
\]
设 $E$ 为直线 $AA'$ 与 $MB'$ 的交点, $F$ 为 $MA'$ 与 $BB'$ 的交点. 我们用三角形 $FEB'$ 与 $FEA'$ 的面积来估计三角形 $EB'C$ 与 $FA'C$ 的面积. 我们关注 $FEA'$. 由于 $FEA'$ 与 $MAA'$ 相似, 而 $FEC$ 与 $BAC$ 相似, 我们算得 $FEA'$ 的高为 $h_1+\frac{h_0h_1}{h_0+2h_1}<2h_1$, 并算得底
\[
|EF|=\frac{bh_{1}}{h_{0}+2h_{1}}<c\,\frac{bh_{1}}{h_{0}}.
\]
总面积被 $\frac{bh_{1}^{2}}{h_{0}}$ 控制. 同样的估计适用于 $FEB'$.

为构造集合 $E$, 我们从一个底为 $b$、高为 $h_{0}$ 的等腰三角形出发, 以高度增量 $\frac{h_0}{2}$ 使它发芽. 这得到两个三角形, 其中每一个再以高度增量 $\frac{h_0}{3}$ 发芽. 如此继续到第 $k$ 步. 我们有 $2^k$ 个三角形, 每个的底为 $b2^{-k}$, 高为 $h_0\sum_{i}\frac{1}{i}\approx h_0\log k$. 设 $E$ 是所有这些三角形的并. $E$ 的面积是多少? 在第 $i$ 步, 我们有 $2^i$ 个三角形; 取定其中一个. 它的高至少为 $h_0\log i$, 底至少为 $b2^{-i}$, 而高度增量为 $\frac{h_0}{i}$. 发芽过程增加的面积增量约为
\[
2\cdot\frac{b2^{-i}}{2}\cdot\frac{(h_0/i)^{2}}{h_0\log i}
=2\,\frac{b\,h_{0}\,2^{-i}}{i^{2}\log i}.
\]
由于有 $2^i$ 个三角形, 总增量为 $2\frac{bh_{0}}{i^{2}\log i}$, 于是 $E$ 的面积被
\[
\frac{bh_{0}}{2}+2\sum_{i}\frac{bh_{0}}{i^{2}\log i}<3bh_{0}
\]
控制.

现在我们对上面得到的 $2^k$ 个三角形构造相应的矩形 $R_j$. 图~\ref{fig:10} 画出了这一想法. 设 $ABC$ 是其中一个发芽后的三角形. 把直线 $AC$ 与 $BC$ 延长为直线 $GAC$ 与 $EBC$. 过点 $A$ 作 $DAC$ 的垂线, 直到它与 $EBC$ 交于 $F$. 线段 $AF$ 是 $R_j$ 的一条边. 另一条边是 $AG$, 其中 $G$ 的选取使得 $AC$ 的长度等于发芽三角形 $ABC$ 的高.

\begin{figure}[htbp]
\centering
\begin{minipage}[b]{0.45\textwidth}
\centering
\begin{tikzpicture}[scale=0.75]
  % 图 9: 发芽
  \draw (0,0) node[below left] {$A$} -- (6,0) node[below right] {$B$} -- (3,4.2) node[above] {$C$} -- cycle;
  \draw (0,0) -- (2.2,5.2) node[above] {$A'$};
  \draw (6,0) -- (1.4,5.2) node[above] {$B'$};
  \draw (2.2,5.2) -- (3,0) node[below] {$M$} -- (1.4,5.2);
  \fill (1.85,2.55) circle (1.5pt) node[below left] {$E$};
  \fill (3.6,2.75) circle (1.5pt) node[below right] {$F$};
  \node at (3,-1.4) {};
\end{tikzpicture}
\caption{Besicovitch 三角形的发芽过程}
\label{fig:9}
\end{minipage}\hfill
\begin{minipage}[b]{0.45\textwidth}
\centering
\begin{tikzpicture}[scale=0.75]
  % 图 10: 矩形构造
  \draw (0.4,1.2) node[left] {$A$} -- (1.6,1.2) node[right] {$B$};
  \draw (0.4,1.2) -- (3.4,6.2) node[above] {$C$};
  \draw (1.6,1.2) -- (3.4,6.2);
  \draw (0.4,1.2) -- (-1.6,-3.4) node[left] {$G$};
  \draw (3.4,6.2) -- (-0.9,-3.2) node[below] {$E$};
  \draw (-1.6,-3.4) -- (0.9,0.9);
  \draw (0.75,0.95) arc (0:-60:0.45);
  \node at (1.35,0.55) {$\theta$};
  \fill (1.15,0.62) circle (1.5pt) node[below right] {$F$};
\end{tikzpicture}
\caption{沿发芽三角形构造细长矩形}
\label{fig:10}
\end{minipage}
\end{figure}

显然, 如此构造的诸 $R_j$ 互不相交, 长度为 $h_0\log k$, 宽度至多为 $b2^{-k}$. 我们还需要宽度的下界估计. 图~\ref{fig:10} 表明, 宽度由发芽三角形的细长程度控制. 若 $\theta$ 是 $AF$ 与 $AB$ 之间的夹角, 则 $\frac{AF}{AB}=\cos\theta$. 现在 $\theta=90-\angle(AB,AC)$, 而 $\angle(AB,AC)\geq 60$, 因为发芽过程增大了角部的角度, 而我们最初是一个等腰三角形. 于是 $\theta<30$ 且 $\cos\theta<\frac{\sqrt{3}AC}{2}=\frac{\sqrt{3}b2^{-k}}{2}$. 则
\[
\sum|R_{j}|=2^{k}|R_{0}|\geq 2^{k}\left(\frac{\sqrt{3}\,b\,2^{-k}}{2}\right)(h_{0}\log k),
\]
且
\[
\frac{\sum|R_{j}|}{|E|}\geq\frac{\sqrt{3}}{2}\cdot\frac{h_{0}b\log k}{3h_{0}b}\geq C\log k,
\]
正如所断言的.

最后, 我们估计
\[
|\tilde{R}_{j}\cap E|\geq\frac{1}{4}|\tilde{R}_{j}|.
\]
参考图~\ref{fig:11}. 显然
\begin{align*}
\bigl|\tilde{R}_{j}\cap E\bigr|&\geq\bigl|\tilde{R}_{j}\cap(ABC)\bigr|=\bigl|ABED\bigr| \\
&=\bigl|AFED\bigr|-\bigl|ABF\bigr|\geq\bigl|AFED\bigr|-\bigl|AFG\bigr|;
\end{align*}
在这个估计中我们用了 $AD=\text{height }ABC\leq AC$. 现在, 由于 $D$ 不等于 $E$, $|AFED|\geq\frac{1}{2}|\tilde{R}_j|$. 最后,
\begin{align*}
|AFG|&=\frac{1}{2}(AF)(FG)=(AF)^{2}\,\frac{\tan\theta}{2} \\
&\leq(AF)^{2}\,\frac{\tan\frac{\pi}{6}}{2}
\leq(AB)(AF)\,\frac{1}{2\sqrt{3}} \\
&=\frac{b2^{-k}}{2\sqrt{3}}(AF)\leq h_{0}\,\frac{\log k}{4}\,(AF)=\frac{1}{4}|\tilde{R}_{j}|.
\end{align*}
于是,
\[
|\tilde{R}_{j}\cap E|\geq\frac{1}{2}|\tilde{R}_{j}|-\frac{1}{4}|\tilde{R}_{j}|.
\]
\end{proof}

\begin{figure}[htbp]
\centering
\begin{tikzpicture}[scale=0.8]
  % 图 11: R̃_j 与三角形 ABC
  \draw (0,0) node[below left] {$A$} -- (1.6,0) node[below right] {$B$}
        -- (2.6,3.4) node[above right] {$C$} -- cycle;
  \draw (0,0) -- (0.5,4.4) node[above] {$A'$};
  \draw (0,0) rectangle (0.5,4.4);
  \fill (0.5,2.9) circle (1.5pt) node[right] {$D$};
  \fill (0.5,4.0) circle (1.5pt) node[right] {$E$};
  \draw (0,0) -- (2.9,4.4);
  \draw (2.6,3.4) -- (2.9,4.4);
  \fill (2.75,4.1) circle (1.5pt) node[above right] {$F$};
  \draw (0,0) -- (2.0,1.2) node[below] {$G$};
\end{tikzpicture}
\caption{矩形与 Kakeya 集合交集的面积下界}
\label{fig:11}
\end{figure}

\begin{theorem}\label{thm:7.5}
\index{乘子!球的乘子问题 (multiplier problem for the ball)}
$T_0$ 在 $L^p(\mathbb{R}^n)$ 上无界, 若 $n>1$ 且 $p\neq 2$.
\end{theorem}

\begin{proof}
设 $T_0$ 有界; 我们可以假设 $p>2$; 我们将证明: 对每个 $k>k_0$,
\[
{}_{p}\|T_{0}\|>C(\log k)^{\frac{1}{2}-\frac{1}{p}}.
\]
取定 $k$, 我们如定理~\ref{thm:7.4} 那样构造 Kakeya 集 $E$ 与矩形 $R_j$. 若 $v_j$ 是 $R_j$ 最长边的方向, 我们令 $f_j=\chi_{R_j}$, 并用定理~\ref{thm:7.2} 得到
\[
\Bigl\|\Bigl(\sum|H_{j}f_{j}|^{2}\Bigr)^{\frac{1}{2}}\Bigr\|_{p}
\leq{}_{p}\|T_{0}\|\,\Bigl\|\Bigl(\sum|f_{j}|^{2}\Bigr)^{\frac{1}{2}}\Bigr\|_{p}.
\]
我们将证明: 对 $x\in\tilde{R}_j$, 有 $|H_jf_j(x)|\geq C$. 由此可得
\begin{align*}
C\sum|R_{j}|&=C\sum|\tilde{R}_{j}|\leq C\sum|\tilde{R}_{j}\cap E|
=C\sum\int_{\tilde{R}_{j}\cap E}\leq\sum\int_{\tilde{R}_{j}\cap E}|H_{j}f_{j}|^{2} \\
&\leq\sum\int_{E}|H_{j}f_{j}|^{2}=\int_{E}\sum|H_{j}f_{j}|^{2}
\leq|E|^{1-\frac{2}{p}}\left(\int_{E}\Bigl(\sum|H_{j}f_{j}|^{2}\Bigr)^{\frac{p}{2}}\right)^{\frac{2}{p}} \\
&\leq|E|^{1-\frac{2}{p}}\,{}_p\|T_{0}\|\Bigl\|\Bigl(\sum|H_{j}f_{j}|^{2}\Bigr)^{\frac{1}{2}}\Bigr\|_{p}^{2} \\
&\leq|E|^{1-\frac{2}{p}}\,{}_p\|T_{0}\|\Bigl\|\Bigl(\sum|f_{j}|^{2}\Bigr)^{\frac{1}{2}}\Bigr\|_{p}^{2} \\
&\leq|E|^{1-\frac{2}{p}}\,{}_p\|T_{0}\|\Bigl(\sum|R_{j}|\Bigr)^{\frac{2}{p}}.
\end{align*}
分离出 ${}_p\|T_0\|$, 得
\[
{}_{p}\|T_{0}\|\geq C\left(\frac{\sum|R_{j}|}{|E|}\right)^{1-\frac{2}{p}}
\geq C(\log k)^{1-\frac{2}{p}}.
\]
为完成证明, 我们需要估计 $H_jf_j$. 我们注意到这些估计在旋转与平移下不变, 于是可以假设 $v_j=(0,1)$, $R_j=(0,a)\times(0,b)$. 则
\[
\chi_{P_{j}}(\xi_{1},\xi_{2})=\frac{1-i\,(i\,\mathrm{sign}\,\xi_{2})}{2},
\]
于是
\[
H_{j}=\frac{I\times I-i\,I\times H}{2}.
\]
由于 $f_{j}=\chi_{R_{j}}$ 是实的,
\[
|H_{j}f_{j}|\geq\frac{1}{2}\bigl|(I\times H)f_{j}\bigr|.
\]
但
\begin{align*}
(I\times H)f_{j}(x_{1},x_{2})
&=\iint\delta(x_{1}-y_{1})\,\chi_{(0,a)}(y_{1})\,\frac{\chi_{(0,b)}(y_{2})}{x_{2}-y_{2}}\,dy_{1}\,dy_{2} \\
&=\chi_{(0,a)}(x_{1})\int\frac{\chi_{(0,b)}(y_{2})}{x_{2}-y_{2}}\,dy_{2}.
\end{align*}
现在 $\tilde{R}_j=(0,a)\times(-b,0)\cup(0,a)\times(b,2b)$, 于是 $(x_1,x_2)\in\tilde{R}_j$ 蕴含 $\chi_{(0,a)}=1$ 且 $|x_2-y_2|<2b$, 从而
\[
|H_{j}f_{j}(x_{1},x_{2})|\geq\left|\int\frac{\chi_{(0,b)}(y_{2})}{2b}\,dy_{2}\right|=\frac{1}{2}. \qedhere
\]
\end{proof}

\begin{corollary}\label{cor:7.6}
$T_0$ 在 $L^p(\mathbb{R}^n)$ 上无界, 若 $n>1$ 且 $p\neq 2$.
\end{corollary}

\begin{proof}
若 $T_0$ 在 $L^p$ 上有界, 我们可以把 $\mu_0=\chi_B$, $B=\{\xi \mid |\xi|<1\}$ 限制到 $\mathbb{R}^n$ 中的平面上. 由定理~\ref{thm:1.23}f), 这些限制几乎处处都是 $L^p(\mathbb{R}^2)$ 的乘子. 现在取垂直于 $\xi_1$ 轴的平面, $\mu_0$ 在这些平面上的限制是一列圆盘. 其中至少有一个必是 $L^p(\mathbb{R}^2)$ 的乘子, 这与定理~\ref{thm:7.5} 矛盾.
\end{proof}

\essayhead{ess:7.7}{另一种证明}

在这一短论中, 我们勾画~\ref*{cor:7.6} 的一个更符合``反例应当长什么样''的直观的证明.

我们预计 $\hat{f}_j$ 是以单位向量 $v_j$ 为中心、大小为 $\delta\times\delta^2$ 的矩形的特征函数; 反例函数 $f_j$ 应当是大小为 $N\times N^2$ 的矩形乘以指数 $e^{2\pi iv_jx}$.

\begin{lemma}\label{lem:7.8}
对每个 $\gamma>0$ 与 $k>k_0$, 存在 $N_k$、集合 $E$ 与 $2^k$ 个互不相交的矩形 $\{R_j\}$, 满足:
\begin{enumerate}[label=\arabic*)]
\item 每个 $R_j$ 的较长边的长度为 $N_k^2$; 较短边被一个常数倍上下界住, 即与 $N_k$ 可比.
\item $|E|\leq C(\log\log N_k)^{-1}\sum|R_j|$.
\item 若 $\bar{R}_j$ 是 $R_j$ 沿其最长方向平移距离 $(1+\gamma)N_k^2$ 所得的矩形, 则 $|\bar{R}_j\cap E|\geq C|R_j|$.
\item $\lim_{k\to\infty}N_k=\infty$.
\end{enumerate}
\end{lemma}

\begin{proof}
证明与定理~\ref{thm:7.4} 几乎相同. 起始三角形被赋予初始高度 $2^{2k}\log k$, 并以第 $i$ 步高度增量 $\frac{2^{2k}\log k}{i}$ 发芽. 经过 $k$ 步, 三角形的高为 $N_k^2=2^{2k}\log^{2}k$. 旧的 $R_j$ 照旧构造, 而新的 $R_j$ 由把旧的 $R_j$ 沿远离 $E$ 的方向平移 $\gamma N_k$ 个单位得到, 这些就是我们所说的 $R_j$. 于是 $\bar{R}_j$ 就是原来的 $\tilde{R}_j$, 从而享有同样的估计.
\end{proof}

\par\vspace{1.2em}
{\noindent{\songfont\bfseries 短论~\ref*{ess:7.7} (续)}\par}
\vspace{0.4em}

为了让计算 manageable, 我们将使用一个更简单的模型; 我们不用正确的
\[
K_{0}(x)=\frac{\cos\bigl(2\pi|x|-\frac{3\pi}{4}\bigr)}{|x|^{\frac{3}{2}}}+O\Bigl(|x|^{-\frac{5}{2}}\Bigr),
\]
而用
\[
K(x)=\frac{e^{i|x|}}{|x|^{\frac{3}{2}}}.
\]
从关于 $K$ 的否定信息过渡到关于 $K_{0}$ 的信息是不可能的, 但对 $K_{0}$ 的完整处理要把 $\cos$ 的加法公式用上许多次, 以至于那些不重要的误差项会像苍蝇一样累积. 我们就做 $K$.

如同定理~\ref{thm:7.4} 的证明, 只需证明 $|K\star f_j(x)|\geq c$ 对 $x\in\bar{R}_j$ 成立. 利用~\ref*{sec:6.1} 开头的分析, 我们可以直观地完成计算. 由于 $f_j=e^{iv_j\cdot y}\chi_{R_j}(y)$ 且 $x\in\bar{R}_j$, $x$ 远离 $f_j$ 的支集, 于是
\[
K\star f_{j}(x)\approx\frac{e^{i|x|}}{|x|^{\frac{3}{2}}}\,\hat{f}_{j}\Bigl(\frac{x}{|x|}\Bigr).
\]
现在 $\hat{f}_j$ 是以 $v_j$ 为中心的矩形的特征函数乘以齐性因子 $N^{3}$, 至少对所有使 $\frac{x}{|x|}$ 落在 $\hat{f}_j$ 的支集内的 $x$ 如此. 这发生在 $\bigl|\frac{x}{|x|}-v_j\bigr|<\frac{1}{N}$ 时, 而当 $x\in\bar{R}_j$ 时这是成立的. 于是
\[
|K\star f_{j}|\approx\frac{N^{3}}{|x|^{\frac{3}{2}}},
\]
但 $|x|\approx R_j\ \text{的最长边}=\gamma N^2$, 故
\[
|K\star f_{j}|\approx\frac{C}{\gamma^{\frac{3}{2}}}.
\]
让这个计算严格化是相当容易的; 我们从~\ref*{prop:6.1} 中估计的一个严格版本开始. 设 $w_j$ 是垂直于 $v_j$ 的单位向量; 若 $x\in\bar{R}_j$, $y\in R_j$, $\gamma>2$, 则
\begin{gather*}
\gamma N^{2}<|x-y|<3\gamma N^{2}, \\
\gamma N^{2}\leq\bigl|(x-y)\cdot v_{j}\bigr|<2\gamma N^{2}, \\
\bigl|\bigl(x-y\bigr)\cdot w_{j}\bigr|<\sqrt{2}\,N.
\end{gather*}
则
\begin{align*}
|x-y|&=\left(\bigl[(x-y)\cdot v_{j}\bigr]^{2}+\bigl[(x-y)\cdot w_{j}\bigr]^{2}\right)^{\frac{1}{2}} \\
&=\bigl|(x-y)\cdot v_{j}\bigr|\left(1+\frac{\bigl[(x-y)\cdot w_{j}\bigr]^{2}}{\bigl[(x-y)\cdot v_{j}\bigr]^{2}}\right)^{\frac{1}{2}}.
\end{align*}
现在对 $\sqrt{1+x}$ 在 $0$ 附近、以及对 $e^{ix}$ 在 $0$ 附近使用 Taylor 展开, 我们得到
\[
e^{i|x-y|}=e^{i(x-y)\cdot v_{j}}+Error,
\]
其中
\begin{align*}
|Error|&\leq C\max\frac{\bigl[(x-y)\cdot w_{j}\bigr]^{2}}{\bigl|(x-y)\cdot v_{j}\bigr|}
\leq\frac{C_2N^{2}}{\gamma N^{2}}=\frac{C}{\gamma}.
\end{align*}
现在我们选取 $E$ 的定向使 $(x-y)\cdot v_j\geq 0$; 则
\[
|K\star f_{j}(x)|
\geq\left|\int\frac{e^{i(x-y)\cdot v_{j}}}{|x-y|^{\frac{3}{2}}}\,e^{iv_{j}\cdot y}\chi_{R_{j}}(y)\,dy\right|
-\int_{R_{j}}\frac{Error}{|x-y|^{\frac{3}{2}}}\,dy.
\]
第一项是
\[
\int_{R_{j}}|x-y|^{-\frac{3}{2}}\,dy\geq\bigl(3\gamma N^{2}\bigr)^{-\frac{3}{2}}|R_{j}|=C\gamma^{-\frac{3}{2}}.
\]
第二项被下式从上方控制:
\[
\frac{C}{\gamma}\int_{R_{j}}|x-y|^{-\frac{3}{2}}\,dy
=\frac{C}{\gamma}\cdot\frac{C}{(2\gamma)^{\frac{3}{2}}}\cdot\frac{|R_{j}|}{N^{3}}
\leq\frac{C}{\gamma^{\frac{5}{2}}}.
\]
若 $\gamma$ 充分大, 第二项是第一项的一半, 于是
\[
|K\star f_{j}|\geq\frac{C}{\gamma^{\frac{3}{2}}}.
\]

\section{评注与推广}
\label{sec:7.3}

\noindent (7.1):\quad 定理~\ref{thm:7.1} 取自 Zygmund~\cite{zygmund}; 这里的证明对混合 $f$ 的实部与虚部的算子 $T$ 不适用. Paley 展示了如何用复的 Gaussian 函数对不令人愉快的 $T$ 建立结果. 由于我们实际上处理的是 $(f_1,f_2,\ldots,f_k)$ 上的 Euclid 范数, 这一整个课题被称为``向量值不等式''.

\noindent (7.2):\quad 定理~\ref{thm:7.2} 归功于 Yves~Meyer, 该结果通常称为``Meyer 引理''. 它可以应用于圆盘之外的集合; 事实上, 若 $C$ 是 $\mathbb{R}^2$ 中法方向只有 $N$ 个的凸集, 那么关于这些方向上 Hilbert 变换的 Meyer 引理便可适用. 这将给出 $N$ 边形特征函数的 $L^p$ 乘子范数的下界. 类似地, Meyer 引理可用于证明: 乘子 $|1-|\xi|^2|^{it}$ 的 $L^p$ 有界性导致对卷积核 $K_j(x)=|x\cdot v_j|^{-1-it}$ 的向量值不等式.

Meyer 引理最深远的应用是在 A.~C\'ordoba 与 R.~Fefferman 的理论中, 它把极大函数与乘子联系起来.

\noindent (7.4):\quad Kakeya 集最初是为解决平面上的一个几何问题而构造的. 考虑长度为 $1$ 的线段 $I$; 设 $E$ 是一个集合. 只要不离开 $E$, 我们可以随意旋转与平移 $I$. 使 $I$ 能完整旋转 $360^\circ$ 的 $E$ 的最小面积是多少? 利用 Kakeya 集可以证明: 对每个 $\epsilon$, $|E|$ 都可以做得小于 $\epsilon$. 精彩的阐述见 de Guzm\'an~\cite{deguzman}. 相关推广在 Cunningham~\cite{cunningham} 中讨论.

Kakeya 集的种种变体可用于分析 $\mathbb{R}^2$ 中其他几何--分析问题, 特别是积分的微分. 问题是: 对 $\mathbb{R}^2$ 上的函数, 基于对方形或球作平均构造一个极大函数; 一维的想法给出弱 $L^1$ 与 $L^p$ 界. 但若试图沿任意方向对矩形作平均, Kakeya 集技术表明上确界是无穷. 参见 de Guzm\'an~\cite{deguzman}.

\noindent (7.4):\quad 这个结果归功于 C.~Fefferman. 显然, 同样的想法可以推广到任何沿一个连续的方向集合具有间断的乘子; 特别地, 推广到径向函数的跳跃间断. 振荡型间断的情形仍未解决.

\noindent (7.7):\quad 这里的替代证明源自 C.~Fefferman~\cite{fefferman-ball} 中的概要; 另见~\cite{kenig-tomas}. 计算技巧的意义在于, 即使 Meyer 引理不适用时它们仍然适用; 见 Ruiz~\cite{cordoba-flp}.

% 第 8 章 The Cordoba Multiplier Theorem (原书第 126-146 页)
\chapter{C\'ordoba 乘子定理}
\label{chap:8}

\section{殆正交性}
\label{sec:8.1}
\index{殆正交性!与曲率 (almost orthogonality and curvature)}

奇性的标准模型提示, 我们需要对 $\phi_{\delta}(\xi)=\phi\bigl(\delta^{-1}(|\xi|-1)\bigr)$ 及其乘子性质给出一个好的估计. 对于短论~\ref*{ess:7.3} 中的反例函数, $T_{\delta}f_j=T_0f_j$, 于是我们预计 ${}_4\|T_{\delta}\|\geq C|\log\delta|^{\frac{1}{4}}$. 若事实上这个下界接近真实的算子范数, 则标准模型提示 $\mu_{\lambda}=\sum 2^{-\lambda j}\phi_{2^{-j}}$, 于是
\[
{}_{4}\|T_{\lambda}\|\leq\sum 2^{-\lambda j}\,{}_4\|T_{2^{-j}}\|\leq C\sum 2^{-\lambda j}\,j^{\frac{1}{4}}<\infty
\]
当 $\lambda>0$ 时成立. 这将证明所有 $T_\lambda$ 的 $L^4$ 有界性, 而对偶与内插便会像图~\ref{fig:12} 所勾画的那样, 在其余各处证明最优结果.

~\ref*{ess:7.3} 中的计算提示, $\phi_{\delta}$ 应当被分解成大小为 $\delta\times\delta^{\frac{1}{2}}$ 的小矩形; 在控制了每一个的范数之后, 我们还须证明 Meyer 引理的某种逆命题, 把各个矩形的作用分离开来. 于是我们从建立正交性结果开始.

\begin{figure}[htbp]
\centering
\begin{tikzpicture}[scale=0.9]
  % 阴影区域 (斜线填充)
  \fill[pattern=north west lines, pattern color=black!70]
    (0,2) -- (1.25,0) -- (3.75,0) -- (5,2) -- (5,4) -- (0,4) -- cycle;
  \draw (0,0) rectangle (5,4);
  \draw[dashed] (0,2) -- (1.25,0);
  \draw[dashed] (3.75,0) -- (5,2);
  \draw (0,2) circle (2.4pt);
  \draw (5,2) circle (2.4pt);
  \draw (1.25,0) circle (2.4pt);
  \draw (3.75,0) circle (2.4pt);
  \fill (1.8,0) circle (2.2pt);
  \fill (3.2,0) circle (2.2pt);
  \node[right] at (5.1,2.6) {$\alpha=\frac{1}{2}$};
  \node[below] at (1.25,-0.15) {$\frac{1}{p}=\frac{1}{4}$};
  \node[below] at (3.75,-0.15) {$\frac{1}{p}=\frac{3}{4}$};
\end{tikzpicture}
\caption{圆周乘子的多尺度矩形分解}
\label{fig:12}
\end{figure}

\begin{theorem}\label{thm:8.1}
设 $I_j$ 是等长、内部互不相交且并集为整个 $\mathbb{R}$ 的区间. 设 $P_j$ 是以 $\chi_{I_j}$ 为乘子的乘子算子. 若 $2\leq p<\infty$, 则
\[
\Bigl\|\Bigl(\sum_{j}|P_{j}f|^{2}\Bigr)^{\frac{1}{2}}\Bigr\|_{p}\leq C\|f\|_{p}.
\]
\end{theorem}

\begin{proof}
区间网格可以平移为形如 $I_{j}=(ja,(j+1)a)$ 的区间; 这只是把 $f$ 改变了一个等距算子. 再伸缩得 $(j,j+1)$, 又经平移得 $\bigl(j-\frac{1}{2},j+\frac{1}{2}\bigr)$. 我们得出: 定理中的范数与区间的大小无关.

在证明 Littlewood--Paley 不等式时, 我们需要: a) 结果的一个光滑版本; b) 为证明 a) 所需的一个乘子定理; c) 一个正交化步骤. 为构造乘子结果, 设 $\psi$ 是光滑函数, 在 $I_0$ 上 $\psi=1$, 在 $\tilde{I}_0$ 之外 $\psi=0$. 设 $T_j$ 是对应于 $\psi(\xi-j)$ 的乘子算子; 注意 $P_jf=P_jT_jf$. 由于 H\"ormander--Mikhlin 定理不适用于 $\sum T_j$, 我们需要一个新窍门, 而它将同时解决我们对正交性的需求. 令
\[
\mu(\xi)=\sum_j e^{2\pi ij\theta}\,\psi(\xi-j);\qquad k=\check{\mu}.
\]
则
\[
\int_{I_{0}}\bigl|k\star f(x)\bigr|^{2}\,d\theta
=\int_{I_{0}}\Bigl|\sum_j e^{2\pi ij\theta}\,T_{j}f(x)\Bigr|^{2}\,d\theta
=\sum_j\bigl|T_{j}f(x)\bigr|^{2},
\]
最后一步用了 $I_0$ 上的 Plancherel 定理. 但问题关乎与 $k$ 作卷积的 $L^{p}$ 行为; 为此我们使用 Poisson 求和公式 (参见 Stein 与 Weiss~\cite{steinweiss} 第~7 章第~2 节).

若 $f\in L^1(\mathbb{R})$, $\theta\in I_0$, 则
\[
\sum_j\hat{f}(j)\,e^{2\pi ij\theta}=\sum_j f(\theta+j);
\]
这里 $\hat{f}(j)=\int e^{-2\pi ijy}f(y)\,dy$. 在我们的情形, 需要 $\hat{f}(j)=\psi(\xi-j)$, 于是令
\[
f(x)=\int e^{2\pi iyx}\,\psi(\xi-y)\,dy=e^{2\pi ix\xi}\,\hat{\psi}(x).
\]
由于 $\psi\in\mathcal{S}$, 故 $f\in L^1$. 则
\[
\sum_j\psi(\xi-j)\,e^{2\pi ij\theta}=\sum_j e^{2\pi i(j+\theta)\xi}\,\hat{\psi}(\theta+j),
\]
且
\[
k(x)=\sum_j\bigl[e^{2\pi i(j+\theta)\xi}\bigr]^{\vee}(x)\,\hat{\psi}(\theta+j)
=\sum_j\delta_{j+\theta}(x)\,\hat{\psi}(j+\theta).
\]
现在, 与 $k$ 作卷积的 $L^{p}$ 算子范数由 $k$ 的全变差范数给出:
\[
\|k\|\leq\sum_j\bigl|\hat{\psi}(j+\theta)\bigr|\,\|\delta_{j+\theta}\|
\leq\sum_j\bigl|\hat{\psi}(j+\theta)\bigr|
\leq C\sum_j\bigl((j+\theta)^{2}+1\bigr)^{-1}\leq C.
\]
这控制了乘子范数; 最后我们需要从 $L^{2}$ 过渡到 $L^{p}$; 我们将使用 H\"older 不等式而非独立性. 则
\[
\int\Bigl(\sum|P_{j}f|^{2}\Bigr)^{\frac{p}{2}}
=\int\Bigl(\sum|P_{j}T_{j}f|^{2}\Bigr)^{\frac{p}{2}}
\leq C\int\Bigl(\sum|T_{j}f|^{2}\Bigr)^{\frac{p}{2}},
\]
最后一步用了命题~\ref{prop:4.12}. 而后者又等于
\begin{align*}
C\int\left(\int_{I_{0}}|k\star f|^{2}\,d\theta\right)^{\frac{p}{2}}dx
&=C\int\bigl(\|k\star f\|_{2,\theta}\bigr)^{p}\,dx
\leq\int\bigl(\|k\star f\|_{p,\theta}\bigr)^{p}\,dx \\
&=C\iint_{I_{0}}|k\star f|^{p}\,d\theta\,dx
=C\int_{I_{0}}\int_{\mathbb{R}}|k\star f|^{p}\,dx\,d\theta \\
&\leq C\|k\|\int_{I_{0}}\int|f|^{p}\,dx\,d\theta=C\|f\|_{p}^{p}. \qedhere
\end{align*}
\end{proof}

\begin{corollary}\label{cor:8.2}
设 $B_j=\{(\xi_1,\xi_2) \mid \xi_2\in I_j\}$, 并设 $S_j$ 是以 $\chi_{B_j}$ 为乘子的算子. 则对 $2\leq p<\infty$,
\[
\Bigl\|\Bigl(\sum|S_{j}f|^{2}\Bigr)^{\frac{1}{2}}\Bigr\|_{p}\leq C\|f\|_{p}.
\]
\end{corollary}

\textbf{证明概要}: 由于 $S_j$ 在 $\xi_1$ 方向上如同恒等算子, 该结果是定理~\ref{thm:8.1} 与恒等算子的乘积. 形式上, 对每个 $x_1$ 定义函数 $g(x_2)=f(x_1,x_2)$, 并注意对几乎所有的 $x_1$, $g\in L^1$ 且 $(S_jf)(x_1,x_2)=(P_jg)(x_2)$. 等等.

\begin{remark}\label{rem:8.3}
诸 $I_j$ 的规则分布对定理~\ref{thm:8.1} 的证明并非本质; 我们真正用到的只是 $\psi_j(\xi)=\psi(\xi-j)$ 以及 $\chi_{I_j}\psi_j=\chi_{I_j}$. 这意味着诸 $I_j$ 可以做一些手脚而定理仍然成立. 当然我们可以每隔一个丢掉一个, 从而假设 $\operatorname{distance}(I_j,I_{j+1})\geq c>0$. 但我们还必须要求诸 $I_j$ 的长度与它们的间距以及彼此之间可比:
\[
c_{1}\leq d(I_{j},I_{j+1})\leq c_{2},
\qquad
c_{1}\leq|I_{j}|\leq c_{2}.
\]
我们心中其实已有一个应用: 令
\[
C_{j}=\bigl\{(r,\theta) \,\big|\, j\leq\delta^{-\frac{1}{2}}\theta\leq j+1\bigr\},
\]
并设 $I_{j}$ 是 $C_{j}$ 在 $\xi_{2}$ 轴上的投影. 则
\[
|I_{j}|\approx\delta^{\frac{1}{2}}\,\Bigl|\sin\Bigl(\frac{2\pi j}{\delta^{\frac{1}{2}}}\Bigr)\Bigr|,
\]
于是诸 $I_j$ 的长度不可比. 若把诸 $I_j$ 分成四组, 每组落在围绕 $\frac{\pi}{4}$ 系直线的某个锥形中, 并对每个锥中的区间用单独的平方函数处理, 那么在每个锥内我们就能保证大小可比.
\end{remark}

\essayhead{ess:8.4}{殆正交性与曲率}

不幸的是, 推论~\ref{cor:8.2} 还不够好; 为得到乘子定理, 我们需要形如
\[
\|T_{\delta}f\|_{4}\leq\Bigl\|\Bigl(\sum|S_{j}T_{\delta}f|^{2}\Bigr)^{\frac{1}{2}}\Bigr\|_{4}
\leq\Bigl\|\Bigl(\sum|P_{j}f|^{2}\Bigr)^{\frac{1}{2}}\Bigr\|_{4}
\]
的不等式; 推论~\ref{cor:8.2} 只允许我们从 $\sum|P_jf|^2$ 走到 $|f|$, 所以我们还需要一些正交性. 为看清需要什么, 设 $C_j$ 表示锥
\[
\{(r,\theta) \mid j\leq\delta^{-\frac{1}{2}}\theta\leq j+1\},
\]
并设 $\mu_j=\chi_{C_j}\phi_\delta$. 如上所述, 我们只处理落在锥 $0\leq\theta\leq\frac{\pi}{4}$ 中的 $C_j$. 若 $T_j$ 是对应于 $\mu_j$ 的算子, 令 $T^0=\sum T_j$. 设 $I_j$ 表示诸 $\mu_j$ 的支集在 $\xi_2$ 轴上的投影, 并设 $S_j$ 是以 $\chi_{I_j}(\xi_2)$ 为乘子的算子. 则 $S_jT_j=T_j$, 且估计
\[
\Bigl\|\sum T_{j}f\Bigr\|_{4}=\Bigl\|\Bigl(\sum|T_{j}f|^{2}\Bigr)^{\frac{1}{2}}\Bigr\|_{4}
\]
将蕴含
\[
\|T^{0}f\|_{4}=\Bigl\|\sum T_{j}f\Bigr\|_{4}
\leq C\Bigl\|\Bigl(\sum|T_{j}f|^{2}\Bigr)^{\frac{1}{2}}\Bigr\|_{4}
=\Bigl\|\Bigl(\sum|S_{j}T_{j}f|^{2}\Bigr)^{\frac{1}{2}}\Bigr\|_{4}.
\]
我们希望把它控制为
\[
C\Bigl\|\Bigl(\sum|S_{j}f|^{2}\Bigr)^{\frac{1}{2}}\Bigr\|_{4},
\]
而后者可以用~\ref*{cor:8.2} 处理. 于是我们需要两个不等式:
\begin{align*}
\Bigl\|\sum T_{j}f\Bigr\|_{4}&\leq C\Bigl\|\Bigl(\sum|T_{j}f|^{2}\Bigr)^{\frac{1}{2}}\Bigr\|_{4}, \\
\Bigl\|\Bigl(\sum|S_{j}T_{j}f|^{2}\Bigr)^{\frac{1}{2}}\Bigr\|_{4}
&\leq C\Bigl\|\Bigl(\sum|S_{j}f|^{2}\Bigr)^{\frac{1}{2}}\Bigr\|_{4}.
\end{align*}
为分析第一个, 我们取四次幂: 我们需要
\[
\int\Bigl|\sum T_{j}f\Bigr|^{4}\leq C\int\Bigl(\sum|T_{j}f|^{2}\Bigr)^{2},
\]
即
\[
\int\Bigl|\Bigl(\sum T_{j}f\Bigr)\Bigl(\sum T_{k}f\Bigr)\Bigr|^{2}
\leq C\int\sum|T_{j}f|^{2}\,|T_{k}f|^{2}.
\]
先把积分搁在一边; 我们是在要求
\[
\Bigl(\sum|a_{j}|\Bigr)^{2}\leq C\sum|a_{j}|^{2}.
\]
这非常奇特; 所有真实的不等式, 比如 $(a+b)^2\geq a^2+b^2$, 方向都相反. 当然 $(a+b)^2\leq 2(a^2+b^2)$ 是对的, 但平方里项数越多它越不成立. 因此积分是关键的: 我们最好指望, 若 $x$ 落在某个 $T_jfT_kf$ 的支集之中, 它就不落在其他的支集之中. 简言之, 我们希望诸 $T_jfT_kf$ 在 $L^2$ 上正交.

当然, 这是 $L^{2}$, 而正交性在 Fourier 变换一侧应当有意义:
\begin{align*}
\|T^{0}f\|_{4}^{4}
&=\int\Bigl|\sum_{j,k}\bigl(\mu_{j}\hat{f}\bigr)\star\bigl(\mu_{k}\hat{f}\bigr)\Bigr|^{2}.
\end{align*}
若 $A_{jk}$ 表示
\[
\bigl(\mu_{j}\hat{f}\bigr)\star\bigl(\mu_{k}\hat{f}\bigr)
\]
的支集, 我们将证明: 对所有的 $\xi$, $\sum\chi_{A_{jk}}(\xi)\leq 9$, 即没有一个 $\xi$ 落在多于 $9$ 个的 $A_{jk}$ 之中. 这是完全不相交的替代品, 蕴含正交性. 它同样好用, 因为
\begin{align*}
\Bigl|\sum\bigl(\mu_{j}\hat{f}\bigr)\star\bigl(\mu_{k}\hat{f}\bigr)\Bigr|
&=\Bigl|\sum\bigl(\mu_{j}\hat{f}\bigr)\star\bigl(\mu_{k}\hat{f}\bigr)\chi_{A_{jk}}\Bigr| \\
&\leq\Bigl(\sum|\chi_{A_{jk}}|^{2}\Bigr)^{\frac{1}{2}}
\Bigl(\sum\bigl|\mu_{j}\hat{f}\star\mu_{k}\hat{f}\bigr|^{2}\Bigr)^{\frac{1}{2}} \\
&\leq 3\Bigl(\sum\bigl|\mu_{j}\hat{f}\star\mu_{k}\hat{f}\bigr|^{2}\Bigr)^{\frac{1}{2}}.
\end{align*}
平方并积分,
\begin{align*}
\|T^{0}f\|_{4}^{4}
&=\int\Bigl|\sum\bigl(\mu_{j}\hat{f}\bigr)\star\bigl(\mu_{k}\hat{f}\bigr)\Bigr|^{2}
\leq 9\int\sum\bigl|\bigl(\mu_{j}\hat{f}\bigr)\star\bigl(\mu_{k}\hat{f}\bigr)\bigr|^{2} \\
&=9\int\sum|T_{j}f|^{2}\,|T_{k}f|^{2}
=9\int\Bigl(\sum|T_{j}f|^{2}\Bigr)^{2}
=9\Bigl\|\Bigl(\sum|T_{j}f|^{2}\Bigr)^{\frac{1}{2}}\Bigr\|_{4}^{4},
\end{align*}
于是分析的第一部分完成. 为证 $\sum\chi_{A_{jk}}\leq 9$, 我们注意到
\[
A_{jk}=\operatorname{support}\,(\check{\mu}_{j}\star\check{\mu}_{k})
=\operatorname{support}\,\check{\mu}_{j}+\operatorname{support}\,\check{\mu}_{k}.
\]
我们面前的任务是一个几何任务: 我们必须分析诸 $A_{jk}$ 的交, 即环面的扇区之和. 实际的几何图景很复杂; 图~\ref{fig:13} 画的是若诸 $\mu_j$ 支在一个圆上而非环形区域上时图景的样子; 在图~\ref{fig:14} 中我们画出一个实际的 $A_{jk}$. 为证明不相交性, 诸 $A_{jk}$ 必须在角向与径向上分开. 若 $v_j$ 是 $\mu_j$ 支集中心处的单位向量, 则 $A_{jk}$ 是关于方向 $v_j+v_k$ 对称的集合. 这样的方向对至多有 $C\delta^{-\frac{1}{2}}$ 个; 沿着其中每一个, $A_{jk}$ 包含在半径为 $2\delta^{\frac{1}{2}}$ 的球中. 我们取球的中心为 $v_j+v_k$. 由于这些向量等距分布, 围绕 $v_j+v_k$ 的一个锥便把诸 $A_{jk}$ 彼此隔开; 我们可以得出结论: $A_{jk}\cap A_{jk'}\neq\varnothing$ 蕴含 $|(j+k)-(j'+k')|<2$. 这就完成了角向的隔离.

对于径向, 注意 $A_{jk}$ 的中心在
\[
\Bigl(1-\frac{\delta}{2}\Bigr)\bigl(\cos\theta_{j}+\cos\theta_{k},\ \sin\theta_{j}+\sin\theta_{k}\bigr);
\]
精确到 $\delta$ 的量级, $\theta_j=j\delta^{\frac{1}{2}}$. 诸 $A_{jk}$ 的中心之间的距离因此是向量
\[
\Bigl(1-\frac{\delta}{2}\Bigr)\bigl(\cos\theta_{j}-\cos\theta_{j-1}+\cos\theta_{k}-\cos\theta_{k+1},\ \ldots\bigr)
\]
的长度, 即
\begin{align*}
2\Bigl(1-\frac{\delta}{2}\Bigr)\Bigl(\sin\Bigl(\frac{\theta_{j}+\theta_{j-1}}{2}\Bigr)\sin\Bigl(\frac{\theta_{j}-\theta_{j-1}}{2}\Bigr)&\\
+\sin\Bigl(\frac{\theta_{k}+\theta_{k+1}}{2}\Bigr)&\sin\Bigl(\frac{\theta_{k}-\theta_{k+1}}{2}\Bigr),\ \cdots\Bigr).
\end{align*}
由于精确到 $\delta$ 的量级有 $\theta_j-\theta_{j-1}=\delta^{\frac{1}{2}}$, 诸中心的间距为 $\delta^{\frac{1}{2}}$, 精确到 $\delta$ 的量级.

因此, 沿方向 $v_j+v_k$ 的诸集合 $A_{jk}$ 构成 $\delta^{-\frac{1}{2}}$ 个集合, 每个包含在半径为 $2\delta^{\frac{1}{2}}$ 的球中, 其中心沿 $v_j+v_k$ 方向的一条直线均匀分布. 于是, 沿一条直线没有一个球能与多于两个其他的球重叠, 而在相邻的直线上至多与三个重叠. 这样, 平面上的任何一点至多属于 $9$ 个的 $A_{jk}$.

现在我们需要理解不等式
\[
\Bigl\|\Bigl(\sum|S_{j}T_{j}f|^{2}\Bigr)^{\frac{1}{2}}\Bigr\|_{4}
\leq\Bigl\|\Bigl(\sum|S_{j}f|^{2}\Bigr)^{\frac{1}{2}}\Bigr\|_{4}.
\]
若回忆 Marcinkiewicz 定理, 即定理~\ref{thm:4.14} 的证明, 我们看到需要形如
\[
\int|S_{j}T_{j}f|^{2}\,g\leq\int|S_{j}f|^{2}\,g^{\star}
\]
的不等式, 其中 $g^{\star}$ 是某种极大函数. 按照第~\ref{chap:7} 章的思路, $S_jT_jf$ 在 $v_j$ 方向上如同一个 Hilbert 变换. 若 $M_j$ 是该方向上的一维极大函数, 我们从第~\ref{chap:3}、\ref{chap:4} 章预计
\[
\int|S_{j}T_{j}f|^{2}\,g\leq\int|S_{j}f|^{2}\,\bigl[M_{j}(g^{s})\bigr]^{\frac{1}{s}},
\]
从而
\begin{align*}
\int\sum|S_{j}T_{j}f|^{2}\,g
&\leq\int\sum|S_{j}f|^{2}\,\sup_{j}\bigl[M_{j}(g^{s})\bigr]^{\frac{1}{s}} \\
&\leq\Bigl\|\Bigl(\sum|S_{j}f|^{2}\Bigr)^{\frac{1}{2}}\Bigr\|_{4}^{2}\,
\Bigl\|\sup_{j}\bigl[M_{j}(g^{s})\bigr]^{\frac{1}{s}}\Bigr\|_{2};
\end{align*}
剩下的唯一问题是控制极大函数. 这一做法确实成功, 但 $S_jT_j$ 实际上对应于 $\mu_j$, 它不如 Hilbert 变换奇异; 若把这一点考虑进去, 极大函数将更容易控制. 直观上, $\mu_j$ 是大小为 $\delta\times\delta^{\frac{1}{2}}$ 的矩形的特征函数; 因此 $T_j$ 的作用应当如同在一个具有对偶大小 $\delta^{-1}\times\delta^{-\frac{1}{2}}$ 的矩形上作平均. 这正是我们需要控制的那类极大函数.

\begin{figure}[htbp]
\centering
\begin{minipage}[b]{0.42\textwidth}
\centering
\begin{tikzpicture}[scale=0.7]
  % 图 13: 圆与扇区
  \draw (0,0) circle (2.2);
  \draw (0,0) circle (2.0);
  \foreach \a in {90,120,150} {
    \draw (\a:2.0) -- (\a:2.2);
  }
  \draw[dashed] (90:2.1) -- (150:2.1);
  \node at (2.9,1.9) {$\mu_j$};
\end{tikzpicture}
\caption{圆环支集的弧段扇区划分}
\label{fig:13}
\end{minipage}\hfill
\begin{minipage}[b]{0.42\textwidth}
\centering
\begin{tikzpicture}[scale=0.7]
  % 图 14: 单个 A_jk
  \draw (0,0) circle (2.2);
  \draw[dashed] (1.2,1.85) -- (1.2,-1.85);
  \draw[dashed] (1.75,1.35) -- (1.75,-1.35);
  \draw (1.2,1.85) -- (1.75,1.35) -- (1.75,-1.35) -- (1.2,-1.85) -- cycle;
  \node at (2.6,0) {$A_{jk}$};
\end{tikzpicture}
\caption{相交区域 $A_{jk}$ 的截口形态}
\label{fig:14}
\end{minipage}
\end{figure}

\section{C\'ordoba 极大函数}
\label{sec:8.2}

\begin{definition}\label{def:8.5}
\index{Cordoba 极大函数 (C'ordoba maximal function)|emph}
取定 $N>1$, $a>0$. 设
\[
\mathcal{B}=\{R \mid R\ \text{是尺寸为}\ a\times aN\ \text{的矩形}\}.
\]
C\'ordoba 极大函数 $\mathcal{C}$ 定义为
\[
\mathcal{C}f(x)=\sup\frac{1}{|R|}\int_{R}|f(y)|\,dy,
\]
其中上确界取遍 $\mathcal{B}$ 中所有包含 $x$ 的矩形.
\end{definition}

\begin{theorem}\label{thm:8.6}
\[
\|\mathcal{C}f\|_{2}\leq C\,(\log 6N)^{\frac{1}{2}}\,\|f\|_{2}.
\]
\end{theorem}

\begin{proof}
诸如 $\bigl\|\sum\chi_{R_j}\bigr\|_2^2=\sum_{j,k}|R_j\cap R_k|$ 的量可以像第~\ref{chap:7} 章那样控制, 这便产生我们需要的对数项. 不幸的是, 构成 $\mathcal{C}$ 的过程是在形如 $R_j$ 的矩形上对 $f$ 作平均; 它并不产生 $\chi_{R_j}$.

C\'ordoba 结果的一个关键新想法是: 通过取伴随, 我们可以把``在 $R_j$ 上作平均''换成``乘以 $\chi_{R_j}$'': 若 $\mathcal{C}f=\int f\chi_R$, 则
\[
\mathcal{C}^{\star}f=\Bigl(\int f\Bigr)\chi_{R}.
\]
当然, 这里有技术细节.

第一个技术困难是: 对每个 $x$, 有一个矩形 $R$ 使
\[
\mathcal{C}f(x)\approx\frac{1}{|R|}\int_{R}f.
\]
$R$ 对 $x$ 的依赖使得计算伴随成为不可能. 因此我们的首要任务是把 $x$ 局部化, 使一个矩形能对许多 $x$ 通用.

我们先把 $\mathcal{B}$ 中矩形允许的角度限制到 $(0,\frac{\pi}{4})$; 相应的极大函数仍记作 $\mathcal{C}$. 我们真正寻找的是 $\mathcal{C}$ 在其上为常数的区域. 但 $\mathcal{C}$ 是以 $|R|^{-1}\chi_R$ 为核的卷积算子的上确界; 如同第~\ref{chap:7} 章那样, $\mathcal{C}$ 的作用如同一个紧支集的卷积核, 因而独立地作用在等于其支集直径的方块上. 由于 $\mathcal{B}$ 中所有矩形的大小为 $a\times aN$, 我们把 $\mathbb{R}^2$ 分解成边长为 $aN$ 的互不相交的方块 $Q$ 的并; 设 $\tilde{Q}$ 是与之同心、边长为 $6aN$ 的同心方块. 若 $\mathcal{C}(f\chi_Q)(x)\neq 0$, 则存在 $R\in\mathcal{B}$ 使 $x\in R$ 且
\[
|R|^{-1}\int_{R}f(y)\,\chi_{Q}(y)\,dy>0.
\]
因此 $y\in R\cap Q$, 且
\[
|x-\operatorname{center}Q|\leq|x-y|+|y-\operatorname{center}Q|
\leq\operatorname{diameter}R+\operatorname{diameter}Q
\leq 2\sqrt{2}\,aN\leq 3aN.
\]
于是 $\operatorname{support}\mathcal{C}(f\chi_Q)\subset\tilde{Q}$. 由于每个 $x\in\mathbb{R}^2$ 属于唯一的一个 $Q$, 且至多属于 $30$ 个的 $\tilde{Q}$,
\begin{align*}
\|\mathcal{C}f\|_{2}^{2}
&=\Bigl\|\mathcal{C}\Bigl(\sum_Q f\chi_{Q}\Bigr)\Bigr\|_{2}^{2}
\leq\Bigl\|\sum_Q\mathcal{C}(f\chi_{Q})\Bigr\|_{2}^{2}
=\Bigl\|\sum_Q\mathcal{C}(f\chi_{Q})\,\chi_{\tilde{Q}}\Bigr\|_{2}^{2} \\
&\leq 30\sum_Q\bigl\|\mathcal{C}(f\chi_{Q})\bigr\|_{2}^{2}
\leq C(\log 3N)\sum_Q\|f\chi_{Q}\|_{2}^{2}
=C(\log 3N)\,\|f\|_{2}^{2}.
\end{align*}
于是, 只要证明
\[
\bigl\|\mathcal{C}(f\chi_{Q})\,\chi_{\tilde{Q}}\bigr\|_{2}\leq C(\log 3N)^{\frac{1}{2}}\,\|f\chi_{Q}\|_{2},
\]
定理即告完成.

为此, 取定一个方块 $Q$ 并把 $f\chi_Q$ 简记为 $f$.

至此, 我们已使 $\mathcal{C}f$ 在 $\tilde{Q}$ 之外为常数 (即零); 现在我们把 $\tilde{Q}$ 切成使 $\mathcal{C}$ 为常数的小块. 由于每个 $R$ 的大小为 $a\times aN$, 且诸 $R$ 的方向落在一个大小为 $\frac{\pi}{4}$ 的角形区域内, $\mathcal{C}$ 在垂直于 $R$ 方向的方向上的作用如同支集大小为 $a$. 把 $\tilde{Q}$ 细分为 $36N^2$ 个大小为 $a\times a$ 的方块 $Q_{ij}$. 我们断言: 对每个 $Q_{ij}$, 存在 $R_{ij}\in\mathcal{B}$ 使得
\[
\mathcal{C}f(x)\leq 2\,|\tilde{R}_{i j}|^{-1}\int_{\tilde{R}_{i j}}|f|
\]
对一切 $x\in Q_{ij}$ 成立. 若然, 我们便可以用
\[
\mathcal{C}f(x)\leq 2\sum_{ij}|\tilde{R}_{i j}|^{-1}\int_{\tilde{R}_{i j}}|f|
\]
控制 $\mathcal{C}f$, 然后计算右端算子的伴随.

若 $R\in\mathcal{B}$ 且 $R\cap Q_{ij}\neq\varnothing$, 则 $Q_{ij}$ 包含在 $\tilde{R}$ 中, 其中 $\tilde{R}$ 是与 $R$ 同心、直径八倍的矩形. 则
\[
|R|^{-1}\int_{R}|f|\leq 64\,|\tilde{R}|^{-1}\int_{\tilde{R}}|f|;
\]
且对 $x\in Q_{ij}$,
\[
\mathcal{C}f(x)\leq\sup_{R\in\mathcal{B}}64\,|\tilde{R}|^{-1}\int_{\tilde{R}}|f|
\leq C\sup_{R\in\mathcal{B}'}|R|^{-1}\int_{R}|f|,
\]
其中
\[
\mathcal{B}'=\{R \mid Q_{ij}\subset R;\ \operatorname{direction}R\in(0,\tfrac{\pi}{4}),\ \operatorname{size}R=8a\times 8aN\}.
\]
但 $R\in\mathcal{B}'$ 的上确界在 $Q_{ij}$ 上为常数, 于是我们可以找到 $R_{ij}\in\mathcal{B}$ 使 $\tilde{R}_{ij}\in\mathcal{B}'$, 且
\[
\frac{1}{2}\mathcal{C}f(x)\leq 64\,|\tilde{R}_{i j}|^{-1}\int_{\tilde{R}_{i j}}|f|
\]
对一切 $x\in Q_{ij}$ 成立.

到目前为止我们所做的, 是把极大函数局部化; 现在我们证明: 对矩形之交的预期估计足以完成证明. 于是, 设已选定 $f$ 并得到了诸 $\tilde{R}_{ij}$. 假设我们能证明
\[
\bigl|\tilde{R}_{i j}\cap\tilde{R}_{i' j'}\bigr|\leq\frac{C\,aN}{|j-j'|+1}.
\]
现在定义算子 $T:L^2(Q)\to L^2(\tilde{Q})$ 为
\[
Tg(x)=\sum_{ij}|\tilde{R}_{i j}|^{-1}\int_{R_{i j}}g(y)\,dy\;\chi_{Q_{i j}}(x).
\]
注意, 对我们取定的 $f$, 有 $\mathcal{C}f(x)\leq 2T(|f|)(x)$, 因此一旦得到 $T$ 的 $L^2$ 界, 定理即得证, 而这可以由 $T^{\ast}$ 的界得出. 但
\begin{align*}
\int T^{\star}h(y)\,g(y)\,dy&=\int h(x)\,Tg(x)\,dx \\
&=\sum_{ij}|\tilde{R}_{i j}|^{-1}\int_{R_{i j}}g(y)\,dy\int h(x)\,\chi_{Q_{i j}}(x)\,dx \\
&=\int\left(\sum_{ij}|\tilde{R}_{i j}|^{-1}\int_{Q_{i j}}h(x)\,dx\;\chi_{\tilde{R}_{i j}}(y)\right)g(y)\,dy.
\end{align*}
于是
\[
T^{\star}h(y)=\sum_{ij}|\tilde{R}_{i j}|^{-1}\int_{Q_{i j}}h(x)\,dx\;\chi_{\tilde{R}_{i j}}(y)
=C N^{-1}\sum_{ij}a_{ij}\,\chi_{\tilde{R}_{i j}}(y),
\]
其中 $a_{ij}=\int_{Q_{ij}}h$. 计算 $L^{2}$ 范数,
\begin{align*}
\int|T^{\star}h|^{2}
&=C N^{-2}\sum_{ij,\,i'j'}a_{ij}\,a_{i'j'}\,\bigl|\tilde{R}_{i j}\cap\tilde{R}_{i' j'}\bigr| \\
&\leq C a^{2}N^{-1}\sum_{ij,\,i'j'}a_{ij}\,a_{i'j'}\,\bigl(|j-j'|+1\bigr)^{-1} \\
&=C a^{2}N^{-1}\sum_{i,i'}\sum_{j}a_{ij}\sum_{j'}a_{i'j'}\,\bigl(|j-j'|+1\bigr)^{-1} \\
&\leq C a^{2}N^{-1}\sum_{i,i'}\left[\sum_{j}a_{ij}^{2}\right]^{\frac{1}{2}}
\left[\sum_{j}\Bigl(\sum_{k}a_{i',j-k}\bigl(|k|+1\bigr)^{-1}\Bigr)^{2}\right]^{\frac{1}{2}} \\
&\leq C a^{2}N^{-1}\log 6N\left(\sum_{i}\left[\sum_{j}a_{ij}^{2}\right]^{\frac{1}{2}}\right)^{2} \\
&\leq C a^{2}\log 6N\sum_{i,j}\left(\frac{1}{|Q_{ij}|}\int_{Q_{ij}}h\right)^{2} \\
&\leq C a^{2}\log 6N\sum_{i,j}\frac{1}{|Q_{ij}|}\int_{Q_{ij}}h^{2}
=C\log 6N\int h^{2}
=C\log 6N\,\|h\|_{2}^{2},
\end{align*}
正如所需.

为完成证明, 我们必须分析诸 $\tilde{R}_{ij}$ 的交; 我们把行编号为 $j$, 列编号为 $i$. 关于 $\tilde{R}_{ij}$ 的精确位置, 我们只知道它与 $Q_{ij}$ 相交. 两个 $\tilde{R}_{ij}$ 的最大可能的交出现在二者平行时; 此时测度如同 $a^2N$. 若两个矩形垂直, 则交如同 $a^2$.

\begin{figure}[htbp]
\centering
\begin{tikzpicture}[scale=0.65]
  % 图 15: 两个相交的长矩形
  \draw[rotate=6] (-3.4,-0.35) rectangle (3.4,0.35);
  \draw[rotate=-6] (-3.4,-0.35) rectangle (3.4,0.35);
  \draw[dashed] (-0.6,-0.75) -- (-0.1,0.5);
  \draw[dashed] (0.1,0.5) -- (0.6,-0.75);
  \node at (1.9,-1.4) {$a$};
  \node at (-2.2,1.3) {$a$};
\end{tikzpicture}
\caption{任意方向两相交矩形的重叠面积控制}
\label{fig:15}
\end{figure}

为控制这个交, 我们把它最大化: 在两个矩形确有相交的约束下, 平移诸 $\tilde{R}_{ij}$ 使二者之间的夹角最小.

我们还用最大的交来控制任何交; 即图~\ref{fig:15} 中所示的交. 按此做法, 这个交不随矩形左右滑动而变, 而由于列由 $i$ 编号, 我们所作的估计与 $i$ 无关. 于是为方便起见, 设 $i=i'$.

图~\ref{fig:15} 中交的面积容易计算; 每条虚线都是其中一个矩形的厚度, 即 $a$, 因为平行四边形的高为 $a$. 平行四边形的宽度满足 $\sin\theta=\frac{a}{N}$, 且
\[
\bigl|\tilde{R}_{i j}\cap\tilde{R}_{i' j'}\bigr|=\frac{a^{2}}{\sin\theta}
\]
(若 $\theta=0$, 这不准确; 修正将在下面给出).

给定条件 $\tilde{R}_{ij}\cap\tilde{R}_{i'j'}\neq\varnothing$ 与 $\tilde{R}_{ij}\cap Q_{i'j'}\neq\varnothing$, 角 $\theta$ 容易算出. 为使诸 $\tilde{R}_{ij}$ 的交最大, 角 $\theta$ 必须最小. 最小的 $\theta$ 出现在 $R_{ij}$ 与 $R_{i'j'}$ 恰好相触时, 如图~\ref{fig:16}.

\begin{figure}[htbp]
\centering
\begin{tikzpicture}[scale=0.65]
  % 图 16: 恰好相触的两个矩形
  \draw[rotate=90] (-3,-0.35) rectangle (3,0.35);
  \draw[rotate=75] (-3,-0.35) rectangle (3,0.35);
  \draw[dashed] (0,0) -- (0,2.6);
  \node at (0.7,2.2) {$\delta$};
  \node at (1.6,0.4) {$a$};
\end{tikzpicture}
\caption{临界相切矩形的夹角与边长关系}
\label{fig:16}
\end{figure}

则
\[
\sin\theta=\cos\Bigl(\frac{\pi}{2}-\theta\Bigr)\approx\frac{a\,|j-j'|}{\delta N}=\frac{|j-j'|}{N}.
\]
于是
\[
\bigl|\tilde{R}_{i j}\cap\tilde{R}_{i' j'}\bigr|\leq\frac{a^{2}N}{|j-j'|}.
\]
若 $j=j'$, 我们可以得到更好的估计, 故一般地,
\[
\bigl|\tilde{R}_{i j}\cap\tilde{R}_{i' j'}\bigr|\leq\frac{a^{2}N}{|j-j'|+1}.
\]
\end{proof}

\begin{remark}\label{rem:8.7}
上面对 $\mathcal{C}$ 的估计是最佳的; 我们现在证明这一事实. 取 $a=1$; 令当 $|x|\leq N$ 时 $f(x)=(1+|x|)^{-1}$, 其余情形 $f=0$. 则 $\|f\|_2\approx(\log N)^{\frac{1}{2}}$. 为估计 $\mathcal{C}f$, 先设 $|x|\leq N$. 取 $R_x$ 以 $\frac{x}{|x|}$ 为长边方向, 并选取它使 $R_x$ 的短边包含 $x$. 则 $|R_x|^{-1}\int_{R_x}f$ 是径向的, 我们在 $x=(|x|,0)$ 处求值:
\[
\frac{1}{N}\int_{-1}^{1}\int_{|x|-N}^{|x|}\frac{1}{1+|y|}\,dy_{1}\,dy_{2}
\geq C N^{-1}\int_{0}^{|x|}\frac{1}{1+y_{1}}\,dy_{1}
\geq C N^{-1}\log|x|.
\]
于是 $\|\mathcal{C}f\|_2^2\geq CN^{-2}\int_0^N\log^2 r\,dr\approx C\log^2 N$.
\end{remark}

\section{C\'ordoba 乘子定理}
\label{sec:8.3}

\begin{notation}\label{not:8.8}
设 $\phi$ 是 $\mathbb{R}^1$ 上满足当 $|x|>1$ 时 $\phi(x)=0$ 的光滑函数. 在 $\mathbb{R}^2$ 上定义 $\phi_\delta$ 为
\[
\phi_{\delta}(\xi)=\phi\bigl(\delta^{-1}(|\xi|-1)\bigr).
\]
设 $S_{\delta}$ 表示以 $\phi_{\delta}$ 为乘子的算子.
\end{notation}

\begin{theorem}\label{thm:8.9}
\index{Cordoba 乘子定理 (C'ordoba multiplier theorem)|emph}
\[
{}_{4}\|S_{\delta}\|\leq C\,|\log\delta|^{\frac{1}{4}}.
\]
\end{theorem}

\begin{proof}
以往那些乘子定理 (比如 Marcinkiewicz 定理) 的套路提示: 我们需要把 $\phi_{\delta}$ 的各个小块隔离开, 然后用一个极大函数控制各小块, 再控制极大函数. 许多部件已经就位; 例如极大函数的控制. 我们从定义小块开始.

直观上, 我们把 $\phi_\delta$ 想成厚度为 $\delta$ 的环形区域的特征函数. 我们希望彼此隔离的小块是不同的方向. 定理~\ref{thm:8.1} 中的函数将完成隔离, 但它不会把我们从 $S_\delta f$ 带到 $\bigl(\sum|S_jS_\delta f|^2\bigr)^{1/2}$; 相反, 它把我们从 $\sum|S_jf|^2$ 带到 $f$. 这意味着我们得自己做方向的隔离.

我们在~\ref*{ess:8.4} 中所做的计算提示, 若选取大小为 $\delta\times\delta^{\frac{1}{2}}$ 的矩形, 圆的曲率将提供隔离. 我们需要光滑地做这个分解; 取 $S^1$ 上的光滑函数 $\phi$, 设当 $|\theta|>\frac{1}{2}$ 时 $\phi(\theta)=0$, 并令
\[
\phi_{j}(\theta)=\phi\Bigl(\delta^{-\frac{1}{2}}\bigl(\theta-2\pi j[\delta^{\frac{1}{2}}]\bigr)\Bigr).
\]
我们再定义 $\mu_j(\xi)=\mu_j(r,\theta)=\phi_j(\theta)\,\phi_\delta(r)$, 并令 $T_j$ 表示相应的算子. 直观上, $\mu_j$ 是以 $j[\delta^{\frac{1}{2}}]$ 为中心、大小为 $\delta\times\delta^{\frac{1}{2}}$ 的矩形的特征函数. 如同~\ref*{ess:8.4} 中的讨论, 为避免不同方向之间的干扰, 我们需要把角度限制在圆的若干弧段上; 这相当于把所有最终估计乘一个因子 $8$. 于是我们可以假设 $0\leq j\leq[\delta^{-\frac{1}{2}}]/8$. 这些是我们的第一组分解.

开始估计:
\begin{align*}
\|S_{\delta}f\|_{4}^{4}
&=\Bigl\|\sum T_{j}f\Bigr\|_{4}^{4}
=\int\Bigl|\Bigl(\sum T_{j}f\Bigr)^{2}\Bigr|^{2} \\
&=\int\Bigl|\sum T_{j}f\,\sum T_{k}f\Bigr|^{2}
=\int\Bigl|\sum_{j,k}\widehat{T_{j}f}\star\widehat{T_{k}f}\Bigr|^{2}.
\end{align*}
由于每个 $x\in\mathbb{R}^2$ 至多属于 $9$ 个的 $\widehat{T}_jf\star\widehat{T}_kf=\mu_j\hat{f}\star\mu_k\hat{f}$ 的支集, 且这与 $\delta$ 无关, 最后的积分被
\begin{align*}
C\int\sum_{j,k}\bigl|\widehat{T_{j} f}\star\widehat{T_{k} f}\bigr|^{2}
&=C\int\sum_{j,k}|T_{j}f\,T_{k}f|^{2}
=C\int\sum_{j,k}|T_{j}f|^{2}\,|T_{k}f|^{2} \\
&=C\int\Bigl(\sum|T_{j}f|^{2}\Bigr)^{2}
=C\Bigl\|\Bigl(\sum|T_{j}f|^{2}\Bigr)^{\frac{1}{2}}\Bigr\|_{4}^{4}
\end{align*}
控制. 我们得到第一个分解:
\[
\|S_{\delta}f\|_{4}\leq C\Bigl\|\Bigl(\sum|T_{j}f|^{2}\Bigr)^{\frac{1}{2}}\Bigr\|_{4}.
\]
现在的任务是借助 $g$-函数从 $\sum|T_jf|^2$ 回到 $f$. 若令 $I_j$ 表示 $\phi_j$ 的支集在 $\xi_2$ 轴上的投影, 我们当然有 $\chi_{I_j}\mu_j=\mu_j$. 但注~\ref{rem:8.3} 对容许的 $I_j$ 提出了一些要求: 它们必须互不相交, 长度上下有界, 且彼此之间的距离有下界. 就目前而言, 我们取的 $\phi_j$ 是光滑的, 所以诸 $I_j$ 相互重叠. 这很容易克服: 把奇数与偶数的块分开处理, 把所有最终估计乘以 $2$. 我们甚至懒得改变记号. 于是诸 $I_j$ 互不相交, 长度为 $C\delta^{\frac{1}{2}}\sin(\frac{\pi}{2}-\theta_j)$, 其中 $\theta_j=2\pi j\delta^{\frac{1}{2}}$. 由于我们聪明地把 $0\leq\theta_j\leq\frac{\pi}{4}$ 限定了, 长度都被 $C\delta^{\frac{1}{2}}$ 控制. 当然, 诸 $I_j$ 之间的间距与缺失的块可比, 后者又与 $C\delta^{\frac{1}{2}}$ 可比. 于是我们可以用~\ref*{cor:8.2} 得到估计
\[
\Bigl\|\Bigl(\sum|S_{j}f|^{2}\Bigr)^{\frac{1}{2}}\Bigr\|_{4}\leq C\|f\|_{4},
\]
如果我们能走到那一步的话.

下一步是用极大函数控制 $T_jf$. 现在 $T_jf=\check{\mu}_j\star f$, 直观上 $\check{\mu}_j\star f$ 恰是在大小为 $\delta^{-1}\times\delta^{-\frac{1}{2}}$ 的矩形上对 $f$ 作平均. 设该矩形记为 $R_j$. 则 $T_jf\approx\frac{1}{|R_j|}\int_{R_j}f$, 且
\begin{align*}
\int|T_{j}f|^{2}\,g
&=\int\left|\int|R_{j}|^{-1}\chi_{R_{j}}(x-y)\,f(y)\,dy\right|^{2}g(x)\,dx \\
&\leq\int\left(\int\chi_{R_{j}}^{2}(x-y)\,dy\right)
\left(|R_{j}|^{-2}\int\chi_{R_{j}}^{2}(x-y)\,|f(y)|^{2}\,dy\right)g(x)\,dx \\
&=\int|R_{j}|\,|R_{j}|^{-2}\iint\chi_{R_{j}}(x-y)\,g(x)\,|f(y)|^{2}\,dy\,dx \\
&=\int|f(y)|^{2}\,|R_{j}|^{-1}\,\chi_{R_{j}}\star g\,dy.
\end{align*}
因此,
\[
\int\sum|T_{j}f|^{2}\,g=\int\sum|T_{j}S_{j}f|^{2}\,g
\leq\sum_j\int|S_{j}f|^{2}\,|R_{j}|^{-1}\chi_{R_{j}}\star g.
\]
若要沿 Marcinkiewicz 定理的证明走, 我们必须使 $|R_j|^{-1}\chi_{R_j}\star g$ 与 $j$ 无关. 我们对所有的 $j$ 取上确界来做到这一点; 这被一个对所有方向、大小为 $\delta^{-1}\times\delta^{-\frac{1}{2}}$ 的矩形所取的上确界控制. 这是一个 C\'ordoba 极大函数, 其中 $a=\delta^{-\frac{1}{2}}$, $N=\delta^{-\frac{1}{2}}$. 于是
\begin{align*}
\int\sum|T_{j}f|^{2}\,g
&\leq\int\sum|S_{j}f|^{2}\,\mathcal{C}(g) \\
&\leq\left(\int\Bigl(\sum|S_{j}f|^{2}\Bigr)^{2}\right)^{\frac{1}{2}}\|\mathcal{C}g\|_{2} \\
&\leq\Bigl\|\Bigl(\sum|S_{j}f|^{2}\Bigr)^{\frac{1}{2}}\Bigr\|_{4}^{2}\,\|\mathcal{C}g\|_{2} \\
&\leq C\|f\|_{4}^{2}\,\bigl(\log\delta^{-\frac{1}{2}}\bigr)^{\frac{1}{2}}\,\|g\|_{2}.
\end{align*}
对所有满足 $\|g\|_2=1$ 的 $g$ 取上确界:
\[
\Bigl\|\Bigl(\sum|T_{j}f|^{2}\Bigr)^{\frac{1}{2}}\Bigr\|_{4}^{2}
\leq C\,|\log\delta|^{\frac{1}{2}}\,\|f\|_{4}^{2},
\]
这就证明了定理.

当然, 还有很多技术细节. 我们并非在与一个矩形的特征函数作卷积. 另一方面, 卷积核属于 $\mathcal{S}$, 所以我们可以用 Schwartz 函数的想法来估计它. 首先, 注意诸 $\mu_j$ 都是同一个函数的旋转, 于是我们可以就其中一个 (记为 $K_0$) 估计卷积核. 直观上,
\[
K_{0}(x_{1},x_{2})=\delta^{\frac{3}{2}}\,\check{\phi}\bigl(\delta x_{1},\,\delta^{\frac{1}{2}}x_{2}\bigr),
\]
利用 $\phi\in\mathcal{S}$ 的想法,
\[
|K_{0}(x_{1},x_{2})|\leq C\,\delta^{\frac{3}{2}}\,|\delta x_{1}|^{-p}\,|\delta^{\frac{1}{2}}x_{2}|^{-q}
\]
对所有的 $p,q\geq 0$ 成立. 这意味着 $K_0$ 衰减得如此之快, 以至于我们可以用矩形特征函数之和来控制它. 令
\[
R_{i}^{0}=\{(x_{1},x_{2}) \mid |x_{1}|<2^{i}\delta^{-1},\ |x_{2}|<2^{i}\delta^{-\frac{1}{2}}\};
\]
我们断言
\[
|K_{0}(x)|\leq C\sum_{i}2^{-i}\,|R_{i}^{0}|^{-1}\chi_{R_{i}^{0}}(x).
\]
为证之, 取 $x\in R_{k+1}$; 设 $x$ 不在 $R_k$ 中. 若 $|x_1|>2^k\delta^{-1}$,
\[
|K_{0}(x)|\leq C\,\delta^{\frac{3}{2}}\,|\delta x_{1}|^{-3}\,|\delta x_{2}|^{0}
\leq C\,\delta^{\frac{3}{2}}\,\bigl|\delta\,2^{k}\delta^{-1}\bigr|^{-3}
=C\,\delta^{\frac{3}{2}}\,2^{-3k},
\]
而
\[
\sum_i 2^{-i}\,|R_{i}^{0}|^{-1}\chi_{R_{i}^{0}}(x)
=\sum_{i\geq k+1}2^{-i}\,\delta^{\frac{3}{2}}\,2^{-2i}
\approx C\,\delta^{\frac{3}{2}}\,2^{-3k}.
\]
证明的其余部分与之前一样:
\[
|T_{j}S_{j}f|\leq\sum_i 2^{-i}\,\Bigl||R_{i}^{j}|^{-1}\chi_{R_{i}^{j}}\star S_{j}f\Bigr|,
\]
且
\[
\Bigl\|\Bigl(\sum_j|T_{j}S_{j}f|^{2}\Bigr)^{\frac{1}{2}}\Bigr\|_{4}
\leq\sum_i 2^{-i}\,\Bigl\|\Bigl(\sum_j\bigl(|R_{i}^{j}|^{-1}\chi_{R_{i}^{j}}\star S_{j}f\bigr)^{2}\Bigr)^{\frac{1}{2}}\Bigr\|_{4}.
\]
由于 $2^{-i}$ 本身就可求和, 只需证明该 $L^{4}$ 范数关于 $i$ 一致有界. 这可以完全照旧去做, 只是矩形 $R_{i}^{j}$ 的尺寸为 $2^{i}\delta^{-1}\times2^{i}\delta^{-\frac{1}{2}}$. 这意味着我们取 C\'ordoba 极大函数, 其中 $a=2^{i}\delta^{-\frac{1}{2}}$, $N=\delta^{-\frac{1}{2}}$. 我们得到的范数 $\log N$ 与 $i$ 无关.

至此证明终于完成, 只差一个技术细节: $\mu_0$ 并非真的是方块的伸缩 $\psi\bigl(\delta^{-1}\xi_1,\delta^{-\frac{1}{2}}\xi_2\bigr)$. 真正的伸缩是在 $(r,\theta)$ 方向上, 这与 Fourier 变换配合得并不好. 现在我们展示如何克服这个问题. 我们的直观估计基于这样一个原理: 矩形的 Fourier 变换还是矩形; 它的精确版本考察 Schwartz 函数的导数. 首先,
\[
|x_{1}K_{0}|\leq\|x_{1}K_{0}\|_{\infty}\leq C\,\|D_{\xi_{1}}\mu_{0}\|_{1}.
\]
我们将证明
\[
\|D_{\xi_{1}}\mu_{0}\|_{1}\leq C\,\delta^{\frac{1}{2}},
\]
由此得到
\[
|K_{0}(x_{1},x_{2})|\leq C\,\delta^{\frac{3}{2}}\,|\delta x_{1}|^{-1},
\]
这正是我们所需要的典型结果.

首先,
\begin{align*}
D_{\xi_{1}}\mu&=\frac{\partial r}{\partial\xi_{1}}\,\frac{\partial\mu}{\partial r}
+\frac{\partial\theta}{\partial\xi_{1}}\,\frac{\partial\mu}{\partial\theta} \\
&=\frac{\xi_{1}}{(\xi_{1}^{2}+\xi_{2}^{2})^{\frac{1}{2}}}\,
\frac{\partial}{\partial r}\phi_{\delta}(r)\,\phi(\theta)
-\frac{\xi_{2}}{(\xi_{1}^{2}+\xi_{2}^{2})^{\frac{1}{2}}}\,
\phi_{\delta}(r)\,\frac{\partial}{\partial\theta}\phi(\theta).
\end{align*}
对 $\xi\in\operatorname{support}\mu_0$, 有 $|\xi_1|\leq 2$, $\xi_1^2+\xi_2^2\geq\frac{1}{2}$, $|\xi_2|\leq 2$, 且当然有
\[
\frac{\partial\phi_{\delta}}{\partial r}=\delta^{-1}\phi'\bigl(\delta^{-1}(r-1)\bigr),
\qquad
\frac{\partial\phi}{\partial\theta}=\delta^{-\frac{1}{2}}\phi'\bigl(\delta^{-\frac{1}{2}}\theta\bigr).
\]
于是,
\[
|D_{\xi_{1}}\mu_{0}|\leq C\,\delta^{-1}\bigl(|\phi'\psi|+\delta^{\frac{1}{2}}\,|\phi\psi'|\bigr)
\leq C\,\delta^{-1}.
\]
最后,
\[
\|D_{\xi_{1}}\mu_{0}\|_{1}\leq C\,\delta^{-1}\,|\operatorname{support}\mu_{0}|
=C\,\delta^{-1}\,\delta^{\frac{3}{2}}.
\]
这个估计以及类似的估计完成了定理的证明.
\end{proof}

\begin{corollary}\label{cor:8.10}
Bochner--Riesz 平均在最优范围内有界.
\end{corollary}

\begin{proof}
我们从一个单位分解 $1=\sum_i\phi_i$ 出发. 设诸 $\phi$ 光滑、径向, 径向支在 $[1-2^{-i},\,1-2^{-i-2}]$ 中, 且 $|D_r^p\phi_i|\leq C2^{ip}$. 令 $\mu_i=\phi_i\mu_\lambda$; 注意 $\mu_\lambda=\sum_i\mu_i$.

直观上, $\mu_i$ 近似是 $\phi_i$ 乘以 $\mu_\lambda$ 在 $\phi_i$ 支集上的最大值, 即 $2^{-\lambda i}$. 因此 $\mu_i$ 的 $L^4$ 范数应当如同 $2^{-\lambda i}\,\bigl|\log 2^{-i}\bigr|^{\frac{1}{4}}$. 若 $\lambda>0$, 这可以对和求和, 于是只要 $\lambda>0$, 所有的 Bochner--Riesz 算子 $T_\lambda$ 都在 $L^4$ 上有界. 内插现在给出最优范围的其余部分.

当然这只是直观; 为给出精确的证明, 我们重新审视~\ref*{thm:8.9} 的证明. 我们注意到, 取 $\delta=2^{-i}$ 时, $\mu_i$ 具有那里用到的 $\phi_\delta$ 的一切性质. 在照搬那个证明之前, 唯一需要验证的是对 $|K_0|$ 的估计. 我们将证明
\[
|K_{0}(x_{1},x_{2})|\leq C\,\delta^{\lambda}\,\delta^{\frac{3}{2}}\,|x_{1}\delta|^{-p}\,|x_{2}\delta^{\frac{1}{2}}|^{-q},
\]
而如同~\ref*{thm:8.9}, 这给出 $L^4$ 范数 $\delta^{\lambda}\,|\log\delta|^{\frac{1}{4}}$, 这正是我们需要的.

对 $K_{0}$ 的估计遵循与~\ref*{thm:8.9} 中相同的模式. 一个典型的估计要求
\[
\left|\frac{\partial}{\partial r}\mu_{i}\right|\leq\delta^{\lambda}\,\delta^{-1}.
\]
但
\[
\left|\frac{\partial}{\partial r}\mu_{i}\right|
\leq\left|\frac{\partial}{\partial r}\phi_{i}\,\mu_{\lambda}\right|
+\left|\frac{\partial}{\partial r}\mu_{\lambda}\,\phi_{i}\right|.
\]
现在 $\bigl|\frac{\partial}{\partial r}\phi_i\bigr|\leq C2^i$, 而若 $r\in\operatorname{support}\phi_i$, 则 $\bigl|\mu_\lambda(r)\bigr|\leq\delta^\lambda$. 类似地,
\[
\left|\frac{\partial}{\partial r}\mu_{\lambda}\right|\leq C\,\bigl|(1-r^{2})^{\lambda-1}\bigr|,
\]
而对 $r\in\operatorname{support}\phi_i$, 有 $1-\delta<r<1-\frac{\delta}{4}$, 故
\[
\left|\frac{\partial}{\partial r}\mu_{\lambda}\right|\leq C\,\delta^{\lambda-1}.
\]
其余的估计照~\ref*{thm:8.9} 的证明进行.
\end{proof}

\section{重访限制性}
\label{sec:8.4}

在本节中, 我们将证明 $\mathbb{R}^2$ 中限制性在最优范围内的 $L^p$--$L^q$ 有界性. 我们从标准模型开始; 我们考察一个环形区域而非单位圆. 注意
\[
\int_{S^{1}}\bigl|\hat{f}(\theta)\bigr|^{q}\,d\theta
=\lim_{\delta\to0}\frac{1}{2\delta}\int_{1-\delta}^{1+\delta}\int_{S^{1}}\bigl|\hat{f}(r,\theta)\bigr|^{q}\,r\,dr\,d\theta,
\]
于是我们可以由环形区域恢复出限制性. 设 $0\leq\phi\leq 1$ 如~\ref*{not:8.8} 中那样, 且 $\chi_{(1-\delta,1+\delta)}(|\xi|)\leq\phi_\delta(\xi)$. 只需证明
\[
\frac{1}{2\delta}\int_{1-\delta}^{1+\delta}\int_{S^{1}}\bigl|\phi_{\delta}(r)\,\hat{f}(r,\theta)\bigr|^{q}\,r\,dr\,d\theta\leq C\|f\|_{p}^{q},
\]
亦即
\[
\bigl\|\phi_{\delta}\hat{f}\bigr\|_{L^{q}(S^{1})}\leq\delta^{\frac{1}{q}}\,\|f\|_{p}.
\]
设 $Rf=\phi_{\delta}\hat{f}$; 则
\[
R^{\star}g=\widehat{\phi_{\delta}\,g}=\check{\phi}_{\delta}\star g=S_{\delta}g.
\]

\begin{theorem}\label{thm:8.11}
\index{离散分数次积分 (discrete fractional integration)}
设 $q=\frac{p'}{3}$. 则
\begin{align*}
\|S_{\delta}g\|_{p'}&\leq C\,\delta^{\frac{1}{q}}\,\|g\|_{q'},\qquad 1\leq p<\frac{4}{3}, \\
\|S_{\delta}g\|_{4}&\leq C\,\delta^{\frac{3}{4}}\,|\log\delta|^{\frac{1}{4}}\,\|g\|_{4}.
\end{align*}
\end{theorem}

\begin{proof}
若 $p\leq\frac{4}{3}$, 则 $r=\frac{p'}{2}\geq 2$, 且
\begin{align*}
\|S_{\delta}g\|_{p'}^{p'}
&=\|S_{\delta}g\cdot S_{\delta}g\|_{r}^{r}
\leq\bigl\|\widehat{S_{\delta}g}\star\widehat{S_{\delta}g}\bigr\|_{r'}^{r}
=\|(\phi_{\delta}\check{g})\star(\phi_{\delta}\check{g})\|_{r'}^{r} \\
&=\Bigl\|\sum_{j,k}(\phi_{\delta}\phi_{j}\check{g})\star(\phi_{\delta}\phi_{k}\check{g})\Bigr\|_{r'}^{r} \\
&\leq C\left(\sum_{j,k}\int\bigl|(\phi_{\delta}\phi_{j}\check{g})\star(\phi_{\delta}\phi_{k}\check{g})\bigr|^{r'}\right)^{\frac{r}{r'}};
\end{align*}
最后一步用了~\ref*{thm:8.9} 证明中的殆正交性. 我们稍后将证明: 若 $1\leq r'\leq\infty$, 则
\[
\bigl\|(\phi_{\delta}\phi_{j}f)\star(\phi_{\delta}\phi_{k}f)\bigr\|_{r'}
\leq C\left[\frac{\delta^{\frac{3}{2}}}{|j-k|+1}\right]^{\frac{1}{r}}\|\phi_{j}f\|_{r'}\,\|\phi_{k}f\|_{r'}.
\]
这给我们一个界
\[
C\delta^{\frac{3}{2}}\left(\sum_{k}\|\phi_{k}f\|_{r'}^{r'}\left[\sum_{j}\frac{\|\phi_{j}f\|_{r'}^{r'}}{(|j-k|+1)^{\frac{r'}{r}}}\right]\right)^{\frac{r}{r'}}.
\]
令 $b_j=\|\phi_jf\|_{r'}^{r'}$. 则我们需要控制
\begin{align*}
&\,C\delta^{\frac{3}{2}}\left(\sum_{k}b_{k}\left[\sum_{j}\frac{b_{j}}{(|j-k|+1)^{\frac{r'}{r}}}\right]\right)^{\frac{r}{r'}} \\
&\leq C\delta^{\frac{3}{2}}\left(\sum_{k}b_{k}^{s}\right)^{\frac{r}{r's}} \\
&\qquad\times\left(\sum_{k}\left[\sum_{j}b_{j}\,(|j-k|+1)^{1-r'}\right]^{s'}\right)^{\frac{r}{r's'}}.
\end{align*}
若 $p<\frac{4}{3}$, 则 $2-r'>0$, 我们选取 $\frac{1}{s'}=\frac{1}{2}-(2-r')$, 并用离散分数次积分理论 (Hardy--Littlewood--P\'olya~\cite{hardy-littlewood-polya}) 得到界
\[
C\delta^{\frac{3}{2}}\left(\sum_{k}b_{k}^{s}\right)^{\frac{2r}{r's}}.
\]
但 $r's=q'$, 于是我们可以用 H\"older 不等式从 $r'$ 过渡到 $q'$ 来估计 $\|\phi_jf\|_{r'}^{r'}$. 这给出
\begin{align*}
C\delta^{\frac{3}{2}}\,\delta^{\frac{3}{2}\left(\frac{q'}{r'}-1\right)\frac{2r}{q'}}\left(\sum_{k}\int|\phi_{k}f|^{q'}\right)^{\frac{2r}{q'}}
&=C\delta^{3}\,\|f\|_{q'}^{2r}.
\end{align*}
这正是我们需要的.

另一方面, 若 $p=\frac{4}{3}$, 则 $r'=r=s=s'=2$, 于是我们可以用估计
\[
\left(\sum_{k}\left[\sum_{j}\frac{b_{j}}{|j-k|+1}\right]^{2}\right)^{\frac{1}{2}}
\leq\left(\sum_{k}(|j-k|+1)^{-1}\right)\left(\sum_{j}b_{j}^{2}\right)^{\frac{1}{2}}.
\]
这给我们一个界
\begin{align*}
C\delta^{\frac{3}{2}}\,|\log\delta|\,\sum_{j}b_{j}^{2}
&=C\delta^{\frac{3}{2}}\,|\log\delta|\sum_{j}\|\phi_{j}f\|_{2}^{4}
\leq C\delta^{\frac{3}{2}}\,|\log\delta|\,\|f\|_{4}^{4}.
\end{align*}
为证明估计
\[
\bigl\|(\phi_{\delta}\phi_{j}f)\star(\phi_{\delta}\phi_{k}f)\bigr\|_{r'}
\leq C\left[\frac{\delta^{\frac{3}{2}}}{|j-k|+1}\right]^{\frac{1}{r}}\|\phi_{j}f\|_{r'}\,\|\phi_{k}f\|_{r'},
\]
我们使用双线性内插. 若 $B(f,g)=f\star g$ 是双线性映射
\[
B:L^{p}(S_{j})\times L^{p}(S_{k})\to L^{p}(S_{j}+S_{k}),
\]
我们将从 $p=1$ 与 $p=\infty$ 对 $B$ 作内插. 则
\begin{align*}
\|(\phi_{\delta}\phi_{j}f)\star(\phi_{\delta}\phi_{k}f)\|_{1}
&\leq\|\phi_{\delta}\phi_{j}f\|_{1}\,\|\phi_{\delta}\phi_{k}f\|_{1} \\
&\leq\|\phi_{\delta}\|_{\infty}^{2}\,\|\phi_{j}f\|_{1}\,\|\phi_{k}f\|_{1}
=\|\phi_{j}f\|_{1}\,\|\phi_{k}f\|_{1}.
\end{align*}
在 $p=\infty$ 的情形,
\[
\|(\phi_{\delta}\phi_{j}f)\star(\phi_{\delta}\phi_{k}f)\|_{\infty}
\leq\|\phi_{\delta}\|_{\infty}^{2}\,\|\phi_{j}f\|_{\infty}\,\|\phi_{k}f\|_{\infty}\,\|\phi_{j}\star\phi_{k}\|_{\infty}.
\]
惯常的``平行四边形面积''技巧给出
\[
\|\phi_{j}\star\phi_{k}\|_{\infty}\leq C\,\frac{\delta^{\frac{3}{2}}}{|j-k|+1},
\]
于是内插证明了结果.
\end{proof}

\begin{remark}\label{rem:8.12}
\begin{enumerate}[label=\alph*)]
\item 已经知道限制性定理在 $p=\frac{4}{3}$ 处失效; 因子 $|\log\delta|$ 的出现解释了这一失效.
\item 容易证明这个估计是最优的; 设 $S_j=\operatorname{support}\phi_j$, 并令 $g\equiv 1$. 则 $\|g\|_4^4=\delta$, 且
\[
\|S_{\delta}\hat{g}\|_{4}^{4}=\int\Bigl|\sum_{j,k}\chi_{S_{j}}\star\chi_{S_{k}}\Bigr|^{2}
\geq\sum_{j,k}\int\bigl|\chi_{S_{j}}\star\chi_{S_{k}}\bigr|^{2}.
\]
但 $\chi_{S_{j}}\star\chi_{S_{k}}\geq C\frac{\delta^{\frac{3}{2}}}{|j-k|+1}$ 在 $\chi_{S_{j}}\star\chi_{S_{k}}$ 的支集的至少一半上成立, 而后者的测度至少为 $C\delta^{\frac{3}{2}}(|j-k|+1)$. 于是
\[
\|S_{\delta}\hat{g}\|_{4}^{4}
\geq C\delta^{3}\,\delta^{\frac{3}{2}}\sum_{j,k}\bigl(|j-k|+1\bigr)^{-1}
=C\delta^{3}\,\delta^{\frac{3}{2}}\,\delta^{-\frac{1}{2}}\,|\log\delta|.
\]
\item 这些计算与我们``除非 $\widehat{d\sigma}\in L^p$, 否则限制性必失效''的直观有关. 如上, 令 $\hat{g}\equiv 1$; 则 $S_{\delta}\hat{g}=\hat{\phi}_{\delta}$, 且上面的估计表明
\[
\|\hat{\phi}_{\delta}\|_{4}\approx C\,\delta^{\frac{3}{4}}\,|\log\delta|^{\frac{1}{4}},
\]
而这个估计是最优的. 由于 $d\sigma=\lim_{\delta\to0}\delta^{-1}\phi_{\delta}$, 我们再次看到 $\widehat{d\sigma}$ 不属于 $L^4$. 这也可以与曲率性质联系起来; 我们需要知道没有一个 $\xi$ 落在多于 $9$ 个的 $S_j+S_k$ 之中, 而且
\[
\bigl|\chi_{S_{j}}\star\chi_{S_{k}}\bigr|\leq\frac{|S_{j}|}{|j-k|+1}.
\]
这两个性质对直线来说都惨遭失败. 例如, 设 $S=[-1,1]\times\{0\}$, $D_j=\bigl[j\delta^{\frac{1}{2}},(j+1)\delta^{\frac{1}{2}}\bigr]\times[-\delta,\delta]$, $-\delta^{-\frac{1}{2}}\leq j\leq\delta^{-\frac{1}{2}}$. 则原点落在所有的 $D_j+D_k$ 之中, 于是重叠是巨大的. 类似地, $\|\chi_{D_j}\star\chi_{D_k}\|_{\infty}=|D_j|$, 于是我们得不到任何随 $j-k$ 的衰减.
\end{enumerate}
\end{remark}

\section{评注与推广}
\label{sec:8.5}

\noindent (定理~\ref{thm:8.1}):\quad 这个结果取自 C\'ordoba~\cite{cordoba-kakeya}, 不过 C\'ordoba 说明它更早便已为人所知. 最近有人证明, 对间距与一致性的要求可以完全去掉; 结果对任意互不相交的区间都成立.

\noindent (注~\ref*{ess:8.4}):\quad 殆正交性最早出现在 C.~Fefferman~\cite{fefferman-spherical} 中; 一个较弱的版本被用于 C.~Fefferman 与 E.~M.~Stein~\cite{fefferman-stronglysingular} 的一个限制性定理. 在后者中, 人们用到这样一个事实: 每个长度与方向恰好对应圆上至多两条弦. $S_{j}+S_{k}$ 中的一个点对应于稍微加厚的圆上的一个方向与一个长度.

\noindent (定理~\ref{thm:8.6}):\quad 这里的证明取自 C\'ordoba~\cite{cordoba-kakeya}. 通过伴随控制极大函数的想法最早由 C\'ordoba 发展; 参见~\cite{deguzman} 的附录. 原始想法可追溯到 Kolmogorov--Seliverstov--Plessner 定理; 参见 Zygmund~\cite{zygmund}.

关于这些极大函数有界性的其他证明, 见 Wainger~\cite{wainger-applications}.

\noindent (定理~\ref{thm:8.9}):\quad 这里的证明取自 C\'ordoba~\cite{cordoba-bochnerriesz}. 对 $T_{\delta}$ 的控制最早由 C\'ordoba~\cite{cordoba-kakeya} 建立, 不过在那种情形必须先控制一个强得多的极大函数, 即``Kakeya 极大函数''. 它取遍所有边长之比为 $N$ 的矩形上的平均. 这些想法可以追溯到 C.~Fefferman 的文章~\cite{fefferman-spherical}.

\noindent (推论~\ref{cor:8.10}):\quad 这个结果最早由 L.~Carleson 与 P.~Sj\"olin~\cite{christ} 证明. 证明的途径是把卷积核分解成小块, 如同~\ref*{thm:6.9} 中那样; 然后以直接的方式遵循~\ref*{sec:6.1} 中的直观. Carleson 与 Sj\"olin 不是证明 $T_{\lambda}$ 的各部分如同限制算子, 而是先证明 Plancherel 型的估计, 再证明 $L^{4}$ 型的估计.

~\ref*{cor:8.10} 的第二个证明由 C.~Fefferman~\cite{fefferman-spherical} 给出; 从历史上看, 它出现在 Kakeya 集被引进之后, 并且是第一篇把 Kakeya 集思想融入正面结果的论文.

\noindent (定理~\ref{thm:8.11}):\quad 当 $p\neq\frac{4}{3}$ 时这个结果的第一个证明归功于 C.~Fefferman 与 E.~M.~Stein~\cite{fefferman-spherical}; C\'ordoba~\cite{cordoba-flp} 给出了一个几何证明. 由环形区域出发的结果是~\cite{tomas2} 中证明的简化版本.


% ---------- 参考文献 ----------
% 模板提供 \makebibliography：\chapter{参考文献} + \printbibliography[heading=none]
% \chapter 自带 \cleardoublepage，此处不再前置，以免多出一张空白页。
% 印刷保持编号章与页眉; 但 PDF 书签要求不带章号: 先手工放置根级书签,
% 再局部令 \addcontentsline 失效, 使该 \chapter 不生成“第九章”自动书签.
% \cleardoublepage 须在书签之前, 使书签锚点落在文献首页; 其后 \chapter
% 自带的清页在空页上是空操作, 不会多出空白页.
\begingroup
\makeatletter
\long\def\addcontentsline#1#2#3{}%
\makeatother
\cleardoublepage
\pdfbookmark[0]{参考文献}{chap:bib}%
\makebibliography
\endgroup

% ---------- 索引 ----------
\printindex

\end{document}
